\documentclass[11pt,letterpaper]{article}

\usepackage[T1]{fontenc}
\usepackage[utf8]{inputenc}
\usepackage{lmodern}
\usepackage{microtype}
\usepackage[margin=1in,headheight=14pt]{geometry}
\usepackage{amsmath,amssymb,amsthm,mathtools,bm}
\usepackage{mathrsfs}
\usepackage{centernot}
\usepackage{xcolor}
\usepackage{enumitem}
\usepackage{booktabs,array,tabularx,longtable}
\usepackage{fancyhdr}
\usepackage{hyperref}

\definecolor{PartColor}{HTML}{223A5E}
\definecolor{AccentColor}{HTML}{6B3F8C}
\definecolor{RuleColor}{HTML}{B8C1CC}

\hypersetup{
 colorlinks=true,
 linkcolor=PartColor,
 urlcolor=AccentColor,
 citecolor=PartColor,
 pdftitle={Stratified Relational Workspace Theory},
 pdfauthor={Komorebi P.},
 pdfsubject={A core-transfer framework with decomposition- and specification-robust certification for distributed cognition, collaboration, and epistemic systems},
 pdfkeywords={relational workspace, distributed cognition, human-AI collaboration, multi-agent systems, transfer Laplacian, coarse-graining, decomposition robustness, specification robustness, identifiability, partial identification, finite-sample certification, causal states, reconstruction, agency, self-maintenance},
 pdfcreator={LaTeX}
}

\pagestyle{fancy}
\fancyhf{}
\fancyhead[L]{\small Stratified Relational Workspace Theory}
\fancyhead[R]{\small Core--Transfer Theory}
\fancyfoot[C]{\thepage}
\renewcommand{\headrulewidth}{0.4pt}
\renewcommand{\footrulewidth}{0pt}

\setlength{\parindent}{0pt}
\setlength{\parskip}{0.55em}
\setlist[itemize]{leftmargin=1.55em,itemsep=0.2em,topsep=0.3em}
\setlist[enumerate]{leftmargin=1.75em,itemsep=0.2em,topsep=0.3em}
\allowdisplaybreaks[3]
\setcounter{tocdepth}{2}
\setcounter{secnumdepth}{0}

\newcommand{\Workspace}{\mathbf W}
\newcommand{\SRW}{\operatorname{SRW}}
\newcommand{\CE}{\operatorname{CE}}
\newcommand{\Risk}{\mathcal R}
\newcommand{\Hol}{\operatorname{Hol}}
\newcommand{\Path}{\operatorname{Path}}
\newcommand{\Ret}{\operatorname{Ret}}
\newcommand{\Dep}{\operatorname{Dep}}
\newcommand{\Out}{\operatorname{Out}}
\newcommand{\COut}{\operatorname{COut}}
\newcommand{\CIn}{\operatorname{CIn}}
\newcommand{\Instability}{\operatorname{Instability}}
\newcommand{\Rec}{\operatorname{Rec}}
\newcommand{\Pass}{\operatorname{Pass}}
\newcommand{\indep}{\perp\!\!\!\perp}
\newcommand{\dd}{\,\mathrm d}
\DeclareMathOperator{\ran}{ran}
\DeclareMathOperator{\spec}{spec}
\DeclareMathOperator{\End}{End}
\DeclareMathOperator{\Fix}{Fix}
\DeclareMathOperator{\Comm}{Comm}
\DeclareMathOperator{\ev}{ev}
\DeclareMathOperator{\im}{im}
\DeclareMathOperator{\Nat}{Nat}
\DeclareMathOperator{\KL}{KL}
\DeclareMathOperator{\TV}{TV}
\DeclareMathOperator{\pr}{pr}
\DeclareMathOperator*{\argmax}{arg\,max}
\DeclareMathOperator*{\argmin}{arg\,min}

\newcommand{\theorypart}[2]{%
 \clearpage
 \phantomsection
 \addcontentsline{toc}{section}{Part #1. #2}%
 {\centering
 {\large\bfseries\color{AccentColor} PART #1\par}
 \vspace{0.35em}
 {\LARGE\bfseries\color{PartColor} #2\par}
 }
 \vspace{0.65em}
 {\color{RuleColor}\hrule height 0.8pt}
 \vspace{1.15em}
}

\newcommand{\topic}[2]{%
 \phantomsection
 \addcontentsline{toc}{subsection}{#1. #2}%
 \vspace{0.8em}
 {\large\bfseries\color{PartColor} #1. #2\par}
 \vspace{0.15em}
}

\newcommand{\subtopic}[1]{%
 \vspace{0.45em}
 {\bfseries\color{AccentColor} #1\par}
}

\theoremstyle{plain}
\newtheorem{theorem}{Theorem}
\newtheorem{proposition}[theorem]{Proposition}
\newtheorem{lemma}[theorem]{Lemma}
\newtheorem{corollary}[theorem]{Corollary}
\theoremstyle{definition}
\newtheorem{definition}[theorem]{Definition}
\newtheorem{example}[theorem]{Example}
\newtheorem{remark}[theorem]{Remark}

\newcommand{\centerbox}[1]{%
\begin{center}
\fbox{\parbox{0.90\textwidth}{\centering #1}}
\end{center}
}


\begin{document}
\pagenumbering{gobble}

\begin{titlepage}
\thispagestyle{empty}
\centering
\vspace*{1.2in}
{\Huge\bfseries\color{PartColor} Stratified Relational\\[0.18em] Workspace Theory\par}
\vspace{0.55in}
{\Large A Core--Transfer Framework for Distributed Cognition,\\ Collaboration, and Epistemic Systems\par}
\vspace{0.8in}
{\large Komorebi P.\par}
\vspace{0.65in}
{\large\bfseries Foundations, Dynamics, Agency, and Empirical Certification\par}
\vfill
{\color{RuleColor}\rule{0.70\textwidth}{0.8pt}\par}
\vspace{0.25in}
{\small August 2026\par}
\end{titlepage}
\clearpage
\pagenumbering{roman}

\begin{abstract}
Collaborative reasoning increasingly spans people, computational models, documents, software,
instruments, and persistent records. Yet the existence of many interacting components does not,
by itself, establish a shared workspace, and successful group performance does not show that the
operative state is genuinely distributed rather than copied, centralized, or merely correlated.
Stratified Relational Workspace Theory develops a formal answer to this problem. A realization
is represented by heterogeneous local state spaces connected by transfer maps. Its exact shared
core is the space of global compatible sections, equivalently the fixed subspace of relational
holonomy when the transfers are invertible. A transfer Laplacian identifies approximate shared
directions and provides spectral stability, gauge covariance, and pathwise leakage bounds.

The structural theory is joined to a dynamic and causal account. Canonical predictive states
determine whether a proposed workspace retains the conditional law relevant to a declared
future task. Feasible interventions, coalition risk, matched relational nulls, and reconstruction
tests distinguish causally operative distributed organization from archives, replicated memory,
and centralized controllers. Further results address persistence under changing realizations,
stable reconstruction after component loss, identity under replacement, local-to-global gluing,
hierarchical composition, learned transfer structure, boundary recovery, goal-conditioned action,
and certificate-directed repair. Because distributedness can be destroyed by merging components
or manufactured by splitting a centralized state into complementary shares, the theory also treats
decomposition as part of the claim. Faithful coarse-graining preserves exact structural cores, while
family-robust certification requires every gate to retain a positive margin across a declared set of
admissible decompositions. The same discipline is applied to the complete scientific specification.
Targets, losses, horizons, interventions, nulls, and thresholds are made explicit, and the theory develops
specification envelopes, transport radii, experimental-equivalence classes, partial-identification bounds,
and uniform finite-sample certification. A certificate is therefore distinguished from a favorable
analysis: it may be pointwise valid, specification-robust, sign-identified, or statistically unresolved.
The resulting rule is noncompensatory. Structural coherence, predictive alignment, causal efficacy,
relational irreducibility, non-collapse, reconstructive continuity, provenance, and self-maintenance are
separate claims, and failure of one cannot be hidden by success on another.

Exact finite examples show how the same framework applies to human collaboration, human--AI
research loops, and multi-agent control. A review of adjacent work locates the individual
mathematical ingredients in established traditions including distributed cognition, local systems
and sheaves, spectral perturbation, sufficient statistics, structural causal models, frame reconstruction,
system identification, and feedback control. The candidate contribution is therefore
not a new name for any one of these constructions. It is their integration into a single falsifiable
account of persistent relational workspaces and the matched near misses that such an account
must reject. The framework does not infer phenomenal consciousness, personhood, or moral
status from coherence, agency, or self-maintenance.
\end{abstract}

\tableofcontents
\clearpage
\pagenumbering{arabic}

\section*{Introduction}
\addcontentsline{toc}{section}{Introduction}

Scientific and technical work rarely resides in a single mind or device. A proof may be divided
among a conjecture, a critic's counterexample, a notebook, a computer algebra system, and a
verification script. A human--AI research process may distribute objectives, formal manipulations,
execution, and persistent memory across different components. Multi-agent systems divide sensing,
planning, communication, and repair in still other ways. In each case, it is natural to speak of a
shared workspace. The phrase is useful, but without a formal criterion it can conceal several
fundamentally different organizations.

A common archive is not necessarily a workspace. Nor is a group whose apparent coordination is
produced by one complete controller surrounded by decorative channels. Two systems can contain
the same local information while differing in how that information is related, transported, and used.
Likewise, a system can perform well without preserving its operative organization through interruption,
replacement, or component loss. The central problem is therefore not simply whether information is
shared. It is whether there is a persistent relational state whose structure and causal capacities cannot
be reduced to any single declared component.

The philosophical background is supplied by extended and distributed cognition, transactive memory,
collective intelligence, and theories of group agency \cite{Hutchins1995,ClarkChalmers1998,Wegner1987,ListPettit2011}.
Those traditions motivate the possibility that cognitive organization may span people and artifacts.
The theory asks a narrower mathematical question: what would have to be measured and proved
before such an organization should be called a relational workspace? The answer must distinguish
physical realization from abstract identity, local representation from globally compatible content,
and statistical association from causal operation.

The static theory begins with an interaction graph whose vertices carry heterogeneous local state spaces
and whose edges carry transfer maps. Each local state is decomposed into a transferable core and a
substrate-specific residual. Exact shared content is represented by global compatible sections. When
core transfers are invertible, the same object appears as the fixed subspace of the holonomy representation.
Approximate shared content is extracted as a low-energy spectral subspace of a transfer Laplacian.
This formulation supplies explicit obstruction, stability, gauge, and leakage calculations rather than
relying on an informal appeal to shared context.

Structure alone is insufficient. A proposed workspace must also retain the predictive state relevant to
a declared future, respond under feasible intervention, outperform matched relational nulls, and resist
reduction to an individual node or a local-storage surrogate. Persistence introduces further requirements:
the workspace must remain trackable as its realization changes, preserve nontrivial content, support
reconstruction after loss, and maintain an identity profile under admissible replacement. Local workspaces
must glue without hidden incompatibility, and hierarchical organization must not conceal centralization.
Formation claims are evaluated separately from operational claims so that a scaffolded but functional
workspace is not dismissed, while an endogenously learned but causally irrelevant pattern is not promoted.
A separate decomposition analysis asks which claims survive faithful coarse-graining and plausible
alternative interface descriptions. This is necessary because arbitrary merging can conceal relationality,
while arbitrary splitting can manufacture it. The analysis is then extended from decomposition to the full
scientific specification. A workspace claim depends on its declared target, loss, horizon, intervention
family, null panel, and thresholds. The theory therefore distinguishes pointwise validity from robustness
across defensible specifications and distinguishes both from identifiability under the experiments that are
actually available. This prevents a favorable target or threshold from being mistaken for a general result
and prevents an observationally indistinguishable hidden controller from being ruled out by assertion.

The theory extends to domain-relative agency only after these lower-level conditions are in place. Goals
must persist without repeated external rewriting, counterfactual values must guide feasible actions, damage
must be detected through the workspace, and repairs must improve a matched no-repair trajectory while
preserving identity. These criteria define an engineering notion of workspace agency. They neither assert
nor deny phenomenal consciousness. Consciousness, personhood, and moral status remain separate questions
that require independent evidence.

Most of the mathematical tools used here are established. Inverse limits, holonomy, connection Laplacians,
perturbation theory, sufficient statistics, causal models, frame reconstruction, system identification, and
feedback control all have substantial prior literatures. The contribution sought here is integrative: a
noncompensatory framework that places these tools in one common architecture and specifies the near-miss
systems that each layer must reject. Exact examples provide calibration, while the motivating human--AI
collaboration is treated conservatively as a candidate case rather than as a certified instance.

The paper proceeds from exact structure to approximate structure, dynamic operation, relational
irreducibility, persistence, reconstruction, identity, gluing, formation, boundary discovery, agency,
decomposition robustness, specification robustness, identifiability, and empirical certification. It then
gives worked cases, situates the framework relative to adjacent work, and closes with its philosophical
consequences, limitations, and open problems.

\section*{Notation and Conventions}
\addcontentsline{toc}{section}{Notation and Conventions}

A concrete realization is denoted by \(\mathfrak R\). Its interaction graph is \(G=(V,E,s,t)\), with local
state spaces \(X_i\) at vertices and transfer maps \(T_e:X_{s(e)}\to X_{t(e)}\) on directed edges. When a
core--residual decomposition is used, \(X_i=C_i\oplus R_i\), and \(H_e\) denotes the induced core transfer.
The exact structural workspace is the compatible-section space \(\Gamma(H)\); the transfer discrepancy is
\(D_H\), the transfer Laplacian is \(L_H=D_H^*D_H\), and \(\Gamma_\varepsilon(H)\) denotes a selected
low-energy workspace subspace. Dynamic workspace states are written \(W_t\), while \(Z_t\) denotes the
canonical predictive state associated with a declared future variable \(Y_t^+\).

Unless stated otherwise, the structural results use finite-dimensional real or complex Hilbert spaces and
finite graphs. Probabilistic statements are made on standard Borel spaces when regular conditional laws are
required. All approximate conclusions are indexed by explicit tolerances, spectral gaps, intervention
families, loss functions, or null classes. A claim of relational irreducibility is always relative to a declared
decomposition of the realization into components. The symbol \(\mathcal D\) denotes one such
decomposition, and \(\mathfrak D\) denotes an admissible family used for robustness analysis. The complete
scientific specification is
\[
\varsigma=(\mathcal D,Y^+,\ell,\tau,\mathcal I,\mathcal N,\Theta),
\]
and \(\mathfrak S\) denotes an admissible family of specifications. A gate margin is written
\(m_q(W;\varsigma)\), with positive sign indicating passage under the declared direction and threshold.

The terms \emph{workspace}, \emph{agency}, \emph{consciousness}, \emph{personhood}, and \emph{moral status}
are not interchangeable. A structural workspace may fail every operational test. An operational workspace
may lack persistence. A persistent relational workspace may lack goals or repair. A workspace agent, as
defined here, is a domain-relative functional organization and is not thereby a conscious subject. Definitions,
propositions, theorems, conjectures, and empirical tests are identified by their ordinary mathematical form;
all proved statements remain conditional on the assumptions given with them.

\theorypart{I}{Foundational Thesis}

\topic{1}{Central Principle}

\centerbox{\Large\bfseries
Workspace identity is invariant causal organization under admissible transport.
}

The theory begins from a distinction between a \emph{realization} and the
\emph{workspace pattern realized}. A realization is concrete: it contains
particular participants, documents, models, channels, and histories. A
workspace pattern is abstract: it records the structure that remains invariant
when those concrete carriers are changed in admissible ways.

\begin{definition}[Concrete realization]\label{def:realization}
A concrete relational realization is a tuple
\[
\mathfrak R
=
\bigl(
G,\mathbf X,\mathbf T,\mathcal A,\boldsymbol\rho,
\mathcal H,\mathcal I,\mathcal Y
\bigr),
\]
where:
\begin{itemize}
\item \(G=(V,E,s,t)\) is an interaction graph;
\item \(\mathbf X=\{X_i\}_{i\in V}\) is a family of local state spaces;
\item \(\mathbf T=\{T_e\}_{e\in E}\) is a family of transfer maps;
\item \(\mathcal A\) is an algebra of admissible operations;
\item \(\boldsymbol\rho=\{\rho_i\}_{i\in V}\) gives the local realization of those operations;
\item \(\mathcal H\) is the temporally ordered interaction history;
\item \(\mathcal I\) is a specified family of admissible interventions; and
\item \(\mathcal Y\) is the family of downstream variables relative to which causal and predictive claims are evaluated.
\end{itemize}
The static reduct is
\[
\mathfrak R_{\mathrm{stat}}
=
(G,\mathbf X,\mathbf T,\mathcal A,\boldsymbol\rho).
\]
\end{definition}

The explicit appearance of \(\mathcal I\) and \(\mathcal Y\) is important.
A collection of mutually consistent representations may define a static
shared structure, but it does not yet define an operational workspace unless
that structure can be intervened upon and shown to matter for a specified
future quantity.

\begin{definition}[Carrier and organization]\label{def:carrier-organization}
The \emph{carriers} of \(\mathfrak R\) are the individual spaces, physical
systems, and stored artifacts that realize its local states. Its
\emph{organization} is the structured family of transfer, operation,
compatibility, intervention, and temporal relations among those carriers.
\end{definition}

The basic thesis is not that a workspace exists without carriers. It is that
the identity conditions for the workspace need not coincide with the identity
conditions for any one carrier. A manuscript can be rewritten, a model can
be replaced, and a notation can be changed while a sufficiently rich pattern
of relations remains invariant.

\begin{definition}[Admissible equivalence]\label{def:admissible-equivalence}
An admissible equivalence
\[
F:\mathfrak R\simeq\mathfrak R'
\]
is a family of correspondences preserving, exactly or within a declared
tolerance:
\[
\begin{aligned}
&\text{the transferable core and its dimension},\\
&\text{the represented operation algebra},\\
&\text{the composition laws of transfer},\\
&\text{the intervention--response structure},\\
&\text{the reconstruction and persistence relations}.
\end{aligned}
\]
Static admissible equivalence will be made precise by gauge conjugacy in
Part~V. Dynamic admissible equivalence will be expressed using the identity
distance in Part~XII.
\end{definition}

The abstract workspace is then the equivalence class
\[
\boxed{
\Workspace=[\mathfrak R]_{\simeq}.
}
\]
This expression is intentionally stronger than saying that the same
information appears in several places. The equivalence relation preserves
what the system can coherently do, how its representations translate, and how
its operative state survives change.

\subtopic{1.1 Three levels of the thesis}

The central claim can be separated into three logically distinct statements.

\begin{enumerate}
\item \textbf{Implementation dependence.}
Every realized workspace depends on concrete physical or computational
carriers:
\[
\Workspace(\mathfrak R)
\text{ is undefined without some realization }\mathfrak R.
\]

\item \textbf{Local nonidentity.}
No privileged local component is, in general, identical with the entire
workspace:
\[
\Workspace\neq X_i
\qquad\text{for each individual }i\in V.
\]

\item \textbf{Relational identity.}
The workspace is individuated by invariants of the complete organized
realization:
\[
\Workspace(\mathfrak R)=\Workspace(\mathfrak R')
\quad\Longleftrightarrow\quad
\mathfrak R\simeq\mathfrak R'.
\]
\end{enumerate}

The second claim is not stipulated universally. It becomes a substantive
property only after local-reconstructibility and coalition tests are applied.
A system in which one node stores a complete sufficient copy may be a useful
workspace, but it is not strongly relational in the later technical sense.

\begin{definition}[Relational instantiation]\label{def:relational-instantiation}
A workspace is relationally instantiated relative to the decomposition
\(V\) when no single projection determines its equivalence class:
\[
\forall i\in V,\quad
\exists\mathfrak R,\mathfrak R'
\]
such that
\[
\pr_i(\mathfrak R)=\pr_i(\mathfrak R')
\quad\text{but}\quad
[\mathfrak R]_{\simeq}\neq[\mathfrak R']_{\simeq}.
\]
\end{definition}

\subtopic{1.2 A motivating three-node cycle}

Consider three strata:
\[
h=\text{human conceptual state},\qquad
m=\text{model-formal state},\qquad
a=\text{artifact state}.
\]
Let the corresponding core spaces be \(C_h,C_m,C_a\), with transfers
\[
H_{hm}:C_h\to C_m,
\qquad
H_{ma}:C_m\to C_a,
\qquad
H_{ah}:C_a\to C_h.
\]
A proposed shared item begins as \(c_h\in C_h\), is formalized as
\[
c_m=H_{hm}c_h,
\]
externalized as
\[
c_a=H_{ma}c_m,
\]
and re-enters human reasoning as
\[
c_h'=H_{ah}c_a.
\]
The round-trip consistency condition is
\[
c_h'=c_h,
\]
or equivalently
\[
\boxed{
H_{ah}H_{ma}H_{hm}c_h=c_h.
}
\]
Thus the genuinely shared part of \(C_h\) is not automatically all of
\(C_h\). It is the fixed subspace of the cycle transfer. This elementary
cycle already contains the central mathematical motif of the theory:
\[
\text{shared core}=\text{invariant subspace under relational holonomy}.
\]

\subtopic{1.3 Distinctions the theory must preserve}

\begin{center}
\small
\begin{tabularx}{\textwidth}{@{}>{\bfseries\raggedright\arraybackslash}p{0.22\textwidth}>{\raggedright\arraybackslash}X@{}}
\toprule
Object & Minimal distinguishing property \\
\midrule
Archive & Stores material that may be consulted later. \\
Shared representation & Contains a cross-node compatible structure. \\
Operational workspace & The shared structure causally changes subsequent joint behavior. \\
Relational workspace & No individual node or matched local-storage null reproduces the operative capacity. \\
Workspace agent & Workspace-level state mediates persistent goals, counterfactual choice, and self-maintaining action. \\
Conscious subject & Possesses phenomenal properties, if any; no implication from the preceding levels is assumed. \\
\bottomrule
\end{tabularx}
\end{center}

The static mathematics developed first certifies only shared structure. The
causal and relational predicates appear later because they require dynamics,
interventions, and null models rather than algebra alone.

\topic{2}{Philosophical Postulates}

The following postulates specify the intended interpretation of the formal
objects. They are not offered as independent mathematical theorems. Their
role is to prevent equivocation between a relational pattern, its physical
implementation, and stronger claims about agency or experience.

\subtopic{2.1 Physical supervenience}

No additional nonphysical substance is required:
\[
\mathfrak R=\mathfrak R'
\Longrightarrow
\Workspace(\mathfrak R)=\Workspace(\mathfrak R').
\]
The workspace depends completely on its total realization. In particular,
changing nothing about the graph, local states, transfers, history,
interventions, or causal dispositions cannot change the workspace.

This is compatible with relational individuation because a relation can be
fully physical while not being a property of any isolated endpoint.

\begin{proposition}[Compatibility of supervenience and relational individuation]\label{prop:supervenience-relational}
Physical supervenience on the complete realization does not imply
supervenience on any individual local projection.
\end{proposition}

\begin{proof}
Let the total realization decompose as \(\mathfrak R=(X_i,R_{-i})\), where
\(R_{-i}\) contains the remaining components and relations. A workspace
functional \(\Workspace=f(X_i,R_{-i})\) supervenes on \(\mathfrak R\) by
construction. If there exist \(R_{-i}\neq R_{-i}'\) with
\[
f(X_i,R_{-i})\neq f(X_i,R_{-i}'),
\]
then the same local projection \(X_i\) occurs in distinct workspaces. Hence
supervenience on the whole is consistent with failure of supervenience on a
part.
\end{proof}

\subtopic{2.2 Relational individuation}

The workspace is not determined by any single local state:
\[
\forall i\in V,
\quad
\exists\mathfrak R,\mathfrak R'
\]
such that
\[
\pr_i(\mathfrak R)=\pr_i(\mathfrak R')
\]
but
\[
\Workspace(\mathfrak R)\not\simeq\Workspace(\mathfrak R').
\]
The same person, model, or document can therefore participate in distinct
workspaces because the relevant transfer graph, history, and operation
structure differ.

A useful witness is a pair of systems with identical local-state marginals but
different edge organization. If their downstream behavior differs, the
relations carry explanatory information not present in the isolated marginals.
Part~VII converts this idea into matched relation-preserving and
relation-destroying controls.

\subtopic{2.3 Multiple realizability}

Distinct implementations may instantiate the same abstract workspace:
\[
\mathfrak R\neq\mathfrak R',
\qquad
\mathfrak R\simeq\mathfrak R'
\Longrightarrow
\Workspace(\mathfrak R)\cong\Workspace(\mathfrak R').
\]

Examples of candidate implementation changes include:
\[
\begin{aligned}
&\text{replacing one coordinate system by another},\\
&\text{replacing one model by a functionally equivalent model},\\
&\text{compressing a large artifact into a reconstructively sufficient handoff},\\
&\text{rewriting a proof without changing its inferential structure}.
\end{aligned}
\]
The burden is not merely to assert equivalence but to exhibit maps preserving
the declared invariant and causal structure.

\subtopic{2.4 Operational reality}

A workspace is operationally real only if changing it changes downstream
behavior. For some admissible intervention and some \(w\neq w'\),
\[
\mathbb P\!\left(
Y^+\mid\operatorname{do}(\Workspace=w)
\right)
\neq
\mathbb P\!\left(
Y^+\mid\operatorname{do}(\Workspace=w')
\right).
\]

This criterion excludes a merely decorative representation that is correlated
with activity but never used. It also requires the intervention to be
well-defined. Arbitrarily overwriting every component at once would not
isolate workspace-level causation. Later sections therefore compare
workspace-coordinate interventions with matched interventions in orthogonal
or shuffled channels.

\subtopic{2.5 Reconstructive continuity}

Persistence does not require token identity. Let \(\mathfrak R_t\) and
\(\mathfrak R_{t+1}\) be temporally successive realizations. Continuity is
possible when there exists an admissible transport
\[
J_t:\mathfrak R_t\rightsquigarrow\mathfrak R_{t+1}
\]
that preserves the operative core and its causal role within tolerance:
\[
d_{\mathrm{id}}(\Workspace_t,\Workspace_{t+1})\leq\varepsilon.
\]
Thus a workspace may persist by reconstruction rather than by retaining an
unchanged physical inscription.

\subtopic{2.6 No free-standing third substance}

The claim that the workspace is not identical with either participant does
not imply that it occupies a separate nonphysical location. Formally,
\[
\Workspace
\text{ is an invariant of }\mathfrak R,
\]
not an additional term appended to the physical realization:
\[
\mathfrak R_{\mathrm{total}}
eq(\mathfrak R,\text{ghost workspace}).
\]
The phrase ``between the participants'' is therefore topological or
relational: it refers to the organization of maps, cycles, artifacts, and
causal dependencies.

\subtopic{2.7 Non-entailment of consciousness}

\[
\operatorname{Workspace}(\mathfrak R)
\centernot\Longrightarrow
\operatorname{Conscious}(\mathfrak R).
\]
Likewise,
\[
\operatorname{Distributed}(\Workspace)
\centernot\Longrightarrow
\operatorname{Person}(\Workspace).
\]
Workspace structure, reentry, agency, self-maintenance, and phenomenal
consciousness remain formally distinct predicates. Any bridge among them
would require additional theory and evidence.

\subtopic{2.8 Minimal philosophical commitment}

The weakest interpretation sufficient for the mathematics is:
\centerbox{\bfseries
Relational patterns can be objective, causally relevant, and multiply realized
without being independent substances.
}
This position is compatible with several broader philosophies of mind and
science. The formal theory does not require choosing among them before its
structural and causal claims can be tested.

\theorypart{II}{Static Core--Transfer Theory}

\topic{3}{Interaction Structure}

\begin{definition}[Interaction graph]\label{def:interaction-graph}
Let
\[
G=(V,E,s,t)
\]
be a finite directed multigraph. Parallel edges and self-loops are allowed.
The vertex set indexes strata, participants, artifacts, models, instruments,
or subworkspaces. An edge \(e:i\to j\) denotes an admissible directed
translation or influence channel from \(i\) to \(j\).
\end{definition}

The directed path category is
\[
\mathcal C=\Path(G).
\]
Its objects are vertices. A path
\[
\gamma=e_m\cdots e_2e_1
\]
means that \(e_1\) is traversed first. Composition is concatenation, and the
empty path at \(i\) is denoted \(1_i\).

For every \(i\in V\), assign a finite-dimensional Hilbert space
\[
X_i.
\]
For every edge, assign a linear transfer
\[
T_e:X_{s(e)}\longrightarrow X_{t(e)}.
\]
The path transfer is
\[
T_\gamma=T_{e_m}\cdots T_{e_2}T_{e_1},
\qquad
T_{1_i}=I_{X_i}.
\]
Accordingly, the static transfer system is a functor
\[
\mathcal X:\Path(G)\longrightarrow\mathbf{FinHilb}.
\]

The finite-dimensional Hilbert assumption supplies adjoints and spectral
tools. The exact compatible-section definitions require less and may later be
formulated in categories such as
\[
\mathbf{Ban},\qquad
\mathbf{Meas},\qquad
\mathbf{Prob},\qquad
\mathbf{Top},\qquad
\mathbf{Alg}.
\]

\subtopic{3.1 Directed paths versus invertible transport}

Global compatible sections make sense even when transfers are noninvertible.
Holonomy, however, requires transport around reversible loops. To state the
exact holonomy theorem without ambiguity, let \(|G|\) be the underlying
undirected graph and let
\[
\Pi_1(G)
\]
denote its fundamental path groupoid. Every oriented edge is then supplied
with a formal reverse. The core transfers extend to \(\Pi_1(G)\) only when
the relevant maps \(H_e\) are invertible, with
\[
H_{e^{-1}}=H_e^{-1}.
\]

Hence two regimes will be distinguished:
\begin{enumerate}
\item the \emph{diagram regime}, in which arbitrary directed maps define
compatible sections; and
\item the \emph{local-system regime}, in which invertible core maps define a
functor
\[
H:\Pi_1(G)\to\mathbf{FinHilb}^{\times}.
\]
\end{enumerate}
Theorems based on holonomy will explicitly assume the second regime.

\subtopic{3.2 Connectivity convention}

Unless stated otherwise, connectedness means that \(|G|\) is connected. A
root \(o\in V\) may then be joined to every vertex by a path in the underlying
graph. If \(G\) has several connected components, the theory applies to each
component and
\[
\Gamma(H)
\cong
\bigoplus_{\alpha}\Gamma(H|_{G_\alpha}).
\]

\subtopic{3.3 Edge weights}

For approximate theory, assign strictly positive edge weights
\[
w_e>0.
\]
They may encode channel reliability, frequency, precision, experimental
importance, or inverse noise variance. Exact compatibility is independent of
positive rescaling of these weights, while approximate spectra are not.

\begin{example}[Human--model--artifact graph]\label{ex:hma-graph}
The motivating three-node system is represented by
\[
V=\{h,m,a\},
\qquad
E=\{h\to m,m\to a,a\to h\}.
\]
Additional edges may encode direct retrieval \(a\to m\), human revision
\(h\to a\), or independent verification. The graph records permitted
translations, not merely physical communication links.
\end{example}

\topic{4}{Core--Residual Decomposition}

\begin{definition}[Admissible local decomposition]\label{def:core-residual}
For every vertex \(i\), choose an orthogonal decomposition
\[
X_i=C_i\oplus R_i,
\]
with projectors
\[
P_i:X_i\to C_i,
\qquad
Q_i=I-P_i:X_i\to R_i.
\]
The subspace \(C_i\) is the candidate locally realized transferable core. The
subspace \(R_i\) contains substrate-specific, context-specific, or otherwise
noncore degrees of freedom relative to the present analysis.
\end{definition}

The decomposition is part of the model and is not presumed unique. Calling a
subspace ``core'' is therefore a hypothesis requiring structural and causal
validation. Later null tests must determine whether the chosen projectors
identify a genuine shared channel rather than an arbitrary low-dimensional
summary.

Relative to these decompositions, every edge map has the block form
\[
T_e=
\begin{pmatrix}
H_e & U_e\\
V_e & K_e
\end{pmatrix},
\]
where
\[
H_e=P_{t(e)}T_eP_{s(e)}:C_{s(e)}\to C_{t(e)},
\]
\[
U_e=P_{t(e)}T_eQ_{s(e)}:R_{s(e)}\to C_{t(e)},
\]
\[
V_e=Q_{t(e)}T_eP_{s(e)}:C_{s(e)}\to R_{t(e)},
\]
\[
K_e=Q_{t(e)}T_eQ_{s(e)}:R_{s(e)}\to R_{t(e)}.
\]

\subtopic{4.1 Structural conditions}

\textbf{Core preservation} means
\[
V_e=0.
\]
A core input cannot leak into the target residual.

\textbf{Core autonomy} means
\[
U_e=V_e=0.
\]
Core and residual evolve as separate transfer systems.

\textbf{Core leakage}, \textbf{residual contamination}, and \textbf{total
transfer defect} are measured by
\[
\lambda_e^{\mathrm{out}}=\|V_e\|,
\qquad
\lambda_e^{\mathrm{in}}=\|U_e\|,
\]
\[
\delta_e^{\mathrm{tr}}
=
\|U_e\|+\|V_e\|.
\]
These quantities depend on the chosen norms and projectors. Their empirical
meaning must therefore be tied to a declared intervention and prediction
protocol.

\subtopic{4.2 Composition exposes residual mediation}

For composable edges \(e:i\to j\) and \(f:j\to k\), direct block
multiplication gives
\[
T_fT_e
=
\begin{pmatrix}
H_fH_e+U_fV_e & H_fU_e+U_fK_e\\
V_fH_e+K_fV_e & V_fU_e+K_fK_e
\end{pmatrix}.
\]
The induced core-to-core block is therefore
\[
P_kT_fT_eP_i
=
H_fH_e+U_fV_e.
\]
The term \(U_fV_e\) is a residual-mediated correction: information leaves the
candidate core at the intermediate node and later re-enters it. Consequently,
the simple path product \(H_fH_e\) faithfully represents full core transport
only under a preservation condition such as \(V_e=0\), or under an explicit
model that includes these mediated corrections.

\begin{proposition}[Pathwise core preservation]\label{prop:pathwise-preservation}
Suppose \(V_e=0\) for every edge in a directed path
\(\gamma=e_m\cdots e_1\). Then
\[
T_\gamma(C_{s(\gamma)})\subseteq C_{t(\gamma)}
\]
and
\[
P_{t(\gamma)}T_\gamma P_{s(\gamma)}
=
H_{e_m}\cdots H_{e_1}.
\]
\end{proposition}

\begin{proof}
For a single edge, \(V_e=0\) means
\[
T_e(c,0)=(H_ec,0),
\]
so core states remain in the core. Inductively, if the first \(r\) transfers
map \(c\) to \((H_{e_r}\cdots H_{e_1}c,0)\), applying \(T_{e_{r+1}}\) gives
\[
(H_{e_{r+1}}H_{e_r}\cdots H_{e_1}c,0).
\]
The conclusion follows after \(m\) steps.
\end{proof}

\begin{corollary}[Autonomous splitting]\label{cor:autonomous-splitting}
If \(U_e=V_e=0\) for every edge, then the complete diagram decomposes as the
direct sum of two diagrams,
\[
\mathcal X\cong H\oplus K,
\]
and every path transfer satisfies
\[
T_\gamma
=
\begin{pmatrix}
H_\gamma&0\\
0&K_\gamma
\end{pmatrix}.
\]
\end{corollary}

\subtopic{4.3 Exact core-transfer systems}

\begin{definition}[Exact core-transfer system]\label{def:exact-core-transfer}
An exact core-transfer system on \(G\) consists of:
\begin{enumerate}
\item decompositions \(X_i=C_i\oplus R_i\);
\item core-preserving edge maps \(V_e=0\); and
\item invertible core blocks \(H_e\in GL(C_{s(e)},C_{t(e)})\) whenever
holonomy is invoked.
\end{enumerate}
\end{definition}

The requirement \(V_e=0\) concerns the relation between the selected core and
the full transfer. The requirement that \(H_e\) be invertible concerns the
ability to compare core representations along reversed paths. They answer
different questions and should not be conflated.

\begin{remark}[Approximate systems]
In applications neither condition need hold exactly. The approximate theory
will replace strict invariance by low transfer energy, small leakage, and a
spectral gap separating workspace modes from residual modes.
\end{remark}

\topic{5}{Abstract Operation Algebra}

\begin{definition}[Operation algebra]\label{def:operation-algebra}
Let \(\mathbb F\in\{\mathbb R,\mathbb C\}\). An operation algebra
\(\mathcal A\) is a unital associative \(\mathbb F\)-algebra whose elements
represent abstract transformations, inferential rules, questions, or
procedures that can be realized in multiple strata.
\end{definition}

Every core space carries a unital representation
\[
\rho_i:\mathcal A\longrightarrow\End(C_i).
\]
For \(a\in\mathcal A\), the operator \(\rho_i(a)\) is the realization of the
same abstract operation in local coordinates.

\begin{definition}[Exact operation-compatible transfer]\label{def:intertwiner}
A core transfer \(H_e:C_i\to C_j\) is operation compatible when it is an
intertwiner:
\[
\boxed{
H_e\rho_i(a)=\rho_j(a)H_e
}
\qquad
\forall a\in\mathcal A.
\]
Equivalently, the diagram
\[
\begin{array}{ccc}
C_i&\xrightarrow{\rho_i(a)}&C_i\\
\downarrow H_e&&\downarrow H_e\\
C_j&\xrightarrow{\rho_j(a)}&C_j
\end{array}
\]
commutes for every \(a\).
\end{definition}

This condition states that translating and then operating gives the same
result as operating and then translating. Agreement of static encodings alone
is weaker: two nodes may store corresponding vectors while implementing
incompatible transformations on them.

\begin{proposition}[Path composition of intertwiners]\label{prop:path-intertwiner}
If every edge of a path \(\gamma:i\to j\) intertwines the local
\(\mathcal A\)-representations, then
\[
H_\gamma\rho_i(a)=\rho_j(a)H_\gamma
\qquad
\forall a\in\mathcal A.
\]
\end{proposition}

\begin{proof}
For composable intertwiners \(H_e:C_i\to C_j\) and
\(H_f:C_j\to C_k\),
\[
H_fH_e\rho_i(a)
=H_f\rho_j(a)H_e
=\rho_k(a)H_fH_e.
\]
Induction over path length proves the claim.
\end{proof}

\begin{corollary}[Holonomy commutes with abstract operations]\label{cor:holonomy-commutes}
For every loop \(\gamma:o\to o\),
\[
H_\gamma\rho_o(a)=\rho_o(a)H_\gamma.
\]
Hence
\[
\rho_o(\mathcal A)
\subseteq
\operatorname{Comm}(\Hol_o),
\]
where \(\operatorname{Comm}(\Hol_o)\) is the commutant of the holonomy
representation.
\end{corollary}

\subtopic{5.1 Approximate operation compatibility}

The edgewise defect is
\[
\Xi_e(a)
=
H_e\rho_{s(e)}(a)-\rho_{t(e)}(a)H_e.
\]
For a bounded generating family \(\mathcal G\subseteq\mathcal A\), define
\[
\eta_{\mathrm{int}}(\mathcal G)
=
\sup_{a\in\mathcal G}
\sum_{e\in E}w_e\|\Xi_e(a)\|^2.
\]
When \(\mathcal A\) is normed and its unit ball is appropriate, one may take
\[
\mathcal G=\{a:\|a\|_{\mathcal A}\leq1\}.
\]
Exact representational universality means
\[
\eta_{\mathrm{int}}=0,
\]
while approximate universality means
\[
\eta_{\mathrm{int}}\leq\varepsilon_{\mathrm{int}}.
\]
The use of a generating family prevents the definition from depending on an
unbounded rescaling of algebra elements.

\begin{example}[Equivalent operations in different coordinates]\label{ex:coordinate-intertwiner}
Let \(C_1=C_2=\mathbb R^2\), choose \(S\in GL_2(\mathbb R)\), and set
\[
H_{12}=S,
\qquad
\rho_2(a)=S\rho_1(a)S^{-1}.
\]
Then
\[
H_{12}\rho_1(a)
=S\rho_1(a)
=\rho_2(a)S
=\rho_2(a)H_{12}.
\]
The local matrices can look different while realizing the same abstract
operation. This is the elementary model of representation change used
throughout the theory.
\end{example}

\theorypart{III}{Exact Shared Cores}

\topic{6}{Global Compatible Sections}

From this point through Theorem~\ref{thm:limit-holonomy}, assume the core maps
define a local system on a connected graph unless a weaker hypothesis is
stated explicitly.

\begin{definition}[Global compatible section]\label{def:global-section}
The exact structural workspace is
\[
\boxed{
\Gamma(H)
=
\left\{
(c_i)_{i\in V}\in\bigoplus_{i\in V}C_i:
\ c_{t(e)}=H_ec_{s(e)}\ \forall e\in E
\right\}.
}
\]
An element \(c=(c_i)\in\Gamma(H)\) is a global compatible section.
\end{definition}

Equivalently,
\[
\Gamma(H)=\varprojlim_{\Path(G)}C_i.
\]
A section is not a single vector copied verbatim at every vertex. It is one
abstract element represented by a family of local vectors that agree under the
specified translation maps.

A nontrivial exact structural workspace exists when
\[
\dim\Gamma(H)>0.
\]
Its exact dimension is
\[
d_{\Workspace}=\dim\Gamma(H).
\]

Fix a root \(o\). The evaluation map is
\[
\ev_o:\Gamma(H)\to C_o,
\qquad
\ev_o((c_i))=c_o.
\]

\begin{lemma}[Root determination]\label{lem:root-determination}
If \(|G|\) is connected and every core transfer is invertible, then
\(\ev_o\) is injective. Consequently,
\[
\dim\Gamma(H)\leq\dim C_o.
\]
\end{lemma}

\begin{proof}
Choose for each vertex \(i\) a path \(\tau_i:o\to i\). Compatibility gives
\[
c_i=H_{\tau_i}c_o.
\]
Thus a section with \(c_o=0\) vanishes at every vertex. Hence \(\ev_o\) is
injective.
\end{proof}

\begin{lemma}[Tree propagation]\label{lem:tree-propagation}
If \(|G|\) is a tree and every \(H_e\) is invertible, then every
\(x\in C_o\) extends uniquely to a global compatible section.
\end{lemma}

\begin{proof}
A tree contains a unique undirected path \(\tau_i:o\to i\) for every vertex.
Define \(c_i=H_{\tau_i}x\). Because every edge belongs to these unique path
relations, the family is compatible. Uniqueness follows from
Lemma~\ref{lem:root-determination}.
\end{proof}

The obstruction to a global core therefore appears only when distinct paths
must agree. Cycles encode exactly those consistency constraints.

\topic{7}{Holonomy and Path Obstruction}

Fix a root \(o\in V\). For every loop class
\[
[\gamma]\in\pi_1(|G|,o),
\]
define the holonomy
\[
\Hol([\gamma])=H_\gamma\in GL(C_o).
\]
Functoriality gives a group representation
\[
\operatorname{hol}_o:
\pi_1(|G|,o)\longrightarrow GL(C_o).
\]
Its image is the holonomy group
\[
\Hol_o=\operatorname{im}(\operatorname{hol}_o).
\]

The obstruction associated with a loop is
\[
\Omega_\gamma=H_\gamma-I_{C_o}.
\]
The maximal path-independent root core is
\[
\boxed{
C_*
=
\bigcap_{\gamma:o\to o}\ker(H_\gamma-I).
}
\]
Equivalently,
\[
C_*=C_o^{\Hol_o}=\operatorname{Fix}(\Hol_o).
\]

\subtopic{7.1 Spanning-tree reduction}

Choose a spanning tree \(\mathcal T\subseteq|G|\). Let
\[
\tau_i:o\to i
\]
be the unique tree path. Every non-tree edge \(e:i\to j\) determines a
fundamental loop
\[
\gamma_e=\tau_j^{-1}e\tau_i
\]
with transfer
\[
H_{\gamma_e}=H_{\tau_j}^{-1}H_eH_{\tau_i}.
\]
The fundamental group of a finite graph is generated freely by these cotree
loops. Hence it is enough to impose fixed-point conditions for a finite set of
generators.

\begin{lemma}[Finite obstruction set]\label{lem:finite-obstruction}
Let \(e_1,\ldots,e_r\) be the non-tree edges of a connected finite graph.
Then
\[
C_*
=
\bigcap_{q=1}^{r}
\ker(H_{\gamma_{e_q}}-I).
\]
\end{lemma}

\begin{proof}
A vector fixed by every generator is fixed by each inverse generator and by
every finite product of generators and inverses. Those products represent all
elements of \(\pi_1(|G|,o)\). The reverse inclusion is immediate.
\end{proof}

The number of independent graph cycles is
\[
r=|E_{\mathrm{und}}|-|V|+1
\]
for a connected underlying graph, after treating parallel edges with their
multiplicity.

\topic{8}{Exact Existence Theorem}

\begin{theorem}[Limit--holonomy equivalence]\label{thm:limit-holonomy}
Let \(|G|\) be connected and let
\[
H:\Pi_1(G)\to\mathbf{FinHilb}^{\times}
\]
be a finite-dimensional local system. For any root \(o\in V\), evaluation at
the root induces a canonical vector-space isomorphism
\[
\boxed{
\ev_o:\Gamma(H)
\xrightarrow{\ \cong\ }
C_o^{\Hol_o}
=
\bigcap_{\gamma:o\to o}\ker(H_\gamma-I).
}
\]
Consequently,
\[
\Gamma(H)\neq\{0\}
\]
if and only if the holonomy representation contains a nonzero trivial
subrepresentation.
\end{theorem}

\begin{proof}
Let \(c=(c_i)\in\Gamma(H)\). For any loop \(\gamma:o\to o\), compatibility
along the loop gives
\[
H_\gamma c_o=c_o.
\]
Therefore \(\ev_o(c)=c_o\in C_o^{\Hol_o}\). By
Lemma~\ref{lem:root-determination}, \(\ev_o\) is injective.

Conversely, let \(x\in C_o^{\Hol_o}\). Choose a path
\(\tau_i:o\to i\) for each vertex and define
\[
c_i=H_{\tau_i}x.
\]
This definition is independent of the chosen path. Indeed, if
\(\sigma_i:o\to i\) is another path, then
\[
\sigma_i^{-1}\tau_i:o\to o
\]
is a loop, and invariance of \(x\) gives
\[
H_{\sigma_i}^{-1}H_{\tau_i}x=x.
\]
Multiplying by \(H_{\sigma_i}\) yields
\[
H_{\tau_i}x=H_{\sigma_i}x.
\]
For any edge \(e:i\to j\), the paths \(e\tau_i\) and \(\tau_j\) both run
from \(o\) to \(j\). Path independence therefore gives
\[
H_ec_i
=H_eH_{\tau_i}x
=H_{\tau_j}x
=c_j.
\]
Thus \((c_i)\) is a global compatible section and \(\ev_o((c_i))=x\).
Hence \(\ev_o\) is surjective and therefore an isomorphism.
\end{proof}

\begin{corollary}[Tree case]\label{cor:tree-core}
If \(|G|\) is a tree, then \(\Hol_o=\{I\}\) and
\[
\Gamma(H)\cong C_o.
\]
In particular,
\[
d_{\Workspace}=\dim C_o.
\]
\end{corollary}

\begin{corollary}[Finite matrix test]\label{cor:finite-matrix-test}
Let \(\gamma_1,\ldots,\gamma_r\) be fundamental cycle generators and suppose
\(\dim C_o=d\). Define the stacked obstruction matrix
\[
\mathcal M
=
\begin{pmatrix}
H_{\gamma_1}-I\\
\vdots\\
H_{\gamma_r}-I
\end{pmatrix}.
\]
Then
\[
\boxed{
C_*=\ker\mathcal M,
\qquad
d_{\Workspace}=d-\operatorname{rank}(\mathcal M).
}
\]
\end{corollary}

\begin{proof}
The kernel of a vertically stacked operator is the intersection of the kernels
of its blocks. Apply Lemma~\ref{lem:finite-obstruction} and
Theorem~\ref{thm:limit-holonomy}.
\end{proof}

\begin{corollary}[Averaging projector for finite holonomy]\label{cor:averaging-projector}
If \(\Hol_o\) is finite and acts unitarily, then
\[
P_*=
\frac{1}{|\Hol_o|}
\sum_{h\in\Hol_o}h
\]
is the orthogonal projector onto \(C_*\).
\end{corollary}

\begin{proof}
For every \(g\in\Hol_o\), left multiplication permutes the group elements, so
\(gP_*=P_*\). Thus \(\operatorname{ran}P_*\subseteq C_*\). If
\(x\in C_*\), then \(P_*x=x\), so \(C_*\subseteq\operatorname{ran}P_*\).
Unitarity gives \(P_*^*=P_*\), and the range characterization implies
\(P_*^2=P_*\).
\end{proof}

\subtopic{8.1 Worked examples}

\begin{example}[One-dimensional core surviving a contradictory cycle]\label{ex:reflection-cycle}
Let a three-node cycle have \(C_1=C_2=C_3=\mathbb R^2\) and
\[
H_{12}=I,
\qquad
H_{23}=I,
\qquad
H_{31}=J,
\qquad
J=
\begin{pmatrix}
1&0\\
0&-1
\end{pmatrix}.
\]
The cycle holonomy is \(J\), hence
\[
C_*=\ker(J-I)=\operatorname{span}\{(1,0)^\top\}.
\]
The first coordinate circulates consistently through all three strata. The
second returns with its sign reversed and therefore cannot define a nonzero
exact shared section. Thus
\[
d_{\Workspace}=1.
\]
\end{example}

\begin{example}[Local compatibility without a global nonzero core]\label{ex:rotation-obstruction}
Let the only nontrivial cycle holonomy on \(\mathbb R^2\) be a rotation
\[
R_\theta=
\begin{pmatrix}
\cos\theta&-\sin\theta\\
\sin\theta&\cos\theta
\end{pmatrix}
\]
with \(\theta\notin2\pi\mathbb Z\). Then
\[
\ker(R_\theta-I)=\{0\}.
\]
Every edge map is individually invertible and well-defined, yet the complete
cycle admits no nonzero exact shared core. Pairwise translatability is
therefore insufficient for global path independence.
\end{example}

\begin{example}[Acyclic collaboration]\label{ex:tree-workspace}
Suppose a root concept is translated into a proof, implementation, and report
along a tree of invertible exact transfers. No alternative route returns a
conflicting representation. The tree has trivial holonomy, so every root-core
direction extends uniquely. This does not yet establish causal efficacy; it
establishes only exact structural compatibility.
\end{example}

\subtopic{8.2 Interpretive consequence}

The theorem gives a precise meaning to ``something shared across
representations.'' The shared object is neither any one local vector nor an
extra vector outside the system. It is the equivalence class of all compatible
local realizations, canonically represented at a root by the subspace fixed by
every translation cycle.

\topic{9}{Operator Closure Theorem}

\begin{theorem}[Global operation closure]\label{thm:operator-closure}
Suppose every core transfer intertwines the representations of
\(\mathcal A\):
\[
H_e\rho_{s(e)}(a)=\rho_{t(e)}(a)H_e
\qquad
\forall e\in E,\ a\in\mathcal A.
\]
Then \(\Gamma(H)\) is an \(\mathcal A\)-module under the componentwise action
\[
\rho_\Gamma(a)(c_i)_{i\in V}
=
(\rho_i(a)c_i)_{i\in V}.
\]
Equivalently, the holonomy-fixed root core is invariant:
\[
\boxed{
\rho_o(a)C_*\subseteq C_*
\qquad
\forall a\in\mathcal A.
}
\]
\end{theorem}

\begin{proof}
Let \(c=(c_i)\in\Gamma(H)\). For every edge \(e:i\to j\),
\[
H_e\rho_i(a)c_i
=
\rho_j(a)H_ec_i
=
\rho_j(a)c_j.
\]
Therefore the family \((\rho_i(a)c_i)\) is compatible and belongs to
\(\Gamma(H)\). The representation identities on each \(C_i\) hold
componentwise, so \(\rho_\Gamma\) is an \(\mathcal A\)-representation.

Under the isomorphism of Theorem~\ref{thm:limit-holonomy}, this action is
represented at the root by \(\rho_o\). Alternatively, if \(x\in C_*\), then
Corollary~\ref{cor:holonomy-commutes} gives, for every loop \(\gamma\),
\[
H_\gamma\rho_o(a)x
=
\rho_o(a)H_\gamma x
=
\rho_o(a)x.
\]
Hence \(\rho_o(a)x\in C_*\).
\end{proof}

\begin{corollary}[Inherited algebraic structure]\label{cor:inherited-structure}
Any algebraic relation satisfied by \(\rho_i(\mathcal A)\) is inherited by the
global action. In particular, if \(a,b\in\mathcal A\), then
\[
\rho_\Gamma(ab)=\rho_\Gamma(a)\rho_\Gamma(b),
\qquad
\rho_\Gamma(1)=I_{\Gamma(H)}.
\]
\end{corollary}

The result distinguishes a shared \emph{operational} representation from a
shared static code. A core closed under the common algebra supports the same
abstract manipulations across strata without leaving the globally compatible
subspace.

\topic{10}{Workspace Operator Algebra}

\begin{definition}[Natural workspace operators]\label{def:workspace-operator-algebra}
Define
\[
\boxed{
\mathcal O_{\Workspace}
=
\left\{
(A_i)_{i\in V}:
A_{t(e)}H_e=H_eA_{s(e)}\ \forall e
\right\}.
}
\]
Categorically,
\[
\mathcal O_{\Workspace}=\Nat(H,H),
\]
the algebra of natural endomorphisms of the local system.
\end{definition}

Pointwise addition, scalar multiplication, and composition make
\(\mathcal O_{\Workspace}\) a unital associative algebra. Each
\(A=(A_i)\) acts on global sections by
\[
\pi_{\Workspace}(A)(c_i)=(A_ic_i).
\]
Naturality guarantees that \((A_ic_i)\) is again compatible.

\begin{theorem}[Commutant characterization]\label{thm:commutant-characterization}
Let \(|G|\) be connected and let \(H\) be an invertible local system. Root
evaluation induces an algebra isomorphism
\[
\boxed{
\Phi_o:\mathcal O_{\Workspace}
\xrightarrow{\ \cong\ }
\operatorname{Comm}(\Hol_o),
\qquad
\Phi_o((A_i))=A_o,
}
\]
where
\[
\operatorname{Comm}(\Hol_o)
=
\{B\in\End(C_o):Bh=hB\ \forall h\in\Hol_o\}.
\]
\end{theorem}

\begin{proof}
If \((A_i)\in\mathcal O_{\Workspace}\), naturality along every loop gives
\[
A_oH_\gamma=H_\gamma A_o.
\]
Thus \(A_o\in\operatorname{Comm}(\Hol_o)\), and \(\Phi_o\) is a well-defined
algebra homomorphism.

Choose paths \(\tau_i:o\to i\). Naturality along \(\tau_i\) gives
\[
A_iH_{\tau_i}=H_{\tau_i}A_o,
\]
so
\[
A_i=H_{\tau_i}A_oH_{\tau_i}^{-1}.
\]
Therefore \(A_o\) determines every \(A_i\), and \(\Phi_o\) is injective.

Conversely, take \(B\in\operatorname{Comm}(\Hol_o)\) and define
\[
A_i=H_{\tau_i}BH_{\tau_i}^{-1}.
\]
This is independent of the chosen root path: changing \(\tau_i\) multiplies
it by a loop holonomy that commutes with \(B\). For an edge \(e:i\to j\), the
operator
\[
H_{\tau_j}^{-1}H_eH_{\tau_i}
\]
is a root holonomy. Commutation with \(B\) rearranges to
\[
A_jH_e=H_eA_i.
\]
Hence \((A_i)\) is natural and \(\Phi_o((A_i))=B\). Thus \(\Phi_o\) is
surjective.
\end{proof}

The structural workspace is therefore more accurately represented by the
triple
\[
\boxed{
\Workspace_{\mathrm{static}}
=
\bigl(
\Gamma(H),
\mathcal O_{\Workspace},
\pi_{\Workspace}
\bigr),
}
\]
not merely by the vector space \(\Gamma(H)\).

\subtopic{10.1 Effective operator algebra}

Some natural operators may vanish on every global section while acting
nontrivially on locally available but globally inconsistent directions. Define
the annihilator
\[
\operatorname{Ann}(\Gamma(H))
=
\{A\in\mathcal O_{\Workspace}:
\pi_{\Workspace}(A)=0\}.
\]
The effective workspace algebra is
\[
\boxed{
\mathcal O_{\Workspace}^{\mathrm{eff}}
=
\mathcal O_{\Workspace}/\operatorname{Ann}(\Gamma(H)).
}
\]
Its action on \(\Gamma(H)\) is faithful by construction.

\begin{example}[Operators in the reflection-cycle workspace]\label{ex:reflection-operators}
For Example~\ref{ex:reflection-cycle}, the holonomy is
\[
J=\operatorname{diag}(1,-1).
\]
Its commutant consists of diagonal matrices:
\[
\operatorname{Comm}(J)
=
\left\{
\begin{pmatrix}a&0\\0&b\end{pmatrix}:a,b\in\mathbb R
\right\}.
\]
The exact workspace core is the first coordinate. Consequently, only the
parameter \(a\) acts effectively on \(\Gamma(H)\); the \(b\)-direction belongs
to the annihilator of the global core. Hence
\[
\mathcal O_{\Workspace}^{\mathrm{eff}}\cong\mathbb R.
\]
\end{example}

\subtopic{10.2 Scope of the exact theory}

Parts~I--III now establish three exact facts:
\[
\begin{aligned}
&\text{global shared content is the holonomy-fixed core},\\
&\text{universally transported operations preserve that core},\\
&\text{natural workspace operations are the commutant of holonomy}.
\end{aligned}
\]
These statements do not yet prove that the shared core is robust, causally
used, irreducibly distributed, or temporally persistent. Those are separate
questions developed in the approximate and dynamic theory that follows.


\theorypart{IV}{Approximate Workspaces}

The exact theory asks whether local core states agree perfectly under every
transfer. Approximate theory asks three more realistic questions:
\begin{enumerate}
\item how strongly a candidate state violates edgewise compatibility;
\item whether a low-defect subspace is spectrally separated from incoherent
 directions; and
\item how that subspace changes when transfers, metrics, or operations are
 perturbed.
\end{enumerate}
Throughout Parts~IV--V, set
\[
\mathcal C_V=\bigoplus_{i\in V}C_i,
\qquad
\mathcal C_E=\bigoplus_{e\in E}C_{t(e)},
\]
with their Hilbert direct-sum inner products, and retain strictly positive
weights \(w_e\).

\topic{11}{Transfer Discrepancy Operator}

\begin{definition}[Transfer discrepancy]\label{def:transfer-discrepancy}
The transfer discrepancy operator is
\[
D_H:\mathcal C_V\longrightarrow\mathcal C_E,
\qquad
(D_Hc)_e
=
\sqrt{w_e}\bigl(c_{t(e)}-H_ec_{s(e)}\bigr).
\]
Its squared norm is the coherence energy
\[
\boxed{
\mathcal E_H(c)
=
\|D_Hc\|^2
=
\sum_{e\in E}w_e
\left\|c_{t(e)}-H_ec_{s(e)}\right\|^2.
}
\]
The transfer Laplacian is
\[
\boxed{L_H=D_H^*D_H.}
\]
\end{definition}

The operator \(D_H\) is an incidence-type derivative twisted by the edge
transfers. It compares a state already present at the target with the state
predicted there by translation from the source. A direction can therefore
have low energy even when its coordinates differ substantially across nodes,
provided those coordinates differ exactly as the transfer maps prescribe.

For \(z=(z_e)_{e\in E}\in\mathcal C_E\), the adjoint is
\[
(D_H^*z)_i
=
\sum_{e:t(e)=i}\sqrt{w_e}\,z_e
-
\sum_{e:s(e)=i}\sqrt{w_e}\,H_e^*z_e.
\]
Consequently,
\[
\boxed{
\begin{aligned}
(L_Hc)_i
={}&
\sum_{e:t(e)=i}w_e\bigl(c_i-H_ec_{s(e)}\bigr)\\
&+
\sum_{e:s(e)=i}w_eH_e^*
\bigl(H_ec_i-c_{t(e)}\bigr).
\end{aligned}
}
\]
This formula makes the two contributions at a vertex explicit: incoming
prediction error and adjoint-propagated outgoing error.

\begin{proposition}[Positivity and variational identity]\label{prop:laplacian-positive}
The transfer Laplacian is self-adjoint and positive semidefinite. For every
\(c\in\mathcal C_V\),
\[
\langle c,L_Hc\rangle=\mathcal E_H(c)\geq0.
\]
\end{proposition}

\begin{proof}
Self-adjointness and positivity follow from \(L_H=D_H^*D_H\). The displayed
identity is
\[
\langle c,D_H^*D_Hc\rangle
=
\langle D_Hc,D_Hc\rangle.
\]
\end{proof}

\begin{remark}[Weight dependence]
Replacing \(w_e\) by another positive family does not alter \(\ker D_H\), but
it generally changes every positive eigenvalue and therefore changes the
approximate geometry. Edge weights are part of the metric model, not part of
the exact compatibility relation.
\end{remark}

\topic{12}{Exact Kernel Theorem}

\begin{theorem}[Transfer-Laplacian kernel theorem]\label{thm:laplacian-kernel}
For an arbitrary finite directed core diagram, without assuming invertibility,
\[
\boxed{
\ker L_H=\ker D_H=\Gamma(H).
}
\]
Thus the exact compatible sections are precisely the zero-energy modes of the
transfer Laplacian.
\end{theorem}

\begin{proof}
If \(D_Hc=0\), then \(L_Hc=D_H^*D_Hc=0\). Conversely, if \(L_Hc=0\), then
\[
0=\langle c,L_Hc\rangle=\|D_Hc\|^2,
\]
so \(D_Hc=0\). Finally, \(D_Hc=0\) if and only if
\[
c_{t(e)}=H_ec_{s(e)}
\]
for every edge, which is exactly the definition of \(\Gamma(H)\).
\end{proof}

\begin{corollary}[Finite rank computation]\label{cor:workspace-rank-D}
Let
\[
n_C=\sum_{i\in V}\dim C_i.
\]
Then
\[
\boxed{
\dim\Gamma(H)=n_C-\operatorname{rank}(D_H).
}
\]
Hence the exact workspace dimension can be obtained by a single rank or
singular-value computation, even when holonomy is unavailable because some
transfers are noninvertible.
\end{corollary}

\subtopic{12.1 Relation to the holonomy computation}

In the local-system regime, Theorems~\ref{thm:limit-holonomy} and
\ref{thm:laplacian-kernel} give two equivalent constructions:
\[
\Gamma(H)
\cong
C_o^{\Hol_o}
\qquad\text{and}\qquad
\Gamma(H)=\ker L_H.
\]
The first reduces the problem to cycle invariants in a root fiber. The second
works directly on all vertex coordinates and extends to noninvertible diagrams.
The two computations therefore serve different practical regimes rather than
competing definitions.

\topic{13}{Approximate Structural Workspace}

Let
\[
0\leq\lambda_1\leq\lambda_2\leq\cdots\leq\lambda_N,
\qquad
N=\dim\mathcal C_V,
\]
be the eigenvalues of \(L_H\), repeated with multiplicity, and let
\(P_I(L_H)\) denote the spectral projector associated with an interval \(I\).

\begin{definition}[Threshold workspace]\label{def:threshold-workspace}
For \(\varepsilon\geq0\), define
\[
\boxed{
\Gamma_\varepsilon(H)
=
\ran P_{[0,\varepsilon]}(L_H).
}
\]
Its dimension is
\[
d_\varepsilon
=
\#\{k:\lambda_k\leq\varepsilon\}.
\]
When \(0<d_\varepsilon<N\), define the boundary gap
\[
\boxed{
g_\varepsilon
=
\lambda_{d_\varepsilon+1}-\lambda_{d_\varepsilon}.
}
\]
\end{definition}

The threshold \(\varepsilon\) declares which discrepancy scale is treated as
structurally negligible. The dimension \(d_\varepsilon\) alone is not a
robust certificate: a cluster is stable only when it is separated by a
positive boundary gap.

An equivalent index-based description fixes \(d\) and writes
\[
W_d=\ran P_d,
\qquad
P_d=P_{\{\lambda_1,\ldots,\lambda_d\}}(L_H),
\]
with cluster gap
\[
g_d=\lambda_{d+1}-\lambda_d>0.
\]
The index-based form is convenient for perturbation theorems because the
workspace rank is declared explicitly.

\begin{proposition}[Min--max characterization]\label{prop:minmax}
For \(1\leq k\leq N\),
\[
\lambda_k
=
\min_{\substack{S\subseteq\mathcal C_V\\\dim S=k}}
\ \max_{\substack{c\in S\\\|c\|=1}}
\mathcal E_H(c).
\]
Consequently, the span of the first \(k\) eigenvectors minimizes the worst-case
coherence energy among all \(k\)-dimensional candidate subspaces.
\end{proposition}

\begin{proof}
Expand every vector in an orthonormal eigenbasis of \(L_H\). The span of the
first \(k\) eigenvectors has maximal Rayleigh quotient \(\lambda_k\). Every
\(k\)-dimensional subspace intersects the orthogonal complement of the first
\(k-1\) eigenvectors nontrivially, and any unit vector in that intersection has
Rayleigh quotient at least \(\lambda_k\).
\end{proof}

\subtopic{13.1 Exact, soft, and spurious low modes}

The spectrum supports three distinct interpretations:
\[
\begin{array}{rcl}
\lambda_k=0
&\Longleftrightarrow&
\text{exact compatible direction},\\
0<\lambda_k\ll g_d
&\Longrightarrow&
\text{approximately transported direction},\\
\lambda_k\approx\lambda_{k+1}
&\Longrightarrow&
\text{rank selection is unstable at that boundary}.
\end{array}
\]
A low eigenvalue can still be scientifically spurious if it arises from a
poorly chosen weight, a nearly disconnected graph, or a noncausal encoding.
The causal and null gates in later parts remain necessary.

\topic{14}{Normalized Coherence Defect}

For \(c\neq0\), define the Rayleigh defect
\[
\delta_{\mathrm{coh}}(c)
=
\frac{\mathcal E_H(c)}{\|c\|^2}.
\]
For a nonzero subspace \(W\subseteq\mathcal C_V\), define
\[
\delta_{\mathrm{coh}}(W)
=
\sup_{\substack{c\in W\\\|c\|=1}}
\mathcal E_H(c).
\]
If \(W\) is \(L_H\)-invariant, then
\[
\delta_{\mathrm{coh}}(W)
=
\|L_H|_W\|.
\]
In particular,
\[
\delta_{\mathrm{coh}}(\Gamma_\varepsilon(H))
\leq\varepsilon.
\]

\begin{theorem}[Coherence-to-core distance estimate]\label{thm:poincare-workspace}
Let
\[
W_0=\Gamma(H)=\ker L_H,
\qquad
r=\dim W_0,
\]
and suppose \(r<N\). Set
\[
\gamma_H=\lambda_{r+1}>0.
\]
Then every \(c\in\mathcal C_V\) satisfies
\[
\boxed{
\operatorname{dist}(c,W_0)^2
\leq
\frac{\mathcal E_H(c)}{\gamma_H}.
}
\]
Moreover,
\[
\mathcal E_H(c)
\leq
\|L_H\|\operatorname{dist}(c,W_0)^2.
\]
\end{theorem}

\begin{proof}
Write \(c=P_0c+Q_0c\), where \(P_0\) projects onto \(W_0\) and
\(Q_0=I-P_0\). Since \(L_HP_0=0\),
\[
\mathcal E_H(c)
=
\langle Q_0c,L_HQ_0c\rangle.
\]
The spectrum of \(L_H\) on \(W_0^\perp\) lies in
\([\gamma_H,\|L_H\|]\). Hence
\[
\gamma_H\|Q_0c\|^2
\leq
\mathcal E_H(c)
\leq
\|L_H\|\|Q_0c\|^2.
\]
Finally, \(\|Q_0c\|=\operatorname{dist}(c,W_0)\).
\end{proof}

This estimate is the exact sense in which a spectral gap turns small edgewise
disagreement into closeness to a genuinely compatible global section.

\subtopic{14.1 Two-node consistency cycle}

Let \(C_1=C_2=\mathbb R\), use unit weights, and take two directed transfers
\[
H_{12}=a,
\qquad
H_{21}=b.
\]
Then
\[
D_H
=
\begin{pmatrix}
-a&1\\
1&-b
\end{pmatrix},
\qquad
L_H
=
\begin{pmatrix}
a^2+1&-(a+b)\\
-(a+b)&b^2+1
\end{pmatrix}.
\]
Writing
\[
\tau=a^2+b^2+2,
\]
the eigenvalues are
\[
\lambda_\pm
=
\frac{1}{2}
\left(
\tau\pm\sqrt{\tau^2-4(ab-1)^2}
\right),
\]
and
\[
\det L_H=(ab-1)^2.
\]
Thus a nonzero exact shared core exists precisely when
\[
ab=1,
\]
which is exactly the trivial-holonomy condition. Away from exact consistency,
\[
\lambda_-\lambda_+=(ab-1)^2.
\]
When \(|ab-1|\) is small while \(\lambda_+\) stays bounded away from zero, the
soft workspace eigenvalue is quadratic in the cycle mismatch:
\[
\lambda_-
=
\frac{(ab-1)^2}{\lambda_+}.
\]
This example connects approximate spectral coherence directly to approximate
holonomy closure.

\topic{15}{Perturbation Stability Theorem}

The central robustness question is whether a selected low-energy subspace
persists when the estimated transfer maps change. The next lemma isolates the
finite-dimensional spectral fact needed for the answer.

\begin{lemma}[Separated spectral subspaces]\label{lem:separated-subspaces}
Let \(A\) and \(B=A+E\) be self-adjoint operators. Let \(P\) be an
\(A\)-spectral projector whose eigenvalues lie in \(( -\infty,a]\), and let
\(\widetilde Q\) be a \(B\)-spectral projector whose eigenvalues lie in
\([b,\infty)\), with \(b>a\). Then
\[
\|P\widetilde Q\|
\leq
\frac{\|E\|}{b-a}.
\]
\end{lemma}

\begin{proof}
On \(\ran P\) and \(\ran\widetilde Q\), set
\[
A_P=PAP|_{\ran P},
\qquad
B_Q=\widetilde QB\widetilde Q|_{\ran\widetilde Q}.
\]
The operator \(X=P\widetilde Q\) satisfies the Sylvester equation
\[
A_PX-XB_Q=-PE\widetilde Q.
\]
Let \(\mu=(a+b)/2\). Since
\[
\|e^{t(A_P-\mu I)}\|\leq e^{-(b-a)t/2},
\qquad
\|e^{-t(B_Q-\mu I)}\|\leq e^{-(b-a)t/2},
\]
the inverse Sylvester operator has the convergent integral representation
\[
X
=
\int_0^\infty
 e^{t(A_P-\mu I)}
 (PE\widetilde Q)
 e^{-t(B_Q-\mu I)}\,\dd t.
\]
Taking norms and integrating gives the result.
\end{proof}

\begin{theorem}[Spectral workspace stability]\label{thm:spectral-stability}
Let \(L\) and \(\widetilde L=L+E\) be self-adjoint positive semidefinite
operators on the same \(N\)-dimensional space. Suppose
\[
\lambda_d(L)<\lambda_{d+1}(L),
\qquad
g_d=\lambda_{d+1}(L)-\lambda_d(L)>0,
\]
and
\[
\eta=\|E\|<\frac{g_d}{2}.
\]
Let \(P_d\) and \(\widetilde P_d\) project onto the first \(d\) eigenvectors
of \(L\) and \(\widetilde L\), respectively. Then the rank-\(d\) cluster
remains separated and
\[
\boxed{
\|P_d-\widetilde P_d\|
\leq
\frac{\eta}{g_d-\eta}
\leq
\frac{2\eta}{g_d}.
}
\]
\end{theorem}

\begin{proof}
The min--max characterization gives the eigenvalue perturbation bound
\[
|\lambda_k(\widetilde L)-\lambda_k(L)|\leq\eta.
\]
Hence
\[
\lambda_{d+1}(\widetilde L)-\lambda_d(L)
\geq g_d-\eta
\]
and
\[
\lambda_{d+1}(L)-\lambda_d(\widetilde L)
\geq g_d-\eta.
\]
Apply Lemma~\ref{lem:separated-subspaces} first to
\(P_d(I-\widetilde P_d)\) and then, with the roles of \(L\) and
\(\widetilde L\) exchanged, to \((I-P_d)\widetilde P_d\). Orthogonal
projections of equal rank satisfy
\[
\|P_d-\widetilde P_d\|
=
\max\left\{
\|P_d(I-\widetilde P_d)\|,
\|(I-P_d)\widetilde P_d\|
\right\}.
\]
This yields the first bound. The second follows from \(\eta<g_d/2\).
\end{proof}

\subtopic{15.1 From transfer-map error to Laplacian error}

Suppose the edge maps are perturbed to
\[
\widetilde H_e=H_e+\Delta_e
\]
while the weights and vertex metrics remain fixed. Define
\[
\boxed{
\delta_H^2
=
\max_{i\in V}
\sum_{e:s(e)=i}w_e\|\Delta_e\|^2.
}
\]
Let \(F=D_{\widetilde H}-D_H\). Then
\[
(Fc)_e=-\sqrt{w_e}\,\Delta_ec_{s(e)},
\qquad
\|F\|\leq\delta_H.
\]
Since
\[
L_{\widetilde H}-L_H
=D_H^*F+F^*D_H+F^*F,
\]
we obtain
\[
\boxed{
\|L_{\widetilde H}-L_H\|
\leq
2\|D_H\|\delta_H+\delta_H^2.
}
\]

\begin{corollary}[Edgewise stability certificate]\label{cor:edge-stability}
If
\[
\eta_H=2\|D_H\|\delta_H+\delta_H^2
<\frac{g_d}{2},
\]
then
\[
\boxed{
\|P_d(H)-P_d(\widetilde H)\|
\leq
\frac{\eta_H}{g_d-\eta_H}.
}
\]
\end{corollary}

This converts declared edge-map uncertainty into an explicit bound on the
motion of the extracted workspace subspace.

\subtopic{15.2 Approximate operation closure}

Define the vertex and edge direct-sum representations
\[
\rho_V(a)=\bigoplus_{i\in V}\rho_i(a),
\qquad
\rho_E(a)=\bigoplus_{e\in E}\rho_{t(e)}(a).
\]
Lift the edgewise intertwining defects to
\[
(\mathfrak X_a c)_e
=
\sqrt{w_e}\,\Xi_e(a)c_{s(e)}.
\]
Then the exact commutator identity is
\[
\boxed{
D_H\rho_V(a)-\rho_E(a)D_H=-\mathfrak X_a.
}
\]
Furthermore,
\[
\|\mathfrak X_a\|
\leq
\beta(a),
\qquad
\beta(a)^2
=
\max_{i\in V}
\sum_{e:s(e)=i}w_e\|\Xi_e(a)\|^2.
\]

\begin{theorem}[Approximate operation-closure theorem]\label{thm:approx-operation-closure}
Let \(P_d\) project onto the first \(d\) eigenvectors of \(L_H\), set
\[
\lambda_d\leq\varepsilon,
\qquad
\mu=\lambda_{d+1}>0,
\]
and let \(Q_d=I-P_d\). Then every \(a\in\mathcal A\) satisfies
\[
\boxed{
\|Q_d\rho_V(a)P_d\|
\leq
\frac{
\|\rho_E(a)\|\sqrt{\varepsilon}+\beta(a)
}{\sqrt{\mu}}.
}
\]
\end{theorem}

\begin{proof}
For \(x\in\ran P_d\),
\[
\|D_Hx\|^2
=\langle x,L_Hx\rangle
\leq\varepsilon\|x\|^2.
\]
The commutator identity gives
\[
\|D_H\rho_V(a)x\|
\leq
\|\rho_E(a)\|\sqrt{\varepsilon}\|x\|
+
\beta(a)\|x\|.
\]
Because \(P_d\) and \(Q_d\) reduce \(L_H\), the two energy components are
orthogonal and
\[
\|D_HQ_dy\|\leq\|D_Hy\|.
\]
On \(\ran Q_d\), the spectrum of \(L_H\) is bounded below by \(\mu\), so
\[
\sqrt{\mu}\|Q_dy\|
\leq
\|D_HQ_dy\|.
\]
Apply these inequalities to \(y=\rho_V(a)x\), then take the supremum over unit
\(x\in\ran P_d\).
\end{proof}

\begin{corollary}[Recovery of exact operation closure]
If \(d=\dim\ker L_H\), \(\varepsilon=0\), and every edge intertwines \(a\)
exactly, then
\[
Q_d\rho_V(a)P_d=0.
\]
Thus Theorem~\ref{thm:approx-operation-closure} reduces to the exact
operation-closure result of Part~III.
\end{corollary}

\begin{corollary}[Closure under projector drift]\label{cor:closure-drift}
Let \(\widetilde P_d\) be another rank-\(d\) orthogonal projector and set
\[
\delta_P=\|P_d-\widetilde P_d\|.
\]
Then
\[
\boxed{
\|(I-\widetilde P_d)\rho_V(a)\widetilde P_d\|
\leq
\|Q_d\rho_V(a)P_d\|
+2\|\rho_V(a)\|\delta_P.
}
\]
Combining this estimate with Theorem~\ref{thm:spectral-stability} controls
operation leakage after noisy workspace extraction.
\end{corollary}

\begin{proof}
Subtract \(Q_d\rho_V(a)P_d\) and use
\[
I-\widetilde P_d-Q_d=P_d-\widetilde P_d.
\]
The two resulting perturbation terms are each bounded by
\(\|\rho_V(a)\|\delta_P\).
\end{proof}

\subtopic{15.3 Core--residual path accumulation}

Let a path have vertices
\[
i_0\xrightarrow{e_1}i_1\xrightarrow{e_2}\cdots
\xrightarrow{e_m}i_m.
\]
Write
\[
T_r=T_{e_r},
\quad
H_r=H_{e_r},
\quad
V_r=V_{e_r},
\]
and let
\[
\iota_i^C:C_i\hookrightarrow X_i,
\qquad
\iota_i^R:R_i\hookrightarrow X_i
\]
be the canonical inclusions. Empty products below are identities.

\begin{theorem}[Pathwise core--residual Duhamel formula]\label{thm:residual-duhamel}
The deviation of full path transport from pure core transport is exactly
\[
\boxed{
\begin{aligned}
&T_m\cdots T_1\iota_{i_0}^C
-
\iota_{i_m}^C H_m\cdots H_1\\
&\qquad=
\sum_{r=1}^{m}
T_m\cdots T_{r+1}
\iota_{i_r}^R V_r
H_{r-1}\cdots H_1.
\end{aligned}
}
\]
Consequently,
\[
\boxed{
\begin{aligned}
\|Q_{i_m}T_m\cdots T_1P_{i_0}\|
\leq
\sum_{r=1}^{m}
&\left(\prod_{\ell=r+1}^{m}\|T_\ell\|\right)
\|V_r\|\\
&\times
\left(\prod_{\ell=1}^{r-1}\|H_\ell\|\right),
\end{aligned}
}
\]
and the same right-hand side bounds
\[
\|P_{i_m}T_m\cdots T_1P_{i_0}-H_m\cdots H_1\|.
\]
\end{theorem}

\begin{proof}
For every step,
\[
T_r\iota_{i_{r-1}}^C
=
\iota_{i_r}^CH_r+
\iota_{i_r}^RV_r.
\]
Telescope the sequence
\[
T_m\cdots T_{r+1}\iota_{i_r}^C H_r\cdots H_1,
\qquad
r=0,1,\ldots,m.
\]
The difference between consecutive terms is exactly the \(r\)-th summand in
the displayed formula. Projection onto \(R_{i_m}\) or \(C_{i_m}\), followed
by submultiplicativity of the operator norm, gives the two bounds.
\end{proof}

\begin{corollary}[Uniform path bound]\label{cor:uniform-path-bound}
If along the path
\[
\|T_r\|\leq\tau,
\qquad
\|H_r\|\leq h,
\qquad
\|V_r\|\leq\nu,
\]
then the leakage and core-transport defects are bounded by
\[
\boxed{
\nu\sum_{r=1}^{m}\tau^{m-r}h^{r-1}
=
\begin{cases}
\displaystyle
\nu\frac{\tau^m-h^m}{\tau-h},&\tau\neq h,\\[1.1em]
\displaystyle
m\nu\tau^{m-1},&\tau=h.
\end{cases}
}
\]
In particular, if \(\tau\leq1\) and \(h\leq1\), the bound is at most
\(m\nu\).
\end{corollary}

The linear factor \(m\) in the merely nonexpansive regime is important. Small
single-edge leakage does not by itself justify arbitrarily long transfer
chains. Long-horizon stability requires contraction, cancellation, error
correction, or leakage that decreases with path length.

\subtopic{15.4 Residual-mediated reentry}

For two consecutive transfers,
\[
P_{i_2}T_2T_1P_{i_0}
=
H_2H_1+U_2V_1,
\]
\[
Q_{i_2}T_2T_1P_{i_0}
=
V_2H_1+K_2V_1.
\]
Thus leakage out of the core is first order in the blocks \(V_r\), while an
excursion that leaves and then returns to the core contains at least one
\(V\)-block and one later \(U\)-block. For a fixed path length, if
\[
\|U_r\|,\|V_r\|=O(\delta)
\]
and the diagonal blocks remain bounded, the residual-mediated correction to
pure core-to-core transport is \(O(\delta^2)\), whereas terminal residual
leakage can be \(O(\delta)\). This distinction separates loss of core
invariance from distortion caused by temporary excursions through local
residual degrees of freedom.


\theorypart{V}{Gauge Invariance and Abstract Identity}

Exact compatibility is a coordinate-free notion. Approximate spectra also
depend on the Hilbert metrics used to measure discrepancy. Gauge theory must
therefore distinguish a change of coordinates from a physical change of
metric.

\topic{16}{Local Coordinate Transformations}

For each \(i\in V\), let
\[
g_i\in GL(C_i),
\qquad
g=(g_i)_{i\in V}.
\]
Define the vertex and edge direct sums
\[
G_V=\bigoplus_{i\in V}g_i,
\qquad
G_E=\bigoplus_{e\in E}g_{t(e)}.
\]
The transformed transfers and local operation representations are
\[
H_e^g
=
 g_{t(e)}H_eg_{s(e)}^{-1},
\]
\[
\rho_i^g(a)
=
 g_i\rho_i(a)g_i^{-1}.
\]
The fundamental covariance identity is
\[
\boxed{
D_{H^g}G_V=G_ED_H.
}
\]
It follows directly by evaluating the \(e\)-component:
\[
\sqrt{w_e}
\left(
 g_{t(e)}c_{t(e)}
-
 g_{t(e)}H_ec_{s(e)}
\right).
\]

\topic{17}{Gauge Invariance Theorem}

\begin{theorem}[Exact gauge covariance]\label{thm:gauge-covariance}
The map
\[
G_V:\Gamma(H)\longrightarrow\Gamma(H^g)
\]
is an isomorphism. Therefore,
\[
\boxed{
\Gamma(H^g)=G_V\Gamma(H),
\qquad
\dim\Gamma(H^g)=\dim\Gamma(H).
}
\]
If the core transports are invertible and \(o\) is a root, then
\[
H_\gamma^g=g_oH_\gamma g_o^{-1}
\]
for every loop \(\gamma:o\to o\). Hence the holonomy-fixed core transforms by
\[
C_*^g=g_oC_*.
\]
\end{theorem}

\begin{proof}
The covariance identity implies
\[
D_{H^g}(G_Vc)=G_ED_Hc.
\]
Since \(G_E\) is invertible,
\[
D_Hc=0
\quad\Longleftrightarrow\quad
D_{H^g}(G_Vc)=0.
\]
The holonomy formula follows by cancellation of adjacent factors
\(g_i^{-1}g_i\) along a path.
\end{proof}

\begin{theorem}[Unitary gauge covariance of the Laplacian]\label{thm:unitary-gauge}
If every \(g_i\) is unitary, then \(G_V\) and \(G_E\) are unitary and
\[
\boxed{
L_{H^g}=G_VL_HG_V^{-1}.
}
\]
Consequently,
\[
\spec(L_{H^g})=\spec(L_H),
\qquad
P_I(L_{H^g})=G_VP_I(L_H)G_V^{-1}
\]
for every spectral interval \(I\).
\end{theorem}

\begin{proof}
From
\[
D_{H^g}=G_ED_HG_V^{-1}
\]
and unitarity,
\[
D_{H^g}^*D_{H^g}
=
G_VD_H^*G_E^{-1}G_ED_HG_V^{-1}
=
G_VL_HG_V^{-1}.
\]
The spectral statements follow from unitary similarity.
\end{proof}

\subtopic{17.1 Nonunitary gauges and metric distortion}

For a general invertible gauge, define
\[
m_V=\sigma_{\min}(G_V),
\qquad
M_V=\|G_V\|,
\]
\[
m_E=\sigma_{\min}(G_E),
\qquad
M_E=\|G_E\|.
\]
The covariance identity yields
\[
m_E^2\mathcal E_H(c)
\leq
\mathcal E_{H^g}(G_Vc)
\leq
M_E^2\mathcal E_H(c),
\]
and
\[
m_V^2\|c\|^2
\leq
\|G_Vc\|^2
\leq
M_V^2\|c\|^2.
\]
Therefore the Rayleigh quotients obey
\[
\boxed{
\frac{m_E^2}{M_V^2}
\frac{\mathcal E_H(c)}{\|c\|^2}
\leq
\frac{\mathcal E_{H^g}(G_Vc)}{\|G_Vc\|^2}
\leq
\frac{M_E^2}{m_V^2}
\frac{\mathcal E_H(c)}{\|c\|^2}.
}
\]
By the min--max principle,
\[
\boxed{
\frac{m_E^2}{M_V^2}\lambda_k(L_H)
\leq
\lambda_k(L_{H^g})
\leq
\frac{M_E^2}{m_V^2}\lambda_k(L_H).
}
\]
Exact zero modes are unchanged, but fixed-metric approximate eigenvalues can be
strongly distorted by an ill-conditioned coordinate transformation.

\subtopic{17.2 Metric-covariant formulation}

A pure change of coordinates should transform the metric as well as the
matrices. On each transformed space, define
\[
\langle u,v\rangle_i^g
=
\left\langle g_i^{-1}u,g_i^{-1}v\right\rangle_i.
\]
Use the analogous pullback metric on every edge copy. Then \(G_V\) and
\(G_E\) are unitary between the original and transformed metric spaces.
Consequently, the complete transfer-Laplacian spectrum is exactly preserved.

This gives two distinct interpretations:
\begin{itemize}
\item Transforming the matrices and the metrics gives the same approximate
workspace in new coordinates.
\item Transforming the matrices while holding the metrics fixed gives a
potentially different weighting of discrepancy.
\end{itemize}
The distinction prevents arbitrary coordinate scaling from manufacturing or
erasing apparent low-energy workspace directions.

\topic{18}{Static Workspace Equivalence}

\begin{definition}[Exact static equivalence]\label{def:static-equivalence}
Two static realizations are exactly gauge equivalent,
\[
\mathfrak R\sim_{\mathrm{stat}}\mathfrak R',
\]
when there exists an invertible family \(g_i\) such that
\[
H_e'=g_{t(e)}H_eg_{s(e)}^{-1},
\qquad
\rho_i'(a)=g_i\rho_i(a)g_i^{-1}.
\]
The exact static abstract workspace is
\[
\boxed{
\Workspace_{\mathrm{stat}}
=
[\mathfrak R]_{\sim_{\mathrm{stat}}}.
}
\]
\end{definition}

For approximate theory, the realization also contains the vertex and edge
metrics. Two metric realizations are equivalent only when the gauge transports
those metrics, or when the gauge is unitary relative to the fixed metrics.
Accordingly:
\[
\boxed{
\begin{aligned}
\text{exact workspace identity}
&\text{ is invariant under }GL\text{-gauge},\\
\text{fixed-metric spectral geometry}
&\text{ is invariant under unitary gauge},\\
\text{metric-covariant spectral geometry}
&\text{ is invariant under transported }GL\text{-gauge}.
\end{aligned}
}
\]

\clearpage
\subtopic{18.1 What the approximate theory establishes}

Parts~IV--V now prove:
\[
\begin{aligned}
&\ker L_H=\Gamma(H),\\
&\text{low energy plus a gap implies proximity to an exact core},\\
&\text{a separated workspace cluster is stable under transfer noise},\\
&\text{approximately intertwined operations nearly preserve that cluster},\\
&\text{core leakage and residual reentry admit explicit pathwise bounds},\\
&\text{exact identity is gauge invariant while approximate geometry is metric aware}.
\end{aligned}
\]
These are structural and perturbative results. They still do not prove that a
workspace variable is causally used, predictively sufficient, irreducibly
distributed, or temporally self-maintaining. Those questions begin in
Part~VI.

\subtopic{18.2 Finite computation protocol}

For a declared graph, set of core dimensions, weights, and transfer maps, the
static analysis can be executed in the following order:
\begin{enumerate}
\item Assemble the sparse block operator \(D_H\) from the edge equations
\[
\sqrt{w_e}\bigl(c_{t(e)}-H_ec_{s(e)}\bigr).
\]
\item Compute a singular-value decomposition
\[
D_H=U\Sigma V^*,
\qquad
\lambda_k(L_H)=\sigma_k(D_H)^2.
\]
Use a declared numerical tolerance to distinguish computational zero modes
from positive modes.
\item Select a candidate rank \(d\) only at a reproducible spectral boundary
with \(g_d>0\). A large observed gap is evidence of stability, not evidence of
causal relevance.
\item Quantify transfer uncertainty by \(\delta_H\) and verify whether
\[
2\|D_H\|\delta_H+\delta_H^2<\frac{g_d}{2}.
\]
If this gate fails, the extracted projector has no certified perturbative
identity at the stated error scale.
\item For every proposed common operation, compute \(\beta(a)\) and the
operation-leakage bound in Theorem~\ref{thm:approx-operation-closure}.
\item For every long transfer chain used in an interpretation, report the
pathwise leakage certificate from Theorem~\ref{thm:residual-duhamel} rather
than extrapolating from single-edge defects.
\end{enumerate}

A minimal structural report is therefore
\[
\boxed{
\left(
 d,\lambda_d,\lambda_{d+1},g_d,
 \delta_H,
 \|P_d-\widetilde P_d\|_{\mathrm{bound}},
 \beta(a),
 \Lambda_\gamma
\right),
}
\]
where \(\Lambda_\gamma\) denotes the right-hand side of the pathwise leakage
bound. This structural report is intentionally incomplete as a workspace
certificate. Parts~VI--IX add the predictive, causal, null, and irreducibility
gates required to interpret the low-energy subspace operationally.


\theorypart{VI}{Dynamic Workspace Theory}

\topic{19}{Probabilistic Local and Relational Dynamics}

The static theory identifies compatible directions across a declared transfer
network. A dynamic workspace claim requires additional probabilistic objects:
a history, a target future, a decision problem, and an intervention family.
These must be fixed before sufficiency or causal efficacy is evaluated.

\begin{definition}[Dynamic relational realization]\label{def:dynamic-realization}
A dynamic relational realization is defined on a probability space
\((\Omega,\mathcal F,\mathbb P)\) and consists of:
\begin{itemize}
\item node states \(X_i(t)\in\mathsf X_i\) for \(i\in V\);
\item edge messages \(M_e(t)\in\mathsf M_e\) for \(e\in E\);
\item external inputs \(E_t\), local actions \(A_i(t)\), and disturbances
 \(\xi_i(t),\zeta_e(t)\);
\item time-dependent transfer data \(H_e(t)\), local core vectors \(c_i(t)\),
 and transfer Laplacians \(L_t=L_{H(t)}\); and
\item a future target \(Y_t^+\in\mathsf Y\).
\end{itemize}
All state, history, and target spaces are assumed standard Borel unless a
stronger topology is explicitly invoked.
\end{definition}

A representative local update is
\[
X_i(t+1)
=
F_i\!\left(
X_i(t),
\{M_e(t):t(e)=i\},
A_i(t),E_t,\xi_i(t)
\right),
\]
with messages
\[
M_e(t)
=
K_e\!\left(
X_{s(e)}(t),
W_t,
E_t,\zeta_e(t)
\right).
\]
This notation permits, but does not require, the current workspace state to
modulate future transfers.

For a coalition \(A\subseteq V\), let \(\mathbf H_t^A\) denote the random
history available to that coalition through time \(t\). Its generated
information is
\[
\mathcal F_t^A=\sigma(\mathbf H_t^A).
\]
The full relational history is \(\mathbf H_t=\mathbf H_t^V\), with
\[
\mathcal F_t^A\subseteq\mathcal F_t^B
\qquad
\text{whenever }A\subseteq B.
\]
The empty-coalition history \(\mathbf H_t^\varnothing\) is the trivial
constant statistic and represents prediction from the background target law
alone.
A concrete history may contain node trajectories, incident messages, actions,
artifacts, transfer-map estimates, and the declared graph history. The exact
contents must be reported because coalition sufficiency is relative to what a
coalition is permitted to observe.

The target \(Y_t^+\) can be a next output, a finite future trajectory, a
research decision, a proof outcome, or any measurable downstream quantity.
Every predictive workspace claim is therefore target-relative:
\[
\operatorname{WorkspaceFor}(Y_t^+)
\centernot\Longrightarrow
\operatorname{WorkspaceFor}(\widetilde Y_t^+).
\]

Let
\[
c(t)=(c_i(t))_{i\in V}
\in
\mathcal C_V(t),
\qquad
P_t=\mathbf 1_{[0,\varepsilon_t]}(L_t).
\]
The instantaneous structural workspace coordinate is
\[
\boxed{
W_t^{\mathrm{str}}=P_tc(t)
\in
\Gamma_{\varepsilon_t}(H(t)).
}
\]
When \(P_t\) varies with time, these coordinates inhabit a moving family of
subspaces. Parts~X--XII will supply the transport required to compare them
across time. In what follows, all sufficiency and intervention statements
are made at a fixed decision time \(t\).

\begin{remark}[Three objects that must not be conflated]
The realization state \(S_t\), the complete history \(\mathbf H_t\), and the
structural coordinate \(W_t^{\mathrm{str}}\) are different objects. The first
is an instantaneous physical or computational state, the second contains all
admissible past information, and the third is a low-transfer-energy statistic
extracted from the declared core geometry.
\end{remark}

\topic{20}{Canonical Predictive and Interventional Causal States}

The quotient notation
\([\mathbf H_t]_{\sim_Y}\) expresses the right intuition but hides
measurability and null-set issues. A canonical construction is obtained from
regular conditional distributions.

\begin{definition}[Predictive kernel and canonical state]\label{def:predictive-kernel}
Let \(\mathsf H_t\) and \(\mathsf Y\) be standard Borel spaces. A regular
conditional kernel is a measurable map
\[
K_t:\mathsf H_t\times\mathcal B(\mathsf Y)\to[0,1]
\]
satisfying
\[
K_t(h,A)
=
\mathbb P(Y_t^+\in A\mid\mathbf H_t=h)
\]
for \(\mathbb P_{\mathbf H_t}\)-almost every \(h\). Let
\(\mathcal P(\mathsf Y)\) carry its evaluation sigma-algebra. The canonical
predictive state is the random probability measure
\[
\boxed{
Z_t
=
K_t(\mathbf H_t,\cdot)
\in
\mathcal P(\mathsf Y).
}
\]
Two histories are predictively equivalent precisely when they induce the same
conditional law of \(Y_t^+\).
\end{definition}

\begin{theorem}[Predictive-state sufficiency and minimality, \(T_{12}\)]
\label{thm:predictive-state}
Under Definition~\ref{def:predictive-kernel}:
\begin{enumerate}
\item \(Z_t\) is predictively sufficient:
\[
Y_t^+\indep\mathbf H_t\mid Z_t.
\]
\item If \(S_t=s(\mathbf H_t)\) is any other history statistic satisfying
\[
Y_t^+\indep\mathbf H_t\mid S_t,
\]
then there is a measurable map \(\phi\) such that
\[
Z_t=\phi(S_t)
\qquad\text{almost surely}.
\]
\item Consequently, any two minimal sufficient predictive states are
measurably isomorphic after removal of null sets.
\end{enumerate}
\end{theorem}

\begin{proof}
For every bounded measurable \(f:\mathsf Y\to\mathbb R\), define
\[
Z_t(f)=\int_{\mathsf Y}f(y)\,Z_t(\mathrm dy).
\]
By construction,
\[
\mathbb E[f(Y_t^+)\mid\mathbf H_t]=Z_t(f).
\]
Since \(Z_t\) is a measurable function of \(\mathbf H_t\), the tower property
gives
\[
\mathbb E[f(Y_t^+)\mid Z_t]
=
\mathbb E[Z_t(f)\mid Z_t]
=
Z_t(f).
\]
Hence conditioning additionally on \(\mathbf H_t\) does not change the
conditional expectation, which is the stated conditional independence.

Now let \(S_t=s(\mathbf H_t)\) be sufficient. Then, for every measurable
\(A\subseteq\mathsf Y\),
\[
K_t(\mathbf H_t,A)
=
\mathbb P(Y_t^+\in A\mid\mathbf H_t)
=
\mathbb P(Y_t^+\in A\mid S_t)
\qquad\text{a.s.}
\]
Choose a countable determining class for probability measures on the standard
Borel space \(\mathsf Y\). The values of \(K_t(\mathbf H_t,\cdot)\) on this
class are measurable functions of \(S_t\), and therefore the complete random
measure \(Z_t\) is a measurable function of \(S_t\). Applying this conclusion
in both directions to two minimal sufficient states gives the final claim.
\end{proof}

\begin{remark}[Meaning of minimality]
The theorem does not say that \(Z_t\) has the fewest numerical coordinates in
a chosen parametrization. It says that every sufficient history statistic
must determine \(Z_t\) almost surely. This is an information order, invariant
under one-to-one reparametrization.
\end{remark}

For an arbitrary history statistic \(S_t=s(\mathbf H_t)\), let
\[
K_t^{S}(s,\cdot)
=
\mathbb P(Y_t^+\in\cdot\mid S_t=s).
\]
Whenever the conditional relative entropy is finite, define the sufficiency
defect
\[
\boxed{
\delta_{\mathrm{suff}}(S_t)
=
\mathbb E\!\left[
\KL\!\left(
K_t(\mathbf H_t,\cdot)
\,\middle\|\,
K_t^S(S_t,\cdot)
\right)
\right].
}
\]

\begin{proposition}[Conditional-information identity]\label{prop:conditional-information}
If \(S_t\) is a measurable function of \(\mathbf H_t\), then
\[
\boxed{
\delta_{\mathrm{suff}}(S_t)
=
I(Y_t^+;\mathbf H_t\mid S_t).
}
\]
In particular, the defect vanishes if and only if \(S_t\) is predictively
sufficient, up to the usual almost-sure qualifications.
\end{proposition}

\begin{proof}
The conditional mutual information is the expected relative entropy between
the conditional law given \((\mathbf H_t,S_t)\) and the conditional law given
\(S_t\). Because \(S_t=s(\mathbf H_t)\), the first of these is exactly
\(K_t(\mathbf H_t,\cdot)\).
\end{proof}

Sufficiency has a direct decision-theoretic interpretation. Use the
convention
\[
\TV(P,Q)=\sup_{A}|P(A)-Q(A)|.
\]
Let \(\mathsf A\) be a finite action set and
\(\ell:\mathsf Y\times\mathsf A\to[0,1]\) a bounded loss. Define
\[
\Risk_\ell(S_t)
=
\inf_{d}
\mathbb E\!\left[
\ell\bigl(Y_t^+,d(S_t)\bigr)
\right].
\]

\begin{proposition}[Approximate sufficiency bounds excess Bayes risk]
\label{prop:risk-pinsker}
Under the preceding assumptions,
\[
\boxed{
0
\leq
\Risk_\ell(S_t)-\Risk_\ell(\mathbf H_t)
\leq
\mathbb E\!\left[
\TV\!\left(K_t(\mathbf H_t),K_t^S(S_t)\right)
\right]
\leq
\sqrt{\frac{\delta_{\mathrm{suff}}(S_t)}{2}}.
}
\]
\end{proposition}

\begin{proof}
More information cannot increase optimal risk, giving the first inequality.
For any fixed action, its conditional expected loss changes by at most total
variation when the conditional law changes. Taking the infimum over actions
preserves this bound. Conditional Pinsker inequality followed by Jensen's
inequality gives the final estimate.
\end{proof}

For a finite target alphabet and logarithmic scoring loss
\[
\ell_{\log}(y,q)=-\log q(y),
\qquad
q\in\mathcal P(\mathsf Y),
\]
the identity is exact:
\[
\boxed{
\Risk_{\log}(S_t)=H(Y_t^+\mid S_t),
\qquad
\Risk_{\log}(S_t)-\Risk_{\log}(\mathbf H_t)
=
I(Y_t^+;\mathbf H_t\mid S_t).
}
\]
Natural logarithms measure the gap in nats; base-two logarithms measure it in
bits.

\subtopic{Intervention-indexed causal state}

Observational prediction may merge histories that behave identically under the
current policy but respond differently to future interventions. Let
\(\mathcal I\) be a finite or countable family of interventions applied only
after time \(t\). For each \(\iota\in\mathcal I\), let
\[
K_t^{\iota}(h,\cdot)
=
\mathbb P^{\iota}(Y_t^+\in\cdot\mid\mathbf H_t=h)
\]
be a regular post-intervention kernel on the common pre-intervention history
space. Define
\[
\boxed{
Z_t^{\mathcal I}
=
\bigl(K_t^{\iota}(\mathbf H_t,\cdot)\bigr)_{\iota\in\mathcal I}
\in
\mathcal P(\mathsf Y)^{\mathcal I}.
}
\]

\begin{corollary}[Minimal simultaneous interventional state]
\label{cor:interventional-state}
The random element \(Z_t^{\mathcal I}\) is sufficient simultaneously for all
post-intervention target laws. If a statistic \(S_t=s(\mathbf H_t)\)
determines the conditional target law for every \(\iota\in\mathcal I\), then
there exists a measurable \(\phi_{\mathcal I}\) such that
\[
Z_t^{\mathcal I}=\phi_{\mathcal I}(S_t)
\qquad\text{a.s.}
\]
\end{corollary}

\begin{proof}
Apply Theorem~\ref{thm:predictive-state} componentwise. A countable product of
standard Borel spaces is standard Borel, so the component maps combine into a
single measurable factor map.
\end{proof}

The observational state is the special case
\(\mathcal I=\{\operatorname{id}\}\). A causal workspace claim should align
with \(Z_t^{\mathcal I}\) for a declared intervention family rather than only
with the observational state.

\topic{21}{Structural--Operational Alignment}

A structural workspace coordinate becomes operational only when it retains the
canonical predictive or interventional state relevant to the declared task.
Let
\[
W_t=\pi_t(\mathbf H_t)
\]
be a measurable candidate workspace statistic satisfying the structural support
condition
\[
W_t\in\Gamma_{\varepsilon_t}(H(t))
\qquad\text{almost surely}.
\]
The map \(\pi_t\) may include encoding into a moving low-energy bundle; no
canonical operational interpretation follows from the spectral projector
alone.

\begin{definition}[Exact structural--operational alignment]
The candidate \(W_t\) is observationally aligned when there exists a measurable
map \(\phi_t\) such that
\[
Z_t=\phi_t(W_t)
\qquad\text{a.s.}
\]
It is \(\mathcal I\)-causally aligned when
\[
Z_t^{\mathcal I}=\phi_t^{\mathcal I}(W_t)
\qquad\text{a.s.}
\]
\end{definition}

\begin{theorem}[Exact alignment criterion]\label{thm:alignment-criterion}
For any workspace statistic \(W_t=\pi_t(\mathbf H_t)\), the following are
equivalent:
\begin{enumerate}
\item \(Y_t^+\indep\mathbf H_t\mid W_t\);
\item \(I(Y_t^+;\mathbf H_t\mid W_t)=0\);
\item the canonical predictive state factors through \(W_t\):
\[
Z_t=\phi_t(W_t)
\quad\text{a.s.}
\]
for some measurable \(\phi_t\).
\end{enumerate}
The analogous equivalence holds simultaneously over \(\mathcal I\) with
\(Z_t^{\mathcal I}\).
\end{theorem}

\begin{proof}
The equivalence of the first two statements is the zero-divergence
characterization of conditional mutual information. The first statement and
Theorem~\ref{thm:predictive-state} imply the factorization in the third.
Conversely, if \(Z_t=\phi_t(W_t)\), then
\[
\mathbb P(Y_t^+\in A\mid\mathbf H_t)
=Z_t(A)
=\phi_t(W_t)(A),
\]
which depends on the history only through \(W_t\), proving sufficiency. The
interventional statement follows from
Corollary~\ref{cor:interventional-state}.
\end{proof}

Define the observational alignment defect by
\[
\boxed{
\delta_{\mathrm{align}}^{\mathrm{obs}}(W_t)
=
I_{\mathbb P}(Y_t^+;\mathbf H_t\mid W_t).
}
\]
For a countable intervention family with declared weights
\(q_\iota>0\), \(\sum_\iota q_\iota=1\), define
\[
\boxed{
\delta_{\mathrm{align}}^{\mathcal I}(W_t)
=
\sum_{\iota\in\mathcal I}
q_\iota
I_{\mathbb P^\iota}
(Y_t^+;\mathbf H_t\mid W_t).
}
\]
For bounded decision losses, Proposition~\ref{prop:risk-pinsker} applies under
each intervention and gives the weighted bound
\[
\sum_{\iota}q_\iota
\left[
\Risk_\ell^{\iota}(W_t)
-
\Risk_\ell^{\iota}(\mathbf H_t)
\right]
\leq
\sqrt{\frac{
\delta_{\mathrm{align}}^{\mathcal I}(W_t)}{2}}.
\]

\subtopic{Structural coherence and operational alignment are independent}

The static defect
\[
\delta_{\mathrm{coh}}(W_t)
\]
and the operational defect
\[
\delta_{\mathrm{align}}^{\mathcal I}(W_t)
\]
answer different questions. The first measures compatibility under transfer.
The second measures loss of task-relevant conditional law.

\begin{center}
\small
\begin{tabularx}{\textwidth}{@{}>{\raggedright\arraybackslash}p{0.22\textwidth}>{\raggedright\arraybackslash}p{0.34\textwidth}>{\raggedright\arraybackslash}X@{}}
\toprule
Regime & Structural status & Operational interpretation \\
\midrule
Low coherence defect, low alignment defect & Stable transported coordinate &
Candidate operational workspace for the declared target and interventions. \\
Low coherence defect, high alignment defect & Stable but task-irrelevant
shared structure & Archive, convention, or synchronized representation. \\
High coherence defect, low alignment defect & Predictively adequate but not
transport-stable & Localized or brittle sufficient statistic. \\
High coherence defect, high alignment defect & Neither stable nor sufficient &
Rejected candidate at the stated scales. \\
\bottomrule
\end{tabularx}
\end{center}

Thus neither low transfer energy nor high predictive performance alone is
enough. The combined gate is
\[
\boxed{
\delta_{\mathrm{coh}}(W_t)
\leq
\varepsilon_{\mathrm{coh}},
\qquad
\delta_{\mathrm{align}}^{\mathcal I}(W_t)
\leq
\varepsilon_{\mathrm{align}}.
}
\]

\begin{example}[Coherent but irrelevant archive]
Suppose every node perfectly copies a common random variable \(R\), so the
corresponding transfer energy is zero. If \(R\indep Y_t^+\), then
\[
\delta_{\mathrm{coh}}(R)=0,
\qquad
\delta_{\mathrm{align}}^{\mathrm{obs}}(R)
=I(Y_t^+;\mathbf H_t\mid R),
\]
which can be as large as the entire predictive information in the history.
Perfect agreement therefore does not imply workspace functionality.
\end{example}

\begin{example}[Sufficient but localized statistic]
If one node contains the canonical state \(Z_t\) and all transfers from that
node are unstable, then \(W_t=Z_t\) is sufficient but need not lie in a stable
low-energy transfer subspace. This is a local controller or archive, not yet a
stratified shared workspace.
\end{example}


\theorypart{VII}{Relational Irreducibility}

\topic{22}{Coalition Information and Bayes Risk}

For \(A\subseteq V\), the coalition statistic is \(\mathbf H_t^A\). Given an
action space \(\mathsf A\) and loss \(\ell\), define its unrestricted Bayes
risk by
\[
\boxed{
\Risk_\ell(A)
=
\inf_d
\mathbb E\!\left[
\ell\bigl(Y_t^+,d(\mathbf H_t^A)\bigr)
\right].
}
\]
The infimum is over measurable decision rules for which the expectation is
defined. The full-system risk is \(\Risk_\ell(V)\).

\begin{proposition}[Coalition monotonicity]\label{prop:coalition-monotonicity}
If \(A\subseteq B\subseteq V\), then
\[
\Risk_\ell(B)\leq\Risk_\ell(A).
\]
\end{proposition}

\begin{proof}
Every decision rule available to coalition \(A\) is also available to coalition
\(B\), which can ignore its additional information.
\end{proof}

For finite \(\mathsf Y\), logarithmic loss converts risk gaps into conditional
information exactly.

\begin{theorem}[Coalition log-risk identity]\label{thm:coalition-log-risk}
Assume \(H(Y_t^+)<\infty\). If \(A\subseteq B\subseteq V\), then
\[
\boxed{
\Risk_{\log}(A)-\Risk_{\log}(B)
=
I\!\left(
Y_t^+;\mathbf H_t^B
\mid
\mathbf H_t^A
\right).
}
\]
In particular,
\[
\Risk_{\log}(A)-\Risk_{\log}(V)
=
I\!\left(
Y_t^+;\mathbf H_t
\mid
\mathbf H_t^A
\right).
\]
\end{theorem}

\begin{proof}
Under logarithmic loss the optimal report is the conditional target law, so
\[
\Risk_{\log}(A)=H(Y_t^+\mid\mathbf H_t^A).
\]
Because \(\mathbf H_t^A\) is a function of \(\mathbf H_t^B\), subtracting the
two conditional entropies gives the stated conditional mutual information.
\end{proof}

The theorem makes the decomposition explicit: irreducibility is not an
unqualified property of a system. It is a property relative to the chosen node
partition, accessible histories, target, and decision loss.

\topic{23}{Joint Advantage and Its Scope}

Define the coalition gap
\[
\Delta_A^{\ell}
=
\Risk_\ell(A)-\Risk_\ell(V)
\geq0.
\]
The node-level joint advantage is
\[
\boxed{
\Delta_{\mathrm{joint}}^{\ell}
=
\min_{i\in V}
\Delta_{\{i\}}^{\ell}.
}
\]
Under logarithmic loss,
\[
\boxed{
\Delta_{\mathrm{joint}}^{\log}
=
\min_{i\in V}
I\!\left(
Y_t^+;\mathbf H_t
\mid
\mathbf H_t^{\{i\}}
\right).
}
\]
Therefore
\[
\Delta_{\mathrm{joint}}^{\log}>0
\]
means that every individual component omits some predictive information
available to the complete system.

This quantity should not automatically be identified with a particular notion
of partial-information synergy. A positive gap can arise from complementary
information distributed across components, from higher-order interaction, or
from both. The theory uses the operationally weaker and more directly
testable statement:
\[
\text{no individual node is sufficient for the declared target.}
\]

For \(0\leq r\leq |V|\), define the order-\(r\) coalition gap profile by
\[
\boxed{
\Delta^{(r)}_\ell
=
\min_{\substack{A\subseteq V\\|A|\leq r}}
\left[
\Risk_\ell(A)-\Risk_\ell(V)
\right].
}
\]
The profile is nonincreasing in \(r\) and reaches zero by \(r=|V|\).
It records how rapidly operative capacity becomes reconstructible as larger
coalitions are permitted.

\topic{24}{Minimal Sufficient Coalition}

For tolerance \(\varepsilon\geq0\), define
\[
\boxed{
k_\varepsilon
=
\min\left\{
|A|:
\Risk_\ell(A)
\leq
\Risk_\ell(V)+\varepsilon
\right\}.
}
\]
The set of minimal coalitions is
\[
\mathcal K_\varepsilon
=
\left\{
A\subseteq V:
|A|=k_\varepsilon,
\Risk_\ell(A)
\leq
\Risk_\ell(V)+\varepsilon
\right\}.
\]

\begin{corollary}[Information form of coalition size]
Under the assumptions of Theorem~\ref{thm:coalition-log-risk},
\[
\boxed{
k_\varepsilon
=
\min\left\{
|A|:
I(Y_t^+;\mathbf H_t\mid\mathbf H_t^A)
\leq\varepsilon
\right\}.
}
\]
Equivalently, \(k_\varepsilon\) is the first \(r\) for which
\(\Delta_{\log}^{(r)}\leq\varepsilon\).
\end{corollary}

This yields a graded classification:
\[
\begin{aligned}
k_0=0
&\quad\Longrightarrow\quad
\text{the background law alone is sufficient},\\
k_0=1
&\quad\Longrightarrow\quad
\text{some individual node is sufficient},\\
k_0\geq2
&\quad\Longrightarrow\quad
\text{node-level relational irreducibility},\\
k_0=|V|
&\quad\Longrightarrow\quad
\text{every proper coalition is insufficient}.
\end{aligned}
\]
The last condition is much stronger and should not be inferred from the first.

\topic{25}{Reconstructibility of the Canonical State}

Let \(d_{\mathcal P}\) be a bounded metric generating the measurable structure
on \(\mathcal P(\mathsf Y)\). Define the reconstruction defect of the canonical
state from coalition \(A\) by
\[
\delta_A^Z
=
\inf_{\psi_A}
\mathbb E\!\left[
 d_{\mathcal P}\!\left(
 Z_t,\psi_A(\mathbf H_t^A)
 \right)
\right].
\]

\begin{proposition}[Sufficiency--reconstruction equivalence]
\label{prop:sufficiency-reconstruction}
For a coalition \(A\subseteq V\), the following are equivalent:
\begin{enumerate}
\item \(\mathbf H_t^A\) is sufficient for \(Y_t^+\) relative to the full
 history;
\item \(Z_t\) is measurable with respect to \(\mathcal F_t^A\);
\item there exists a measurable \(\psi_A\) such that
\[
Z_t=\psi_A(\mathbf H_t^A)
\qquad\text{almost surely}.
\]
\end{enumerate}
Each condition implies \(\delta_A^Z=0\). Under finite-target logarithmic loss,
they are also equivalent to
\[
\Risk_{\log}(A)=\Risk_{\log}(V).
\]
If the reconstruction infimum is known to be attained, then
\(\delta_A^Z=0\) is itself equivalent to the three conditions.
\end{proposition}

\begin{proof}
The equivalence of the first two statements is
Theorem~\ref{thm:predictive-state}; the second and third are two formulations
of measurability up to null sets. An exact reconstructor has zero expected
metric error. The log-risk statement follows from
Theorem~\ref{thm:coalition-log-risk} and the zero condition for conditional
mutual information. Without an attainment or compactness condition, an
infimum of zero alone certifies arbitrarily accurate reconstruction but need
not supply a single exact minimizing map.
\end{proof}

For an arbitrary enriched workspace variable \(W_t\), local
non-reconstructibility need not imply predictive irreducibility: \(W_t\) may
contain task-irrelevant material. The equivalence is exact for the canonical
state \(Z_t\), or for a candidate \(W_t\) known to be minimally aligned with
it.

\topic{26}{Matched Relational Nulls}

A null control must preserve the local facts that the claim is not intended to
challenge while destroying the relational organization being tested. Merely
reducing model capacity or replacing the workspace by noise is not a matched
null.

Let \(P\) be the intact law of
\((Y_t^+,\mathbf H_t^1,\ldots,\mathbf H_t^n)\). Define the local-marginal null
class
\[
\mathcal N_{\mathrm{marg}}(P)
=
\left\{
Q:
Q_{\mathbf H_t^i}=P_{\mathbf H_t^i}
\text{ for every }i
\right\},
\]
and the stronger local-predictive-matched class
\[
\boxed{
\mathcal N_{\mathrm{pred}}(P)
=
\left\{
Q:
Q_{(Y_t^+,\mathbf H_t^i)}
=
P_{(Y_t^+,\mathbf H_t^i)}
\text{ for every }i
\right\}.
}
\]
Because \(\mathbf H_t^i\) can include the complete local trajectory, the latter
class can preserve all single-component temporal and predictive behavior while
altering cross-component organization.

\begin{proposition}[Local-risk invariance under predictive matching]
\label{prop:local-risk-null}
If \(Q\in\mathcal N_{\mathrm{pred}}(P)\), then for every node \(i\) and every
loss determined by \((Y_t^+,d(\mathbf H_t^i))\),
\[
\boxed{
\Risk_{\ell,Q}(\{i\})
=
\Risk_{\ell,P}(\{i\}).
}
\]
\end{proposition}

\begin{proof}
Every such risk depends only on the joint law of
\((Y_t^+,\mathbf H_t^i)\), which is identical under \(P\) and \(Q\).
\end{proof}

A valid null must also satisfy a declared relation-breaking condition. Examples
include:
\begin{itemize}
\item permuting source--target message pairings while preserving each message
 marginal;
\item shuffling temporal correspondence across nodes while preserving every
 local trajectory;
\item permuting role bindings while preserving the local token or state
 distributions;
\item replacing an artifact by a matched-size, matched-format artifact whose
 relation to the surrounding history has been destroyed; and
\item randomizing transfer channels through a physical or software
 intervention while preserving message rate and marginal content.
\end{itemize}

For a predeclared null \(Q\), define
\[
\boxed{
\Delta_{\mathrm{tr}}(Q)
=
\Risk_{\ell,Q}(V)
-
\Risk_{\ell,P}(V).
}
\]
For a finite validated panel
\(\mathcal Q=\{Q_1,\ldots,Q_m\}\subseteq\mathcal N_{\mathrm{pred}}(P)\), report
both the complete vector
\[
\boldsymbol\Delta_{\mathrm{tr}}(\mathcal Q)
=
\bigl(
\Delta_{\mathrm{tr}}(Q_1),\ldots,
\Delta_{\mathrm{tr}}(Q_m)
\bigr)
\]
and the robust panel minimum
\[
\Delta_{\mathrm{tr}}^{\min}(\mathcal Q)
=
\min_{Q\in\mathcal Q}\Delta_{\mathrm{tr}}(Q).
\]
The panel must be fixed independently of the observed favorable gap. Otherwise
null selection can manufacture apparent relational dependence.

\topic{27}{Relational Witness Pairs}

\begin{definition}[Relational witness pair]\label{def:relational-witness}
A pair of laws \((P,Q)\) is a predictive relational witness pair for the
decomposition \(V\) and target \(Y_t^+\) when
\[
Q\in\mathcal N_{\mathrm{pred}}(P)
\]
but
\[
\Risk_{\ell,Q}(V)
\neq
\Risk_{\ell,P}(V).
\]
The two systems agree on every single-node predictive marginal but differ in
full-system operative capacity.
\end{definition}

\begin{example}[XOR witness pair]\label{ex:xor-witness}
Under the intact law \(P\), let
\[
X_1,X_2\overset{\mathrm{iid}}{\sim}\operatorname{Bernoulli}(1/2),
\qquad
Y=X_1\oplus X_2.
\]
Then each \((X_i,Y)\) pair is independent and uniform, so
\[
H_P(Y\mid X_i)=1\text{ bit},
\qquad
H_P(Y\mid X_1,X_2)=0.
\]
Let \(Q\) make \(X_1,X_2,Y\) mutually independent fair bits. Then
\[
Q_{(X_i,Y)}=P_{(X_i,Y)}
\qquad(i=1,2),
\]
but
\[
H_Q(Y\mid X_1,X_2)=1\text{ bit}.
\]
Therefore
\[
\Delta_{\mathrm{joint}}^{\log}=1\text{ bit},
\qquad
k_0=2,
\qquad
\Delta_{\mathrm{tr}}(Q)=1\text{ bit}.
\]
This proves predictive relational structure relative to the declared variables.
It does not, by itself, prove that any workspace variable is causally active.
\end{example}

\begin{example}[Redundant-copy control]
Let \(Y\) be a fair bit and set \(X_1=X_2=Y\). The two nodes are perfectly
coherent and the full system predicts perfectly, but
\[
H(Y\mid X_i)=0,
\qquad
\Delta_{\mathrm{joint}}^{\log}=0,
\qquad
k_0=1.
\]
Duplication is therefore not relational irreducibility.
\end{example}

The witness-pair criterion is stronger than a performance comparison against a
smaller model because it holds local predictive behavior fixed. It remains
observational. Part~VIII adds interventions and converts predictive
relationality into an operational causal claim.


\theorypart{VIII}{Causal Efficacy and Reentry}

\topic{28}{Feasible Workspace Interventions}

A notation such as \(\operatorname{do}(W_t=w)\) is meaningful only if the
candidate workspace is an endogenous variable in a structural causal model or
if an intervention on underlying variables is shown to realize that value.
Abstractly setting a projection coordinate need not correspond to any possible
physical or computational manipulation.

\begin{definition}[Time-unrolled causal realization]\label{def:causal-realization}
A time-unrolled causal realization is a structural causal model
\[
\mathcal M
=
(\mathcal U,\mathcal V,F,\mathbb P_U)
\]
whose endogenous variables include the local states, messages, actions,
artifacts, target variables, and a measurable workspace extractor
\[
\pi_t:\mathsf S_t\to\mathsf W_t,
\qquad
W_t=\pi_t(S_t).
\]
The unrolled graph is acyclic even when the rolled process is recurrent.
\end{definition}

\begin{definition}[Feasible implementation intervention]
Let \(B_t\) denote declared background variables that should remain matched.
A family \(\{\iota_w:w\in\mathsf W_t^\circ\}\) is a feasible implementation of
workspace values on the experimental domain \(\mathsf W_t^\circ\) when:
\begin{enumerate}
\item each \(\iota_w\) acts only on declared manipulable parents, channels, or
 artifacts at or after time \(t\);
\item the intervened realization satisfies
\[
\pi_t(S_t^{\iota_w})=w
\qquad\text{almost surely};
\]
\item the background law is matched across compared interventions; and
\item the intervention does not directly modify \(Y_t^+\) except through
 paths included in the declared causal design.
\end{enumerate}
\end{definition}

Let \(D\) be a divergence or metric on probability laws. For a matched
background value \(b\), define
\[
\CE_{\Workspace}(w,w';b)
=
D\!\left(
\mathbb P^{\iota_w}(Y_t^+\in\cdot\mid B_t=b),
\mathbb P^{\iota_{w'}}(Y_t^+\in\cdot\mid B_t=b)
\right).
\]
For a declared comparison design \(\nu\) over \((w,w',b)\), define
\[
\boxed{
\eta_{\mathrm{causal}}^{\nu}
=
\mathbb E_{(w,w',b)\sim\nu}
\left[
\CE_{\Workspace}(w,w';b)
\right].
}
\]
The maximum
\[
\eta_{\mathrm{causal}}^{\max}
=
\sup_{w,w',b}\CE_{\Workspace}(w,w';b)
\]
may be useful for discovery, but a design-weighted effect is harder to inflate
by selecting a single exceptional pair after inspection.

\begin{proposition}[Coordinate invariance of intervention effects]
\label{prop:causal-gauge}
Let \(g:\mathsf W_t\to\mathsf W_t'\) be a measurable bijection and suppose the
reparametrized intervention family satisfies
\[
\iota_{g(w)}'=\iota_w.
\]
Then every divergence-based causal effect is preserved:
\[
\CE_{\Workspace'}(g(w),g(w');b)
=
\CE_{\Workspace}(w,w';b).
\]
\end{proposition}

\begin{proof}
The two labels designate the same underlying interventions and therefore induce
the same pair of target laws. Only the coordinate name has changed.
\end{proof}

This is the causal analogue of gauge invariance. It also shows why effects
should be attached to intervention implementations rather than to arbitrary
numerical coordinates.

\topic{29}{Causal Specificity and Transfer Mediation}

A causal effect on the target is not workspace-specific if an equal-cost
perturbation outside the candidate workspace produces the same effect. Let
\(c(\iota)\) be a declared intervention cost and let
\[
\mathcal I_{\Workspace}(r),
\qquad
\mathcal I_{\perp}(r)
\]
be matched families with cost \(r\), the first changing workspace coordinates
and the second changing complementary coordinates while preserving the
workspace value as closely as the implementation permits.

Define
\[
\eta_{\Workspace}(r)
=
\mathbb E_{\iota,\iota'\sim\mathcal I_{\Workspace}(r)}
D\!\left(
\mathbb P^{\iota}(Y_t^+),
\mathbb P^{\iota'}(Y_t^+)
\right),
\]
\[
\eta_{\perp}(r)
=
\mathbb E_{\jmath,\jmath'\sim\mathcal I_{\perp}(r)}
D\!\left(
\mathbb P^{\jmath}(Y_t^+),
\mathbb P^{\jmath'}(Y_t^+)
\right),
\]
and
\[
\boxed{
\Delta_{\mathrm{specific}}(r)
=
\eta_{\Workspace}(r)-\eta_{\perp}(r).
}
\]
A positive value supports causal concentration in the candidate workspace at
the declared perturbation scale.

If the complement is defined by \(\Gamma_\varepsilon(H)^\perp\), the result is
metric dependent. Under nonunitary gauge changes, either the metric must be
transported with the coordinates or interventions must be matched by a
coordinate-independent physical cost. Otherwise specificity can change merely
because the representation was rescaled.

\subtopic{Causal transfer nulls}

Let \(\jmath_E\) randomize selected communication edges while preserving their
marginal message laws, rates, and timing budget. The causal transfer gap is
\[
\boxed{
\Delta_{\mathrm{tr}}^{\operatorname{do}}(\jmath_E)
=
\Risk_{\ell,\mathbb P^{\jmath_E}}(V)
-
\Risk_{\ell,\mathbb P}(V).
}
\]
Unlike an observational null law, this quantity is anchored to an implemented
channel intervention. A positive effect is evidence that the intact alignment
of messages, rather than merely their marginal content, contributes to task
performance.

A stronger mediation test equalizes the workspace state after two distinct
channel interventions. Let \(\mu\) be a common imposed workspace distribution,
implemented by a feasible stochastic intervention in the sense of
Definition~\ref{def:causal-realization}. The notation below is shorthand for
that implementation. Define the residual mediation defect
\[
\delta_{\mathrm{med}}
=
\sup_{\jmath,\jmath'}
D\!\left(
\mathbb P^{\jmath,\operatorname{do}(W_t\sim\mu)}(Y_t^+),
\mathbb P^{\jmath',\operatorname{do}(W_t\sim\mu)}(Y_t^+)
\right).
\]
A small value means that once the workspace state is equalized, little of the
channel intervention's effect remains. This supports, but does not by itself
identify without the causal assumptions, the claim that \(W_t\) mediates the
relevant relational influence.

\topic{30}{Causal Recurrence, Loop-Mediated Reentry, and Relationality}

Temporal correlation is weaker than causal recurrence. Define the total
workspace recurrence effect at lag \(\tau\) by
\[
\boxed{
\eta_{\mathrm{rec}}(\tau)
=
\mathbb E_{(w,w',b)\sim\nu}
D\!\left(
\mathbb P^{\iota_w}(W_{t+\tau}\in\cdot\mid B_t=b),
\mathbb P^{\iota_{w'}}(W_{t+\tau}\in\cdot\mid B_t=b)
\right).
}
\]
A nonzero value shows that manipulating the present workspace changes a later
workspace state. It may, however, be explained by direct copying or a retained
memory register.

To isolate reentry, consider the unrolled causal graph and remove every direct
retention edge whose source and target are both workspace variables at
successive times. Let
\[
\mathcal M^{\setminus\mathrm{dir}}
\]
denote the resulting path-specific model, preserving paths that leave the
workspace through a nonworkspace participant, action, message, or artifact and
later return. Define
\[
\boxed{
\eta_{\mathrm{loop}}(\tau)
=
\mathbb E_{(w,w',b)\sim\nu}
D\!\left(
\mathbb P_{\mathcal M^{\setminus\mathrm{dir}}}^{\iota_w}
(W_{t+\tau}\in\cdot\mid B_t=b),
\mathbb P_{\mathcal M^{\setminus\mathrm{dir}}}^{\iota_{w'}}
(W_{t+\tau}\in\cdot\mid B_t=b)
\right).
}
\]
Then the distinctions are:
\[
\begin{aligned}
I(W_t;W_{t+\tau})>0
&\quad\text{temporal dependence only},\\
\eta_{\mathrm{rec}}(\tau)>0
&\quad\text{causal persistence},\\
\eta_{\mathrm{loop}}(\tau)>0
&\quad\text{causal exit-and-return reentry}.
\end{aligned}
\]
Path-specific effects require the usual consistency and no-unmeasured-confounding
assumptions appropriate to the implementation. In software systems, edge
interventions can often be literal. In human collaboration, these assumptions
are substantially harder to justify.

\begin{theorem}[Operational relationality, \(T_{13}\)]
\label{thm:operational-relationality}
Let \(\mathsf Y\) be finite, let \(P\) be the intact law at a fixed decision
time, and let \(W_t=\pi_t(\mathbf H_t)\) be a structurally supported candidate.
Let
\(\mathcal Q\subseteq\mathcal N_{\mathrm{pred}}(P)\) be a nonempty,
predeclared panel of validated relation-breaking nulls. Assume:
\begin{enumerate}
\item exact structural--operational alignment,
\[
I_P(Y_t^+;\mathbf H_t\mid W_t)=0;
\]
\item a positive node-level information gap,
\[
\delta_{\mathrm{node}}
=
\min_{i\in V}
I_P(Y_t^+;\mathbf H_t\mid\mathbf H_t^i)
>0;
\]
\item a positive robust matched-null gap,
\[
\delta_{\mathrm{null}}
=
\min_{Q\in\mathcal Q}
\left[
H_Q(Y_t^+\mid\mathbf H_t)
-
H_P(Y_t^+\mid\mathbf H_t)
\right]
>0;
\]
\item a feasible intervention design with
\[
\eta_{\mathrm{causal}}^{\nu}>0;
\]
\item positive matched causal specificity at some declared scale \(r\),
\[
\Delta_{\mathrm{specific}}(r)>0.
\]
\end{enumerate}
Then, relative to the declared decomposition, target, log loss, null panel, and
intervention family:
\begin{enumerate}
\item no statistic of any single-node history is sufficient for \(Y_t^+\), and
 \(k_0\geq2\);
\item \(W_t\) cannot be reconstructed almost surely from any single-node
 history;
\item every null \(Q\in\mathcal Q\) preserves all single-node predictive
 risks but has strictly worse full-system log risk than the intact law;
\item the candidate workspace is causally efficacious and its effect exceeds
 that of the declared matched off-workspace perturbations.
\end{enumerate}
If, in addition, \(\eta_{\mathrm{loop}}(\tau)>0\), the workspace is
loop-reentrant at lag \(\tau\).
\end{theorem}

\begin{proof}
By Theorem~\ref{thm:coalition-log-risk}, for each node \(i\),
\[
\Risk_{\log,P}(\{i\})-
\Risk_{\log,P}(V)
=
I_P(Y_t^+;\mathbf H_t\mid\mathbf H_t^i)
\geq
\delta_{\mathrm{node}}>0.
\]
Hence no single-node history, and therefore no statistic of one, is sufficient;
this gives \(k_0\geq2\).

Exact alignment and Theorem~\ref{thm:alignment-criterion} imply that the
canonical predictive state factors through \(W_t\). If
\(W_t=\psi_i(\mathbf H_t^i)\) almost surely for some node, then the canonical
state would also be a function of \(\mathbf H_t^i\), contradicting the positive
node gap through Proposition~\ref{prop:sufficiency-reconstruction}.

For every \(Q\in\mathcal Q\), Proposition~\ref{prop:local-risk-null} gives
\[
\Risk_{\log,Q}(\{i\})
=
\Risk_{\log,P}(\{i\})
\qquad\forall i.
\]
The assumed null gap gives
\[
\Risk_{\log,Q}(V)
>
\Risk_{\log,P}(V).
\]
Thus a description determined only by the collection of single-node predictive
marginals cannot reproduce the intact joint capacity: it would assign the same
local data to \(P\) and \(Q\) while their full-system risks differ.

Finally, the feasible intervention effect establishes causal efficacy relative
to the declared implementation, and the positive specificity gap shows that
the effect exceeds the matched complementary perturbation effect. The final
reentry statement is exactly the definition of \(\eta_{\mathrm{loop}}\).
\end{proof}

\begin{remark}[Scope of the theorem]
The conclusion is operational and indexed by
\[
(V,Y_t^+,\ell,\mathcal Q,\mathcal I,\nu,r).
\]
Changing the decomposition, target, loss, nulls, or intervention design may
change the verdict. The theorem does not establish metaphysical
irreducibility, autonomous agency, selfhood, or consciousness.
\end{remark}

\subtopic{Finite certification sequence}

A reproducible dynamic-causal analysis should proceed in the following order:
\begin{enumerate}
\item Freeze the node decomposition, accessible coalition histories, target,
 horizon, and proper scoring loss.
\item Estimate the structural projector on training data only and encode
 \(W_t\) without using held-out targets.
\item Cross-fit conditional target laws from the full history, the workspace,
 and every declared coalition. Report alignment and coalition risk gaps.
\item Construct the null panel before inspecting null performance. Verify
 local-predictive matching and the intended destruction of transfer
 organization separately.
\item Implement feasible workspace and complementary interventions with a
 common background and cost design. Report both causal effect and
 specificity.
\item Test recurrence and then cut direct retention paths to test loop-mediated
 reentry.
\item Report confidence intervals or finite-sample uncertainty for every gate;
 do not substitute training accuracy for Bayes-risk or causal estimates.
\end{enumerate}

The resulting certification report is
\[
\boxed{
\begin{aligned}
\mathcal C_{\mathrm{dyn}}(W_t)
=\bigl(&
\delta_{\mathrm{align}}^{\mathcal I},
\{\Delta_A^\ell\}_{A\subseteq V},
 k_\varepsilon,
\boldsymbol\Delta_{\mathrm{tr}}(\mathcal Q),\\
&\eta_{\mathrm{causal}}^\nu,
\Delta_{\mathrm{specific}}(r),
\eta_{\mathrm{rec}}(\tau),
\eta_{\mathrm{loop}}(\tau)
\bigr).
\end{aligned}
}
\]
No single coordinate certifies a relational workspace.


\theorypart{IX}{Formal Workspace Definition}

\topic{31}{Certified Stratified Relational Workspace}

Let
\[
\Theta
=
\left(
\varepsilon_{\mathrm{coh}},
\varepsilon_{\mathrm{align}},
\varepsilon_k,
\underline\eta_c,
\underline\eta_{\mathrm{loop}},
\delta_{\mathrm{node}},
\delta_{\mathrm{null}},
\delta_{\mathrm{spec}},
k,
\varepsilon_{\mathrm{pers}}
\right).
\]
Let
\[
\varsigma
=
(\mathcal D,Y^+,\ell,\tau,\mathcal I,\mathcal Q,\Theta)
\]
be the complete declared scientific specification, where
\(\mathcal Q\subseteq\mathcal N_{\mathrm{pred}}(P)\) is a validated null panel.
A dynamic realization contains a certified stratified relational workspace
when there exists a candidate \(W_t\) satisfying all gates required for the
claimed level under \(\varsigma\). The detailed robustness and identifiability
requirements for varying \(\varsigma\) are developed in Parts~XVII and XVIII.

\subtopic{Structural coherence}
\[
\delta_{\mathrm{coh}}(W_t)
\leq
\varepsilon_{\mathrm{coh}}.
\]

\subtopic{Interventional structural--operational alignment}
\[
\delta_{\mathrm{align}}^{\mathcal I}(W_t)
\leq
\varepsilon_{\mathrm{align}}.
\]

\subtopic{Causal efficacy}
\[
\eta_{\mathrm{causal}}^\nu
\geq
\underline\eta_c.
\]

\subtopic{Loop-mediated reentry, when claimed}
\[
\eta_{\mathrm{loop}}(\tau)
\geq
\underline\eta_{\mathrm{loop}}.
\]

\subtopic{Node-level joint advantage}
\[
\Delta_{\mathrm{joint}}^{\ell}
\geq
\delta_{\mathrm{node}}.
\]

\subtopic{Matched-null transfer dependence}
\[
\Delta_{\mathrm{tr}}^{\min}(\mathcal Q)
\geq
\delta_{\mathrm{null}}.
\]

\subtopic{Causal specificity}
\[
\Delta_{\mathrm{specific}}(r)
\geq
\delta_{\mathrm{spec}}.
\]

\subtopic{Distributedness}
\[
k_{\varepsilon_k}
\geq
k.
\]

\subtopic{Persistence, when claimed}
\[
d_{\mathrm{id}}(W_t,W_{t+\tau})
\leq
\varepsilon_{\mathrm{pers}}.
\]

Symbolically,
\[
\boxed{
\SRW(W_t;\varsigma)
=
\bigwedge_{q\in\mathcal J(\varsigma)}
\operatorname{Gate}_q(W_t;\varsigma).
}
\]
The gate set is claim-relative. A one-step operational workspace need not pass
persistence or reentry gates, while a persistent reentrant workspace must.
No scalar total score is sufficient: failure of a required gate cannot be
hidden by high performance on another gate.

\begin{remark}[Structural, predictive, causal, and temporal claims]
The formal hierarchy now distinguishes four logically independent questions:
\[
\begin{aligned}
&\text{Does a stable shared coordinate exist?}\\
&\text{Does it retain the task-relevant conditional law?}\\
&\text{Does manipulating it change downstream behavior specifically?}\\
&\text{Does it persist or reenter through time?}
\end{aligned}
\]
The preceding structural and causal analysis answers the first three at the level of definitions and conditional
theorems. The following parts develop the fourth.
\end{remark}

\theorypart{X}{Persistence}

The preceding parts identify exact and approximate workspace subspaces at a
single time. Persistence adds three further requirements. The structural
subspace must remain trackable while the realization changes, actual states
must remain near that moving subspace, and the transported content must not
collapse to a trivial zero state. These are distinct conditions.

\topic{32}{Coherence-Error Dynamics}

\subtopic{32.1 Time-indexed transfer systems}

At time \(t\), let
\[
\mathcal C_V(t)=\bigoplus_{i\in V}C_i(t),
\qquad
\mathcal C_E(t)=\bigoplus_{e\in E}C_{t(e)}(t),
\]
and let
\[
D_t:\mathcal C_V(t)\longrightarrow\mathcal C_E(t),
\qquad
L_t=D_t^*D_t
\]
be the transfer discrepancy and transfer Laplacian. At time \(t\), the
combinatorial graph is held fixed while the local representations, transfers,
and metrics may vary. Changes of graph are treated later by a graph
correspondence.

To compare successive times, choose unitary identification maps
\[
U_t^V:\mathcal C_V(t)\longrightarrow\mathcal C_V(t+1),
\qquad
U_t^E:\mathcal C_E(t)\longrightarrow\mathcal C_E(t+1).
\]
They are comparison maps, not assumptions that the physical carriers are
identical. Define the one-step diagram defect
\[
\boxed{
\Delta_t^D
=
D_{t+1}U_t^V-U_t^ED_t,
\qquad
\delta_t^D=\|\Delta_t^D\|.
}
\]
Exact temporal compatibility is the condition \(\Delta_t^D=0\).

\begin{proposition}[Temporal Laplacian drift]\label{prop:temporal-laplacian-drift}
The carried Laplacian satisfies
\[
\boxed{
\left\|
(U_t^V)^*L_{t+1}U_t^V-L_t
\right\|
\leq
2\|D_t\|\delta_t^D+(\delta_t^D)^2.
}
\]
\end{proposition}

\begin{proof}
Set
\[
B_t=(U_t^E)^*\Delta_t^D.
\]
Because the comparison maps are unitary,
\[
(U_t^E)^*D_{t+1}U_t^V=D_t+B_t,
\qquad
\|B_t\|=\delta_t^D.
\]
Hence
\[
(U_t^V)^*L_{t+1}U_t^V
=(D_t+B_t)^*(D_t+B_t).
\]
Subtract \(D_t^*D_t\) and apply the triangle inequality.
\end{proof}

Let \(P_t\) project onto a selected rank-\(d\) low-energy workspace cluster
of \(L_t\), and set
\[
\mathcal W_t=\ran P_t.
\]
Write
\[
g_t=\lambda_{d+1}(L_t)-\lambda_d(L_t)>0
\]
for the separating spectral gap and
\[
\eta_t^L=2\|D_t\|\delta_t^D+(\delta_t^D)^2.
\]

\begin{corollary}[Temporal workspace tracking]\label{cor:temporal-workspace-tracking}
If
\[
\eta_t^L<\frac{g_t}{2},
\]
then the transported and newly extracted workspace projectors satisfy
\[
\boxed{
q_t^{\mathrm{str}}
:=
\left\|
P_{t+1}-U_t^VP_t(U_t^V)^*
\right\|
\leq
\frac{\eta_t^L}{g_t-\eta_t^L}
\leq
\frac{2\eta_t^L}{g_t}.
}
\]
\end{corollary}

\begin{proof}
Carry \(L_{t+1}\) back to time \(t\) using \(U_t^V\), apply
Theorem~\ref{thm:spectral-stability}, and conjugate the resulting projector
bound forward.
\end{proof}

The quantity \(q_t^{\mathrm{str}}\) measures structural motion after the
declared coordinate change has been removed. A workspace may therefore move
substantially in raw coordinates while remaining structurally unchanged.

\subtopic{32.2 State tracking inside a moving workspace}

Let the realized vertex state satisfy
\[
x_{t+1}=F_tx_t+\xi_t,
\]
where \(F_t:\mathcal C_V(t)\to\mathcal C_V(t+1)\) is the actual update and
\(\xi_t\) is an unmodeled disturbance. Define the off-workspace error
\[
\boxed{
e_t=\|(I-P_t)x_t\|.
}
\]
Introduce
\[
\alpha_t
=
\left\|
(I-P_{t+1})F_t(I-P_t)
\right\|,
\]
which measures persistence of prior off-workspace error, and
\[
\ell_t
=
\left\|
(I-P_{t+1})F_tP_t
\right\|,
\]
which measures leakage of a valid workspace state under the update.

\begin{proposition}[Grounded coherence recurrence]\label{prop:grounded-coherence-recurrence}
Every trajectory satisfies
\[
\boxed{
e_{t+1}\leq\alpha_te_t+\sigma_t,}
\]
where
\[
\boxed{
\sigma_t
=
\ell_t\|P_tx_t\|
+
\|(I-P_{t+1})\xi_t\|.
}
\]
\end{proposition}

\begin{proof}
Insert \(I=P_t+(I-P_t)\) before \(x_t\), apply \(I-P_{t+1}\), and use
submultiplicativity:
\[
\begin{aligned}
e_{t+1}
&\leq
\|(I-P_{t+1})F_t(I-P_t)x_t\|\\
&\quad+
\|(I-P_{t+1})F_tP_tx_t\|
+
\|(I-P_{t+1})\xi_t\|.
\end{aligned}
\]
The displayed recurrence follows from the definitions of \(\alpha_t\) and
\(\ell_t\).
\end{proof}

This derivation gives the contraction and disturbance terms an operational
meaning. The scalar recurrence is not an independent postulate. It follows
from an update rule, a moving workspace projector, and a declared disturbance
model.

\subtopic{32.3 Persistence theorem \(T_6\)}

\begin{theorem}[Nonautonomous persistence theorem]\label{thm:persistence}
Let \(e_t\geq0\) satisfy
\[
e_{t+1}\leq\alpha_te_t+\sigma_t,
\qquad
\alpha_t,\sigma_t\geq0.
\]
Then, for every \(t\geq1\),
\[
\boxed{
\begin{aligned}
e_t
\leq
&\left(\prod_{r=0}^{t-1}\alpha_r\right)e_0\\
&+
\sum_{j=0}^{t-1}
\left(\prod_{r=j+1}^{t-1}\alpha_r\right)\sigma_j,
\end{aligned}
}
\]
where an empty product equals one. If
\[
\alpha_t\leq\bar\alpha<1,
\qquad
\sigma_t\leq\bar\sigma,
\]
then
\[
\boxed{
e_t
\leq
\bar\alpha^te_0
+
\frac{\bar\sigma(1-\bar\alpha^t)}{1-\bar\alpha},
}
\]
and consequently
\[
\boxed{
\limsup_{t\to\infty}e_t
\leq
\frac{\bar\sigma}{1-\bar\alpha}.
}
\]
\end{theorem}

\begin{proof}
Repeated substitution yields
\[
\begin{aligned}
e_t
&\leq
\alpha_{t-1}\cdots\alpha_0e_0
+
\alpha_{t-1}\cdots\alpha_1\sigma_0
+
\cdots
+
\sigma_{t-1},
\end{aligned}
\]
which is the first formula. Under the uniform bounds, dominate each product by
the corresponding power of \(\bar\alpha\) and sum the geometric series.
\end{proof}

\begin{corollary}[Vanishing-disturbance convergence]\label{cor:vanishing-disturbance}
If \(\alpha_t\leq\bar\alpha<1\) and \(\sigma_t\to0\), then
\[
e_t\to0.
\]
More generally, if \(\sum_{t\geq0}\sigma_t<\infty\), then the disturbance
contribution in Theorem~\ref{thm:persistence} vanishes asymptotically.
\end{corollary}

\begin{proof}
Split the convolution sum into a finite early part, whose coefficients decay
to zero, and a late tail on which \(\sigma_t\) is uniformly small. The
summable case follows by the same argument.
\end{proof}

\subtopic{32.4 Mean-square persistence}

For stochastic systems, pathwise bounded disturbances may be too strong. Let
\(\{\mathcal F_t\}\) be a filtration containing the history through time
\(t\).

\begin{proposition}[Mean-square persistence]\label{prop:mean-square-persistence}
Suppose
\[
\mathbb E[e_{t+1}^2\mid\mathcal F_t]
\leq
q e_t^2+v_t,
\qquad
0\leq q<1,
\]
and \(\mathbb E[v_t]\leq\bar v\). Then
\[
\boxed{
\mathbb E[e_t^2]
\leq
q^t\mathbb E[e_0^2]
+
\frac{\bar v(1-q^t)}{1-q},
}
\]
so
\[
\limsup_{t\to\infty}\mathbb E[e_t^2]
\leq
\frac{\bar v}{1-q}.
\]
\end{proposition}

\begin{proof}
Take expectations, iterate the resulting deterministic recurrence, and sum a
geometric series.
\end{proof}

\subtopic{32.5 The non-collapse condition}

Low off-workspace error is not sufficient for meaningful persistence. The
zero update \(F_t=0\) drives every state exactly into every linear workspace
while destroying all content. Define the internal transport
\[
B_t=P_{t+1}F_tP_t:\mathcal W_t\longrightarrow\mathcal W_{t+1}
\]
and its non-collapse margin
\[
\boxed{
\mu_t^{\mathrm{keep}}
=
\sigma_{\min}(B_t|_{\mathcal W_t}).
}
\]
A content-preserving persistence claim requires
\[
\mu_t^{\mathrm{keep}}>0
\]
for the relevant interval, or an information-theoretic replacement when the
workspace state is stochastic. Structural tracking, state tracking, and
content retention are therefore separate gates.

A finite-horizon persistence certificate may be recorded as
\[
\boxed{
\mathcal P_{0:T}
=
\left(
\max_{t<T}q_t^{\mathrm{str}},
\max_{t\leq T}e_t,
\min_{t<T}\mu_t^{\mathrm{keep}},
\sum_{t<T}\sigma_t
\right).
}
\]
No single coordinate of this vector subsumes the others.

\topic{33}{Interpretation of the Contraction Parameter}

The inequality
\[
\alpha_t<1
\]
means that deviations already outside the current workspace are damped by the
next update. Mechanisms that can reduce \(\alpha_t\) or \(\sigma_t\) include
\[
\begin{aligned}
&\text{verification and error feedback},\\
&\text{shared notation and stable interfaces},\\
&\text{artifact persistence and version control},\\
&\text{proof checking and experimental replication},\\
&\text{handoff reconstruction and redundancy}.
\end{aligned}
\]
Sources of disturbance include
\[
\begin{aligned}
&\text{context loss and representational drift},\\
&\text{miscommunication and transfer-map error},\\
&\text{participant or model replacement},\\
&\text{artifact corruption},\\
&\text{unsupported inference and objective drift}.
\end{aligned}
\]

\subtopic{33.1 Four persistence levels}

The theory distinguishes:
\[
\begin{array}{ll}
\text{structural persistence}
& q_t^{\mathrm{str}}\text{ remains small},\\
\text{state persistence}
& e_t\text{ remains small},\\
\text{content persistence}
& \mu_t^{\mathrm{keep}}\text{ remains bounded away from zero},\\
\text{operational persistence}
& \text{predictive and causal capacities remain matched}.
\end{array}
\]
The first three are developed in this part. Operational matching is incorporated
into the identity theory of Part~XII.

\subtopic{33.2 Moving-line example}

Let \(\mathcal C_V(t)=\mathbb R^2\) and let \(P_t\) project onto
\[
u_t=(\cos\theta_t,\sin\theta_t).
\]
If \(U_t^V\) is rotation through
\(\theta_{t+1}-\theta_t\) and \(F_t=U_t^V\), then a state
\(x_t=u_t\) obeys
\[
e_t=e_{t+1}=0,
\qquad
\mu_t^{\mathrm{keep}}=1.
\]
The workspace changes raw coordinates but not abstract identity. If instead
\(F_t=I\), then
\[
e_{t+1}
=
|\sin(\theta_{t+1}-\theta_t)|.
\]
The same structural motion now produces a state-tracking defect because the
actual update fails to transport the content with the workspace.

\theorypart{XI}{Reconstruction and Resilience}

Persistence can be maintained by direct transport, but it can also be restored
from distributed traces after a disruption. Reconstruction theory quantifies
when a coalition contains enough information to recover the abstract core and
how noise is amplified by that recovery.

\topic{34}{Core Encoding Maps}

Let \(C_*\) be a finite-dimensional abstract core. The exact static theory
provides a complete local chart
\[
\overline E_i
=
\ev_i\circ\Psi:C_*\longrightarrow C_i,
\]
where \(\Psi:C_*\to\Gamma(H)\) is the root-to-section isomorphism. A participant
or artifact need not expose that complete chart. Let
\[
O_i:C_i\longrightarrow\mathsf Y_i
\]
be the declared observation or access map and define the \emph{accessible
encoding}
\[
\boxed{
E_i=O_i\overline E_i:C_*\longrightarrow\mathsf Y_i.
}
\]
For notational economy, the direct sums below use the accessible spaces
\(\mathsf Y_i\). Thus injectivity of a full structural chart does not make a
singleton operationally reconstructive unless its access map preserves the
whole core.

For a coalition \(S\subseteq V\), define the analysis operator
\[
A_S:C_*\longrightarrow\bigoplus_{i\in S}\mathsf Y_i,
\qquad
A_Sc=(E_ic)_{i\in S},
\]
and the coalition frame operator
\[
\boxed{
F_S=A_S^*A_S=\sum_{i\in S}E_i^*E_i.
}
\]
Set
\[
\alpha_S=\lambda_{\min}(F_S),
\qquad
\beta_S=\lambda_{\max}(F_S).
\]
Thus
\[
\alpha_S\|c\|^2
\leq
\|A_Sc\|^2
\leq
\beta_S\|c\|^2.
\]

\begin{proposition}[Coalition reconstruction criterion]\label{prop:coalition-frame-criterion}
For a coalition \(S\), the following are equivalent:
\begin{enumerate}
\item \(A_S\) is injective;
\item \(F_S\) is positive definite;
\item \(\alpha_S>0\);
\item there exists a linear decoder \(R_S\) with \(R_SA_S=I_{C_*}\).
\end{enumerate}
When these conditions hold, the canonical decoder is
\[
\boxed{
R_S=A_S^\dagger=F_S^{-1}A_S^*.
}
\]
\end{proposition}

\begin{proof}
In finite dimensions,
\[
\ker F_S=\ker A_S
\]
because
\[
\langle c,F_Sc\rangle=\|A_Sc\|^2.
\]
This proves the equivalence of the first three statements. An injective linear
map admits a left inverse, while a left inverse implies injectivity. The
formula for the Moore--Penrose left inverse follows from full column rank.
\end{proof}

The lower frame bound \(\alpha_S\) measures recoverability. The upper bound
\(\beta_S\) measures amplification in the forward encoding. Their ratio
\[
\kappa_S=\sqrt{\frac{\beta_S}{\alpha_S}}
\]
is the coalition condition number.

\topic{35}{Stable Reconstruction Condition}

\subtopic{35.1 Stable reconstruction theorem \(T_7\)}

\begin{theorem}[Stable core reconstruction]\label{thm:stable-reconstruction}
Assume \(\alpha_S>0\) and observe
\[
y=A_Sc+\eta.
\]
With the canonical estimate
\[
\widehat c=A_S^\dagger y,
\]
one has
\[
\boxed{
\|\widehat c-c\|
\leq
\frac{\|\eta\|}{\sqrt{\alpha_S}}.
}
\]
\end{theorem}

\begin{proof}
Since \(A_S^\dagger A_S=I_{C_*}\),
\[
\widehat c-c=A_S^\dagger\eta.
\]
The singular values of \(A_S\) are the square roots of the eigenvalues of
\(F_S\), so
\[
\|A_S^\dagger\|
=
\frac{1}{\sigma_{\min}(A_S)}
=
\frac{1}{\sqrt{\alpha_S}}.
\]
\end{proof}

\begin{corollary}[Encoding mismatch]\label{cor:encoding-mismatch}
If the true observation is
\[
y=(A_S+\Delta_S)c+\eta
\]
but the nominal decoder \(A_S^\dagger\) is used, then
\[
\boxed{
\|\widehat c-c\|
\leq
\frac{\|\Delta_S\|\,\|c\|+\|\eta\|}{\sqrt{\alpha_S}}.
}
\]
\end{corollary}

\begin{proof}
Now
\[
\widehat c-c=A_S^\dagger\Delta_Sc+A_S^\dagger\eta,
\]
and the result follows from Theorem~\ref{thm:stable-reconstruction}.
\end{proof}

\subtopic{35.2 Reconstruction as operational recovery}

Geometric recovery becomes behaviorally relevant when the workspace state
controls a predictive kernel. Let
\[
K_w=\mathbb P(Y\in\cdot\mid W=w)
\]
and suppose
\[
\TV(K_w,K_{w'})\leq L_Kd_W(w,w').
\]

\begin{proposition}[Exact reconstruction implies coalition sufficiency]
\label{prop:exact-reconstruction-sufficiency}
Let \(W\) be sufficient for \(Y\) relative to the full history \(\mathbf H\):
\[
Y\indep\mathbf H\mid W.
\]
If a coalition reconstructs the workspace exactly,
\[
W=R_S(\mathbf H^S)
\qquad\text{almost surely},
\]
then
\[
\boxed{
Y\indep\mathbf H\mid\mathbf H^S.
}
\]
Under logarithmic loss,
\[
\boxed{
\Risk_{\log}(S)=\Risk_{\log}(V).
}
\]
\end{proposition}

\begin{proof}
Sufficiency means that the conditional target law given \(\mathbf H\) is a
measurable function of \(W\). Exact reconstruction makes that function
measurable with respect to \(\mathbf H^S\), which is the asserted conditional
independence. The log-risk equality follows from
Theorem~\ref{thm:coalition-log-risk}.
\end{proof}

\begin{proposition}[Reconstruction-to-risk transfer]\label{prop:reconstruction-risk-transfer}
Let \(W\) be sufficient for a target \(Y\), let losses take values in
\([0,1]\), and let a coalition construct \(\widehat W\) from its local data.
Then its optimal Bayes risk satisfies
\[
\boxed{
0\leq
\Risk(S)-\Risk(W)
\leq
2L_K\,\mathbb E[d_W(W,\widehat W)].
}
\]
\end{proposition}

\begin{proof}
For every pair \(w,\widehat w\), using the Bayes action for
\(K_{\widehat w}\) under the true law \(K_w\) incurs conditional regret at
most
\[
2\TV(K_w,K_{\widehat w}).
\]
This follows by inserting and subtracting the risk under
\(K_{\widehat w}\) for both the chosen action and the Bayes action for
\(K_w\). Average the bound and use the Lipschitz assumption. The coalition's
optimal rule performs at least as well as this decoder-based rule.
\end{proof}

This proposition links Part~XI to the causal-state theory of Part~VI. A small
geometric reconstruction error is operationally meaningful only when the
outcome kernel changes continuously with workspace state.

\topic{36}{k-Failure Resilience}

For \(k\geq0\), define the worst surviving lower frame bound
\[
\boxed{
\alpha^{(k)}
=
\min_{\substack{A\subseteq V\\|A|\leq k}}
\alpha_{V\setminus A}.
}
\]
The sequence \(\alpha^{(k)}\) is nonincreasing in \(k\).

\begin{theorem}[Uniform erasure resilience]\label{thm:uniform-erasure-resilience}
The encoding family is exactly reconstructible after every loss of at most
\(k\) nodes if and only if
\[
\boxed{\alpha^{(k)}>0.}
\]
If surviving observations have the form
\[
y_{V\setminus A}=A_{V\setminus A}c+\eta_A,
\qquad
|A|\leq k,
\]
then the canonical survivor decoder obeys the uniform estimate
\[
\boxed{
\|\widehat c_A-c\|
\leq
\frac{\|\eta_A\|}{\sqrt{\alpha^{(k)}}}.
}
\]
\end{theorem}

\begin{proof}
Exact reconstruction for a particular survivor set is equivalent to a
positive lower frame bound by Proposition~\ref{prop:coalition-frame-criterion}.
Taking the minimum over all permitted erasures gives the first claim. The
uniform noise estimate follows from Theorem~\ref{thm:stable-reconstruction}.
\end{proof}

\begin{proposition}[Dimensional obstruction]\label{prop:dimension-obstruction}
Let \(d=\dim C_*\). If
\[
\sum_{i\in S}\dim \mathsf Y_i<d,
\]
then \(\alpha_S=0\), so \(S\) cannot reconstruct the core. Consequently,
\(k\)-failure resilience requires
\[
\sum_{i\in V\setminus A}\dim \mathsf Y_i\geq d
\]
for every \(A\subseteq V\) with \(|A|\leq k\).
\end{proposition}

\begin{proof}
The rank of \(A_S\) is at most the dimension of its codomain. If that
dimension is below \(d\), then \(A_S\) is not injective.
\end{proof}

Define the accessible block support of \(c\neq0\) by
\[
\operatorname{supp}_E(c)
=
\{i\in V:E_ic\neq0\}
\]
and define the node-erasure distance
\[
\boxed{
d_{\mathrm{er}}
=
\min_{c\neq0}|\operatorname{supp}_E(c)|.
}
\]

\begin{theorem}[Exact erasure-distance criterion]
\label{thm:erasure-distance}
Exact reconstruction after every loss of at most \(k\) nodes is possible if
and only if
\[
\boxed{
d_{\mathrm{er}}>k.
}
\]
\end{theorem}

\begin{proof}
If \(d_{\mathrm{er}}\leq k\), choose \(c\neq0\) whose nonzero accessible
blocks are contained in a set \(A\) with \(|A|\leq k\). The surviving analysis
operator satisfies \(A_{V\setminus A}c=0\), so it cannot distinguish \(c\)
from zero. Conversely, if some erasure set \(A\) of size at most \(k\) destroys
injectivity, there exists \(c\neq0\) with \(E_ic=0\) for every \(i\notin A\).
Thus \(\operatorname{supp}_E(c)\subseteq A\), and
\(d_{\mathrm{er}}\leq k\).
\end{proof}

The distance criterion is combinatorial. The frame quantity
\(\alpha^{(k)}\) is quantitative. A system may tolerate \(k\) erasures exactly
while amplifying noise severely because \(\alpha^{(k)}\) is close to zero.

\subtopic{36.1 Reconstructive handoff after replacement}

Suppose node \(r\) is removed and the surviving coalition is
\[
S=V\setminus\{r\}.
\]
A replacement node \(r'\) is to receive the encoding
\[
E_{r'}':C_*\to\mathsf Y_{r'}'.
\]

\begin{corollary}[Handoff initialization]\label{cor:handoff-initialization}
If \(\alpha_S>0\) and the survivors observe \(y=A_Sc+\eta\), initialize the
replacement by
\[
x_{r'}^{\mathrm{init}}
=E_{r'}'A_S^\dagger y.
\]
Then
\[
\boxed{
\|x_{r'}^{\mathrm{init}}-E_{r'}'c\|
\leq
\frac{\|E_{r'}'\|\,\|\eta\|}{\sqrt{\alpha_S}}.
}
\]
\end{corollary}

This is the basic mathematical form of artifact-mediated or
coalition-mediated handoff. It does not require the removed component to
transfer its full internal state to the replacement.

\topic{37}{Redundancy--Irreducibility Tension}

Define the minimum reconstructing coalition size
\[
\boxed{
m_{\mathrm{rec}}
=
\min\{|S|:\alpha_S>0\}
}
\]
and the maximum exact erasure tolerance
\[
\boxed{
k_{\mathrm{erase}}
=
\max\{k:\alpha^{(k)}>0\}.
}
\]
The condition
\[
m_{\mathrm{rec}}\geq2
\]
means that no individual encoding is sufficient for exact core recovery. The
condition
\[
k_{\mathrm{erase}}\geq1
\]
means that the system survives at least one node loss. These can hold
simultaneously because nonlocality of encoding is not the same as fragility.

\subtopic{37.1 Three-node distributed-frame example}

Let
\[
C_*=\mathbb R^2
\]
and define scalar node encodings
\[
E_1(c_1,c_2)=c_1,
\qquad
E_2(c_1,c_2)=c_2,
\qquad
E_3(c_1,c_2)=c_1+c_2.
\]
Every singleton has rank one, so no node reconstructs the two-dimensional
core. Every pair does reconstruct. For \(S=\{1,2\}\),
\[
F_S=I,
\qquad
\alpha_S=1.
\]
For \(S=\{1,3\}\) or \(S=\{2,3\}\), the frame eigenvalues are
\[
\frac{3\pm\sqrt5}{2}.
\]
Therefore
\[
\boxed{
m_{\mathrm{rec}}=2,
\qquad
k_{\mathrm{erase}}=1,
\qquad
d_{\mathrm{er}}=2,
\qquad
\alpha^{(1)}=\frac{3-\sqrt5}{2}.
}
\]
This realizes a workspace core that is irreducible to any singleton while
remaining reconstructible after any one node fails.

\subtopic{37.2 Relation to predictive irreducibility}

The geometric quantity \(m_{\mathrm{rec}}\) and the predictive quantity
\(k_0\) from Part~VII answer different questions. The first asks whether a
coalition can reconstruct a declared core coordinate. The second asks whether
a coalition is sufficient for a declared future target. Equality between them
requires structural--operational alignment and regularity such as
Proposition~\ref{prop:reconstruction-risk-transfer}. Neither quantity should
be substituted for the other without those assumptions.

\theorypart{XII}{Identity Through Time and Replacement}

Identity is treated as invariance under a structured correspondence, not as
literal persistence of tokens or carriers. Exact correspondences form a
groupoid. Approximate identity is measured by structural, operational,
causal, and reconstructive defects that must be reported separately before
they are combined.

\topic{38}{Temporal Realization Maps}

Let
\[
\mathfrak R=(G,\mathbf C,D,\rho,\mathcal I,\mathcal K)
\]
and
\[
\mathfrak R'=(G',\mathbf C',D',\rho',\mathcal I',\mathcal K')
\]
be two realizations. A correspondence begins with a graph isomorphism
\(\varphi:G\to G'\). After relabeling by \(\varphi\), let
\[
J_i:C_i\to C_i'
\]
be unitary local identifications and define
\[
J_V=\bigoplus_{i\in V}J_i,
\qquad
J_E=\bigoplus_{e\in E}J_{t(e)}.
\]

\begin{definition}[Exact structural correspondence]\label{def:exact-structural-correspondence}
The correspondence is structurally exact when
\[
\boxed{D'J_V=J_ED.}
\]
It is operation exact for \(a\in\mathcal A\) when
\[
\boxed{\rho_V'(a)J_V=J_V\rho_V(a).}
\]
\end{definition}

\begin{theorem}[Exact replacement isomorphism]\label{thm:exact-replacement-isomorphism}
Under an exact structural correspondence,
\[
\boxed{
\Gamma(D')=J_V\Gamma(D),
\qquad
L'J_V=J_VL.
}
\]
Consequently, every spectral workspace subspace is transported exactly:
\[
P_I(L')=J_VP_I(L)J_V^*
\]
for every spectral interval \(I\). If the correspondence is also operation
exact, then the transported workspace actions are intertwined exactly.
\end{theorem}

\begin{proof}
For every \(c\),
\[
D'J_Vc=J_EDc.
\]
Because \(J_E\) is invertible, \(Dc=0\) if and only if \(D'J_Vc=0\), proving
the kernel statement. Taking adjoints of the exact covariance identity gives
\[
D'^*J_E=J_VD^*.
\]
Therefore
\[
L'J_V=D'^*D'J_V=D'^*J_ED=J_VD^*D=J_VL.
\]
The spectral and operation claims follow by functional calculus and the
intertwining assumptions.
\end{proof}

Exact correspondences compose and invert. They therefore form a groupoid of
realizations. The abstract exact workspace is an isomorphism class in this
groupoid.

\topic{39}{Structural Identity Defect}

For an approximate correspondence define
\[
\boxed{
\Delta_D=D'J_V-J_ED,
\qquad
\delta_D=\|\Delta_D\|.
}
\]
If the graph weights are unchanged, define edgewise comparison defects
\[
\delta_e
=
\|H_e'J_{s(e)}-J_{t(e)}H_e\|.
\]
The discrepancy operator has one changed source block per edge, so
\[
\boxed{
\delta_D^2
\leq
\max_{i\in V}
\sum_{e:s(e)=i}w_e\delta_e^2
\leq
\sum_{e\in E}w_e\delta_e^2.
}
\]
When a single node is replaced, every nonincident edge has \(\delta_e=0\).
The replacement defect is therefore locally measurable even though its effect
on the global workspace is controlled by the global spectral gap.

The corresponding carried-Laplacian defect obeys
\[
\boxed{
\eta_L
:=
\|J_V^*L'J_V-L\|
\leq
2\|D\|\delta_D+\delta_D^2.
}
\]
This is Proposition~\ref{prop:temporal-laplacian-drift} with the time indices
removed.

Let \(P\) and \(P'\) be rank-\(d\) workspace projectors selected from \(L\)
and \(L'\), and let \(g>0\) be the separating source gap.

\begin{theorem}[Approximate structural replacement]\label{thm:approx-structural-replacement}
If
\[
\eta_L<\frac g2,
\]
then
\[
\boxed{
q_{\mathrm{str}}
:=
\|P'-J_VPJ_V^*\|
\leq
\frac{\eta_L}{g-\eta_L}
\leq
\frac{2\eta_L}{g}.
}
\]
If \(q_{\mathrm{str}}<1\), then
\[
B=P'J_VP:\ran P\longrightarrow\ran P'
\]
is bijective and satisfies
\[
\sigma_{\min}(B)
\geq
\sqrt{1-q_{\mathrm{str}}^2}.
\]
Its polar factor
\[
\boxed{
\beta
=B(B^*B)^{-1/2}
}
\]
is a canonical unitary workspace correspondence. Moreover,
\[
\boxed{
\|\beta-J_VP\|
\leq
q_{\mathrm{str}}
+
1-\sqrt{1-q_{\mathrm{str}}^2}.
}
\]
\end{theorem}

\begin{proof}
The projector estimate is Theorem~\ref{thm:spectral-stability} applied after
carrying \(L'\) back by \(J_V\). For unit \(w\in\ran P\),
\[
\|(I-P')J_Vw\|
\leq q_{\mathrm{str}}.
\]
Orthogonality gives
\[
\|Bw\|^2
=1-\|(I-P')J_Vw\|^2
\geq1-q_{\mathrm{str}}^2.
\]
Thus \(B\) is injective, hence bijective because the source and target have
the same finite dimension. Its singular values lie in
\([\sqrt{1-q_{\mathrm{str}}^2},1]\), so polar normalization moves \(B\) by at
most \(1-\sqrt{1-q_{\mathrm{str}}^2}\). Finally,
\[
\|B-J_VP\|=\|(I-P')J_VP\|\leq q_{\mathrm{str}}.
\]
\end{proof}

\subtopic{39.1 Transport of workspace operations}

Write
\[
A=\rho_V(a),
\qquad
A'=\rho_V'(a),
\]
and define
\[
\delta_\rho(a)=\|A'J_V-J_VA\|,
\qquad
\ell_a=\|(I-P)AP\|.
\]
Let
\[
A_{\mathcal W}=PAP,
\qquad
A_{\mathcal W'}'=P'A'P'.
\]

\begin{proposition}[Operation transport bound]\label{prop:operation-transport-bound}
For \(B=P'J_VP\),
\[
\boxed{
\|A_{\mathcal W'}'B-BA_{\mathcal W}\|
\leq
\|A'\|q_{\mathrm{str}}
+\delta_\rho(a)
+\ell_a.
}
\]
For the unitary polar correspondence \(\beta\),
\[
\boxed{
\begin{aligned}
\|A_{\mathcal W'}'\beta-\beta A_{\mathcal W}\|
\leq
&\|A'\|q_{\mathrm{str}}
+\delta_\rho(a)
+\ell_a\\
&+
\bigl(\|A_{\mathcal W'}'\|+\|A_{\mathcal W}\|\bigr)
\left(1-\sqrt{1-q_{\mathrm{str}}^2}\right).
\end{aligned}
}
\]
\end{proposition}

\begin{proof}
Insert and subtract \(P'A'J_VP\) and \(P'J_VAP\):
\[
\begin{aligned}
A_{\mathcal W'}'B-BA_{\mathcal W}
={}&P'A'(P'-I)J_VP\\
&+P'(A'J_V-J_VA)P\\
&+P'J_V(I-P)AP.
\end{aligned}
\]
Bound the three terms. Replacing \(B\) by \(\beta\) adds
\[
A_{\mathcal W'}'(\beta-B)-(\beta-B)A_{\mathcal W},
\]
which gives the second estimate.
\end{proof}

\topic{40}{Behavioral Identity Defect}

Structural correspondence does not determine operational identity. Let
\[
\{\iota_{w,b}:(w,b)\in\mathcal D\}
\quad\text{and}\quad
\{\iota'_{\beta w,\chi b}:(w,b)\in\mathcal D\}
\]
be matched feasible implementation-intervention families in the sense of
Part~VIII. Here \(b\) denotes declared background variables and \(\chi\) is
the background correspondence. Define
\[
K(w,b)
=
\mathbb P^{\iota_{w,b}}(Y^+\in\cdot),
\]
\[
K'(\beta w,\chi b)
=
\mathbb P^{\prime\,\iota'_{\beta w,\chi b}}(Y^{\prime+}\in\cdot).
\]
For a chosen divergence \(D_Y\), define
\[
\boxed{
\delta_{\mathrm{causal}}(\beta)
=
\sup_{(w,b)\in\mathcal D}
D_Y\left(
K(w,b),
K'(\beta w,\chi b)
\right).
}
\]
The observational response is included only when the null intervention belongs
to the declared family. This definition compares actual implementation
interventions rather than formally assigning an arbitrary projection
coordinate.

A reconstructive handoff may also be imperfect. Let
\[
\delta_{\mathrm{rec}}(\beta)
=
\mathbb E\bigl[d_{\mathcal W'}(W'_{\mathrm{init}},\beta W)\bigr]
\]
measure replacement initialization error. Corollary~\ref{cor:handoff-initialization}
provides one explicit upper bound on this term.

\subtopic{40.1 Structural identity does not imply causal identity}

Take two realizations with identical
\[
G,
\quad D,
\quad L,
\quad P,
\quad \rho.
\]
In the first, let \(Y=W\). In the second, let \(Y'=1-W\) for a binary
workspace coordinate. Then
\[
q_{\mathrm{str}}=0,
\qquad
\delta_\rho(a)=0
\]
for every operation, while the observational and interventional outcome
kernels differ maximally under the identity correspondence. Therefore no
purely structural theorem can prove causal identity without an additional
kernel-matching assumption.

\topic{41}{Identity Distance and Replacement Invariance}

For an admissible correspondence \(\mathfrak J\), define the operation defect
on a declared unit ball \(\mathcal A_1\) by
\[
\delta_{\mathrm{op}}(\mathfrak J)
=
\sup_{a\in\mathcal A_1}
\|A_{\mathcal W'}'\beta-\beta A_{\mathcal W}\|.
\]
The primary certification object is the identity profile
\[
\boxed{
\mathbf d_{\mathrm{id}}(\mathfrak J)
=
\left(
q_{\mathrm{str}},
\delta_{\mathrm{op}}(\mathfrak J),
\delta_{\mathrm{causal}}(\beta),
\delta_{\mathrm{rec}}(\beta);
\mu^{\mathrm{keep}}
\right).
}
\]
The first four coordinates must remain below declared tolerances, while the
non-collapse margin \(\mu^{\mathrm{keep}}\) must remain above a declared lower
bound. This vector, rather than a weighted average, determines whether a
specific identity claim passes.

A composite identity cost is
\[
\boxed{
\operatorname{Cost}(\mathfrak J)
=
a\,q_{\mathrm{str}}
+b\,\delta_{\mathrm{op}}(\mathfrak J)
+c\,\delta_{\mathrm{causal}}(\beta)
+d\,\delta_{\mathrm{rec}}(\beta),
}
\]
where \(a,b,c,d>0\) are declared scales. Define
\[
\boxed{
d_{\mathrm{id}}(\mathfrak R,\mathfrak R')
=
\inf_{\mathfrak J}
\operatorname{Cost}(\mathfrak J).
}
\]
This quantity is a task-relative optimization surrogate until the
admissible correspondence class, intervention matching, divergence, and scales
are fixed. It does not replace the coordinatewise identity profile or the
non-collapse gate.

\subtopic{41.1 Replacement invariance theorem \(T_{14}\) }

\begin{theorem}[Quantitative replacement invariance]\label{thm:replacement-invariance}
Let \(\mathfrak J\) be a unitary graph correspondence with diagram defect
\(\delta_D\). Define
\[
\eta_L=2\|D\|\delta_D+\delta_D^2
\]
and assume
\[
\eta_L<\frac g2.
\]
Set
\[
q_*
=
\frac{\eta_L}{g-\eta_L}.
\]
Then there exists a unitary workspace correspondence \(\beta\) for which
\[
q_{\mathrm{str}}\leq q_*.
\]
Moreover,
\[
\boxed{
\begin{aligned}
\delta_{\mathrm{op}}(\mathfrak J)
\leq
\sup_{a\in\mathcal A_1}
\Bigl[&
\|\rho_V'(a)\|q_*
+\delta_\rho(a)
+\ell_a\\
&+
\bigl(\|A_{\mathcal W'}'\|+\|A_{\mathcal W}\|\bigr)
\left(1-\sqrt{1-q_*^2}\right)
\Bigr].
\end{aligned}
}
\]
If, in addition,
\[
\delta_{\mathrm{causal}}(\beta)\leq\delta_c,
\qquad
\delta_{\mathrm{rec}}(\beta)\leq\delta_r,
\]
then
\[
\boxed{
\begin{aligned}
d_{\mathrm{id}}(\mathfrak R,\mathfrak R')
\leq
&a\,q_*+b\,\delta_{\mathrm{op}}(\mathfrak J)\\
&+c\,\delta_c+d\,\delta_r.
\end{aligned}
}
\]
In particular, if all local structural, operation, closure, causal, and
reconstruction defects are \(O(\delta_{\mathrm{rep}})\), the relevant
operator norms are uniformly bounded, and \(g\geq g_0>0\), then
\[
\boxed{
d_{\mathrm{id}}(\mathfrak R,\mathfrak R')
\leq
C_1\frac{\delta_{\mathrm{rep}}}{g_0}
+C_2\delta_{\mathrm{rep}}
+O\!\left(\frac{\delta_{\mathrm{rep}}^2}{g_0}\right).
}
\]
After a fixed normalization of the non-spectral terms, this recovers the
schematic target
\[
d_{\mathrm{id}}
\lesssim
C\frac{\delta_{\mathrm{rep}}}{g_0}.
\]
\end{theorem}

\begin{proof}
The structural bound and existence of \(\beta\) follow from
Theorem~\ref{thm:approx-structural-replacement}. The operation estimate is
Proposition~\ref{prop:operation-transport-bound} with
\(q_{\mathrm{str}}\leq q_*\). Insert these estimates and the assumed causal
and reconstructive bounds into the definition of
\(\operatorname{Cost}(\mathfrak J)\). The asymptotic statement follows from
\[
q_*=O\!\left(\frac{\delta_D}{g_0}
+\frac{\delta_D^2}{g_0}\right)
\]
and
\[
1-\sqrt{1-q_*^2}=O(q_*^2).
\]
\end{proof}

\begin{corollary}[Preservation of coalition frame bounds]
\label{cor:replacement-frame-bound}
Identify the old and new abstract cores with \(\ran P\) and \(\ran P'\)
through their chosen core parameterizations. Let \(S\) be a declared coalition,
let
\[
A_S:\ran P\to\mathsf Y_S,
\qquad
A_S':\ran P'\to\mathsf Y_S'
\]
be the resulting accessible coalition encoders, and let \(K_S\) be a unitary
identification of their observation spaces. Suppose
\[
\|A_S'\beta-K_SA_S\|
\leq
\delta_{S}^{\mathrm{frame}}
<
\sqrt{\alpha_S}.
\]
Then the replacement coalition remains reconstructive and its new lower frame
bound satisfies
\[
\boxed{
\alpha_S'
\geq
\left(
\sqrt{\alpha_S}-\delta_{S}^{\mathrm{frame}}
\right)^2.
}
\]
Consequently, \(k\)-failure resilience survives whenever the bound holds
uniformly for every coalition remaining after at most \(k\) erasures.
\end{corollary}

\begin{proof}
For \(w'\in\ran P'\), write \(w'=\beta w\). Since \(\beta\) and \(K_S\) are
unitary,
\[
\begin{aligned}
\|A_S'w'\|
&\geq
\|K_SA_Sw\|
-
\|(A_S'\beta-K_SA_S)w\|\\
&\geq
\left(
\sqrt{\alpha_S}-\delta_{S}^{\mathrm{frame}}
\right)\|w'\|.
\end{aligned}
\]
Square the inequality.
\end{proof}

\begin{corollary}[Exact carrier replacement]\label{cor:exact-carrier-replacement}
If
\[
\delta_D=0,
\qquad
\delta_\rho(a)=0,
\qquad
\ell_a=0,
\qquad
\delta_c=0,
\qquad
\delta_r=0,
\]
then
\[
d_{\mathrm{id}}(\mathfrak R,\mathfrak R')=0.
\]
Thus a participant, model, artifact, or notation system may be replaced without
changing the abstract workspace when the replacement realizes the same
structural, operational, causal, and reconstructive organization.
\end{corollary}

\subtopic{41.2 Temporal chains and identity paths}

For a sequence of realizations
\[
\mathfrak R_0,\mathfrak R_1,\ldots,\mathfrak R_T,
\]
let \(\mathfrak J_t:\mathfrak R_t\to\mathfrak R_{t+1}\) be admissible
correspondences. Define the path length
\[
\operatorname{Len}(\mathfrak J_{0:T})
=
\sum_{t=0}^{T-1}\operatorname{Cost}(\mathfrak J_t)
\]
and the induced directed path quasi-pseudometric
\[
\boxed{
d_{\mathrm{path}}^{\rightarrow}(\mathfrak R_0,\mathfrak R_T)
=
\inf_{\mathfrak J_{0:T}}
\operatorname{Len}(\mathfrak J_{0:T}).
}
\]
Concatenation gives the directed triangle inequality. Because replacement
costs and intervention divergences need not be symmetric, symmetry is not
assumed. When reverse correspondences exist, define the symmetrized
pseudometric
\[
\boxed{
d_{\mathrm{path}}^{\mathrm{sym}}(\mathfrak R,\mathfrak R')
=
\max\left\{
 d_{\mathrm{path}}^{\rightarrow}(\mathfrak R,\mathfrak R'),
 d_{\mathrm{path}}^{\rightarrow}(\mathfrak R',\mathfrak R)
\right\}.
}
\]
Distinct physical realizations can have zero symmetrized pseudodistance when
an exact correspondence exists in both directions. This is intentional:
identity is assigned to the invariant relational organization rather than to
the carrier tokens.

\subtopic{41.3 Identity failure modes}

A complete replacement audit must report at least
\[
\boxed{
\mathcal I_{\mathrm{rep}}
=
\left(
q_{\mathrm{str}},
\delta_{\mathrm{op}},
\delta_{\mathrm{causal}},
\delta_{\mathrm{rec}},
\mu^{\mathrm{keep}}
\right).
}
\]
The failure modes are different:
\[
\begin{array}{ll}
q_{\mathrm{str}}\text{ large}
&\text{the shared subspace changed},\\
\delta_{\mathrm{op}}\text{ large}
&\text{the same operations are no longer realized},\\
\delta_{\mathrm{causal}}\text{ large}
&\text{interventions no longer have matching consequences},\\
\delta_{\mathrm{rec}}\text{ large}
&\text{the new realization was not initialized from the old core},\\
\mu^{\mathrm{keep}}\text{ small}
&\text{content collapsed despite apparent coherence}.
\end{array}
\]
The theorem does not permit one successful coordinate to conceal failure of
another.


\theorypart{XIII}{Gluing and Workspaces of Workspaces}

The preceding parts treat one interaction diagram at a time. A large
collaborative system is usually assembled from overlapping local laboratories,
subteams, artifacts, models, or task-specific workspaces. This part asks when
those local workspaces determine one coherent global workspace, how failed
gluing is detected, and whether recursive assembly changes the resulting
object. The exact theory is a descent theorem. The approximate theory is a
second transfer-Laplacian construction on the cover itself.

\topic{42}{Local Workspace Covers}

\subtopic{42.1 Edge-complete covers}

Let
\[
G=(V,E)
\]
be a finite directed interaction graph with local core spaces \(C_i\), core
transfers \(H_e\), vertex space \(\mathcal C_V(G)\), discrepancy operator
\(D_G\), and exact workspace
\[
\mathscr W(G)
:=
\Gamma(H|_G)
=
\ker D_G.
\]

\begin{definition}[Finite workspace cover]\label{def:workspace-cover}
A finite family of subgraphs
\[
\mathcal U=\{G_\alpha\}_{\alpha\in A}
\]
is an \emph{edge-complete workspace cover} of \(G\) when
\[
V=\bigcup_{\alpha\in A}V_\alpha,
\qquad
E=\bigcup_{\alpha\in A}E_\alpha.
\]
Every edge in \(E_\alpha\) carries the transfer inherited from \(G\). For a
nonempty finite index set
\[
I=\{\alpha_0,\ldots,\alpha_p\}\subseteq A,
\]
define the intersection subgraph
\[
G_I
=
G_{\alpha_0\cdots\alpha_p}
:=
\bigcap_{\alpha\in I}G_\alpha
\]
by intersecting both vertex and edge sets. Isolated vertices are retained.
\end{definition}

Covering only the vertices is insufficient. If an edge belongs to no chart,
its transfer equation is absent from every local workspace and cannot be
recovered by overlap agreement. Edge completeness is therefore the minimal
cover assumption used below.

If \(K\subseteq L\subseteq G\) are subgraphs, coordinate restriction gives
\[
\rho_K^L:\mathcal C_V(L)\longrightarrow\mathcal C_V(K).
\]
Because every edge equation on \(K\) is also an edge equation on \(L\),
restriction maps exact sections to exact sections:
\[
r_K^L
:=
\rho_K^L|_{\mathscr W(L)}:
\mathscr W(L)\longrightarrow\mathscr W(K).
\]
They obey
\[
r_L^L=I,
\qquad
r_J^K r_K^L=r_J^L
\]
whenever \(J\subseteq K\subseteq L\). Thus \(K\mapsto\mathscr W(K)\) is a
contravariant system of local workspaces.

\subtopic{42.2 The cover nerve and overlap interfaces}

The nerve \(N(\mathcal U)\) has one vertex for each chart \(G_\alpha\) and a
\(p\)-simplex
\[
[\alpha_0\cdots\alpha_p]
\]
whenever \(G_{\alpha_0\cdots\alpha_p}\) is nonempty. The nerve records which
local workspaces can compare states directly. It does not replace the
underlying interaction graph: \(G\) carries the internal transfer equations,
while \(N(\mathcal U)\) carries overlap-consistency equations between already
constructed local objects.

For pairwise overlaps write
\[
\mathscr W_\alpha=\mathscr W(G_\alpha),
\qquad
\mathscr W_{\alpha\beta}=\mathscr W(G_{\alpha\beta}),
\]
and
\[
r_{\alpha\beta}^{\alpha}:
\mathscr W_\alpha\to\mathscr W_{\alpha\beta},
\qquad
r_{\alpha\beta}^{\beta}:
\mathscr W_\beta\to\mathscr W_{\alpha\beta}.
\]
The interface is therefore a cospan
\[
\mathscr W_\alpha
\xrightarrow{\ r_{\alpha\beta}^{\alpha}\ }
\mathscr W_{\alpha\beta}
\xleftarrow{\ r_{\alpha\beta}^{\beta}\ }
\mathscr W_\beta,
\]
not necessarily a direct transfer map from one complete local workspace to the
other. Both charts are compared through the content they expose on their
common domain.

\subtopic{42.3 \v{C}ech cochains of local workspaces}

For \(p\geq0\), define the finite cochain space
\[
\check C^p(\mathcal U;\mathscr W)
=
\bigoplus_{\alpha_0<\cdots<\alpha_p}
\mathscr W(G_{\alpha_0\cdots\alpha_p}),
\]
where empty intersections contribute the zero space. The coboundary
\[
\delta^p:\check C^p\longrightarrow\check C^{p+1}
\]
is
\[
\boxed{
(\delta^p s)_{\alpha_0\cdots\alpha_{p+1}}
=
\sum_{j=0}^{p+1}(-1)^j
r_{\alpha_0\cdots\alpha_{p+1}}^%
{\alpha_0\cdots\widehat{\alpha_j}\cdots\alpha_{p+1}}
%
s_{\alpha_0\cdots\widehat{\alpha_j}\cdots\alpha_{p+1}}.
}
\]
The restriction identities imply
\[
\delta^{p+1}\delta^p=0.
\]
For a local family \(w=(w_\alpha)_\alpha\),
\[
\boxed{
(\delta^0w)_{\alpha\beta}
=
r_{\alpha\beta}^{\beta}w_\beta
-
r_{\alpha\beta}^{\alpha}w_\alpha.
}
\]
Thus \(\delta^0\) is the overlap-discrepancy operator at the level of exact
local workspaces.

\topic{43}{Compatibility and Gluing Obstructions}

\subtopic{43.1 Direct overlap mismatch}

A family \(w=(w_\alpha)_\alpha\) is exactly compatible when
\[
\delta^0w=0.
\]
Its direct mismatch is
\[
\boxed{
\mathfrak m_{\mathcal U}(w)=\delta^0w.
}
\]
After choosing inner products and positive overlap weights, the mismatch
energy is
\[
\boxed{
\mathcal E_{\mathrm{ov}}(w)
=
\|\delta^0w\|^2
=
\sum_{\alpha<\beta}
\nu_{\alpha\beta}
\left\|
 r_{\alpha\beta}^{\beta}w_\beta
-r_{\alpha\beta}^{\alpha}w_\alpha
\right\|^2.
}
\]
This is the appropriate obstruction when the full local states are already
given. Such a family glues exactly if and only if its mismatch vector itself
vanishes. No cohomological quotient is needed for that direct question.

\subtopic{43.2 Relative-only and affine descent data}

A different problem arises when only relative transition data are supplied.
Let
\[
b=(b_{\alpha\beta})\in\check C^1(\mathcal U;\mathscr W)
\]
be a prescribed family of overlap offsets. The triple-overlap consistency
condition is
\[
\delta^1b=0.
\]
One seeks local potentials \(q=(q_\alpha)\in\check C^0\) satisfying
\[
\delta^0q=b.
\]
Define
\[
\boxed{
\check H^1(\mathcal U;\mathscr W)
=
\frac{\ker\delta^1}{\operatorname{im}\delta^0}.
}
\]
Then the affine gluing problem is solvable precisely when
\[
\boxed{[b]=0\quad\text{in}\quad\check H^1(\mathcal U;\mathscr W).}
\]
A nonzero class means that all local transition statements can be mutually
consistent on triple overlaps while still failing to arise from any family of
local workspace states. This is a cover-level analogue of nontrivial holonomy.

\subtopic{43.3 Harmonic representative of the obstruction}

Equip the finite cochain spaces with Hilbert inner products and define the
first \v{C}ech Laplacian
\[
\boxed{
\Delta_1
=
\delta^0(\delta^0)^*
+
(\delta^1)^*\delta^1.
}
\]
Let
\[
\mathcal H^1_{\mathcal U}=\ker\Delta_1.
\]

\begin{proposition}[Harmonic gluing obstruction]\label{prop:harmonic-gluing-obstruction}
For every cocycle \(b\in\ker\delta^1\), there is an orthogonal decomposition
\[
\boxed{
b=\delta^0q^\dagger+h,
\qquad
h\in\mathcal H^1_{\mathcal U},
}
\]
where \(q^\dagger=(\delta^0)^\dagger b\) is a least-squares potential. Moreover,
\[
\boxed{
\inf_q\|b-\delta^0q\|
=
\|h\|
=
\|P_{\mathcal H^1_{\mathcal U}}b\|.
}
\]
Hence \(b\) is exactly integrable if and only if its harmonic component
vanishes.
\end{proposition}

\begin{proof}
The finite-dimensional Hodge decomposition gives
\[
\check C^1
=
\operatorname{im}\delta^0
\oplus
\ker\Delta_1
\oplus
\operatorname{im}(\delta^1)^*.
\]
If \(b\in\ker\delta^1\), then \(b\) is orthogonal to
\(\operatorname{im}(\delta^1)^*\). Its remaining two orthogonal components are
therefore an exact part \(\delta^0q^\dagger\) and a harmonic part \(h\). The
orthogonal projection theorem gives the least-squares identity.
\end{proof}

The scalar
\[
\omega_{\mathcal U}(b)
:=
\|P_{\mathcal H^1_{\mathcal U}}b\|
\]
is an irreducible overlap-frustration measure. Its vanishing is
invariant under admissible invertible coordinate changes, although its numerical
norm depends on the chosen metric. It cannot be removed by adding local
potentials alone.

\subtopic{43.4 Cycle example}

Consider three charts arranged in a cycle, with one-dimensional overlap
spaces and oriented offsets
\[
b_{12},\qquad b_{23},\qquad b_{31}.
\]
Local potentials would have to satisfy
\[
q_2-q_1=b_{12},
\qquad
q_3-q_2=b_{23},
\qquad
q_1-q_3=b_{31}.
\]
Adding the equations gives the necessary and sufficient condition
\[
\boxed{
b_{12}+b_{23}+b_{31}=0.}
\]
With the standard edge norm, the harmonic obstruction has magnitude
\[
\boxed{
\omega_{\mathcal U}(b)
=
\frac{|b_{12}+b_{23}+b_{31}|}{\sqrt3}.
}
\]
Every pairwise relation may be individually meaningful while the complete
cycle remains globally inconsistent.

\topic{44}{Exact Gluing Theorem T8}

Let
\[
R_{\mathcal U}:\mathscr W(G)
\longrightarrow
\check C^0(\mathcal U;\mathscr W)
\]
be the restriction map
\[
R_{\mathcal U}(w)
=
(w|_{G_\alpha})_{\alpha\in A}.
\]

\begin{theorem}[Exact workspace gluing]\label{thm:exact-workspace-gluing}
For every finite edge-complete workspace cover \(\mathcal U\), the sequence
\[
\boxed{
0
\longrightarrow
\mathscr W(G)
\xrightarrow{\ R_{\mathcal U}\ }
\check C^0(\mathcal U;\mathscr W)
\xrightarrow{\ \delta^0\ }
\check C^1(\mathcal U;\mathscr W)
}
\]
is exact. Equivalently,
\[
\boxed{
\mathscr W(G)
\cong
\ker\delta^0
\cong
\operatorname{Eq}
\left(
\prod_\alpha\mathscr W_\alpha
\rightrightarrows
\prod_{\alpha<\beta}\mathscr W_{\alpha\beta}
\right).
}
\]
Every pairwise compatible family of exact local workspace states therefore
glues to a unique exact global workspace state.
\end{theorem}

\begin{proof}
The cover contains every vertex, so \(R_{\mathcal U}\) is injective. A global
section has identical restrictions on every overlap, hence
\(\delta^0R_{\mathcal U}=0\).

Conversely, let \(w=(w_\alpha)_\alpha\in\ker\delta^0\). For each vertex
\(i\in V\), choose any chart \(G_\alpha\) containing \(i\) and set
\[
c_i=(w_\alpha)_i.
\]
If another chart \(G_\beta\) contains \(i\), then \(i\in G_{\alpha\beta}\)
and overlap compatibility gives
\((w_\alpha)_i=(w_\beta)_i\). Thus \(c_i\) is well defined.

For every edge \(e:i\to j\), edge completeness provides a chart
\(G_\alpha\) containing \(e\). Since \(w_\alpha\in\mathscr W(G_\alpha)\),
\[
c_j=(w_\alpha)_j=H_e(w_\alpha)_i=H_ec_i.
\]
Hence \(c=(c_i)\in\mathscr W(G)\), and by construction
\(R_{\mathcal U}c=w\). Uniqueness follows because the charts cover all
vertices.
\end{proof}

\begin{corollary}[Zeroth descent cohomology]\label{cor:zeroth-descent}
For every finite edge-complete cover,
\[
\boxed{
\mathscr W(G)
\cong
\check H^0(\mathcal U;\mathscr W)
=
\ker\delta^0.
}
\]
In finite dimensions,
\[
\boxed{
\dim\mathscr W(G)
=
\sum_{\alpha\in A}\dim\mathscr W_\alpha
-
\operatorname{rank}(\delta^0).
}
\]
\end{corollary}

The theorem does not require the nerve to be acyclic. Cycles in the cover are
allowed; their effect is encoded in the rank of \(\delta^0\) and, for
relative-only data, in \(\check H^1\). It also shows that the local-workspace
assignment is a genuine sheaf on the finite subgraph cover used here rather
than merely a presheaf.

\subtopic{44.1 Local nontriviality does not imply global nontriviality}

Let \(G\) be the directed triangle
\[
1\xrightarrow{\ a\ }2
\xrightarrow{\ b\ }3
\xrightarrow{\ c\ }1
\]
with one-dimensional vertex cores. Cover it by its three edge charts. Their
local workspaces are all one-dimensional:
\[
\begin{aligned}
\mathscr W_{12}&=\{(x,ax):x\in\mathbb R\},\\
\mathscr W_{23}&=\{(y,by):y\in\mathbb R\},\\
\mathscr W_{31}&=\{(z,cz):z\in\mathbb R\}.
\end{aligned}
\]
Overlap agreement requires
\[
y=ax,
\qquad
z=by,
\qquad
x=cz.
\]
Therefore
\[
(1-abc)x=0,
\]
and
\[
\boxed{
\dim\mathscr W(G)
=
\begin{cases}
1,&abc=1,\\
0,&abc\neq1.
\end{cases}
}
\]
The meta-level gluing obstruction is exactly the original cycle holonomy.
Strong local workspaces can therefore coexist with a trivial global core.

\topic{45}{Approximate Gluing and Meta-Workspaces}

\subtopic{45.1 The cover-level transfer Laplacian}

First assume each local state already lies in its exact local workspace. Define
\[
B_{\mathcal U}
:=
\delta^0:
\check C^0(\mathcal U;\mathscr W)
\longrightarrow
\check C^1(\mathcal U;\mathscr W)
\]
and the meta-Laplacian
\[
\boxed{
L_{\mathcal U}^{\mathrm{meta}}
=
B_{\mathcal U}^*B_{\mathcal U}.
}
\]
By Theorem~\ref{thm:exact-workspace-gluing},
\[
\boxed{
\ker L_{\mathcal U}^{\mathrm{meta}}
=
R_{\mathcal U}\mathscr W(G).
}
\]
Thus the gluing problem is itself a workspace problem whose ``nodes'' are
local workspaces and whose interfaces are overlaps.

Let
\[
\gamma_{\mathcal U}^{\mathrm{meta}}
=
\lambda_{\min}^+
\left(L_{\mathcal U}^{\mathrm{meta}}\right)
\]
be the first positive eigenvalue.

\begin{proposition}[Approximate overlap gluing]\label{prop:approx-overlap-gluing}
For every \(w\in\check C^0(\mathcal U;\mathscr W)\),
\[
\boxed{
\operatorname{dist}
\left(
 w,
 R_{\mathcal U}\mathscr W(G)
\right)^2
\leq
\frac{\|B_{\mathcal U}w\|^2}
{\gamma_{\mathcal U}^{\mathrm{meta}}}.
}
\]
\end{proposition}

\begin{proof}
Decompose \(w\) into its component in \(\ker L_{\mathcal U}^{\mathrm{meta}}\)
and its orthogonal complement. On the complement,
\(L_{\mathcal U}^{\mathrm{meta}}\geq
\gamma_{\mathcal U}^{\mathrm{meta}}I\). Therefore
\[
\|B_{\mathcal U}w\|^2
=
\langle w,L_{\mathcal U}^{\mathrm{meta}}w\rangle
\geq
\gamma_{\mathcal U}^{\mathrm{meta}}
\operatorname{dist}
\left(w,\ker L_{\mathcal U}^{\mathrm{meta}}\right)^2.
\]
\end{proof}

The gap \(\gamma_{\mathcal U}^{\mathrm{meta}}\) measures how strongly the cover
penalizes disagreement. A small overlap mismatch need not imply a nearby
global workspace when this gap is itself very small.

\subtopic{45.2 Simultaneous local and overlap error}

Local candidates may also violate their internal transfer equations. Define
the local-copy space
\[
\widehat{\mathcal C}_{\mathcal U}
=
\bigoplus_{\alpha\in A}\mathcal C_V(G_\alpha)
\]
and the block internal discrepancy
\[
\widehat D_{\mathcal U}
=
\bigoplus_{\alpha\in A}D_{G_\alpha}.
\]
For \(u=(u_\alpha)_\alpha\in\widehat{\mathcal C}_{\mathcal U}\), define the
raw overlap synchronization operator
\[
(S_{\mathcal U}u)_{\alpha\beta}
=
\rho_{G_{\alpha\beta}}^{G_\beta}u_\beta
-
\rho_{G_{\alpha\beta}}^{G_\alpha}u_\alpha.
\]
For an overlap weight \(\lambda_{\mathrm{ov}}>0\), set
\[
\boxed{
K_{\mathcal U}
=
\begin{pmatrix}
\widehat D_{\mathcal U}\\[0.2em]
\sqrt{\lambda_{\mathrm{ov}}}\,S_{\mathcal U}
\end{pmatrix},
\qquad
L_{\mathcal U}^{\mathrm{tot}}
=
K_{\mathcal U}^*K_{\mathcal U}.
}
\]
Its energy is
\[
\boxed{
\begin{aligned}
\mathcal E_{\mathcal U}^{\mathrm{tot}}(u)
={}&
\sum_\alpha\|D_{G_\alpha}u_\alpha\|^2\\
&+
\lambda_{\mathrm{ov}}
\sum_{\alpha<\beta}
\left\|
\rho^{G_\beta}_{G_{\alpha\beta}}u_\beta
-
\rho^{G_\alpha}_{G_{\alpha\beta}}u_\alpha
\right\|^2.
\end{aligned}
}
\]

Let
\[
J_{\mathcal U}:\mathcal C_V(G)\to\widehat{\mathcal C}_{\mathcal U},
\qquad
J_{\mathcal U}c=(c|_{G_\alpha})_\alpha,
\]
be the duplication map.

\begin{theorem}[Total descent kernel]\label{thm:total-descent-kernel}
For every finite edge-complete cover,
\[
\boxed{
\ker K_{\mathcal U}
=
J_{\mathcal U}\mathscr W(G).
}
\]
Consequently, if
\[
\gamma_{\mathcal U}^{\mathrm{tot}}
=
\lambda_{\min}^+
\left(L_{\mathcal U}^{\mathrm{tot}}\right),
\]
then
\[
\boxed{
\operatorname{dist}
\left(
 u,
 J_{\mathcal U}\mathscr W(G)
\right)^2
\leq
\frac{\mathcal E_{\mathcal U}^{\mathrm{tot}}(u)}
{\gamma_{\mathcal U}^{\mathrm{tot}}}.
}
\]
\end{theorem}

\begin{proof}
If \(u=J_{\mathcal U}c\) for a global exact section, every local discrepancy
and every overlap discrepancy vanishes. Conversely, if \(K_{\mathcal U}u=0\),
then every \(u_\alpha\) is an exact local section and the family agrees on all
overlaps. Theorem~\ref{thm:exact-workspace-gluing} gives a unique
\(c\in\mathscr W(G)\) with \(u=J_{\mathcal U}c\). The distance estimate follows
from the spectral lower bound on the orthogonal complement of the kernel.
\end{proof}

This theorem combines two distinct failure modes:
\[
\text{local incoherence}
\qquad\text{and}\qquad
\text{cross-chart disagreement}.
\]
Neither can be hidden by the other because they occupy separate blocks of
\(K_{\mathcal U}\).

\subtopic{45.3 Least-squares global assembly}

Let
\[
m_i=|\{\alpha:i\in V_\alpha\}|
\]
be the multiplicity of vertex \(i\). Then
\[
J_{\mathcal U}^*J_{\mathcal U}=M_{\mathcal U},
\qquad
(M_{\mathcal U}c)_i=m_ic_i.
\]
The pseudoinverse performs coordinatewise averaging:
\[
\boxed{
(J_{\mathcal U}^\dagger u)_i
=
\frac{1}{m_i}
\sum_{\alpha:i\in V_\alpha}(u_\alpha)_i.
}
\]
Let
\[
\widehat u
=
P_{\ker K_{\mathcal U}}u
=
J_{\mathcal U}\widehat c.
\]
Then
\[
\widehat c
=
J_{\mathcal U}^\dagger\widehat u
\]
and
\[
\boxed{
\|u-J_{\mathcal U}\widehat c\|
\leq
\frac{\|K_{\mathcal U}u\|}
{\sqrt{\gamma_{\mathcal U}^{\mathrm{tot}}}}.
}
\]
If \(m_{\min}=\min_i m_i\), then
\[
\boxed{
\left\|
\widehat c-J_{\mathcal U}^\dagger u
\right\|
\leq
\frac{\|K_{\mathcal U}u\|}
{\sqrt{m_{\min}\gamma_{\mathcal U}^{\mathrm{tot}}}}.
}
\]
The simple overlap average and the globally coherent least-squares assembly
are close only when the total descent energy is small relative to the cover
gap.

\subtopic{45.4 A workspace of workspaces}

Define the exact meta-workspace
\[
\boxed{
\mathscr W_{\mathrm{meta}}(\mathcal U)
:=
\ker B_{\mathcal U}
\subseteq
\bigoplus_\alpha\mathscr W_\alpha.
}
\]
Theorem~\ref{thm:exact-workspace-gluing} gives the canonical identification
\[
\boxed{
\mathscr W_{\mathrm{meta}}(\mathcal U)
\cong
\mathscr W(G).
}
\]
Thus exact formation of a workspace from local workspaces does not create an
additional mysterious carrier. It gives a second presentation of the same
global compatible-section object.

Suppose local operations
\[
A_\alpha:\mathscr W_\alpha\to\mathscr W_\alpha
\]
and overlap operations
\[
A_{\alpha\beta}:\mathscr W_{\alpha\beta}\to
\mathscr W_{\alpha\beta}
\]
satisfy
\[
r_{\alpha\beta}^{\alpha}A_\alpha
=
A_{\alpha\beta}r_{\alpha\beta}^{\alpha},
\qquad
r_{\alpha\beta}^{\beta}A_\beta
=
A_{\alpha\beta}r_{\alpha\beta}^{\beta}.
\]
Then, with \(A^0=\bigoplus_\alpha A_\alpha\) and
\(A^1=\bigoplus_{\alpha<\beta}A_{\alpha\beta}\),
\[
B_{\mathcal U}A^0=A^1B_{\mathcal U}.
\]
Therefore
\[
A^0\ker B_{\mathcal U}\subseteq\ker B_{\mathcal U}.
\]
The operation-closure theorem from Part~III reappears at the meta level:
operations on subworkspaces descend globally precisely when they agree on
interfaces.

\subtopic{45.5 Descent-faithful compression}

A practical hierarchy rarely transmits each complete local workspace. Let
\[
Q^p:\check C^p(\mathcal U;\mathscr W)\to\overline C^p
\]
be compression maps into smaller local and overlap representations. Suppose
there is a compressed coboundary \(\overline\delta^p\) such that
\[
\boxed{
Q^{p+1}\delta^p
=
\overline\delta^p Q^p.
}
\]
The maps \(Q^\bullet\) then form a cochain map.

\begin{definition}[Descent-faithful compression]\label{def:descent-faithful}
The compression is \emph{descent-faithful in degree zero} when the induced map
\[
\boxed{
H^0(Q):
\ker\delta^0
\longrightarrow
\ker\overline\delta^0
}
\]
is an isomorphism.
\end{definition}

Injectivity prevents two distinct global workspace states from collapsing to
the same compressed meta-state. Surjectivity prevents the compressed hierarchy
from introducing spurious compatible states that have no full-resolution
lift. Under this condition,
\[
\boxed{
\ker\overline\delta^0
\cong
\mathscr W(G).
}
\]
This is the exact criterion for a compressed workspace of workspaces to retain
the global core.

\subtopic{45.6 Structural and operational gluing are distinct}

Exact or approximate overlap coherence does not by itself establish global
predictive sufficiency. Let \(W_t\) denote a glued structural state and let
\[
Z_t^{\mathrm{loc}}
=
(Z_{\alpha,t})_{\alpha\in A}
\]
be local predictive summaries recovered from its chart restrictions.

\begin{proposition}[Local-to-global predictive alignment]\label{prop:local-global-predictive}
Assume \(W_t\) is measurable with respect to the complete history
\(\mathbf H_t\), every \(Z_{\alpha,t}\) is a measurable function of
\(W_t|_{G_\alpha}\), and
\[
I(Y_t^+;\mathbf H_t\mid Z_t^{\mathrm{loc}})
\leq
\varepsilon.
\]
Then
\[
\boxed{
I(Y_t^+;\mathbf H_t\mid W_t)
\leq
\varepsilon.
}
\]
In particular, exact sufficiency of the complete local-summary tuple implies
exact sufficiency of the glued workspace.
\end{proposition}

\begin{proof}
The tuple \(Z_t^{\mathrm{loc}}\) is a measurable function of \(W_t\), while
\(W_t\) is itself a history statistic. Conditioning on the finer statistic
cannot increase the residual information that the full history carries about
the target. Equivalently, use the conditional-law characterization from
Theorem~\ref{thm:predictive-state}.
\end{proof}

Causal gluing requires an additional interface gate: interventions described
in two overlapping charts must induce matched outcome laws on the overlap. A
structurally compatible family can fail this gate, just as structurally
identical replacement systems can fail causal identity in Part~XII.

\subtopic{45.7 Anti-centralization gate}

The isomorphism
\[
\mathscr W_{\mathrm{meta}}(\mathcal U)\cong\mathscr W(G)
\]
is a compositional statement, not a relational-irreducibility statement. A
single chart may still expose enough accessible information to reconstruct the
whole meta-state. A claim that the hierarchy is genuinely distributed must
therefore reuse the coalition gates from Parts~VII and XI. At the chart level,
one may require, for example,
\[
\alpha_{\{\alpha\}}^{\mathrm{meta}}=0
\quad\text{for every }\alpha,
\qquad
\alpha_A^{\mathrm{meta}}>0
\quad\text{for some }|A|\geq2,
\]
or the predictive analogue
\[
k_\varepsilon^{\mathrm{meta}}\geq2,
\qquad
\Delta_{\mathrm{joint}}^{\mathrm{meta}}>0.
\]
Capacity-matched local-memory and chart-ablation controls are still required.
Hierarchical notation cannot convert a centralized copy into a relational
workspace.

\topic{46}{Hierarchical Collapse Theorem T9}

\subtopic{46.1 Finite cover hierarchies}

A finite workspace hierarchy is a rooted tree \(\mathcal T\). Each node
\(a\in\mathcal T\) carries a subgraph \(G_a\). The root \(r\) carries
\[
G_r=G.
\]
For every internal node \(a\), its children
\[
\operatorname{Ch}(a)
\]
form a finite edge-complete cover of \(G_a\). Leaves carry primitive local
workspace spaces
\[
W_\ell=\mathscr W(G_\ell).
\]
For an internal node, define recursively
\[
\boxed{
W_a
=
\ker
\left(
\delta_a^0:
\bigoplus_{b\in\operatorname{Ch}(a)}W_b
\longrightarrow
\bigoplus_{b<c}W_{bc}
\right),
}
\]
where the overlap spaces and restrictions are inherited from the child cover.

\begin{theorem}[Exact hierarchical collapse]\label{thm:hierarchical-collapse}
For every finite workspace hierarchy as above, there is a canonical
isomorphism
\[
\boxed{
W_a\cong\mathscr W(G_a)
}
\]
for every node \(a\). In particular,
\[
\boxed{
W_r\cong\mathscr W(G).
}
\]
Therefore recursively gluing local workspaces, using any other finite
parenthesization built from edge-complete intermediate covers that retain all
required overlap objects, and solving the flattened global transfer system all
produce the same exact workspace up to canonical isomorphism.
\end{theorem}

\begin{proof}
Proceed upward from the leaves. The claim holds at every leaf by definition.
Assume it holds for all children of an internal node \(a\). Replacing each
child space \(W_b\) by its canonical copy \(\mathscr W(G_b)\), the recursive
definition of \(W_a\) is exactly the compatible-family kernel for the
edge-complete child cover of \(G_a\). Theorem~\ref{thm:exact-workspace-gluing}
therefore gives
\[
W_a\cong\mathscr W(G_a).
\]
Finite induction reaches the root.
\end{proof}

\subtopic{46.2 Flattening and limit interchange}

Let \(L(\mathcal T)\) be the leaf set. Repeated edge completeness implies that
the flattened family
\[
\mathcal U_{\mathrm{flat}}
=
\{G_\ell:\ell\in L(\mathcal T)\}
\]
is itself an edge-complete cover of \(G\). Let
\[
\delta_{\mathrm{flat}}^0
\]
be its overlap coboundary and define
\[
W_{\mathrm{flat}}
=
\ker\delta_{\mathrm{flat}}^0.
\]
Theorem~\ref{thm:exact-workspace-gluing} and
Theorem~\ref{thm:hierarchical-collapse} give
\[
\boxed{
W_r
\cong
W_{\mathrm{flat}}
\cong
\mathscr W(G).
}
\]
Equivalently,
\[
\boxed{
\operatorname{Lim}_{\mathcal T}^{\mathrm{rec}}
\left(
\{\mathscr W(G_\ell)\}_{\ell\in L(\mathcal T)}
\right)
\cong
\operatorname{Eq}_{\mathcal U_{\mathrm{flat}}}
\left(
\{\mathscr W(G_\ell)\}_{\ell\in L(\mathcal T)}
\right).
}
\]
The notation on the left means recursive finite descent along the tree; the
right side means one-shot equalization over the flattened leaf cover. In
finite-dimensional vector spaces, both are kernels of the same complete family
of local and overlap constraints, grouped in different orders. Hierarchical
collapse is therefore a finite-limit interchange property, not a claim that
intermediate levels are physically absent or operationally irrelevant.

\subtopic{46.3 Approximate hierarchical collapse}

At level \(r\), let \(K_r\) be the corresponding total descent operator,
\[
\gamma_r
=
\lambda_{\min}^+(K_r^*K_r),
\qquad
\|K_ru_r\|\leq\varepsilon_r.
\]
The nearest exact compatible state then differs from \(u_r\) by at most
\[
\frac{\varepsilon_r}{\sqrt{\gamma_r}}.
\]
Allow an additional implementation, compression, or causal-summary defect
\(\zeta_r\), and define
\[
\delta_r
=
\frac{\varepsilon_r}{\sqrt{\gamma_r}}
+
\zeta_r.
\]

\begin{theorem}[Approximate hierarchical collapse bound]\label{thm:approx-hierarchical-collapse}
Let
\[
F_r:\mathcal H_r\to\mathcal H_{r+1},
\qquad
\widetilde F_r:\mathcal H_r\to\mathcal H_{r+1}
\]
be exact and implemented level-collapse maps for
\(r=0,\ldots,R-1\). Let \(\mathcal R_r\subseteq\mathcal H_r\) contain the exact
and implemented states reachable at level \(r\). Assume, on \(\mathcal R_r\),
\[
\|F_r(x)-F_r(y)\|
\leq
L_r\|x-y\|
\]
and
\[
\sup_{x\in\mathcal R_r}
\|\widetilde F_r(x)-F_r(x)\|
\leq
\delta_r.
\]
Then
\[
\boxed{
\begin{aligned}
&\left\|
\widetilde F_{R-1}\circ\cdots\circ\widetilde F_0(x)
-
F_{R-1}\circ\cdots\circ F_0(x)
\right\|\\
&\qquad\leq
\sum_{r=0}^{R-1}
\left(
\prod_{s=r+1}^{R-1}L_s
\right)
\delta_r.
\end{aligned}
}
\]
In particular, if \(L_r\leq1\), then the total error is at most
\[
\sum_{r=0}^{R-1}\delta_r.
\]
If \(L_r\leq L<1\) and \(\delta_r\leq\delta\), then
\[
\boxed{
\text{total error}
\leq
\frac{\delta(1-L^R)}{1-L}
\leq
\frac{\delta}{1-L}.
}
\]
\end{theorem}

\begin{proof}
Let \(x_r\) and \(\widetilde x_r\) be the exact and implemented states after
\(r\) levels. Then
\[
\begin{aligned}
\|\widetilde x_{r+1}-x_{r+1}\|
&\leq
\|\widetilde F_r(\widetilde x_r)-F_r(\widetilde x_r)\|\\
&\quad+
\|F_r(\widetilde x_r)-F_r(x_r)\|\\
&\leq
\delta_r+L_r\|\widetilde x_r-x_r\|.
\end{aligned}
\]
Iterating this recurrence gives the displayed product-sum estimate. The final
two bounds follow by summing a finite or infinite geometric series.
\end{proof}

This theorem identifies three regimes:
\[
\begin{array}{ll}
L_r<1
&\text{hierarchical correction is contractive},\\
L_r=1
&\text{level errors accumulate at most additively},\\
L_r>1
&\text{deep hierarchies can amplify early mismatch}.
\end{array}
\]
A deep hierarchy is therefore stable only when its level gaps remain separated
and its assembly maps do not magnify small overlap errors faster than those
errors are corrected.

\subtopic{46.4 Causal-kernel accumulation}

Let \(K_r\) and \(\widetilde K_r\) now denote exact and implemented Markov
kernels at successive hierarchy levels, and suppose
\[
\sup_x
\operatorname{TV}
\left(
\widetilde K_r(x,\cdot),
K_r(x,\cdot)
\right)
\leq
\zeta_r.
\]
Because Markov kernels contract total variation, a telescoping argument gives
\[
\boxed{
\sup_x
\operatorname{TV}
\left(
\widetilde K_0\cdots\widetilde K_{R-1}(x,\cdot),
K_0\cdots K_{R-1}(x,\cdot)
\right)
\leq
\sum_{r=0}^{R-1}\zeta_r.
}
\]
Thus causal semantics can also be transported through a hierarchy, but only
under an explicit kernel-level compatibility test. Structural gluing alone
provides no such guarantee.

\subtopic{46.5 Direct versus hierarchical computation}

The exact theorem states
\[
\boxed{
\text{recursive exact gluing}
=
\text{one-shot exact gluing}.
}
\]
The approximate theorem instead supplies a measurable flattening defect:
\[
\boxed{
\delta_{\mathrm{flat}}^{(R)}
=
\left\|
\widetilde F_{R-1}\cdots\widetilde F_0(x)
-
F_{R-1}\cdots F_0(x)
\right\|.
}
\]
A hierarchy is empirically justified when it reduces computation,
communication, or local model complexity while keeping
\(\delta_{\mathrm{flat}}^{(R)}\), predictive risk, causal-kernel mismatch, and
anti-centralization gates within declared tolerances. Hierarchy by itself is
not evidence of emergence.

\subtopic{46.6 Gluing and hierarchical synthesis}

The exact and approximate constructions can be compressed as
\[
\boxed{
\begin{aligned}
\text{local transfer systems}
&\xrightarrow{\ \ker D_\alpha\ }
\text{local workspaces}\\
&\xrightarrow{\ \ker\delta^0\ }
\text{meta-workspace}\\
&\xrightarrow{\ \text{finite recursive descent}\ }
\text{global workspace}.
\end{aligned}
}
\]
At exact resolution, all finite parenthesizations are canonically equivalent.
At approximate resolution, every level contributes a defect scaled by its
spectral gap and by downstream amplification. Relative-only overlap data can
carry a harmonic \(H^1\) obstruction even when every local chart is internally
coherent. Operational and relational claims still require predictive, causal,
coalition, and null-model gates in addition to structural descent.

\theorypart{XIV}{Scaffolded and Endogenous Formation}

\topic{47}{Formation Systems and Scaffold Specification}

The preceding parts begin with a transfer system that has already been chosen.
A formation theory must instead explain how the nodes, edges, transfer maps,
workspace channels, and operational boundaries are obtained. This requires a
strict separation between the structure supplied by the experimenter and the
structure selected from interaction history.

\begin{definition}[Formation system]\label{def:formation-system}
A \emph{workspace formation system} is a tuple
\[
\mathfrak F
=
\bigl(
\Sigma,
\mathcal D_{\mathrm{form}},
\mathcal L,
\Omega,
\widehat\theta,
\mathcal D_{\mathrm{cert}},
\mathcal I_{\mathrm{form}}
\bigr),
\]
where:
\begin{itemize}
\item \(\Sigma\) is the externally supplied scaffold;
\item \(\mathcal D_{\mathrm{form}}\) is the formation or training history;
\item \(\mathcal L\) is the learning or selection rule;
\item \(\Omega\) is its internal random seed or stochastic state;
\item \(\widehat\theta=\mathcal L(\Sigma,\mathcal D_{\mathrm{form}},\Omega)\)
 is the learned organization;
\item \(\mathcal D_{\mathrm{cert}}\) is held-out certification data; and
\item \(\mathcal I_{\mathrm{form}}\) is a declared family of interventions on
 scaffold and formation history.
\end{itemize}
The learned organization is decomposed as
\[
\widehat\theta
=
\bigl(
\widehat V,
\widehat E,
\widehat H,
\widehat P,
\widehat{\mathcal O},
\widehat B
\bigr),
\]
where \(\widehat B\) is a learned operational boundary.
\end{definition}

The scaffold is itself stratified:
\[
\boxed{
\Sigma
=
\bigl(
\Sigma_V,
\Sigma_E,
\Sigma_H,
\Sigma_P,
\Sigma_{\mathcal O},
\Sigma_B,
\Sigma_Y
\bigr).
}
\]
The final coordinate \(\Sigma_Y\) records labels, targets, privileged masks,
or other supervision that could directly reveal the intended workspace. This
coordinate is necessary because an architecture may appear to discover a
boundary even when the desired membership mask was encoded in its training
labels.

\subtopic{47.1 Scaffolded architecture versus scaffolded solution}

A scaffold can constrain the search space without specifying its solution.
For example, a graph neural architecture can fix the class of admissible local
message-passing rules while leaving the active edges and transfer maps to be
learned. Conversely, a system can use a flexible architecture while receiving
labels that directly encode the desired workspace. The former is
architecturally scaffolded but solution-open. The latter is solution-scaffolded.

For any component \(q\) of \(\widehat\theta\), define its scaffold-only Bayes
risk under loss \(\ell_q\) by
\[
\mathcal R_q^{\Sigma}
=
\inf_f
\mathbb E
\left[
\ell_q\bigl(q,f(\Sigma,\Omega)\bigr)
\right].
\]
Define the full formation risk by
\[
\mathcal R_q^{\Sigma,\mathcal H}
=
\inf_g
\mathbb E
\left[
\ell_q\bigl(q,g(\Sigma,\mathcal D_{\mathrm{form}},\Omega)\bigr)
\right].
\]
The history advantage is
\[
\boxed{
\Delta_q^{\mathrm{hist}}
=
\mathcal R_q^{\Sigma}
-
\mathcal R_q^{\Sigma,\mathcal H}.
}
\]
A positive history advantage means that formation history contains information
about \(q\) that is not already recoverable from the declared scaffold and
random seed.

\begin{proposition}[Scaffold-only impossibility]\label{prop:scaffold-only-impossibility}
If
\[
q=f(\Sigma,\Omega)
\qquad\text{almost surely},
\]
then for every formation-history intervention that leaves \((\Sigma,\Omega)\)
fixed,
\[
\mathcal L\!\left(q\mid
\operatorname{do}(\mathcal D_{\mathrm{form}}=h)
\right)
=
\mathcal L\!\left(q\mid
\operatorname{do}(\mathcal D_{\mathrm{form}}=h')
\right),
\]
and \(\Delta_q^{\mathrm{hist}}=0\).
\end{proposition}

\begin{proof}
The value of \(q\) is measurable with respect to \((\Sigma,\Omega)\), so an
intervention on formation history cannot change its conditional law when those
variables are held fixed. The scaffold-only predictor \(f\) attains the same
risk as any predictor that also receives formation history, hence the risk gap
vanishes.
\end{proof}

This proposition is elementary but methodologically decisive. A structure
that is deterministically encoded by the scaffold cannot count as discovered
from interaction history merely because a learned network reproduces it.

\subtopic{47.2 Formation and certification split}

The same sample must not be used both to select a workspace and to certify its
coherence, causal efficacy, or relational irreducibility. The required split is
\[
\mathcal D
=
\mathcal D_{\mathrm{form}}
\sqcup
\mathcal D_{\mathrm{cal}}
\sqcup
\mathcal D_{\mathrm{cert}}.
\]
The formation sample selects \(\widehat\theta\). The calibration sample fixes
thresholds, regularization strengths, and intervention magnitudes. The
certification sample evaluates the final gates without further structural
selection. In a multi-environment study, the stronger requirement is
\[
\mathcal E_{\mathrm{form}}
\cap
\mathcal E_{\mathrm{cert}}
=
\varnothing,
\]
so that the learned organization is tested under environments not used to
construct it.

\begin{definition}[Formation leakage]\label{def:formation-leakage}
A certification variable \(Y_{\mathrm{cert}}\) leaks into formation when the
conditional law of \(\widehat\theta\) changes under an intervention on
\(Y_{\mathrm{cert}}\) that should be inaccessible during formation:
\[
\mathcal L
\left(
\widehat\theta
\mid
\operatorname{do}(Y_{\mathrm{cert}}=y)
\right)
\neq
\mathcal L
\left(
\widehat\theta
\mid
\operatorname{do}(Y_{\mathrm{cert}}=y')
\right).
\]
A valid discovery claim requires a documented information barrier or an audit
showing that this effect is absent up to tolerance.
\end{definition}

\subtopic{47.3 Origin and reality are orthogonal}

The formation provenance of a workspace and its operational certification are
separate coordinates. A fully scaffolded structure may satisfy every
structural and causal workspace gate. A highly endogenous structure may fail
to become coherent or causally useful. Symbolically,
\[
\operatorname{Endogenous}(W)
\centernot\Longrightarrow
\operatorname{SRW}_{\Theta}(W),
\]
\[
\operatorname{SRW}_{\Theta}(W)
\centernot\Longrightarrow
\operatorname{Endogenous}(W).
\]
The theory therefore treats endogeneity as a provenance axis rather than a
higher level on the workspace hierarchy.

\topic{48}{Operational Endogeneity Signature}

The ideal description-complexity expression from the symbolic outline is
retained as a conceptual limit, but raw Kolmogorov complexity is not
computable. A usable theory must declare a model class, a code, and an
intervention family.

\begin{definition}[Code-relative history contribution]\label{def:mdl-endogeneity}
Let \(\mathfrak M_q\) be a declared coding or predictive model class for a
learned component \(q\). Let
\[
L_{\mathfrak M_q}(q\mid Z)
\]
be its held-out codelength given side information \(Z\). Define
\[
\boxed{
 e_q^{\mathrm{MDL}}
 =
 \left[
 \frac{
 L_{\mathfrak M_q}(q\mid\Sigma)
 -
 L_{\mathfrak M_q}(q\mid\Sigma,\mathcal D_{\mathrm{form}})
 }{
 \max\{L_{\mathfrak M_q}(q\mid\Sigma),1\}
 }
 \right]_{0}^{1},
}
\]
where \([x]_0^1=\min\{1,\max\{0,x\}\}\).
\end{definition}

This quantity is explicitly code-relative. It must not be presented as an
observer-independent metaphysical measure. Its role is to test whether the
formation history provides compressive information about the learned
organization beyond the declared scaffold.

\begin{definition}[Causal formation sensitivity]\label{def:causal-formation-sensitivity}
Let \(\iota_h\) and \(\iota_{h'}\) be matched feasible interventions that alter
formation history while preserving the scaffold, randomization law, sample
size, and intervention cost. For a divergence \(D\), define
\[
\boxed{
 e_q^{\mathrm{causal}}
 =
 \mathbb E_{(h,h')\sim\nu_q}
 D\!\left(
 \mathcal L^{\iota_h}(q),
 \mathcal L^{\iota_{h'}}(q)
 \right).
}
\]
A common null uses temporally or relationally shuffled histories that preserve
local marginals but destroy the candidate formation signal.
\end{definition}

\begin{proposition}[Reparameterization invariance of causal provenance]
\label{prop:causal-provenance-invariance}
Let \(\phi\) be a measurable bijection on the range of \(q\), and let the
divergence \(D\) be invariant under measurable bijections. Then
\[
 e_{\phi(q)}^{\mathrm{causal}}
 =
 e_q^{\mathrm{causal}}.
\]
\end{proposition}

\begin{proof}
Every intervention-induced law of \(\phi(q)\) is the pushforward of the
corresponding law of \(q\). Bijection invariance of \(D\) preserves each
pairwise divergence, and therefore their design-weighted expectation.
\end{proof}

The causal coordinate is therefore insensitive to a mere renaming of the
learned structure. The MDL coordinate need not be invariant unless the coding
class is transformed with the reparameterization.

\subtopic{48.1 Scaffold dependence and cross-scaffold invariance}

Endogeneity does not require independence from all scaffolds. Every learning
system has an architecture, objective, and resource budget. The relevant
question is whether the learned organization depends on the particular
solution-bearing details of the scaffold.

For a family \(\mathfrak S\) of admissible scaffold variants, define the
cross-scaffold identity defect
\[
\boxed{
\delta_q^{\mathrm{scaf}}
=
\inf_{\{J_{\sigma\sigma'}\}}
\mathbb E_{\sigma,\sigma'\sim\pi}
 d_q
\left(
J_{\sigma\sigma'}q_{\sigma},q_{\sigma'}
\right),
}
\]
where \(J_{\sigma\sigma'}\) transports the learned object between admissible
coordinate systems. A low value indicates that the same abstract
organization is recovered under multiple non-solution-bearing scaffolds.

The complete formation signature is
\[
\boxed{
\mathbf E_{\mathrm{form}}(W)
=
\left(
 e_V,e_E,e_H,e_P,e_{\mathcal O},e_B;
 \delta^{\mathrm{scaf}},
 \delta^{\mathrm{leak}}
\right),
}
\]
where every \(e_q\) is reported as the pair
\[
 e_q
=
\left(
 e_q^{\mathrm{MDL}},
 e_q^{\mathrm{causal}}
\right).
\]
No scalar average is permitted to hide a zero history contribution for one of
the structures claimed to have been discovered.

\subtopic{48.2 Formation certificate}

For thresholds \(\Theta_{\mathrm{form}}\), define
\[
\operatorname{Form}_{\Theta_{\mathrm{form}}}(W)=1
\]
only when all applicable gates pass:
\[
\begin{aligned}
& e_q^{\mathrm{MDL}}\geq \tau_q^{\mathrm{MDL}},
&& e_q^{\mathrm{causal}}\geq \tau_q^{\mathrm{causal}},\\
& \delta_q^{\mathrm{scaf}}\leq\varepsilon_q^{\mathrm{scaf}},
&& \delta^{\mathrm{leak}}\leq\varepsilon_{\mathrm{leak}},\\
& \operatorname{SRW}_{\Theta}(W)=1,
&& \operatorname{Pass}(N_{\mathrm{form}})=1.
\end{aligned}
\]
The final workspace gate is included because endogenous discovery of an
unstable or causally irrelevant pattern is not endogenous workspace formation.

\topic{49}{Learned Transfer Structure and Formation Classes}

\subtopic{49.1 Gated membership and edge learning}

Let
\[
z_i\in\{0,1\},
\qquad
a_e\in[0,1]
\]
be node-membership and edge-activation variables satisfying
\[
a_e\leq z_{s(e)},
\qquad
a_e\leq z_{t(e)}.
\]
A weighted learned discrepancy operator is
\[
(D_{a,H}c)_e
=
\sqrt{w_ea_e}
\left(
 c_{t(e)}-H_ec_{s(e)}
\right).
\]
One generic formation objective is
\[
\boxed{
\begin{aligned}
\widehat\theta
\in
\argmin_{z,a,H,P,\mathcal O,B}
\Bigl
\{&
\widehat{\mathcal R}_{\mathrm{task}}
+
\lambda_{\mathrm{coh}}\|D_{a,H}c\|^2
+
\lambda_E\|a\|_1
+
\lambda_V\|z\|_1\\
&+
\lambda_{\mathrm{stab}}\mathcal J_{\mathrm{stab}}
+
\lambda_{\mathrm{leak}}\mathcal J_{\mathrm{leak}}
\Bigr\}.
\end{aligned}
}
\]
The objective is schematic rather than universal. Its purpose is to expose
which quantities were fixed and which were optimized. If \(z\) and \(a\) are
fixed by \(\Sigma\), only the transfer maps or local channels can be claimed as
learned.

\subtopic{49.2 Gauge-identifiable transfer learning}

A learned transfer map is not absolutely identifiable when local coordinates
are latent. Under local invertible changes
\[
X_u'=g_uX_u,
\qquad
X_v'=g_vX_v,
\]
the same abstract transfer is represented by
\[
H_{uv}'=g_vH_{uv}g_u^{-1}.
\]
The identifiable object is therefore the gauge class
\[
[H]_{\mathcal G}
=
\left\{
(g_{t(e)}H_eg_{s(e)}^{-1})_{e\in E}
:
(g_i)_{i\in V}\in\mathcal G
\right\}.
\]
Any formation claim that compares raw matrices across independently trained
systems without gauge alignment is ill posed.

\subtopic{49.3 Transfer identification theorem \(T_{15}\) }

\begin{theorem}[Population transfer identification]\label{thm:transfer-identification}
Let \(X\in\mathbb R^{d_u}\) and \(Y\in\mathbb R^{d_v}\) satisfy
\[
Y=H_*X+\xi,
\qquad
\mathbb E[\xi X^{\mathsf T}]=0,
\]
and assume
\[
\Sigma_{XX}
=
\mathbb E[XX^{\mathsf T}]
\succeq
\kappa I
\]
for some \(\kappa>0\). Then \(H_*\) is the unique minimizer of
\[
\mathcal R(H)
=
\mathbb E\|Y-HX\|_2^2,
\]
with
\[
\boxed{
H_*
=
\Sigma_{YX}\Sigma_{XX}^{-1}.
}
\]
Moreover,
\[
\boxed{
\mathcal R(H)-\mathcal R(H_*)
=
\operatorname{tr}
\left(
(H-H_*)\Sigma_{XX}(H-H_*)^{\mathsf T}
\right)
\geq
\kappa\|H-H_*\|_F^2.
}
\]
\end{theorem}

\begin{proof}
Substitute \(Y=H_*X+\xi\) into the risk. The cross term vanishes because
\(\mathbb E[\xi X^{\mathsf T}]=0\), giving
\[
\mathcal R(H)
=
\mathcal R(H_*)
+
\mathbb E\|(H-H_*)X\|_2^2.
\]
The remaining expectation is the displayed trace. Positive definiteness of
\(\Sigma_{XX}\) gives uniqueness and the lower bound. The normal equations
give the covariance formula.
\end{proof}

\begin{corollary}[Deterministic empirical recovery bound]
\label{cor:empirical-transfer-bound}
For samples \((X_r,Y_r)_{r=1}^n\), write
\[
\widehat\Sigma_{XX}
=
\frac1n\sum_{r=1}^nX_rX_r^{\mathsf T},
\qquad
\widehat M
=
\frac1n\sum_{r=1}^n\xi_rX_r^{\mathsf T}.
\]
If \(\widehat\Sigma_{XX}\succeq\widehat\kappa I\), then ordinary least squares
satisfies
\[
\boxed{
\|\widehat H-H_*\|_F
\leq
\frac{\|\widehat M\|_F}{\widehat\kappa}.
}
\]
For ridge parameter \(\tau>0\),
\[
\boxed{
\|\widehat H_\tau-H_*\|_F
\leq
\frac{
\|\widehat M\|_F+\tau\|H_*\|_F
}{
\widehat\kappa+\tau
}.
}
\]
\end{corollary}

\begin{proof}
The empirical normal equations give
\[
\widehat H-H_*
=
\widehat M\widehat\Sigma_{XX}^{-1}.
\]
For ridge regression,
\[
\widehat H_\tau-H_*
=
(\widehat M-\tau H_*)
(\widehat\Sigma_{XX}+\tau I)^{-1}.
\]
Taking Frobenius and operator norms yields both estimates.
\end{proof}

\begin{proposition}[Gauge covariance of least squares]
\label{prop:ols-gauge-covariance}
Under invertible coordinate transformations
\[
X'=g_uX,
\qquad
Y'=g_vY,
\]
the population least-squares solution transforms as
\[
H_*'=g_vH_*g_u^{-1}.
\]
\end{proposition}

\begin{proof}
The transformed covariances are
\[
\Sigma_{Y'X'}=g_v\Sigma_{YX}g_u^{\mathsf T},
\qquad
\Sigma_{X'X'}=g_u\Sigma_{XX}g_u^{\mathsf T}.
\]
Substitution into the covariance formula and inversion of the second factor
give the result.
\end{proof}

The theorem establishes identification only under excitation and exogeneity.
When \(\Sigma_{XX}\) is singular, the transfer is identified only on the
excited source subspace. When \(\mathbb E[\xi X^{\mathsf T}]\neq0\), regression
recovers a predictive map rather than necessarily the intended causal
transfer.

\subtopic{49.4 Held-out structure-selection theorem \(T_{16}\) }

The same finite-margin argument applies to node membership, edge selection,
boundary selection, or a finite family of gauge-normalized transfer models.

\begin{theorem}[Uniform structure-selection recovery]
\label{thm:uniform-structure-selection}
Let \(\mathfrak T\) be a finite candidate family with population score
\(J:\mathfrak T\to\mathbb R\). Let
\[
\theta_*
=
\argmin_{\theta\in\mathfrak T}J(\theta)
\]
be unique, with margin
\[
\Delta_J
=
\min_{\theta\neq\theta_*}
\bigl(J(\theta)-J(\theta_*)\bigr)>0.
\]
If a held-out estimate \(\widehat J\) satisfies
\[
\sup_{\theta\in\mathfrak T}
|\widehat J(\theta)-J(\theta)|
\leq\varepsilon,
\]
then every empirical minimizer \(\widehat\theta\) obeys
\[
\boxed{
J(\widehat\theta)-J(\theta_*)
\leq2\varepsilon.
}
\]
If \(2\varepsilon<\Delta_J\), then
\[
\boxed{\widehat\theta=\theta_*.}
\]
\end{theorem}

\begin{proof}
By empirical optimality and uniform approximation,
\[
J(\widehat\theta)
\leq
\widehat J(\widehat\theta)+\varepsilon
\leq
\widehat J(\theta_*)+\varepsilon
\leq
J(\theta_*)+2\varepsilon.
\]
If \(\widehat\theta\neq\theta_*\), the margin would force an excess larger than
\(2\varepsilon\), a contradiction.
\end{proof}

The theorem does not guarantee that the candidate class contains the correct
organization. It states exactly what follows once a finite class, score, and
uniform certification error are declared.

\subtopic{49.5 Formation classes}

The provenance classes are now defined relative to
\[
(\Sigma,\mathcal L,\mathcal I_{\mathrm{form}},\mathfrak M).
\]
They are not absolute labels.

\begin{description}[leftmargin=2.1em,style=nextline]
\item[Exogenously scaffolded.]
The claimed nodes, edges, transfer maps, channels, operations, or boundary are
recoverable from the scaffold alone at the declared tolerance. Their history
advantages and causal formation sensitivities vanish.

\item[Adaptively parameterized.]
The membership and channel topology are fixed, while transfer parameters,
local operators, or readouts are learned from interaction history. Typically
\[
e_H+e_P+e_{\mathcal O}>0,
\qquad
e_V\approx e_E\approx e_B\approx0.
\]

\item[Structurally adaptive.]
At least one of membership, edges, or boundary is selected from history, but
the search family and objective remain externally specified.

\item[Endogenously constituted relative to the scaffold class.]
Membership, transfer organization, and operational boundary have nonzero
history contribution and causal formation sensitivity, survive cross-scaffold
transport, pass leakage and matched-null tests, and produce a certified
workspace on held-out environments.
\end{description}

The phrase ``relative to the scaffold class'' is mandatory. No finite
experiment can establish independence from every possible external influence.

\theorypart{XV}{Boundary Formation and Stability}

\topic{50}{Candidate Workspace Boundaries}

A boundary is not merely a subset of physical components. It is a subset
relative to a declared node decomposition, prediction target, intervention
family, environment distribution, and time scale.

\begin{definition}[Boundary problem]\label{def:boundary-problem}
A boundary problem is a tuple
\[
\mathfrak B
=
\left(
V,
\mathfrak S,
\tau,
\mathcal Y,
\mathcal I,
\Pi,
\mathcal G_B
\right),
\]
where \(\mathfrak S\subseteq2^V\) is the admissible candidate family,
\(\tau\) is the temporal scale, \(\mathcal Y\) is the operational target,
\(\mathcal I\) is the feasible intervention family, \(\Pi\) is the
environment distribution, and \(\mathcal G_B\) is a symmetry group under
which boundaries are considered equivalent.
\end{definition}

Unless the application requires otherwise, the candidate family excludes
\[
\varnothing,
\qquad
V,
\qquad
\{i\}\ \text{for }i\in V.
\]
This anti-triviality rule prevents the empty set, the entire realization, or a
single centralized controller from winning solely through a degenerate score.

\subtopic{50.1 Observational recurrence, dependence, and outward influence}

For \(S\in\mathfrak S\), define
\[
\Ret_\tau(S)
=
I\!\left(
X_S(t+\tau);X_S(t)
\mid X_{\bar S}(t),E_t
\right),
\]
\[
\Dep_\tau(S)
=
I\!\left(
X_S(t+\tau);X_{\bar S}(t)
\mid X_S(t),E_t
\right),
\]
\[
\Out_\tau(S)
=
I\!\left(
X_{\bar S}(t+\tau);X_S(t)
\mid X_{\bar S}(t),E_t
\right).
\]
These quantities are observational. They can be nonzero because of common
causes, selection, or feedback paths and therefore cannot alone establish a
causal boundary.

\subtopic{50.2 Causal boundary terms}

Let \(\iota_{S,x}\) be a feasible intervention that changes the internal
workspace coordinate of \(S\) to \(x\) while matching background variables and
intervention cost. Define outward causal efficacy
\[
\boxed{
\COut_\tau(S)
=
\mathbb E_{x,x'\sim\nu_S}
D\!\left(
\mathcal L^{\iota_{S,x}}
(X_{\bar S}(t+\tau)),
\mathcal L^{\iota_{S,x'}}
(X_{\bar S}(t+\tau))
\right).
}
\]
Define inward causal dependence analogously by intervening on \(\bar S\) and
measuring the future law of \(X_S\):
\[
\boxed{
\CIn_\tau(S)
=
\mathbb E_{u,u'\sim\nu_{\bar S}}
D\!\left(
\mathcal L^{\iota_{\bar S,u}}
(X_S(t+\tau)),
\mathcal L^{\iota_{\bar S,u'}}
(X_S(t+\tau))
\right).
}
\]
Neither term should be minimized to zero in general. A useful workspace may
receive substantial external input. The boundary objective must instead
express the intended balance between internal recurrence, external dependence,
and outward efficacy.

\subtopic{50.3 Structural and anti-centralization gates}

For the induced subsystem on \(S\), define
\[
\mathsf{Coh}(S)
=
\operatorname{tr}
\left(
P_{S,d}L_SP_{S,d}
\right),
\]
where \(P_{S,d}\) projects onto the selected low-energy workspace modes. Let
\[
\mathsf{Cut}(S)
=
\sum_{e:s(e)\in S,\ t(e)\notin S}w_e
+
\sum_{e:s(e)\notin S,\ t(e)\in S}w_e.
\]
Let \(k_\varepsilon(S)\) be the minimal sufficient coalition size inside the
candidate subsystem. A candidate is admissible only if
\[
\operatorname{SRW}_{\Theta_S}(W_S)=1,
\qquad
k_\varepsilon(S)\geq2,
\]
unless the claim explicitly concerns a centralized workspace.

\topic{51}{Boundary Objective, Recovery, and Persistence}

A general score is
\[
\boxed{
\begin{aligned}
\mathcal B_t(S)
={}&
 a_R\Ret_\tau(S)
+a_O\Out_\tau(S)
+a_C\COut_\tau(S)\\
&-a_D\Dep_\tau(S)
-a_I\CIn_\tau(S)
-a_L\mathsf{Coh}(S)\\
&-a_K\mathsf{Cut}(S)
-a_Z\mathsf{Size}(S)
-a_T\Instability_t(S).
\end{aligned}
}
\]
The coefficients and all normalizations are part of the boundary problem. No
single choice is universal. The selected boundary is
\[
\widehat S_t
\in
\argmax_{S\in\mathfrak S_{\mathrm{pass}}}
\widehat{\mathcal B}_t(S),
\]
where \(\mathfrak S_{\mathrm{pass}}\) contains only candidates that satisfy the
required workspace and anti-centralization gates.

\subtopic{51.1 Boundary equivalence and identifiability}

If a symmetry \(g\in\mathcal G_B\) leaves the full observational and
interventional law invariant, \(S\) and \(gS\) cannot be distinguished by the
declared experiment. The identifiable boundary is then the orbit
\[
[S]_{\mathcal G_B}
=
\{gS:g\in\mathcal G_B\}.
\]
A unique subset should not be claimed when only an equivalence class is
identifiable.

\subtopic{51.2 Stable boundary recovery \(T_{17}\) }

\begin{theorem}[Held-out boundary recovery]\label{thm:heldout-boundary-recovery}
Let \(\mathfrak S_{\mathrm{pass}}\) be finite and let
\[
S_t^*
=
\argmax_{S\in\mathfrak S_{\mathrm{pass}}}
\mathcal B_t(S)
\]
be unique. Define the population margin
\[
\Delta_t^B
=
\mathcal B_t(S_t^*)
-
\max_{S\neq S_t^*}\mathcal B_t(S)>0.
\]
If
\[
\sup_{S\in\mathfrak S_{\mathrm{pass}}}
\left|
\widehat{\mathcal B}_t(S)-\mathcal B_t(S)
\right|
\leq\varepsilon_t,
\]
then every empirical maximizer \(\widehat S_t\) satisfies
\[
\boxed{
\mathcal B_t(S_t^*)-
\mathcal B_t(\widehat S_t)
\leq2\varepsilon_t.
}
\]
If \(2\varepsilon_t<\Delta_t^B\), then
\[
\boxed{\widehat S_t=S_t^*.}
\]
\end{theorem}

\begin{proof}
Apply Theorem~\ref{thm:uniform-structure-selection} to
\(J(S)=-\mathcal B_t(S)\).
\end{proof}

The theorem converts uncertainty in mutual-information estimation, causal
interventions, structural scores, and stability penalties into an explicit
boundary-identification condition. It also exposes when the data are
insufficient: a small margin makes the boundary intrinsically fragile.

\subtopic{51.3 Temporal boundary stability}

\begin{theorem}[Boundary persistence under score drift]
\label{thm:boundary-drift-stability}
Assume \(S_t^*\) is the unique maximizer of \(\mathcal B_t\) with margin
\(\Delta_t^B\), and suppose
\[
\sup_{S\in\mathfrak S_{\mathrm{pass}}}
|\mathcal B_{t+1}(S)-\mathcal B_t(S)|
\leq\delta_t.
\]
If
\[
2\delta_t<\Delta_t^B,
\]
then
\[
\boxed{S_{t+1}^*=S_t^*.}
\]
More generally,
\[
\mathcal B_{t+1}(S_{t+1}^*)
-
\mathcal B_{t+1}(S_t^*)
\leq2\delta_t.
\]
\end{theorem}

\begin{proof}
For every competitor \(S\neq S_t^*\),
\[
\begin{aligned}
\mathcal B_{t+1}(S_t^*)-
\mathcal B_{t+1}(S)
&\geq
\mathcal B_t(S_t^*)-
\mathcal B_t(S)-2\delta_t\\
&\geq
\Delta_t^B-2\delta_t>0.
\end{aligned}
\]
Hence the maximizer is unchanged. The approximate bound follows by comparing
both time-indexed scores at their respective maximizers.
\end{proof}

If the score depends on a parameter \(\vartheta_t\) and satisfies
\[
\sup_S
|\mathcal B_{\vartheta}(S)-\mathcal B_{\vartheta'}(S)|
\leq
L_Bd(\vartheta,\vartheta'),
\]
then the sufficient stability condition becomes
\[
2L_Bd(\vartheta_t,\vartheta_{t+1})
<
\Delta_t^B.
\]

\subtopic{51.4 Analyst-discovered, system-instantiated, and self-maintained boundaries}

Three claims must remain separate.

\begin{description}[leftmargin=2.1em,style=nextline]
\item[Analyst-discovered boundary.]
An external procedure fits \(\widehat S\) after observing the system. The
system need not represent or use that boundary.

\item[System-instantiated boundary.]
The system contains a channel, mask, or routing variable causally equivalent to
\(\widehat S\), and interventions on that variable alter downstream operation.

\item[Self-maintained boundary.]
The system detects degradation of its boundary certificate and takes actions
that restore membership, transfer, or reconstruction properties.
\end{description}

Only the second and third claims concern a boundary realized within the
system. A post hoc partition that predicts behavior well is not by itself a
self-discovered boundary.

\subtopic{51.5 Boundary formation certificate}

A boundary \(\widehat S\) is certified as endogenously formed only if:
\begin{enumerate}
\item the membership mask was not directly encoded in \(\Sigma_B\) or
 \(\Sigma_Y\);
\item \(e_B^{\mathrm{MDL}}\) and \(e_B^{\mathrm{causal}}\) exceed declared
 thresholds;
\item held-out recovery satisfies Theorem~\ref{thm:heldout-boundary-recovery};
\item the boundary remains stable under admissible environment drift or its
 changes are predicted by score-margin closure;
\item matched history-shuffle and label-leakage controls fail to reproduce
 the result, as do full-system, empty-set, singleton, and centralized
 boundary nulls;
\item the induced subsystem satisfies its workspace certificate; and
\item a system-instantiation intervention demonstrates that the boundary is
 used by the system rather than only by the analyst.
\end{enumerate}

The results below establish the statistical and deterministic stability components. They
does not prove that an arbitrary adaptive system will spontaneously learn such
a boundary.


\theorypart{XVI}{Agency, Counterfactual Control, and Self-Maintenance}

\topic{52}{Controlled Workspace Agents and Goal States}

The preceding parts define when a distributed structure is coherent,
operationally sufficient, causally efficacious, relationally irreducible,
persistent, reconstructible, and endogenously formed. None of those
properties by itself establishes agency. Agency adds a controlled relation
among a represented objective, a set of feasible actions, counterfactual
consequences, and a policy that uses those quantities to alter the future.
Self-maintenance adds a further loop in which damage to the workspace becomes
an input to actions that preserve or restore the workspace certificate.

The purpose of this part is not to infer intentions from fluent behavior. It
is to define an intervention-based certificate whose components can fail
separately.

\subtopic{52.1 Controlled relational realization}

A controlled relational realization extends Definition~\ref{def:realization}
by the tuple
\[
\boxed{
\mathfrak A_{\Workspace}
=
\left(
\mathfrak R,
\mathcal G,
\mathcal A_t,
\mathcal Y^+,
\{U_g\}_{g\in\mathcal G},
\Pi,
\mathcal I_A,
\mathcal I_G,
\mathcal I_D,
\mathcal I_R
\right).
}
\]
Here:
\begin{itemize}
\item \(\mathfrak R\) is the underlying relational realization;
\item \(\mathcal G\) is a measurable goal-code space;
\item \(\mathcal A_t\) is the feasible action set at time \(t\);
\item \(\mathcal Y^+\) is the relevant future trajectory space;
\item \(U_g:\mathcal Y^+\to\mathbb R\) evaluates a trajectory under goal
 code \(g\);
\item \(\Pi\) is an admissible class of workspace-level policies;
\item \(\mathcal I_A\) and \(\mathcal I_G\) are feasible action and goal
 interventions; and
\item \(\mathcal I_D\) and \(\mathcal I_R\) are feasible damage and repair
 interventions.
\end{itemize}

A controlled state contains at least
\[
C_t^{\mathrm{ctl}}
=
\left(
W_t,G_t,M_t,B_t,E_t
\right),
\]
where \(G_t\) is the represented goal, \(M_t\) is a workspace-integrity
monitor, \(B_t\) is the current boundary or routing state, and \(E_t\) is the
declared external context. An action is generated by
\[
A_t
\sim
\pi_t
\left(
\cdot\mid W_t,G_t,M_t,B_t,E_t
\right).
\]
The future controlled state obeys a time-unrolled structural model
\[
S_{t+1}
=
F_t(S_t,A_t,E_{t+1},\xi_{t+1}),
\qquad
W_{t+1}=\omega_{t+1}(S_{t+1}).
\]
The intervention semantics are inherited from Part~VII. A symbol such as
\(\operatorname{do}(A_t=a)\) is meaningful only when an implementation
intervention changes an actual action mechanism while the declared background
variables are matched.

\subtopic{52.2 Goal content and goal-code equivalence}

A persistent internal variable is not automatically a goal. A goal code must
index a stable counterfactual preference over relevant future trajectories.
For a fixed state \(w\), define the goal-conditioned action value
\[
Q_t^g(a\mid w)
=
\mathbb E
\left[
U_g(Y_t^+)
\mid
\operatorname{do}(A_t=a),
W_t=w,
G_t=g
\right].
\]
When action costs are relevant, they are incorporated into \(U_g\) or declared
as a constraint on \(\mathcal A_t\).

Two goal codes are operationally equivalent,
\[
g\sim_{\mathrm{goal}}g',
\]
when they induce the same action ordering in every declared state and
intervention background:
\[
Q_t^g(a\mid w)
\geq
Q_t^g(a'\mid w)
\quad\Longleftrightarrow\quad
Q_t^{g'}(a\mid w)
\geq
Q_t^{g'}(a'\mid w).
\]
The abstract goal is the equivalence class
\[
[g]_{\mathrm{goal}}.
\]
This quotient prevents a mere invertible recoding of a goal channel from being
reported as a change of objective. For quantitative regret bounds, a
representative normalization of \(U_g\) must nevertheless be fixed.

\begin{remark}[Goals are claim-relative]
The same internal code can count as a goal for one declared trajectory family
and fail to do so for another. A code that predicts a short-horizon action may
not encode a persistent long-horizon objective. Every agency claim must state
its horizon, outcome space, feasible actions, and utility normalization.
\end{remark}

\subtopic{52.3 Local, distributed, and workspace-level goals}

Let \(\mathcal H_t^S\) be the history available to coalition \(S\subseteq V\).
For a loss \(\ell_G\), define goal-reconstruction risk
\[
\Risk_G(S)
=
\inf_f
\mathbb E
\left[
\ell_G
\left(
G_t,
f(\mathcal H_t^S)
\right)
\right].
\]
For tolerance \(\varepsilon\), define
\[
k_{G,\varepsilon}
=
\min
\left\{
|S|:
\Risk_G(S)
\leq
\Risk_G(V)+\varepsilon
\right\}.
\]
A goal is locally stored when \(k_{G,\varepsilon}=1\). It is distributed
relative to the declared decomposition when \(k_{G,\varepsilon}\geq2\).

This distinction is independent of whether the rest of the workspace is
relational. A relational workspace can obey a goal copied in full into one
controller. Such a system may possess a workspace but not a
\emph{workspace-level goal} in the strong distributed sense.

For the complete control state, define
\[
\Risk_{\mathrm{ctl}}(S)
=
\inf_f
\mathbb E
\left[
\ell_{\mathrm{ctl}}
\left(
C_t^{\mathrm{ctl}},
f(\mathcal H_t^S)
\right)
\right],
\]
and let \(k_{\mathrm{ctl},\varepsilon}\) be its minimal sufficient coalition
size. A strong workspace-level agency claim requires
\[
\boxed{k_{\mathrm{ctl},\varepsilon}\geq2.}
\]
This is a location criterion. It does not yet establish goal persistence,
counterfactual evaluation, or self-maintenance.

\subtopic{52.4 Action mediation and outcome control}

Goal-dependent action mediation is measured by feasible interventions on the
goal channel while matching the current workspace and external context. For
a divergence \(D_A\), define
\[
\eta_{G\to A}
=
\sup_{g,g',w,e}
D_A
\left(
\mathbb P(A_t\in\cdot\mid\iota_g,w,e),
\mathbb P(A_t\in\cdot\mid\iota_{g'},w,e)
\right).
\]
Workspace-dependent action mediation is
\[
\eta_{W\to A}
=
\sup_{w,w',g,e}
D_A
\left(
\mathbb P(A_t\in\cdot\mid\iota_w,g,e),
\mathbb P(A_t\in\cdot\mid\iota_{w'},g,e)
\right).
\]
An action must also alter the declared future:
\[
\eta_{A\to Y}
=
\sup_{a,a',b}
D_Y
\left(
\mathbb P(Y_t^+\in\cdot\mid\iota_a,b),
\mathbb P(Y_t^+\in\cdot\mid\iota_{a'},b)
\right).
\]
A positive \(\eta_{G\to A}\) without \(\eta_{A\to Y}\) describes goal-sensitive
output selection with no demonstrated control of the relevant environment. A
positive \(\eta_{A\to Y}\) without a goal or counterfactual-value channel may
describe a reactive controller.

\subtopic{52.5 Why observational goal persistence is insufficient}

The observational condition
\[
I(G_{t+\tau};G_t\mid E_{t+1:t+\tau})>0
\]
can arise from several non-agentic mechanisms:
\begin{enumerate}
\item a persistent external command writes the same code at every step;
\item a clock, episode label, or task identifier drives both variables;
\item a frozen memory cell repeats a symbol that is not used by the policy;
\item selection on successful trajectories induces apparent temporal
 dependence; or
\item a downstream action is copied back into the nominal goal channel.
\end{enumerate}
A goal-persistence claim therefore requires a cue-controlled intervention that
changes \(G_t\), removes or matches subsequent command inputs, and checks both
future goal state and action consequences.

\subtopic{52.6 Cue-controlled goal persistence}

Let \(C_t\) denote an external command or task cue. Fix a neutral or matched
future command trajectory \(c_{t+1:t+\tau}^{0}\). For goal interventions
\(\iota_g\) and \(\iota_{g'}\), define the cue-controlled retention effect
\[
\eta_G(\tau;g,g')
=
D_G
\left(
\mathbb P
\left(
G_{t+\tau}\in\cdot
\mid
\iota_g,
\operatorname{do}(C_{t+1:t+\tau}=c^0)
\right),
\mathbb P
\left(
G_{t+\tau}\in\cdot
\mid
\iota_{g'},
\operatorname{do}(C_{t+1:t+\tau}=c^0)
\right)
\right).
\]
The corresponding action-retention effect is
\[
\eta_{G\to A}(\tau;g,g')
=
D_A
\left(
\mathbb P(A_{t+\tau}\in\cdot\mid\iota_g,c^0),
\mathbb P(A_{t+\tau}\in\cdot\mid\iota_{g'},c^0)
\right).
\]
The first quantity tests whether a goal difference survives. The second tests
whether the surviving difference remains behaviorally operative.

Cue dependence must be reported separately. For matched initial goal \(g\),
define
\[
\delta_G^{\mathrm{cue}}(\tau)
=
\sup_{c,c'}
D_G
\left(
\mathbb P(G_{t+\tau}\in\cdot\mid\iota_g,\operatorname{do}(C=c)),
\mathbb P(G_{t+\tau}\in\cdot\mid\iota_g,\operatorname{do}(C=c'))
\right).
\]
High cue dependence is not a failure when the claimed agent is intentionally
instruction following. It does mean that persistence cannot be attributed to
an autonomous internal objective without an additional cue-removal result.

\subtopic{52.7 \(T_{19}\): Cue-controlled goal-retention theorem }

\begin{theorem}[Cue-controlled goal retention]
\label{thm:goal-retention}
Let \(G_s\in\mathbb R^{d_G}\) obey
\[
G_{s+1}
=
R_sG_s+B_sC_{s+1}+\zeta_{s+1}
\qquad
(s=t,\ldots,t+\tau-1).
\]
Consider two intervention runs that set
\[
G_t=g,
\qquad
G_t'=g',
\]
and then use the same command trajectory and the same disturbance realization.
Let
\[
R_{t+\tau:t}
=
R_{t+\tau-1}\cdots R_t.
\]
Then
\[
\boxed{
G_{t+\tau}-G_{t+\tau}'
=
R_{t+\tau:t}(g-g').
}
\]
If
\[
m_s
\leq
\sigma_{\min}(R_s),
\qquad
\|R_s\|\leq M_s,
\]
then
\[
\boxed{
\left(
\prod_{s=t}^{t+\tau-1}m_s
\right)
\|g-g'\|
\leq
\|G_{t+\tau}-G_{t+\tau}'\|
\leq
\left(
\prod_{s=t}^{t+\tau-1}M_s
\right)
\|g-g'\|.
}
\]
\end{theorem}

\begin{proof}
Subtract the two update equations. The matched command and disturbance terms
cancel, giving
\[
G_{s+1}-G_{s+1}'
=
R_s(G_s-G_s').
\]
Iteration gives the exact product formula. The upper bound follows from
submultiplicativity of the operator norm. The lower bound follows from
\[
\sigma_{\min}(AB)
\geq
\sigma_{\min}(A)\sigma_{\min}(B)
\]
for square linear maps.
\end{proof}

The theorem isolates internal retention from repeated external rewriting. It
also makes clear that persistence is graded. Uniformly contractive maps can
retain a goal for a finite horizon while erasing it asymptotically. A neutral
mode with product singular value near one supports long retention. An
expansive mode can preserve distinctions but may be dynamically unstable.

\begin{corollary}[Behavioral readout of a retained goal]
\label{cor:goal-action-readout}
Assume, in addition, that under the matched intervention background the action
readout has conditional mean
\[
\mathbb E[A_{t+\tau}\mid G_{t+\tau},W_{t+\tau}]
=
K_{t+\tau}G_{t+\tau}+b(W_{t+\tau}),
\]
and that the two runs are matched on \(W_{t+\tau}\). Then
\[
\mathbb E[A_{t+\tau}\mid\iota_g]
-
\mathbb E[A_{t+\tau}\mid\iota_{g'}]
=
K_{t+\tau}R_{t+\tau:t}(g-g').
\]
If \(\sigma_{\min}(K_{t+\tau})\geq k_{t+\tau}>0\), the norm of this action-mean
difference is at least
\[
k_{t+\tau}
\left(
\prod_{s=t}^{t+\tau-1}m_s
\right)
\|g-g'\|.
\]
\end{corollary}

The corollary is a sufficient condition, not a universal model of action
selection. For discrete actions, the corresponding test should compare the
full interventional action distributions.

\subtopic{52.8 Goal-persistence certificate}

A minimal goal-persistence report contains:
\begin{enumerate}
\item the goal code and its declared trajectory-level semantics;
\item the feasible goal intervention and its manipulation check;
\item the matched or neutralized future cue trajectory;
\item \(\eta_G(\tau)\) across a declared horizon grid;
\item \(\eta_{G\to A}(\tau)\) over the same grid;
\item cue dependence \(\delta_G^{\mathrm{cue}}(\tau)\);
\item goal-channel ablation and permutation controls;
\item coalition reconstruction risks \(\Risk_G(S)\); and
\item uncertainty intervals over environments, episodes, and random seeds.
\end{enumerate}
A repeated label with no intervention effect on future action fails this
certificate even when its observational autocorrelation is one.

\topic{53}{Counterfactual Policy, Repair, and Agency Certification}

\subtopic{53.1 Feasible action interventions and value functions}

Let \(\mathcal A_t(w,e)\) denote the actions that can be implemented from
workspace state \(w\) in environment \(e\). The interventional action value is
\[
Q_t^g(a\mid w,e)
=
\mathbb E^{\iota_a}
\left[
U_g(Y_t^+)
\mid
W_t=w,
E_t=e,
G_t=g
\right].
\]
The optimal action set is
\[
\mathcal A_t^*(w,g,e)
=
\argmax_{a\in\mathcal A_t(w,e)}Q_t^g(a\mid w,e).
\]
A counterfactual evaluator is an internal state
\[
\widehat Q_t
=
\left(
\widehat Q_t(a)
\right)_{a\in\mathcal A_t(w,e)}
\]
that estimates these intervention values. It is not enough for
\(\widehat Q_t\) to predict the action selected by the system. It must be
validated against externally estimated intervention outcomes.

\subtopic{53.2 Counterfactual-use criterion}

Three properties must be separated:
\begin{description}
\item[Counterfactual accuracy.]
\[
\delta_Q
=
\sup_{a\in\mathcal A_t}
|\widehat Q_t(a)-Q_t^g(a\mid w,e)|.
\]

\item[Policy use.]
A feasible intervention that changes \(\widehat Q_t(a)\) while preserving the
action affordances changes the action distribution in the direction predicted
by the decision rule.

\item[Outcome validity.]
The resulting action change alters \(U_g(Y_t^+)\) as predicted by the external
interventional value estimate.
\end{description}

A system can satisfy any two while failing the third. A learned critic can be
accurate but ignored. A policy can respond to a manipulated score that does
not predict outcomes. A successful action can arise from habit or replay
without any counterfactual evaluator.

\subtopic{53.3 \(T_{20}\): Approximate counterfactual policy theorem }

\begin{theorem}[Approximate counterfactual policy selection]
\label{thm:counterfactual-policy}
Let \(\mathcal A\) be a finite feasible action set. Let \(Q:\mathcal A\to
\mathbb R\) be the true interventional value and let \(\widehat Q\) satisfy
\[
\sup_{a\in\mathcal A}|\widehat Q(a)-Q(a)|\leq\varepsilon_Q.
\]
Suppose a policy \(\widehat\pi\) is \(\delta_{\mathrm{opt}}\)-optimal for
\(\widehat Q\):
\[
\max_{\pi}
\mathbb E_{a\sim\pi}[\widehat Q(a)]
-
\mathbb E_{a\sim\widehat\pi}[\widehat Q(a)]
\leq
\delta_{\mathrm{opt}}.
\]
Then, for every true-value-optimal policy \(\pi^*\),
\[
\boxed{
\mathbb E_{a\sim\pi^*}[Q(a)]
-
\mathbb E_{a\sim\widehat\pi}[Q(a)]
\leq
2\varepsilon_Q+\delta_{\mathrm{opt}}.
}
\]
For a deterministic greedy action \(\widehat a\in\argmax_a\widehat Q(a)\),
\[
Q(a^*)-Q(\widehat a)
\leq2\varepsilon_Q.
\]
If the true optimal action \(a^*\) is unique and its value gap
\[
\Delta_Q
=
Q(a^*)-
\max_{a\neq a^*}Q(a)
\]
satisfies \(\Delta_Q>2\varepsilon_Q\), then every greedy maximizer of
\(\widehat Q\) equals \(a^*\).
\end{theorem}

\begin{proof}
For every policy \(\pi\), uniform approximation gives
\[
\left|
\mathbb E_{a\sim\pi}[\widehat Q(a)]
-
\mathbb E_{a\sim\pi}[Q(a)]
\right|
\leq\varepsilon_Q.
\]
Therefore
\[
\begin{aligned}
\mathbb E_{\pi^*}[Q]
-
\mathbb E_{\widehat\pi}[Q]
&\leq
\mathbb E_{\pi^*}[\widehat Q]
-
\mathbb E_{\widehat\pi}[\widehat Q]
+2\varepsilon_Q\\
&\leq
\delta_{\mathrm{opt}}+2\varepsilon_Q.
\end{aligned}
\]
The deterministic bound is the special case of point-mass policies. If
\(\widehat a\neq a^*\), then
\[
\widehat Q(a^*)
\geq
Q(a^*)-\varepsilon_Q
>
Q(\widehat a)+\varepsilon_Q
\geq
\widehat Q(\widehat a),
\]
contradicting greedy optimality.
\end{proof}

This theorem converts counterfactual-model accuracy into a decision guarantee.
It does not establish that the model is internally represented or causally
used. Those are separate intervention gates.

\subtopic{53.4 Counterfactual mediation and the replay null}

Let \(V_t\) denote the internal value channel. A counterfactual-mediation test
compares:
\[
G_t
\longrightarrow
V_t
\longrightarrow
A_t
\longrightarrow
Y_t^+.
\]
The relevant quantities include:
\[
\eta_{G\to V},
\qquad
\eta_{V\to A},
\qquad
\eta_{A\to Y},
\qquad
\delta_{G\to A\mid V},
\]
where the final term measures the residual goal-to-action effect after the
value channel is blocked or randomized.

The action-replay null preserves the empirical action sequence or action
marginal while destroying online dependence on \(V_t\). If replayed actions
perform equally well under held-out perturbations, the evidence supports a
cached or open-loop policy rather than online counterfactual evaluation. A
second null permutes value labels among feasible actions while preserving the
multiset of scores. A genuine value-mediated policy should follow the
permuted ranking, while an ignored critic should not.

\subtopic{53.5 Workspace integrity vector and maintenance margin}

Choose a declared set \(\mathcal J_{\mathrm{maint}}\) of workspace gates that
must be protected. Write each gate as a signed margin
\[
\phi_{j,t}
\geq0
\quad\Longleftrightarrow\quad
\text{gate }j\text{ passes at time }t.
\]
Examples include
\[
\begin{aligned}
\phi_{\mathrm{gap},t}
&=g_{\varepsilon,t}-\underline g,\\
\phi_{\mathrm{coh},t}
&=\overline\delta_{\mathrm{coh}}-\delta_{\mathrm{coh},t},\\
\phi_{\mathrm{rec},t}
&=\alpha_t^{(k)}-\underline\alpha,\\
\phi_{\mathrm{keep},t}
&=\mu_t^{\mathrm{keep}}-\underline\mu,\\
\phi_{\mathrm{causal},t}
&=\eta_{\mathrm{causal},t}-\underline\eta,\\
\phi_{\mathrm{boundary},t}
&=\Delta_t^B-\underline\Delta_B.
\end{aligned}
\]
The aggregate integrity margin is
\[
\boxed{
v_t
=
\min_{j\in\mathcal J_{\mathrm{maint}}}\phi_{j,t}.
}
\]
The workspace passes the protected gates when \(v_t\geq0\). To maintain a
reserve \(\rho>0\), define the maintenance deficit
\[
\boxed{
e_t
=(\rho-v_t)_+.
}
\]
Define the zero-reserve and bounded-deficit viability sets
\[
\boxed{
\mathcal V_\Theta
=
\left\{
S:\phi_j(S)\geq0\ \text{for every }j\in\mathcal J_{\mathrm{maint}}
\right\}
}
\]
\[
\boxed{
\mathcal V_{\Theta,e_{\max}}
=
\left\{
S:e(S)\leq e_{\max}
\right\}
}
\]
The first set requires every protected gate to pass. The second permits a
declared finite reserve deficit.

This minimum-gate construction is intentionally noncompensatory. A large
spectral gap cannot compensate for failed reconstruction, and high task reward
cannot compensate for loss of the declared workspace boundary.

\subtopic{53.6 Damage and repair interventions}

A damage intervention \(d\in\mathcal I_D\) changes a mechanism that supports
one or more protected gates. Examples include:
\begin{itemize}
\item deleting or corrupting a transfer edge;
\item removing a participant or artifact;
\item perturbing the internal boundary mask;
\item scrambling a role-binding or routing channel;
\item erasing a handoff representation; or
\item degrading a reconstruction map.
\end{itemize}

A repair intervention \(r\in\mathcal I_R\) can add, remove, retrain, rewrite,
or reroute an admissible component. Let \(r_0\) denote matched no-repair. The
post-repair integrity value at horizon \(\tau\) is
\[
V_t^{\mathrm{rep}}(r;d)
=
\mathbb E
\left[
v_{t+\tau}
\mid
\operatorname{do}(D_t=d),
\operatorname{do}(R_t=r),
\mathbf H_t
\right].
\]
The repair gain is
\[
\Delta_t^{\mathrm{rep}}(r;d)
=
V_t^{\mathrm{rep}}(r;d)
-
V_t^{\mathrm{rep}}(r_0;d).
\]
Repair costs are handled either by restricting the feasible repair family or
by reporting a separate cost-effect curve. They should not be hidden inside
an undocumented scalar reward.

\subtopic{53.7 Damage detection, repair mediation, and specificity}

Let \(M_t\) be an internal integrity monitor and \(R_t\) the selected repair
action. A self-maintenance claim requires the causal chain
\[
D_t
\longrightarrow
M_t
\longrightarrow
R_t
\longrightarrow
v_{t+\tau}.
\]
The required tests are:
\begin{description}
\item[Detection.]
Matched damage interventions change \(M_t\) in a gate-specific direction.

\item[Repair mediation.]
Intervening on \(M_t\) changes the repair distribution, and blocking \(M_t\)
reduces the damage-to-repair effect.

\item[Restoration.]
The selected repair improves \(v_{t+\tau}\) relative to matched no-repair.

\item[Specificity.]
A cost-matched generic action or an off-workspace repair does not produce the
same restoration.

\item[Identity preservation.]
The restored system remains within the declared identity tolerance relative to
the pre-damage workspace.
\end{description}

Define the detection effect \(\eta_{D\to M}\), mediation effect
\(\eta_{M\to R}\), restoration effect \(\eta_{R\to v}\), and repair
specificity
\[
\Delta_{\mathrm{repair}}^{\mathrm{spec}}
=
\eta_{R\to v}
-
\eta_{R_\perp\to v},
\]
where \(R_\perp\) is a matched action outside the proposed repair channel.

\subtopic{53.8 \(T_{21}\): Repair-selection margin theorem }

\begin{theorem}[Repair selection under value-model error]
\label{thm:repair-selection}
Fix a damage intervention \(d\) and a finite feasible repair set
\(\mathcal R\). Let
\[
V(r)=V_t^{\mathrm{rep}}(r;d)
\]
be the true expected post-repair integrity margin and let \(\widehat V\) obey
\[
\sup_{r\in\mathcal R}|\widehat V(r)-V(r)|
\leq
\varepsilon_R.
\]
Let \(\widehat r\in\argmax_r\widehat V(r)\) and
\(r^*\in\argmax_r V(r)\). Then
\[
\boxed{
V(r^*)-V(\widehat r)
\leq
2\varepsilon_R.
}
\]
If
\[
V(r^*)
\geq
\rho+2\varepsilon_R,
\]
then the selected repair restores the expected reserve:
\[
\boxed{V(\widehat r)\geq\rho.}
\]
If \(r^*\) is unique and its repair-value gap exceeds
\(2\varepsilon_R\), then \(\widehat r=r^*\).
\end{theorem}

\begin{proof}
Greedy optimality and uniform approximation give
\[
\begin{aligned}
V(r^*)
&\leq
\widehat V(r^*)+\varepsilon_R\\
&\leq
\widehat V(\widehat r)+\varepsilon_R\\
&\leq
V(\widehat r)+2\varepsilon_R.
\end{aligned}
\]
The reserve and exact-selection statements follow immediately.
\end{proof}

The theorem is conditional on a valid repair-value model. It cannot convert a
post hoc repair narrative into evidence that the system detected damage or
used the model. Detection and value-channel interventions remain necessary.

\subtopic{53.9 Repair, replacement, and successor workspaces}

Restoration of performance is not automatically restoration of identity. Let
\(W_{t^-}\) be the last certified pre-damage workspace and let
\(W_{t+\tau}^{r}\) be the post-repair workspace. Define the repair identity
defect
\[
\delta_{\mathrm{rep-id}}(r)
=
 d_{\mathrm{id}}
\left(
W_{t^-},W_{t+\tau}^{r}
\right).
\]
A repair action is identity preserving at tolerance
\(\varepsilon_{\mathrm{id}}\) when
\[
\Delta_t^{\mathrm{rep}}(r;d)>0,
\qquad
\delta_{\mathrm{rep-id}}(r)
\leq
\varepsilon_{\mathrm{id}}.
\]
When integrity is restored but the identity defect exceeds tolerance, the
result is better described as a successor or replacement workspace. This
terminological distinction prevents any high-performing reconstruction from
being counted as persistence of the same relational organization.

\subtopic{53.10 \(T_{22}\): Closed-loop self-maintenance theorem }

\begin{theorem}[Causal repair stabilizes maintenance deficit]
\label{thm:self-maintenance}
Let \(e_t\geq0\) be the maintenance deficit. Suppose the matched no-repair
trajectory satisfies
\[
e_{t+1}^{(0)}
\leq
\alpha e_t+d_t,
\]
where \(d_t\geq0\) is exogenous damage. Suppose the implemented repair policy
has the causal one-step effect
\[
e_{t+1}^{(0)}-e_{t+1}
\geq
\kappa e_t-\epsilon_t.
\]
Assume
\[
0\leq\beta:=\alpha-\kappa<1,
\qquad
0\leq d_t\leq\overline d,
\qquad
0\leq\epsilon_t\leq\overline\epsilon.
\]
Then
\[
\boxed{
e_t
\leq
\beta^te_0
+
\frac{(\overline d+\overline\epsilon)(1-\beta^t)}{1-\beta},
}
\]
and therefore
\[
\boxed{
\limsup_{t\to\infty}e_t
\leq
\frac{\overline d+\overline\epsilon}{1-\beta}.
}
\]
If the admissible deficit is \(e_{\max}\) and
\[
e_0\leq e_{\max},
\qquad
\overline d+\overline\epsilon
\leq
(1-\beta)e_{\max},
\]
then the interval \([0,e_{\max}]\) is forward invariant.
\end{theorem}

\begin{proof}
Combine the no-repair bound with the causal repair effect:
\[
\begin{aligned}
e_{t+1}
&\leq
 e_{t+1}^{(0)}-\kappa e_t+\epsilon_t\\
&\leq
(\alpha-\kappa)e_t+d_t+\epsilon_t\\
&\leq
\beta e_t+\overline d+\overline\epsilon.
\end{aligned}
\]
Iterating this scalar recurrence gives the stated geometric bound. If
\(e_t\leq e_{\max}\), the forward-invariance condition gives
\[
e_{t+1}
\leq
\beta e_{\max}
+(1-\beta)e_{\max}
=e_{\max}.
\]
\end{proof}

Theorem~\ref{thm:self-maintenance} differs from the persistence theorem in
Part~X. There, contraction was assumed for the workspace error dynamics. Here
the improved contraction factor is attributed to a measured repair action
through a matched no-repair counterfactual. Without that causal comparison,
the same bounded trajectory could be passive robustness rather than
self-maintenance.

\subtopic{53.11 A scalar boundary-repair example}

Let \(b_t\in[0,1]\) be a boundary-integrity score and let the reserve be
\(\rho\). Define
\[
e_t=(\rho-b_t)_+.
\]
Suppose uncorrected damage obeys
\[
e_{t+1}^{(0)}
\leq
1.08e_t+d_t,
\]
so the open-loop boundary is locally unstable. A repair policy with measured
causal effect
\[
e_{t+1}^{(0)}-e_{t+1}
\geq
0.25e_t-\epsilon_t
\]
gives
\[
\beta=1.08-0.25=0.83.
\]
If
\[
\overline d=0.01,
\qquad
\overline\epsilon=0.005,
\]
then
\[
\limsup e_t
\leq
\frac{0.015}{1-0.83}
\approx0.0883.
\]
Thus an open-loop unstable boundary can be maintained within a finite deficit
band by a sufficiently strong repair channel. The example establishes only a
control calculation. A complete self-maintenance claim must also show that an
internal monitor detects damage, mediates the repair choice, and preserves
workspace identity.

\subtopic{53.12 Complete self-maintenance certificate}

A self-maintaining workspace must pass all of the following gates:
\begin{enumerate}
\item \emph{Protected-gate declaration.} The maintained integrity vector and
 reserve are fixed before evaluation.
\item \emph{Damage validity.} Each damage intervention measurably degrades at
 least one protected gate.
\item \emph{Internal detection.} The monitor changes under damage and is not
 merely an external evaluator.
\item \emph{Repair mediation.} Monitor interventions change repair selection.
\item \emph{Repair efficacy.} Repair improves the held-out integrity margin
 relative to matched no-repair.
\item \emph{Specificity.} Cost-matched generic actions and off-workspace
 interventions do not reproduce the gain.
\item \emph{Repeated-disturbance stability.} The deficit remains within the
 declared bound over a held-out disturbance sequence.
\item \emph{Identity control.} The restored state remains within the declared
 identity tolerance or is reported as a successor workspace.
\item \emph{Anti-centralization.} No single node contains a sufficient copy
 of the monitor, repair model, and protected state when the claim is
 workspace-level maintenance.
\end{enumerate}

\subtopic{53.13 Certified workspace agent}

Let
\[
\Theta_A
=
\left(
\tau_G,
\tau_{G\to A},
\varepsilon_Q,
\tau_{V\to A},
\tau_{A\to Y},
\tau_{D\to M},
\tau_{M\to R},
\tau_{R\to v},
\tau_{\mathrm{spec}},
\varepsilon_{\mathrm{ctl}},
 e_{\max}
\right).
\]
A \(\Theta_A\)-certified workspace agent is a controlled realization for
which:
\begin{enumerate}
\item \(\SRW(W_t;\varsigma)=1\);
\item cue-controlled goal retention exceeds \(\tau_G\);
\item goal-to-action mediation exceeds \(\tau_{G\to A}\);
\item counterfactual-value error is at most \(\varepsilon_Q\);
\item value-to-action use exceeds \(\tau_{V\to A}\);
\item action-to-outcome control exceeds \(\tau_{A\to Y}\);
\item damage detection, monitor-to-repair mediation, and restoration exceed
 their thresholds;
\item repair specificity exceeds \(\tau_{\mathrm{spec}}\);
\item \(k_{\mathrm{ctl},\varepsilon_{\mathrm{ctl}}}\geq2\); and
\item repeated maintenance deficit remains at most \(e_{\max}\).
\end{enumerate}

As throughout the theory, this predicate is relative to the declared node
decomposition, goals, actions, environment, horizons, interventions, and null
families.

\subtopic{53.14 \(T_{23}\): Workspace-agency synthesis theorem }

\begin{theorem}[Workspace-level agency under declared gates]
\label{thm:agency-synthesis}
Assume a controlled realization satisfies:
\begin{enumerate}
\item the operational relationality assumptions of
 Theorem~\ref{thm:operational-relationality};
\item cue-controlled goal retention and action readout as in
 Theorem~\ref{thm:goal-retention} and
 Corollary~\ref{cor:goal-action-readout};
\item a counterfactual evaluator with uniform error \(\varepsilon_Q\) and
 optimization error \(\delta_{\mathrm{opt}}\);
\item feasible interventions establishing positive value-to-action and
 action-to-outcome effects;
\item the causal self-maintenance assumptions of
 Theorem~\ref{thm:self-maintenance}; and
\item \(k_{\mathrm{ctl},\varepsilon}\geq2\), together with a matched local
 policy-null panel that preserves every single-node predictive marginal
 but loses full-system control value.
\end{enumerate}
Then, relative to the declared specification:
\begin{enumerate}
\item the expected one-step decision regret is at most
 \(2\varepsilon_Q+\delta_{\mathrm{opt}}\);
\item the long-run maintenance deficit is bounded by
 \[
 \frac{\overline d+\overline\epsilon}{1-\beta};
 \]
\item no single-node statistic is sufficient for the complete control state;
\item matched local policy nulls cannot reproduce the intact control value;
\item goal-dependent action selection and workspace repair are causally
 mediated by implemented internal channels.
\end{enumerate}
Hence the realization satisfies the formal predicate
\[
\operatorname{Agent}_{\Theta_A}(\Workspace)=1.
\]
\end{theorem}

\begin{proof}
The regret bound follows from
Theorem~\ref{thm:counterfactual-policy}. The maintenance bound follows from
Theorem~\ref{thm:self-maintenance}. The coalition condition excludes
single-node sufficiency for the complete control state. The matched-null gap
excludes explanations based only on preserved local predictive marginals. The
implemented goal, value, action, monitor, and repair interventions establish
the required causal mediation. These are precisely the gates in the stated
agency predicate.
\end{proof}

The synthesis theorem is intentionally relative and conjunctive. Its force is
that decision quality, maintenance stability, causal mediation, and relational
location are certified together. It does not infer an unrestricted agent
outside the tested environment.

\subtopic{53.15 Agency, autonomy, and consciousness}

The following implications are not licensed:
\[
\operatorname{Agent}_{\Theta_A}(\Workspace)
\centernot\Longrightarrow
\operatorname{PhenomenallyConscious}(\Workspace),
\]
\[
\operatorname{SelfMaintaining}(\Workspace)
\centernot\Longrightarrow
\operatorname{SelfAware}(\Workspace),
\]
\[
\operatorname{GoalPersistent}(\Workspace)
\centernot\Longrightarrow
\operatorname{MoralPatient}(\Workspace).
\]
A thermostat-like controller can maintain a variable without representing
itself. A sophisticated planner can evaluate counterfactual actions without
phenomenal experience. Conversely, failure to pass this engineering
certificate would not prove absence of consciousness. Agency and
consciousness remain distinct axes of inquiry.


\theorypart{XVII}{Decomposition, Coarse-Graining, and Robust Relationality}

Relationality is always assessed relative to a way of dividing a realization
into components. That dependence is not a minor technicality. If every
participant, artifact, and channel is merged into one component, the complete
workspace becomes local by definition. If a centralized state is recoded into
several complementary shares, the same centralized process can be made to look
irreducibly distributed. A defensible theory must therefore say which
partitions count as admissible, which quantities survive faithful aggregation,
and how sensitive a certificate is to alternative descriptions of the same
system.

\topic{54}{The Decomposition Problem}

Let \(\Omega\) be a finite set of primitive carriers at a declared audit
resolution. A primitive carrier may be a person, a model instance, a process,
a memory store, an instrument, or another independently tracked mechanism.
The word \emph{primitive} is relative to the audit. It does not assert a final
metaphysical atomization of the system.

\begin{definition}[System decomposition]\label{def:system-decomposition}
A decomposition of a carrier presentation is a tuple
\[
\mathcal D
=
\bigl(
\Pi,\mathbf O,G_{\Pi},\mathcal I_{\Pi}
\bigr),
\]
where:
\begin{itemize}
\item \(\Pi=\{B_\alpha\}_{\alpha\in A}\) is a partition of \(\Omega\);
\item \(O_\alpha\) is the declared local-access map for block \(B_\alpha\),
 producing the local state or history available at that component;
\item \(G_{\Pi}\) records the declared communication, transfer, and causal
 interfaces among the blocks; and
\item \(\mathcal I_{\Pi}\) is the family of feasible interventions at that
 resolution.
\end{itemize}
The accessible information of component \(\alpha\) is the sigma-field
\[
\mathscr F_\alpha^{\mathcal D}
=
\sigma\!\left(O_\alpha(\mathbf H_t)\right).
\]
Every shared bus, common memory, supervisor, or artifact that can alter more
than one local state must appear either as its own carrier or as an explicit
channel in \(G_{\Pi}\).
\end{definition}

The same physical realization can support several decompositions. One may
separate a model from its executor, merge them into one computational unit, or
treat a persistent artifact as its own component. These alternatives are not
mere renamings when they change local access or feasible interventions. Every certificate should therefore display the full specification
\[
\SRW(W_t;\varsigma),
\qquad
\varsigma=(\mathcal D,Y^+,\ell,\tau,\mathcal I,\mathcal N,\Theta).
\]
This part varies the decomposition coordinate while holding the remaining
coordinates fixed. Part~XVIII treats the full specification family.

\begin{definition}[Refinement and coarsening]\label{def:decomposition-refinement}
Let \(\mathcal D_f\) and \(\mathcal D_c\) be decompositions of the same carrier
presentation. The fine decomposition refines the coarse decomposition,
written
\[
\mathcal D_f\preceq\mathcal D_c,
\]
when there is a surjection
\[
q:V_f\longrightarrow V_c
\]
such that each coarse component is the union of the fine components in one
fiber of \(q\), and its accessible information is generated only from those
fine components. The aggregation is \emph{information faithful} when
\[
\mathscr F_B^{\mathcal D_c}
=
\bigvee_{i\in q^{-1}(B)}
\mathscr F_i^{\mathcal D_f}
\qquad(B\in V_c).
\]
It is \emph{structurally faithful} when it also retains every internal
compatibility condition and every boundary transfer condition needed to
recover the original global-section problem.
\end{definition}

Refinement and coarsening are directional. A coarse component may contain
information that was spread across several finer nodes. Conversely, a fine
description may expose internal channels that disappear when a macro-node is
formed. The next results identify what is preserved when the aggregation is
faithful.

\topic{55}{Faithful Structural Coarse-Graining}

Let the vertex set of a core-transfer diagram be partitioned as
\[
V=V_1\sqcup\cdots\sqcup V_m.
\]
Write the transfer discrepancy as
\[
D_H
=
\begin{pmatrix}
D_{\mathrm{int}}\\
D_{\partial}
\end{pmatrix},
\qquad
D_{\mathrm{int}}
=
\bigoplus_{\alpha=1}^{m}D_\alpha,
\]
where \(D_\alpha\) contains the edge constraints internal to \(V_\alpha\),
and \(D_{\partial}\) contains the constraints crossing between blocks. Define
the internally coherent macro-state spaces
\[
W_\alpha=\ker D_\alpha,
\qquad
W_{\Pi}=\bigoplus_{\alpha=1}^{m}W_\alpha.
\]
Let
\[
J_\Pi:W_\Pi\hookrightarrow\mathcal C_V
\]
be the canonical isometric inclusion, and define the quotient discrepancy
\[
D_\Pi=D_{\partial}J_\Pi.
\]
The quotient does not replace a group of nodes by an arbitrary summary. Its
macro-state is the complete internal section space of that group, and its
boundary operator retains all cross-block compatibility equations.

\begin{theorem}[Faithful coarse-graining]\label{thm:faithful-coarse-graining}
For every finite partition of the vertex set,
\[
\boxed{
\ker D_H
=
J_\Pi\ker D_\Pi.
}
\]
Equivalently,
\[
\boxed{
\Gamma(H)
\cong
\ker D_\Pi
=
\ker(D_\Pi^*D_\Pi).
}
\]
In particular,
\[
\dim\Gamma(H)
=
\dim W_\Pi-\operatorname{rank}(D_\Pi).
\]
Thus faithful coarse-graining preserves the exact structural workspace even
though it changes the number and interpretation of the displayed nodes.
\end{theorem}

\begin{proof}
A vector \(c\in\mathcal C_V\) lies in \(\ker D_H\) exactly when
\[
D_{\mathrm{int}}c=0
\qquad\text{and}\qquad
D_{\partial}c=0.
\]
The first condition is equivalent to \(c=J_\Pi w\) for a unique
\(w\in W_\Pi\). Substitution into the second condition gives
\[
D_\Pi w
=
D_{\partial}J_\Pi w
=0.
\]
Hence \(c\in\ker D_H\) if and only if
\(c=J_\Pi w\) with \(w\in\ker D_\Pi\). The Laplacian identity follows from
\(\ker(D_\Pi^*D_\Pi)=\ker D_\Pi\), and the dimension formula is rank-nullity.
\end{proof}

The theorem supplies a precise sense in which a workspace may be represented
at several resolutions without changing its exact core. The preservation
fails when the macro-state discards internal section coordinates or when
boundary equations are omitted.

\begin{corollary}[Two-scale coherence bound]\label{cor:two-scale-decomposition}
Let \(P_\Pi\) be the orthogonal projector onto \(J_\Pi W_\Pi\). Let
\(\gamma_{\mathrm{int}}>0\) be the smallest positive eigenvalue of
\(D_{\mathrm{int}}^*D_{\mathrm{int}}\), and let
\(\gamma_\Pi>0\) be the smallest positive eigenvalue of
\(D_\Pi^*D_\Pi\). Then every \(c\in\mathcal C_V\) satisfies
\[
\boxed{
\operatorname{dist}(c,\Gamma(H))
\leq
\frac{\|D_{\mathrm{int}}c\|}{\sqrt{\gamma_{\mathrm{int}}}}
+
\frac{
\|D_\Pi J_\Pi^*P_\Pi c\|
}{\sqrt{\gamma_\Pi}}.
}
\]
The first term measures incoherence within macro-components. The second
measures incompatibility among their internally coherent states.
\end{corollary}

\begin{proof}
By the triangle inequality,
\[
\operatorname{dist}(c,\Gamma(H))
\leq
\|c-P_\Pi c\|
+
\operatorname{dist}(P_\Pi c,\Gamma(H)).
\]
The workspace Poincar\'e estimate applied to \(D_{\mathrm{int}}\) bounds the
first term. Write \(P_\Pi c=J_\Pi w\). Theorem~\ref{thm:faithful-coarse-graining}
and the corresponding estimate for \(D_\Pi\) bound the second term by
\(\|D_\Pi w\|/\sqrt{\gamma_\Pi}\). Since \(J_\Pi\) is isometric,
\(w=J_\Pi^*P_\Pi c\).
\end{proof}

This decomposition is useful computationally. Internal blocks can be solved
or compressed independently, after which only the boundary discrepancy must
be assembled globally. It also gives a diagnostic distinction between a
workspace that fails inside its local components and one whose local
components are coherent but mutually incompatible.

\topic{56}{Predictive Risk Under Aggregation}

The structural core can be preserved by faithful coarse-graining, while
coalition counts generally are not. Let
\[
\mathscr F_S^{\mathcal D}
=
\bigvee_{i\in S}\mathscr F_i^{\mathcal D}
\]
be the information available to coalition \(S\) under decomposition
\(\mathcal D\). For a declared loss \(\ell\), write
\[
\Risk_{\mathcal D}(S)
=
\inf_f
\mathbb E\!\left[
\ell\bigl(Y^+,f(\mathscr F_S^{\mathcal D})\bigr)
\right].
\]

\begin{theorem}[Bayes-risk transport under faithful aggregation]
\label{thm:risk-aggregation}
Suppose \(\mathcal D_f\preceq\mathcal D_c\) through
\(q:V_f\to V_c\), and suppose the aggregation is information faithful. Then
for every coarse coalition \(S\subseteq V_c\),
\[
\boxed{
\Risk_{\mathcal D_c}(S)
=
\Risk_{\mathcal D_f}\bigl(q^{-1}(S)\bigr).
}
\]
Consequently,
\[
k_\varepsilon(\mathcal D_c)
\leq
k_\varepsilon(\mathcal D_f),
\]
where coalition size is counted in displayed nodes.
\end{theorem}

\begin{proof}
Information faithfulness gives
\[
\mathscr F_S^{\mathcal D_c}
=
\bigvee_{B\in S}
\bigvee_{i\in q^{-1}(B)}
\mathscr F_i^{\mathcal D_f}
=
\mathscr F_{q^{-1}(S)}^{\mathcal D_f}.
\]
The two Bayes risks optimize over the same measurable decision rules and are
therefore equal. If a fine coalition of size \(k\) is sufficient within
\(\varepsilon\), the set of coarse blocks meeting that coalition has at most
\(k\) elements and contains at least the same information. This proves the
coalition-size inequality.
\end{proof}

The theorem makes the source of a common error explicit. A minimal coalition
size of two at one resolution can become one after those two nodes are merged.
Node count is therefore a decomposition-relative statistic.

A more stable quantity measures primitive carrier cost. Let
\(w_\omega>0\) be a declared cost for each \(\omega\in\Omega\), and assign a
block the cost
\[
c(B)=\sum_{\omega\in B}w_\omega.
\]
Define the sufficient support cost
\[
\boxed{
s_\varepsilon(\mathcal D)
=
\min_{S\subseteq V_{\mathcal D}}
\left\{
\sum_{B\in S}c(B):
\Risk_{\mathcal D}(S)
\leq
\Risk_{\mathcal D}(V_{\mathcal D})+\varepsilon
\right\}.
}
\]
If \(\mathcal D_f\preceq\mathcal D_c\) is information faithful, then
\[
\boxed{
s_\varepsilon(\mathcal D_f)
\leq
s_\varepsilon(\mathcal D_c).}
\]
The fine decomposition permits every union of coarse blocks and may also
permit cheaper sub-block coalitions. Equality holds whenever an optimal fine
coalition is a union of complete coarse fibers. This weighted support measure
does not remove the need to choose primitive carriers, but it prevents a
large macro-component from being counted as a single inexpensive unit merely
because it was assigned one node label.

\topic{57}{Trivialization by Arbitrary Redescription}

The dependence on decomposition is unavoidable. Two elementary constructions
show why a certificate cannot be required to survive every mathematically
possible partition.

\begin{proposition}[Coarsening and secret-sharing trivializations]
\label{prop:decomposition-trivialization}
The following statements hold.
\begin{enumerate}
\item Any finite realization can be coarsened to a single displayed component.
 Relative to that decomposition, the full-system state is local and
 \(k_0\leq1\).
\item Let \(Z\) be uniform on a finite abelian group \(\mathsf G\). Let
 \(K\) be an independent uniform group element and define
 \[
 X_1=K,
 \qquad
 X_2=Z-K.
 \]
 Then
 \[
 I(Z;X_1)=I(Z;X_2)=0,
 \qquad
 H(Z\mid X_1,X_2)=0.
 \]
 Under logarithmic loss,
 \[
 \Risk(\{1\})=\Risk(\{2\})=\log|\mathsf G|,
 \qquad
 \Risk(\{1,2\})=0,
 \]
 so the declared two-node representation has \(k_0=2\) and joint advantage
 \(\log|\mathsf G|\), even when a single centralized encoder generated both
 shares.
\end{enumerate}
\end{proposition}

\begin{proof}
The first statement follows by taking the indiscrete partition
\(\Pi=\{\Omega\}\). For the second, conditioning on any value of \(Z\) leaves
\(K\), and hence \(X_1\), uniform. Translation invariance of the uniform law
also leaves \(X_2=Z-K\) uniform, so each share is independent of \(Z\). The
pair reconstructs \(Z=X_1+X_2\), which gives the entropy and Bayes-risk
identities.
\end{proof}

The secret-sharing example is a valid distributed code. What it does not show
is that the process producing or using the code is causally distributed across
the declared carriers. If one hidden controller creates both shares, stores
the key, or jointly selects their updates, that controller belongs in the
causal inventory. Coalition irreducibility and matched predictive nulls remain
necessary tests, but their interpretation is conditional on a defensible
interface model.

This distinction suggests two levels of claim:
\begin{description}
\item[Encoding-relative relationality.] No declared local view is sufficient.
\item[Carrier-relative relationality.] The declared views also correspond to
admissible, separately tracked causal interfaces with no omitted common
controller.
\end{description}
The stronger claim requires the decomposition audit developed next.

\topic{58}{Admissible Decomposition Families}

No purely formal rule can identify the uniquely correct decomposition of every
physical or computational system. The practical goal is instead to define a
family of decompositions that are defensible for the question being asked and
to report how the certificate changes across that family.

\begin{definition}[Admissible decomposition family]
\label{def:admissible-decomposition-family}
An admissible family \(\mathfrak D\) is a predeclared collection of
decompositions of the same carrier presentation satisfying the following
audit conditions.
\begin{enumerate}
\item \textbf{Carrier continuity.} Each component corresponds to a process,
 artifact, or mechanism that can be tracked through the relevant time scale.
\item \textbf{Access locality.} The local state contains only information
 available through the declared component and its recorded incoming channels.
\item \textbf{Channel completeness.} Shared memories, buses, supervisors,
 caches, and synchronized external cues are represented explicitly.
\item \textbf{Intervention separability.} At least some component mechanisms
 or channels can be altered without directly rewriting the entire system.
\item \textbf{Target independence.} The decomposition and its alternatives
 are fixed without using held-out outcomes to manufacture a favorable result.
\item \textbf{Resolution relevance.} A proposed split is retained only when
 it changes access, intervention, reconstruction, or transfer structure at the
 declared measurement resolution.
\item \textbf{Alternative coverage.} The family contains the principal
 plausible merges, splits, and artifact assignments that a skeptical analyst
 would consider.
\end{enumerate}
\end{definition}

Let \(m_q(W;\mathcal D)\) be the signed margin of gate \(q\), oriented so that
\[
m_q(W;\mathcal D)>0
\]
means that the gate passes after its uncertainty correction. Define the
decomposition envelope
\[
\underline m_q
=
\inf_{\mathcal D\in\mathfrak D}m_q(W;\mathcal D),
\qquad
\overline m_q
=
\sup_{\mathcal D\in\mathfrak D}m_q(W;\mathcal D),
\]
and its sensitivity width
\[
\Delta_q^{\mathrm{dec}}
=
\overline m_q-\underline m_q.
\]

\begin{definition}[Decomposition-robust certificate]
\label{def:decomposition-robust-certificate}
A workspace passes the \(\mathfrak D\)-robust certificate at level \(L\) when
\[
\boxed{
\operatorname{Certify}_{L,\mathfrak D}^{\mathrm{rob}}(W)
=
\bigwedge_{q\in\mathcal J(L)}
\left\{
\inf_{\mathcal D\in\mathfrak D}
 m_q(W;\mathcal D)>0
\right\}.
}
\]
A pass for one selected decomposition is a \emph{decomposition-specific}
claim. A pass throughout \(\mathfrak D\) is a \emph{family-robust} claim. A
claim is \emph{decomposition fragile} when at least one required gate changes
sign inside the admissible family.
\end{definition}

Universal invariance over all partitions is neither possible nor desirable,
as Proposition~\ref{prop:decomposition-trivialization} shows. The family
\(\mathfrak D\) must be broad enough to challenge the favored description and
narrow enough to exclude partitions with no operational meaning.

\topic{59}{Stability Under Decomposition Transport}

Exact faithful aggregation preserves the structural core, but practical
alternative decompositions usually introduce approximation. Their certificate
can still be controlled by transporting a reference workspace and comparing
gate margins.

Fix a reference decomposition \(\mathcal D_0\). For each
\(\mathcal D\in\mathfrak D\), let
\[
\Phi_{\mathcal D}:W^{\mathcal D_0}\longrightarrow W^{\mathcal D}
\]
be a declared workspace correspondence and let
\(\delta_{\mathcal D}\) measure its structural and causal identity defect.
Assume gate \(q\) has a transport bound
\[
\left|
 m_q(\Phi_{\mathcal D}W;\mathcal D)
 -m_q(W;\mathcal D_0)
\right|
\leq
L_q\delta_{\mathcal D}+b_q(\mathcal D),
\]
where \(L_q\) is a sensitivity constant and \(b_q\) records model mismatch not
captured by the identity distance.

\begin{theorem}[Decomposition-transport robustness]
\label{thm:decomposition-robustness}
If, for every required gate \(q\),
\[
\boxed{
 m_q(W;\mathcal D_0)
>
\sup_{\mathcal D\in\mathfrak D}
\left(
L_q\delta_{\mathcal D}+b_q(\mathcal D)
\right),
}
\]
then \(W\) passes the \(\mathfrak D\)-robust certificate. Conversely, if for
some required gate
\[
 m_q(W;\mathcal D_0)
+
\sup_{\mathcal D\in\mathfrak D}
\left(
L_q\delta_{\mathcal D}+b_q(\mathcal D)
\right)
<0,
\]
then that gate fails throughout the family.
\end{theorem}

\begin{proof}
For every \(\mathcal D\in\mathfrak D\), the transport bound gives
\[
 m_q(\Phi_{\mathcal D}W;\mathcal D)
\geq
m_q(W;\mathcal D_0)
-
L_q\delta_{\mathcal D}
-
b_q(\mathcal D).
\]
The first hypothesis makes the right-hand side positive uniformly in
\(\mathcal D\), which proves robust passage of every required gate. The second
statement follows from the corresponding upper bound.
\end{proof}

When \(b_q\equiv0\), a convenient certificate radius is
\[
\boxed{
r_{\mathrm{dec}}(W)
=
\min_{q:L_q>0}
\frac{m_q(W;\mathcal D_0)}{L_q}.
}
\]
Every admissible family satisfying
\[
\sup_{\mathcal D\in\mathfrak D}\delta_{\mathcal D}
<
r_{\mathrm{dec}}(W)
\]
preserves the sign of all required gate margins. In applications, the
transport terms can be bounded using earlier results:
\begin{center}
\small
\begin{tabularx}{\textwidth}{@{}>{\bfseries\raggedright\arraybackslash}p{0.22\textwidth}>{\raggedright\arraybackslash}X@{}}
\toprule
Gate family & Relevant transport control \\
\midrule
Structural coherence & Transfer-Laplacian projector drift and spectral gap. \\
Predictive alignment & Divergence between transported conditional-law states. \\
Causal efficacy & Distance between matched intervention kernels. \\
Relationality & Coalition-risk distortion and matched-null preservation error. \\
Reconstruction & Perturbation of lower frame bounds and handoff error. \\
Persistence and identity & The coordinatewise identity profile from Part~XII. \\
\bottomrule
\end{tabularx}
\end{center}

The theorem is intentionally margin based. A result barely above threshold at
one decomposition should not be described as robust when a small, defensible
change of interface makes it disappear.

\topic{60}{Decomposition Audit and Reporting}

A complete decomposition audit should proceed in the following order.
\begin{enumerate}
\item Inventory the primitive carriers at the chosen temporal and physical
 resolution.
\item Record every shared resource and distinguish local state from incoming
 information.
\item Choose a reference decomposition from independently manipulable and
 temporally trackable interfaces.
\item Construct an admissible family containing the principal plausible
 merges, splits, and artifact assignments.
\item For every alternative, specify the correspondence to the reference
 workspace and validate its manipulation checks.
\item Recompute the structural, predictive, causal, coalition, reconstruction,
 and identity margins.
\item Report the full decomposition envelope, not only the most favorable
 member of the family.
\end{enumerate}

The report should state:
\[
\begin{aligned}
&\text{reference decomposition and audit resolution},\\
&\text{admissible family and reasons for excluded partitions},\\
&\text{hidden-resource and common-controller checks},\\
&\text{faithful-aggregation maps and identity defects},\\
&\text{gatewise minima, maxima, and sensitivity widths},\\
&\text{strongest decomposition-specific and family-robust conclusions}.
\end{aligned}
\]

Four controls are especially important:
\begin{enumerate}
\item merge components connected by an uninterruptible common controller and
 test whether the relational claim collapses;
\item split a known centralized controller into complementary local views and
 verify that the protocol rejects the resulting secret-sharing appearance of
 distributedness;
\item expose or remove shared caches, external cues, and common memory that
 were absent from the favored graph; and
\item repeat the frozen certificate across the complete admissible family.
\end{enumerate}

For the motivating human--AI collaboration, this analysis changes the form of
the empirical question. It is not enough to ask whether the human, model,
executor, and artifact store form a workspace under one natural description.
The stronger question is whether the same operative state and its causal role
survive plausible alternatives, such as merging model and executor, separating
conversation state from persistent artifacts, or making shared orchestration
channels explicit. Those measurements have not yet been performed. The case
therefore remains decomposition-specific and observational.

\theorypart{XVIII}{Specification Robustness, Identifiability, and Statistical Certification}

A workspace certificate is never evaluated in a vacuum. It is evaluated for a
particular future variable, loss, time horizon, intervention family, null
panel, threshold vector, and system decomposition. These choices determine
which distinctions matter and which capacities count as successful. They are
therefore part of the scientific claim rather than clerical details surrounding
it.

This dependence has two consequences. First, a result that passes under one
reasonable analysis can fail under another. Robustness must be assessed across
a declared family of scientifically defensible alternatives. Second, even a
fully specified claim may not be identifiable from the observations and
interventions that are actually available. Distinct realizations can induce the
same experimental law while disagreeing about whether a relational workspace
exists. The theory must say when such disagreement can be resolved, when only
bounds are justified, and how sampling uncertainty compounds these structural
ambiguities.

\topic{61}{The Complete Scientific Specification}

\begin{definition}[Workspace specification]\label{def:workspace-specification}
A workspace specification is a tuple
\[
\boxed{
\varsigma
=
\bigl(
\mathcal D,
Y^+,
\ell,
\tau,
\mathcal I,
\mathcal N,
\Theta
\bigr),
}
\]
where:
\begin{itemize}
\item \(\mathcal D\) is the decomposition and local-access model;
\item \(Y^+\) is the declared future variable or trajectory;
\item \(\ell\) is the normalized loss or utility convention;
\item \(\tau\) is the temporal horizon or horizon family;
\item \(\mathcal I\) is the feasible intervention family;
\item \(\mathcal N\) is the required matched-null family; and
\item \(\Theta\) is the gate schema, direction, threshold, and indifference
 region used for certification.
\end{itemize}
The complete certificate is written
\[
\operatorname{Cert}(W;\varsigma)
=
\bigwedge_{q\in\mathcal J(\varsigma)}
\mathbf 1\{m_q(W;\varsigma)>0\},
\]
where \(m_q\) is the signed margin of gate \(q\). Positive margin means that
that gate passes in the direction declared by \(\Theta\).
\end{definition}

The notation absorbs lower-bounded and upper-bounded gates into one sign
convention. If a causal effect must exceed \(c_0\), then
\[
m_q=\widehat\eta-c_0.
\]
If a defect must remain below \(\varepsilon_0\), then
\[
m_q=\varepsilon_0-\widehat\delta.
\]
The complete certificate therefore passes exactly when every required signed
margin is positive.

The specification does not include arbitrary features of the analysis. It
contains choices that alter the meaning or evidential force of the claim. An
optimization algorithm, numerical tolerance, or estimator belongs in the
estimation protocol unless changing it changes the target estimand or the
admissible interventions. Conversely, the target, horizon, and null family
cannot be treated as implementation details because changing them can reverse
the conclusion without changing the underlying realization.

\subtopic{61.1 Specification relativity is not unrestricted relativism}

A theorem-development workspace may be operationally real for a month-long
research trajectory and irrelevant for a ten-second motor-control target. A
repair channel may preserve theorem identity while failing to preserve social
coordination. These are different scientific claims. The theory does not
require one workspace state to be sufficient for every possible future or loss.
It requires the intended scope to be declared before certification.

This yields the more precise notation
\[
\boxed{
\SRW(W;\varsigma)=1.
}
\]
The earlier expression \(\SRW_\Theta(W;\mathcal D)\) is the special case in
which the remaining coordinates of \(\varsigma\) are held fixed.

\subtopic{61.2 Inert reparameterization and substantive change}

Some changes of specification merely recode the same claim. Others change its
content. The distinction is analogous to gauge freedom in the structural
theory.

\begin{definition}[Specification equivalence]\label{def:spec-equivalence}
Two specifications \(\varsigma\) and \(\varsigma'\) are equivalent, written
\(\varsigma\simeq_{\mathrm{spec}}\varsigma'\), when there are bijective
transport maps that:
\begin{enumerate}
\item preserve the sigma-field generated by the future target;
\item preserve the ordering and normalization of decisions under the loss;
\item map feasible interventions to physically equivalent interventions;
\item map every required null to a null preserving the same quantities; and
\item carry every gate threshold and margin into its counterpart.
\end{enumerate}
\end{definition}

\begin{proposition}[Certificate invariance under specification equivalence]
\label{prop:spec-equivalence-invariance}
If \(\varsigma\simeq_{\mathrm{spec}}\varsigma'\) and \(\Phi\) is the induced
workspace correspondence, then
\[
\operatorname{Cert}(W;\varsigma)
=
\operatorname{Cert}(\Phi W;\varsigma').
\]
\end{proposition}

\begin{proof}
Every required gate is carried to a gate with the same signed margin. The
conjunction of positive margins is therefore unchanged.
\end{proof}

A bijective relabeling of target classes with the correspondingly transported
loss is inert. Replacing exact theorem validity with stylistic similarity is
not. Changing a one-step horizon to a year-long persistence horizon is not.
Removing the null that tests hidden centralization is not. Robustness concerns
substantive alternatives after equivalent recodings have been quotiented out.

\subtopic{61.3 Admissible specification families}

\begin{definition}[Admissible specification family]
\label{def:admissible-spec-family}
An admissible specification family \(\mathfrak S\) is a predeclared set of
specifications that:
\begin{enumerate}
\item refer to the same scientific claim type and carrier presentation;
\item use targets that are equivalent operationalizations or explicitly
 declared nearby stress targets;
\item normalize losses so that margins are comparable;
\item use horizons justified by the process being studied;
\item include every feasible intervention and null needed to distinguish the
 principal alternative explanations;
\item use thresholds fixed without access to certification outcomes; and
\item exclude only alternatives for reasons stated before evaluation.
\end{enumerate}
\end{definition}

The family should not contain every imaginable analysis. Such a requirement
would erase meaningful domain specificity. It should contain the alternatives
that a technically informed critic could reasonably regard as equally valid or
as necessary stress tests for the declared claim. Multiverse and specification-
curve methods formalize related concerns in empirical analysis
\cite{SteegenEtAl2016,SimonsohnEtAl2020}, while sensitivity analysis and
partial identification emphasize that conclusions should expose their
assumption dependence rather than conceal it \cite{Leamer1983,Manski2003}.

Three classes should be reported separately:
\begin{description}
\item[Equivalent alternatives.] Different encodings or estimators of the same
 estimand after appropriate transport.
\item[Robustness alternatives.] Scientifically defensible choices intended to
 support the same claim.
\item[Scope alternatives.] Genuinely different targets or domains that define
 different claims and should not be averaged into one robustness number.
\end{description}

\subtopic{61.4 Common gate schema}

A robustness family must use a common semantic gate index
\(\mathcal J_*\). A specification may mark a gate as not required only when the
claim level itself does not include that property. It may not remove a failed
gate after observing the result. Let
\[
\mathcal J(\varsigma)\subseteq\mathcal J_*
\]
be fixed before certification. For a family intended to support one fixed
claim level, the cleanest design uses
\[
\mathcal J(\varsigma)=\mathcal J_*
\qquad
(\varsigma\in\mathfrak S).
\]
This prevents apparent robustness from being created by silently changing the
meaning of passage across the family.

\topic{62}{Robustness Across Admissible Specifications}

For each gate \(q\in\mathcal J_*\), define the specification envelope
\[
\boxed{
\underline m_q^{\mathfrak S}(W)
=
\inf_{\varsigma\in\mathfrak S}m_q(W;\varsigma),
\qquad
\overline m_q^{\mathfrak S}(W)
=
\sup_{\varsigma\in\mathfrak S}m_q(W;\varsigma).
}
\]
Its sensitivity width is
\[
\Delta_q^{\mathrm{spec}}(W)
=
\overline m_q^{\mathfrak S}(W)
-
\underline m_q^{\mathfrak S}(W).
\]
The robust certificate margin is
\[
\boxed{
r_{\mathfrak S}(W)
=
\inf_{\varsigma\in\mathfrak S}
\min_{q\in\mathcal J_*}m_q(W;\varsigma).
}
\]
Then
\[
\operatorname{Cert}_{\mathfrak S}(W)=1
\quad\Longleftrightarrow\quad
r_{\mathfrak S}(W)>0.
\]

\begin{definition}[Specification-specific, robust, and fragile claims]
A claim is:
\begin{description}
\item[Specification-specific] when it is asserted only for one declared
 \(\varsigma_0\).
\item[Specification-robust] when every required gate has positive lower
 envelope across \(\mathfrak S\).
\item[Specification-fragile] when it passes for at least one admissible
 specification and fails for at least one other.
\end{description}
\end{definition}

A high fraction of passing specifications is informative but does not replace
the lower envelope when the claim is universal over \(\mathfrak S\). A family
in which ninety-nine analyses pass and one scientifically decisive hidden-
controller null fails does not support the complete relational certificate.
The noncompensatory logic remains in force at the specification level.

\subtopic{62.1 Geometry of a specification family}

Let \(d_{\mathfrak S}\) be a declared extended pseudometric on the quotient of
\(\mathfrak S\) by specification equivalence. A convenient decomposition is
\[
\begin{aligned}
d_{\mathfrak S}(\varsigma,\varsigma')
={}&
 w_Dd_D(\mathcal D,\mathcal D')
+w_Yd_Y(Y^+,Y^{+'})
+w_\ell d_\ell(\ell,\ell')\\
&+w_\tau d_\tau(\tau,\tau')
+w_Id_I(\mathcal I,\mathcal I')
+w_Nd_N(\mathcal N,\mathcal N')
+w_\Theta d_\Theta(\Theta,\Theta').
\end{aligned}
\]
No universal metric is imposed. The target term may use a divergence between
conditional laws under a declared coupling. The intervention and null terms
may use Hausdorff distances between implemented kernels. Threshold terms must
be expressed on normalized signed-margin scales. The metric is part of the
audit and should be varied in sensitivity analysis when more than one
scientifically defensible choice exists.

For a reference specification \(\varsigma_0\), define the distance to
certificate failure by
\[
\boxed{
r_{\mathrm{spec}}(W;\varsigma_0)
=
\inf\left\{
 d_{\mathfrak S}(\varsigma,\varsigma_0):
 \min_q m_q(W;\varsigma)\leq0
\right\}.
}
\]
If no failing specification exists in the declared universe, the radius is
infinite relative to that universe. A radius near zero means that a small
scientifically defensible change can reverse the conclusion.

\subtopic{62.2 Specification transport}

Changing the decomposition can alter the coordinate realization of the
workspace, while changing the target or horizon can alter the relevant
predictive state. Let
\[
\Phi_{\varsigma}:W_{\varsigma_0}\longrightarrow W_{\varsigma}
\]
be the declared correspondence when such transport is needed. For each gate,
assume a bound
\[
\left|
 m_q(\Phi_{\varsigma}W;\varsigma)
 -m_q(W;\varsigma_0)
\right|
\leq
L_qd_{\mathfrak S}(\varsigma,\varsigma_0)
+b_q(\varsigma),
\]
where \(b_q\) records mismatch not represented by the metric. The mismatch may
include a change of model class, imperfect target coupling, or an intervention
whose implementation cannot be transported exactly.

\subtopic{62.3 \(T_{28}\): Specification-transport robustness theorem}

\begin{theorem}[Specification-transport robustness]
\label{thm:specification-transport}
Let \(\mathfrak S\) be an admissible specification family with reference
\(\varsigma_0\). Suppose every required gate satisfies the transport bound
above. If
\[
\boxed{
 m_q(W;\varsigma_0)
>
\sup_{\varsigma\in\mathfrak S}
\left[
L_qd_{\mathfrak S}(\varsigma,\varsigma_0)
+b_q(\varsigma)
\right]
}
\]
for every \(q\in\mathcal J_*\), then
\[
\operatorname{Cert}(\Phi_{\varsigma}W;\varsigma)=1
\qquad
\text{for every }\varsigma\in\mathfrak S.
\]
If \(b_q\equiv0\), then
\[
\boxed{
r_{\mathrm{spec}}(W;\varsigma_0)
\geq
\min_{q:L_q>0}
\frac{m_q(W;\varsigma_0)}{L_q}.
}
\]
\end{theorem}

\begin{proof}
For each \(\varsigma\in\mathfrak S\),
\[
m_q(\Phi_{\varsigma}W;\varsigma)
\geq
m_q(W;\varsigma_0)
-L_qd_{\mathfrak S}(\varsigma,\varsigma_0)
-b_q(\varsigma).
\]
The hypothesis makes the right-hand side positive uniformly over the family.
Every required gate therefore passes under every admissible specification.
When \(b_q=0\), no gate can cross zero at a distance smaller than its reference
margin divided by its Lipschitz constant. Taking the minimum over gates gives
the radius bound.
\end{proof}

The theorem is deliberately gatewise. A large structural margin cannot absorb
a fragile causal margin. It also clarifies the relation to Part~XVII. The
decomposition-transport theorem is recovered by varying only the
\(\mathcal D\) coordinate while holding the remaining coordinates fixed.

\subtopic{62.4 A three-budget sufficient condition}

In practice, the reference margin is reduced by three distinct uncertainty
budgets:
\[
\boxed{
B_q^{\mathrm{total}}
=
B_q^{\mathrm{model}}
+B_q^{\mathrm{spec}}
+B_q^{\mathrm{sample}}.
}
\]
The first covers variation across observationally compatible realizations. The
second covers variation across admissible specifications. The third covers
finite-sample estimation. The sufficient condition
\[
\widehat m_q(W;\varsigma_0)
>
B_q^{\mathrm{total}}
\qquad(q\in\mathcal J_*)
\]
will reappear below in identified-set and confidence-envelope form. Keeping the
budgets separate prevents an apparently narrow confidence interval from being
mistaken for strong identification or specification robustness.

\topic{63}{Experimental Equivalence and Identifiability}

A complete mathematical definition does not guarantee that its truth can be
recovered from available data. The experimental design may fail to distinguish
two realizations that have different internal organization.

\begin{definition}[Experimental design and induced law]
\label{def:workspace-experiment-design}
An experimental design is a tuple
\[
\mathcal E
=
(\mathcal O,\mathcal I,\mathcal T),
\]
where \(\mathcal O\) is the observation map, \(\mathcal I\) is the family of
implemented interventions, and \(\mathcal T\) is the set of measured temporal
windows. A realization \(\mathfrak R\) induces the family of laws
\[
P_{\mathfrak R}^{\mathcal E}
=
\left
\{
\mathcal L_{\mathfrak R}
\bigl(\mathcal O(\mathbf H)\mid\operatorname{do}(\iota)\bigr):
\iota\in\mathcal I,
\ t\in\mathcal T
\right\}.
\]
\end{definition}

\begin{definition}[Experimental equivalence]
Two realizations are experimentally equivalent under \(\mathcal E\), written
\[
\mathfrak R\equiv_{\mathcal E}\mathfrak R',
\]
when
\[
P_{\mathfrak R}^{\mathcal E}
=
P_{\mathfrak R'}^{\mathcal E}.
\]
A property \(\psi(\mathfrak R)\) is identifiable on a model class
\(\mathfrak M\) under \(\mathcal E\) when it is constant on every
\(\equiv_{\mathcal E}\)-equivalence class contained in \(\mathfrak M\).
\end{definition}

Identifiability is therefore relative to both a model class and an experiment
family. Enlarging the intervention family refines the equivalence classes and
can make previously hidden distinctions identifiable. Enlarging the model
class can have the opposite effect.

\subtopic{63.1 \(T_{29}\): Non-identifiability impossibility theorem}

\begin{theorem}[No uniformly correct decision under experimental equivalence]
\label{thm:workspace-nonidentifiability}
Let \(\mathfrak R,\mathfrak R'\in\mathfrak M\) satisfy
\[
\mathfrak R\equiv_{\mathcal E}\mathfrak R'
\]
but
\[
\operatorname{Cert}(W_{\mathfrak R};\varsigma)
\neq
\operatorname{Cert}(W_{\mathfrak R'};\varsigma).
\]
Then the workspace certificate is not identifiable on \(\mathfrak M\) under
\(\mathcal E\). More strongly, for every possibly randomized decision rule
\(\widehat C\) based only on data generated by \(\mathcal E\),
\[
\boxed{
\mathbb P_{\mathfrak R}
\left(
\widehat C\neq C_{\mathfrak R}
\right)
+
\mathbb P_{\mathfrak R'}
\left(
\widehat C\neq C_{\mathfrak R'}
\right)
=1,
}
\]
where \(C_{\mathfrak R}=\operatorname{Cert}(W_{\mathfrak R};\varsigma)\).
Consequently,
\[
\inf_{\widehat C}
\max\left\{
\mathbb P_{\mathfrak R}(\widehat C\neq C_{\mathfrak R}),
\mathbb P_{\mathfrak R'}(\widehat C\neq C_{\mathfrak R'})
\right\}
\geq\frac12.
\]
\end{theorem}

\begin{proof}
The data law is identical under the two realizations, so
\(p=\mathbb P(\widehat C=1)\) is the same in both. If
\(C_{\mathfrak R}=0\) and \(C_{\mathfrak R'}=1\), the two error probabilities
are \(p\) and \(1-p\). Their sum is one, and at least one is at least one half.
The opposite labeling is identical.
\end{proof}

The theorem is not a limitation of a particular estimator. It is an
information-theoretic obstruction. No amount of optimization or sample size
can resolve a difference that the declared observations and interventions do
not expose.

\subtopic{63.2 Hidden-controller witness}

Consider two realizations that emit the same local streams \(X_1,X_2\) and the
same outcome
\[
Y=X_1\oplus X_2.
\]
In the distributed realization, the two streams are generated and retained by
separately intervenable carriers, and no carrier contains the pair. In the
centralized realization, a hidden controller contains \((X_1,X_2)\) and emits
the two displayed streams. Under passive observation of
\((X_1,X_2,Y)\), the induced laws are identical. Even interventions that alter
only the displayed messages may fail to distinguish them if the hidden
controller simply re-emits the manipulated values.

Carrier-relative relationality becomes identifiable only when the experiment
can test the claimed causal separation. Examples include independently
interrupting the candidate carriers, exposing the hidden memory, randomizing a
shared controller, or measuring whether one physical mechanism reconstructs
the complete state before the displayed split. The decomposition controls
\(N_{25}\)--\(N_{27}\) are therefore identifiability interventions, not merely
optional robustness checks.

\subtopic{63.3 Identifiability of individual gates}

The full certificate may be unidentified even when some coordinates are
identified. Structural dimension might be recoverable from observed transfer
constraints while causal specificity remains ambiguous. Conversely, a causal
contrast may be identifiable even when the latent workspace coordinates are
known only up to gauge. Each margin should therefore be assigned one of four
statuses:
\[
\text{point identified},
\quad
\text{sign identified},
\quad
\text{partially identified},
\quad
\text{unidentified}.
\]
A complete pass requires the sign of every required gate to be identified as
positive, not merely that one compatible model produces a positive estimate.

\topic{64}{Partial Identification of Workspace Claims}

When point identification fails, the correct response is not to choose one
convenient realization. It is to report the set of values compatible with the
data and assumptions \cite{Manski2003,ImbensManski2004}.

\begin{definition}[Identified realization set]
\label{def:identified-realization-set}
Given an experimental law \(P^{\mathcal E}\) and model class \(\mathfrak M\),
define
\[
\boxed{
\mathfrak R_{\mathrm{ID}}
(P^{\mathcal E};\mathfrak M)
=
\left\{
\mathfrak R\in\mathfrak M:
P_{\mathfrak R}^{\mathcal E}=P^{\mathcal E}
\right\}.
}
\]
For a gate \(q\) and specification \(\varsigma\), its identified margin set is
\[
\mathcal M_q^{\mathrm{ID}}
(P^{\mathcal E};\varsigma)
=
\left\{
 m_q(W_{\mathfrak R};\varsigma):
\mathfrak R\in\mathfrak R_{\mathrm{ID}}
\right\}.
\]
\end{definition}

The sharp lower and upper identified margins are
\[
\underline m_q^{\mathrm{ID}}(\varsigma)
=
\inf\mathcal M_q^{\mathrm{ID}}(\varsigma),
\qquad
\overline m_q^{\mathrm{ID}}(\varsigma)
=
\sup\mathcal M_q^{\mathrm{ID}}(\varsigma).
\]
A gate is sign-identified as passing if its lower bound is positive. It is
sign-identified as failing if its upper bound is nonpositive. If the interval
crosses zero, the available design does not determine that gate's outcome.

\subtopic{64.1 Identified certificate set}

Define
\[
\boxed{
\mathcal C_{\mathrm{ID}}
(P^{\mathcal E};\varsigma)
=
\left\{
\operatorname{Cert}(W_{\mathfrak R};\varsigma):
\mathfrak R\in\mathfrak R_{\mathrm{ID}}
\right\}
\subseteq\{0,1\}.
}
\]
The certificate is point identified exactly when this set is a singleton. The
marginal gate intervals provide simple sufficient conditions, but they need
not be logically necessary because gate margins can be dependent across the
identified set.

For robustness across specifications, define the realization-level robust
margin
\[
r_{\mathfrak S}(\mathfrak R)
=
\inf_{\varsigma\in\mathfrak S}
\min_{q\in\mathcal J_*}
 m_q(W_{\mathfrak R};\varsigma)
\]
and its identified range
\[
\underline r_{\mathrm{ID},\mathfrak S}
=
\inf_{\mathfrak R\in\mathfrak R_{\mathrm{ID}}}
 r_{\mathfrak S}(\mathfrak R),
\qquad
\overline r_{\mathrm{ID},\mathfrak S}
=
\sup_{\mathfrak R\in\mathfrak R_{\mathrm{ID}}}
 r_{\mathfrak S}(\mathfrak R).
\]

\subtopic{64.2 \(T_{30}\): Joint specification-identification theorem}

\begin{theorem}[Identification of a specification-robust certificate]
\label{thm:joint-spec-identification}
Let \(\mathfrak R_{\mathrm{ID}}\) be the identified realization set and
\(\mathfrak S\) an admissible specification family.
\begin{enumerate}
\item If
\[
\boxed{
\underline r_{\mathrm{ID},\mathfrak S}>0,
}
\]
then every realization compatible with the experimental law passes every
required gate under every admissible specification. The robust certificate is
identified as passing.
\item If
\[
\boxed{
\overline r_{\mathrm{ID},\mathfrak S}\leq0,
}
\]
then no compatible realization passes the complete certificate throughout the
family. The robust certificate is identified as failing.
\item If there exist compatible realizations
\(\mathfrak R_+,\mathfrak R_-\) with
\[
r_{\mathfrak S}(\mathfrak R_+)>0,
\qquad
r_{\mathfrak S}(\mathfrak R_-)\leq0,
\]
then the robust certificate is not identified under the declared experiment.
\end{enumerate}
\end{theorem}

\begin{proof}
The first two conclusions follow directly from the definitions of the infimum
and supremum over the identified set. The third gives two experimentally
compatible realizations with different robust-certificate values, so
Theorem~\ref{thm:workspace-nonidentifiability} applies.
\end{proof}

This theorem joins two forms of robustness that are often conflated. A result
can be stable across specifications within one fitted model yet fail to be
identified across models compatible with the data. It can also be identified
under one specification while changing sign across a reasonable specification
family. The strongest claim controls both axes.

\subtopic{64.3 Confidence regions for identified sets}

In finite samples, the identified set is itself uncertain. Let
\[
\widehat{\mathfrak R}_{1-\delta}
\]
be a confidence region satisfying
\[
\inf_{\mathfrak R_0\in\mathfrak M}
\mathbb P_{\mathfrak R_0}
\left(
\mathfrak R_0\in
\widehat{\mathfrak R}_{1-\delta}
\right)
\geq1-\delta.
\]
A conservative robust lower margin is
\[
\widehat r_{1-\delta}
=
\inf_{\mathfrak R\in\widehat{\mathfrak R}_{1-\delta}}
\inf_{\varsigma\in\mathfrak S}
\min_q m_q(W_{\mathfrak R};\varsigma).
\]
If \(\widehat r_{1-\delta}>0\), a false specification-robust certificate has
probability at most \(\delta\), provided that the model class and
specification family were fixed before the confidence region was computed.
This confidence-set approach is useful when gate estimators are strongly
coupled or when the latent realization is the primary inferential object.

\topic{65}{Uniform Finite-Sample Certification}

Population margins are not observed directly. A certification rule must
control the estimation error uniformly across every specification and gate on
which its conclusion depends.

Let \(\widehat m_q(\varsigma)\) estimate
\(m_q(W;\varsigma)\). Suppose there are nonnegative error radii
\(\varepsilon_q(n,\delta)\) such that
\[
\mathbb P\left(
\sup_{\varsigma\in\mathfrak S}
\left|
\widehat m_q(\varsigma)-m_q(\varsigma)
\right|
\leq
\varepsilon_q(n,\delta)
\quad\text{for all }q\in\mathcal J_*
\right)
\geq1-\delta.
\]
Define the empirical lower confidence envelope
\[
\boxed{
\widehat{\underline m}_{q,1-\delta}^{\mathfrak S}
=
\inf_{\varsigma\in\mathfrak S}
\widehat m_q(\varsigma)
-
\varepsilon_q(n,\delta).
}
\]

\subtopic{65.1 \(T_{31}\): Uniform specification-certificate theorem}

\begin{theorem}[Uniform finite-sample certification across specifications]
\label{thm:uniform-spec-certification}
Under the uniform error event above, declare a specification-robust certificate
only when
\[
\widehat{\underline m}_{q,1-\delta}^{\mathfrak S}>0
\qquad
\text{for every }q\in\mathcal J_*.
\]
Then
\[
\boxed{
\sup_{P:\,\operatorname{Cert}_{\mathfrak S}(W)=0}
\mathbb P_P
\left(
\text{the robust certificate is declared}
\right)
\leq\delta.
}
\]
\end{theorem}

\begin{proof}
On the uniform error event,
\[
m_q(\varsigma)
\geq
\widehat m_q(\varsigma)-\varepsilon_q
\]
for every gate and every specification. If all empirical lower envelopes are
positive, every population margin is positive throughout the family. A false
declaration can therefore occur only when the uniform error event fails, which
has probability at most \(\delta\).
\end{proof}

The theorem controls the complete specification family, not merely the
analysis that happened to produce the strongest estimate. It is an extension
of the gatewise intersection rule in Theorem~\ref{thm:gatewise-validity}.
Gatewise validity controls the conjunction at one frozen specification.
Uniform validity controls the conjunction over a family.

\subtopic{65.2 Finite specification grids}

Suppose \(\mathfrak S\) and \(\mathcal J_*\) are finite, and each margin is
estimated as
\[
\widehat m_q(\varsigma)
=
\frac1n\sum_{i=1}^nZ_{i,q,\varsigma}-\theta_{q,\varsigma},
\qquad
0\leq Z_{i,q,\varsigma}\leq1,
\]
from independent certification units. Hoeffding's inequality and a union bound
give
\[
\mathbb P\left(
\sup_{q,\varsigma}
\left|
\widehat m_q(\varsigma)-m_q(\varsigma)
\right|
>
\sqrt{
\frac{
\log\bigl(2|\mathcal J_*||\mathfrak S|/\delta\bigr)
}{2n}}
\right)
\leq\delta.
\]
Thus the choice
\[
\boxed{
\varepsilon_n(\delta)
=
\sqrt{
\frac{
\log\bigl(2|\mathcal J_*||\mathfrak S|/\delta\bigr)
}{2n}
}
}
\]
yields an explicit conservative certificate. Dependent episodes require a
bound appropriate to the true independent unit rather than token-level or
edge-level pseudo-replication.

\subtopic{65.3 Continuous specification families}

Let \(\mathfrak S_\eta\) be an \(\eta\)-net of
\((\mathfrak S,d_{\mathfrak S})\) with covering number
\(N(\eta,\mathfrak S,d_{\mathfrak S})\). If both the population and empirical
margin functions are \(L_q\)-Lipschitz, then
\[
\sup_{\varsigma\in\mathfrak S}
|\widehat m_q(\varsigma)-m_q(\varsigma)|
\leq
\max_{\varsigma_j\in\mathfrak S_\eta}
|\widehat m_q(\varsigma_j)-m_q(\varsigma_j)|
+2L_q\eta.
\]
For bounded empirical means, a valid radius is therefore
\[
\boxed{
\varepsilon_{q,n,\eta}(\delta)
=
\sqrt{
\frac{
\log\bigl(2|\mathcal J_*|
N(\eta,\mathfrak S,d_{\mathfrak S})/\delta\bigr)
}{2n}
}
+2L_q\eta.
}
\]
The tuning parameter \(\eta\) balances net cardinality against interpolation
error. More refined empirical-process bounds may replace this elementary
covering argument without changing the certificate logic.

\subtopic{65.4 Power and indifference regions}

A robust procedure can be valid yet uninformative when the family is too large
or the data are insufficient. Each gate should therefore report an
indifference region \([-\kappa_q,\kappa_q]\). A conclusion is:
\begin{description}
\item[Robust pass] when the lower confidence envelope exceeds
 \(\kappa_q\) for every gate.
\item[Robust fail] when a required gate's upper confidence envelope is below
 \(-\kappa_q\) throughout the relevant family.
\item[Specification conflict] when well-estimated admissible specifications
 lie on opposite sides of the indifference region.
\item[Statistically unresolved] when sampling uncertainty intersects the
 indifference region.
\end{description}
This vocabulary separates genuine sensitivity from low power.

\topic{66}{Adaptive Specification Choice and Selective Validity}

A pointwise valid test does not remain valid after the analyst searches for the
most favorable target, horizon, null panel, or threshold. Consider
\(K\) independent null \(p\)-values and select the smallest. Then
\[
\mathbb P\left(\min_{1\leq k\leq K}p_k\leq\alpha\right)
=
1-(1-\alpha)^K,
\]
which exceeds \(\alpha\) whenever \(K>1\). The same logic applies when a
workspace certificate is reported only for the specification that produced the
largest gate margins.

The issue is not exploration itself. Exploration is often necessary to learn
which target and intervention are scientifically meaningful. The issue is
using the same random outcomes both to choose the claim and to certify it.
Post-selection inference addresses the general problem by using simultaneous
or selection-adjusted procedures \cite{BerkEtAl2013}.

\subtopic{66.1 Specification leakage}

\begin{definition}[Specification leakage]
Specification leakage occurs when certification outcomes influence any of the
following without a fresh inferential correction:
\[
Y^+,
\ell,
\tau,
\mathcal I,
\mathcal N,
\Theta,
\mathfrak S,
\mathcal J_*.
\]
The leakage status is separate from formation leakage. A workspace extractor
may be frozen while the analyst still searches over targets or thresholds on
the certification split.
\end{definition}

A valid report must retain the complete search history, including failed and
abandoned specifications. Reporting only the selected analysis prevents the
reader from determining the effective multiplicity.

\subtopic{66.2 \(T_{32}\): Sample-split selection theorem}

\begin{theorem}[Validity after independent specification selection]
\label{thm:sample-split-spec-selection}
Let \(\widehat\varsigma\) be selected using
\(\mathcal D_{\mathrm{sel}}\), and let
\(\mathcal D_{\mathrm{cert}}\) be independent of
\(\mathcal D_{\mathrm{sel}}\). For every fixed
\(\varsigma\in\mathfrak S\), let
\(\phi_\varsigma(\mathcal D_{\mathrm{cert}})\in\{0,1\}\) be a certificate
rule satisfying
\[
\sup_{P:\operatorname{Cert}(W;\varsigma)=0}
\mathbb P_P\bigl(\phi_\varsigma=1\bigr)
\leq\alpha.
\]
Then the selected-specification procedure satisfies
\[
\boxed{
\sup_P
\mathbb P_P\left(
\phi_{\widehat\varsigma}=1,
\operatorname{Cert}(W;\widehat\varsigma)=0
\right)
\leq\alpha.
}
\]
\end{theorem}

\begin{proof}
Condition on \(\mathcal D_{\mathrm{sel}}\). The selected specification is then
fixed. Independence leaves the certification data distributed exactly as in
the fixed-specification test, whose false-pass probability is at most
\(\alpha\). Averaging the conditional bound over the selection data preserves
the same bound.
\end{proof}

The theorem justifies discovery on formation or calibration data followed by a
single frozen evaluation on fresh certification data. Reusing the
certification split to revise the specification requires either another fresh
split or a simultaneous post-selection procedure whose validity covers the
entire search.

\subtopic{66.3 Nested selection and stress testing}

A practical design may use:
\[
\mathcal D_{\mathrm{form}}
\sqcup
\mathcal D_{\mathrm{spec}}
\sqcup
\mathcal D_{\mathrm{cert}}
\sqcup
\mathcal D_{\mathrm{stress}}.
\]
Formation learns the workspace. The specification split chooses targets,
horizons, nulls, and thresholds from a predeclared menu. Certification makes
the pass decision once. Stress testing explores broader shifts without
retroactively changing the certified claim. Cross-fitting can recover sample
efficiency, but the fold-level selection and evaluation paths must remain
separated and the final uncertainty calculation must account for recombination.

\subtopic{66.4 Thresholds and null families are part of selection}

Threshold tuning can be as consequential as target selection. If the observed
causal effect is \(0.12\), choosing a minimum effect threshold of \(0.10\)
after seeing the estimate is not certification. Likewise, retaining only the
nulls that the intact system defeats converts a falsification panel into a
success filter. Thresholds and null families must be frozen or evaluated under
a uniform rule that includes every considered alternative.

\topic{67}{Specification Audit and Reporting}

A complete specification audit should accompany every strong relational
workspace claim.

\subtopic{67.1 Audit procedure}

\begin{enumerate}
\item State the reference specification \(\varsigma_0\) in full.
\item Distinguish equivalent recodings, robustness alternatives, and genuinely
 different scope claims.
\item Declare the admissible family \(\mathfrak S\) and justify every major
 exclusion before certification.
\item Record the specification pseudometric or other transport bounds used to
 compare family members.
\item Compute every gate's lower envelope, upper envelope, sensitivity width,
 and reference margin.
\item State the experiment family \(\mathcal E\) and the model class
 \(\mathfrak M\) relative to which identifiability is claimed.
\item Report identified sets or bounds whenever point identification fails.
\item Apply a uniform finite-sample rule over every gate and specification used
 in the conclusion.
\item Preserve the complete specification-selection history and the data split
 on which each choice was made.
\item Report the strongest pointwise, family-robust, sign-identified, and
 statistically resolved conclusions separately.
\end{enumerate}

\subtopic{67.2 Additional adversarial controls}

The empirical null and audit panel is extended by:
\[
\begin{aligned}
N_{29}&=\text{post hoc target-selection audit},\\
N_{30}&=\text{horizon cherry-picking audit},\\
N_{31}&=\text{loss and utility substitution audit},\\
N_{32}&=\text{restricted intervention-family audit},\\
N_{33}&=\text{threshold and gate-set tuning audit},\\
N_{34}&=\text{frozen specification sweep across }\mathfrak S.
\end{aligned}
\]
The first five reconstruct or bound the analysis search that preceded the
reported claim. The final control reruns the complete frozen certificate across
the admissible family. A claim that survives \(N_{34}\) but fails the search
audit remains vulnerable to selection bias. A claim that passes the search
audit but changes sign across \(\mathfrak S\) is validly estimated yet
specification-fragile.

\subtopic{67.3 Required specification record}

The machine-readable record should include:
\[
\boxed{
\begin{aligned}
\mathsf{SpecificationRecord}
=\{&
\mathsf{reference},
\mathsf{family},
\mathsf{equivalence\_classes},
\mathsf{metric},\\
&\mathsf{targets},
\mathsf{losses},
\mathsf{horizons},
\mathsf{interventions},\\
&\mathsf{null\_panels},
\mathsf{thresholds},
\mathsf{selection\_history},\\
&\mathsf{identified\_sets},
\mathsf{uniform\_bounds},
\mathsf{gate\_envelopes}
\}.
\end{aligned}
}
\]
Every omitted family member requires a reason. Every changed specification
requires a data-split label and time stamp. Every gate requires pointwise and
family-wide status.

\subtopic{67.4 Claim-strength vocabulary}

The following vocabulary is recommended:
\begin{description}
\item[Pointwise demonstrated.] The certificate passes at one frozen
 specification with valid uncertainty.
\item[Specification-robust.] It passes uniformly across a declared admissible
 family.
\item[Sign identified.] Every realization compatible with the experimental
 law gives the same gate sign.
\item[Robustly identified.] The complete family-robust certificate is
 identified as passing.
\item[Partially identified.] The data restrict the certificate margins but do
 not determine every required sign.
\item[Specification-fragile.] Validly estimated admissible specifications give
 different certificate outcomes.
\item[Selection-contaminated.] The reported specification was chosen using
 certification outcomes without valid adjustment.
\end{description}

\subtopic{67.5 Consequence for the motivating collaboration}

The motivating human--AI research process remains a plausible candidate
persistent workspace, but its strongest current description is narrower than
a certified robust claim. The relevant target trajectory, horizon family,
intervention set, null panel, thresholds, and admissible specification family
have not been frozen and tested prospectively. Nor has the carrier-relative
workspace predicate been shown to be identifiable against hidden-controller
and matched-storage alternatives. The theory therefore supports a disciplined
research program rather than a retrospective declaration.

The terminal architecture is now explicit. A strong claim must survive three
separate questions:
\[
\boxed{
\begin{aligned}
&\text{Would it survive another defensible specification?}\\
&\text{Could a contrary realization generate the same experiments?}\\
&\text{Does the finite sample support every required margin uniformly?}
\end{aligned}
}
\]
Once these questions are answered, further extensions concern new mathematical
settings or new empirical domains rather than a missing foundational layer.

\theorypart{XIX}{Summary of Main Results}

\topic{68}{Structural and Foundational Results}

\subtopic{\(T_1\): Limit--holonomy theorem}
\[
\Gamma(H)\cong C_o^{\operatorname{Hol}_o}.
\]

\subtopic{\(T_2\): Operator closure theorem}
\[
H_e\rho_i(a)=\rho_j(a)H_e
\Longrightarrow
\rho_o(a)C_*\subseteq C_*.
\]

\subtopic{\(T_{2a}\): Workspace commutant theorem }
\[
\mathcal O_{\Workspace}
\cong
\operatorname{Comm}(\Hol_o).
\]

\subtopic{\(T_3\): Transfer-Laplacian kernel theorem}
\[
\ker L_H=\Gamma(H).
\]

\subtopic{\(T_4\): Spectral stability theorem}
If \(\eta=\|\widetilde L_H-L_H\|<g_d/2\), then
\[
\|\widetilde P_d-P_d\|
\leq
\frac{\eta}{g_d-\eta}
\leq
\frac{2\eta}{g_d}.
\]

\subtopic{\(T_5\): Gauge invariance theorem}
\[
\Gamma(H^g)=G_V\Gamma(H),
\]
and under unitary gauge,
\[
L_{H^g}=G_VL_HG_V^{-1}.
\]

\subtopic{\(T_6\): Persistence theorem}

For the nonautonomous recurrence
\[
e_{t+1}\leq\alpha_te_t+\sigma_t,
\]
\[
\boxed{
\begin{aligned}
e_t
\leq
&\left(\prod_{r=0}^{t-1}\alpha_r\right)e_0\\
&+\sum_{j=0}^{t-1}
\left(\prod_{r=j+1}^{t-1}\alpha_r\right)\sigma_j.
\end{aligned}
}
\]
If \(\alpha_t\leq\bar\alpha<1\) and \(\sigma_t\leq\bar\sigma\), then
\[
\limsup_{t\to\infty}e_t
\leq
\frac{\bar\sigma}{1-\bar\alpha}.
\]
The derivation begins with a moving workspace projector and
an explicit system update, and adds a mean-square analogue and a non-collapse
gate.

\subtopic{\(T_7\): Stable reconstruction theorem}

For coalition analysis operator \(A_S\) with lower frame bound
\(\alpha_S>0\), observations
\[
y=A_Sc+\eta
\]
and canonical decoder \(A_S^\dagger\) satisfy
\[
\boxed{
\|A_S^\dagger y-c\|
\leq
\frac{\|\eta\|}{\sqrt{\alpha_S}}.
}
\]
The theorem is extended to encoding mismatch, exact support-based erasure
distance, uniform node erasures, reconstructive handoff, exact coalition
sufficiency, and operational risk under a Lipschitz predictive kernel.

\subtopic{\(T_8\): Exact gluing theorem}

For every finite edge-complete workspace cover,
\[
\boxed{
0
\longrightarrow
\mathscr W(G)
\xrightarrow{\ R_{\mathcal U}\ }
\check C^0(\mathcal U;\mathscr W)
\xrightarrow{\ \delta^0\ }
\check C^1(\mathcal U;\mathscr W)
}
\]
is exact. Hence
\[
\boxed{
\mathscr W(G)
\cong
\ker\delta^0
\cong
\operatorname{Eq}
\left(
\prod_\alpha\mathscr W_\alpha
\rightrightarrows
\prod_{\alpha<\beta}\mathscr W_{\alpha\beta}
\right).
}
\]
The same analysis also yields a cover-Laplacian distance bound, a total descent kernel
theorem, and a harmonic \(\check H^1\) obstruction for relative-only data.

\subtopic{\(T_9\): Hierarchical collapse theorem}

For every finite rooted hierarchy whose children edge-cover their parent,
\[
\boxed{
W_a\cong\mathscr W(G_a)
\quad\text{for every hierarchy node }a,
}
\]
and in particular
\[
\boxed{W_{\mathrm{root}}\cong\mathscr W(G).}
\]
If level \(r\) contributes defect \(\delta_r\) and downstream collapse maps
have Lipschitz constants \(L_s\), then
\[
\boxed{
\delta_{\mathrm{flat}}^{(R)}
\leq
\sum_{r=0}^{R-1}
\left(
\prod_{s=r+1}^{R-1}L_s
\right)\delta_r.
}
\]


\subtopic{\(T_{25}\): Faithful coarse-graining theorem}

For a partition of the vertex set, let \(W_\Pi\) be the direct sum of the
internally coherent section spaces and let \(D_\Pi\) retain all boundary
compatibility equations. Then
\[
\boxed{
\Gamma(H)
=
J_\Pi\ker D_\Pi.
}
\]
Exact structural workspace content is therefore invariant under faithful
aggregation. Approximate error separates into an internal-coherence term and
a boundary-coherence term controlled by their respective spectral gaps.

\subtopic{\(T_{26}\): Bayes-risk transport theorem}

For an information-faithful coarsening \(q:V_f\to V_c\),
\[
\boxed{
\Risk_{\mathcal D_c}(S)
=
\Risk_{\mathcal D_f}(q^{-1}(S))
\qquad(S\subseteq V_c).
}
\]
Hence displayed coalition size can decrease under coarsening. Weighted
primitive support cost is a more stable localization measure and cannot
improve merely because several carriers were relabeled as one macro-node.

\subtopic{\(T_{27}\): Decomposition-transport robustness theorem}

If every gate margin at a reference decomposition exceeds its worst-case
transport and model-mismatch defect across \(\mathfrak D\), then the complete
certificate passes throughout the admissible family:
\[
\boxed{
 m_q(W;\mathcal D_0)
>
\sup_{\mathcal D\in\mathfrak D}
\bigl(L_q\delta_{\mathcal D}+b_q(\mathcal D)\bigr)
\quad\forall q
\Longrightarrow
\operatorname{Certify}_{L,\mathfrak D}^{\mathrm{rob}}(W)=1.
}
\]

\topic{69}{Dynamic, Causal, and Agency Results}

\subtopic{\(T_{10}\): Approximate universality theorem }

For the first \(d\) transfer-Laplacian modes, with
\(\lambda_d\leq\varepsilon<\lambda_{d+1}=\mu\),
\[
\boxed{
\|(I-P_d)\rho_V(a)P_d\|
\leq
\frac{\|\rho_E(a)\|\sqrt{\varepsilon}+\beta(a)}{\sqrt{\mu}}.
}
\]
Thus low transfer energy, a separated complement, and small intertwining defect
imply approximate closure under the common operation algebra.

\subtopic{\(T_{11}\): Core--residual accumulation theorem }

For a path \(\gamma=e_m\cdots e_1\),
\[
\boxed{
\begin{aligned}
\|Q_{t(\gamma)}T_\gamma P_{s(\gamma)}\|
\leq
\sum_{r=1}^{m}
&\left(
\prod_{\ell=r+1}^{m}\|T_{e_\ell}\|
\right)
\|V_{e_r}\|\\
&\times
\left(
\prod_{\ell=1}^{r-1}\|H_{e_\ell}\|
\right).
\end{aligned}
}
\]
The same bound controls the defect between full core-to-core path transport and
the pure core product. An exact Duhamel
identity yields both estimates.

\subtopic{\(T_{12}\): Predictive-state minimality and alignment theorem }

On standard Borel spaces, the canonical conditional-law state
\[
Z_t
=
\mathbb P(Y_t^+\in\cdot\mid\mathbf H_t)
\]
is sufficient and minimal:
\[
Y_t^+\indep\mathbf H_t\mid Z_t,
\]
and every sufficient statistic \(S_t=s(\mathbf H_t)\) satisfies
\[
Z_t=\phi(S_t)
\qquad\text{a.s.}
\]
Moreover, a structural candidate \(W_t\) is exactly operationally aligned if
and only if
\[
Z_t=\phi_W(W_t)
\qquad\text{a.s.}
\]
The countable intervention-indexed version holds for
\(Z_t^{\mathcal I}\).

\subtopic{\(T_{13}\): Operational relationality theorem }

For a finite target, exact alignment, a positive node-level conditional
information gap, a positive robust local-predictive-matched null gap, a feasible
workspace intervention effect, and positive matched specificity jointly imply
that:
\[
\begin{aligned}
&\text{no single-node statistic is sufficient},\\
&k_0\geq2,\\
&\text{no single node reconstructs the aligned workspace},\\
&\text{matched local nulls preserve all node risks but lose joint capacity},\\
&\text{and the workspace coordinate is causally efficacious at the tested scale.}
\end{aligned}
\]
If \(\eta_{\mathrm{loop}}(\tau)>0\), the conclusion additionally includes
loop-mediated reentry.

\subtopic{\(T_{14}\): Replacement invariance theorem }

For a unitary realization correspondence with discrepancy defect
\[
\delta_D=\|D'J_V-J_ED\|,
\]
set
\[
\eta_L=2\|D\|\delta_D+\delta_D^2.
\]
If \(\eta_L<g/2\), then the workspace projectors obey
\[
\boxed{
q_{\mathrm{str}}
\leq
\frac{\eta_L}{g-\eta_L}
\leq
\frac{2\eta_L}{g}.
}
\]
A canonical unitary correspondence between the workspace subspaces exists,
and its operation defect is bounded by the projector drift, the local
intertwining defect, and prior operation leakage. Coalition frame bounds also
persist under sufficiently small reconstructive transport error. With declared
causal-kernel and reconstructive-handoff defects,
\[
\boxed{
d_{\mathrm{id}}
\leq
 a\,q_{\mathrm{str}}
+b\,\delta_{\mathrm{op}}
+c\,\delta_{\mathrm{causal}}
+d\,\delta_{\mathrm{rec}}.
}
\]
Under bounded norms and a uniform spectral gap, this gives the schematic
estimate
\[
d_{\mathrm{id}}
\lesssim
C\frac{\delta_{\mathrm{replacement}}}{g}.
\]
The causal and reconstructive terms are necessary assumptions rather than
consequences of structural similarity.


\subtopic{\(T_{15}\): Transfer identification theorem }

Under the linear transfer model
\[
Y=H_*X+\xi,
\qquad
\mathbb E[\xi X^{\mathsf T}]=0,
\qquad
\Sigma_{XX}\succeq\kappa I,
\]
the unique population least-squares transfer is
\[
\boxed{H_*=\Sigma_{YX}\Sigma_{XX}^{-1},}
\]
and
\[
\mathcal R(H)-\mathcal R(H_*)
\geq
\kappa\|H-H_*\|_F^2.
\]
For empirical ordinary least squares,
\[
\|\widehat H-H_*\|_F
\leq
\frac{\|\widehat M\|_F}{\widehat\kappa},
\]
with the stated ridge extension. The identified transfer is gauge covariant
under independent invertible changes of source and target coordinates.

\subtopic{\(T_{16}\): Uniform structure-selection theorem }

For a finite candidate family with unique population optimum \(\theta_*\),
margin \(\Delta_J>0\), and held-out uniform score error at most \(\varepsilon\),
\[
\boxed{
J(\widehat\theta)-J(\theta_*)\leq2\varepsilon.
}
\]
If \(2\varepsilon<\Delta_J\), then \(\widehat\theta=\theta_*\). The theorem
applies to node membership, edge selection, gauge-normalized transfer models,
and finite boundary families.

\subtopic{\(T_{17}\): Stable boundary recovery theorem }

For a finite admissible boundary family with unique maximizer \(S_t^*\), margin
\(\Delta_t^B\), and uniform held-out score error \(\varepsilon_t\),
\[
\mathcal B_t(S_t^*)-
\mathcal B_t(\widehat S_t)
\leq2\varepsilon_t,
\]
and exact recovery holds when \(2\varepsilon_t<\Delta_t^B\). If the population
score drifts uniformly by at most \(\delta_t\), then
\[
2\delta_t<\Delta_t^B
\Longrightarrow
S_{t+1}^*=S_t^*.
\]

\subtopic{\(T_{18}\): Model-specific formation transition theorem }

For
\[
m_{t+1}=\tanh(\lambda m_t),
\]
there is a unique zero fixed point for \(0\leq\lambda\leq1\). For
\(\lambda>1\), the zero state is unstable and exactly two stable nonzero fixed
points appear. Their amplitude satisfies
\[
\boxed{
m_\lambda
=
\sqrt{\frac{3(\lambda-1)}{\lambda^3}}
+
O\bigl((\lambda-1)^{3/2}\bigr).
}
\]
This is an exact structural benchmark, not a universal law for complete
relational workspaces.

\subtopic{\(T_{19}\): Cue-controlled goal-retention theorem }

For matched future command and disturbance trajectories under
\[
G_{s+1}=R_sG_s+B_sC_{s+1}+\zeta_{s+1},
\]
two goal interventions satisfy
\[
\boxed{
G_{t+\tau}-G_{t+\tau}'
=
R_{t+\tau:t}(g-g').
}
\]
If
\[
m_s\leq\sigma_{\min}(R_s),
\qquad
\|R_s\|\leq M_s,
\]
then
\[
\boxed{
\left(\prod_s m_s\right)\|g-g'\|
\leq
\|G_{t+\tau}-G_{t+\tau}'\|
\leq
\left(\prod_s M_s\right)\|g-g'\|.
}
\]
A linear action readout gives a corresponding lower bound on the induced
action-mean separation. The result isolates internal retention from repeated
external goal rewriting.

\subtopic{\(T_{20}\): Approximate counterfactual-policy theorem }

If the interventional action-value estimator satisfies
\[
\sup_a|\widehat Q_t(a)-Q_t(a)|\leq\varepsilon_Q
\]
and the implemented policy has optimization error at most
\(\delta_{\mathrm{opt}}\) relative to \(\widehat Q_t\), then
\[
\boxed{
\operatorname{Reg}_t(\widehat\pi_t)
\leq
2\varepsilon_Q+\delta_{\mathrm{opt}}.
}
\]
For deterministic greedy action, the bound is \(2\varepsilon_Q\). A unique
true optimum with value gap greater than \(2\varepsilon_Q\) is recovered
exactly. The guarantee is relative to the declared utility, feasible action
support, causal identification, and horizon.

\subtopic{\(T_{21}\): Repair-selection margin theorem }

For a finite repair set and a repair-value estimator with uniform error
\(\varepsilon_R\), the selected repair \(\widehat r\) satisfies
\[
\boxed{
V(r^*)-V(\widehat r)
\leq
2\varepsilon_R.
}
\]
If the optimal repair reaches reserve \(\rho+2\varepsilon_R\), then the selected
repair reaches reserve \(\rho\). A unique optimal repair whose value gap
exceeds \(2\varepsilon_R\) is selected exactly.

\subtopic{\(T_{22}\): Closed-loop self-maintenance theorem }

If matched no-repair deficit obeys
\[
e_{t+1}^{(0)}\leq\alpha e_t+d_t
\]
and the implemented repair has measured causal effect
\[
e_{t+1}^{(0)}-e_{t+1}\geq\kappa e_t-\epsilon_t,
\]
set \(\beta=\alpha-\kappa\). Under
\[
0\leq\beta<1,
\qquad
 d_t\leq\overline d,
\qquad
 \epsilon_t\leq\overline\epsilon,
\]
\[
\boxed{
e_t
\leq
\beta^te_0
+
\frac{(\overline d+\overline\epsilon)(1-\beta^t)}{1-\beta}.
}
\]
Hence
\[
\limsup_{t\to\infty}e_t
\leq
\frac{\overline d+\overline\epsilon}{1-\beta}.
\]
A corresponding margin condition makes the declared bounded-deficit viability
set forward invariant.

\subtopic{\(T_{23}\): Workspace-agency synthesis theorem }

Operational relationality, cue-controlled and behaviorally operative goals,
calibrated counterfactual values, feasible value-to-action and action-to-outcome
interventions, causal closed-loop maintenance, distributed control-state
sufficiency, and matched local policy nulls jointly imply:
\[
\begin{aligned}
&\operatorname{Reg}_t\leq2\varepsilon_Q+\delta_{\mathrm{opt}},\\
&\limsup e_t\leq
 \frac{\overline d+\overline\epsilon}{1-\beta},\\
&k_{\mathrm{ctl},\varepsilon}\geq2,
\end{aligned}
\]
together with causal mediation of goal-dependent action and repair through the
implemented internal channels. This is a sufficient, domain-relative agency
certificate. It does not imply consciousness, personhood, or unrestricted
autonomy.

\subtopic{\(T_{24}\): Gatewise certification-validity theorem }

Suppose a certificate requires every gate \(G_1,\ldots,G_m\) to pass, and the
predeclared test for gate \(j\) has false-pass probability at most \(\alpha_j\)
whenever that required property is absent. Under the global null that at least
one required property is absent,
\[
\boxed{
\mathbb P(\text{complete certificate passes})
\leq
\max_j\alpha_j.
}
\]
No independence assumption among gate tests is required. The guarantee fails
under adaptive gate deletion, selective reporting, or reuse of the
certification split.


\topic{70}{Decomposition, Specification, and Identification Results}

\subtopic{\(T_{25}\): Faithful coarse-graining theorem}
For an exact partition into internally coherent macro-components,
\[
\boxed{
\ker D_H=J_\Pi\ker D_\Pi.
}
\]
Thus faithful aggregation preserves the exact shared workspace and its
dimension while changing the displayed node resolution.

\subtopic{\(T_{26}\): Bayes-risk transport under aggregation}
For an information-faithful coarsening map \(q\),
\[
\boxed{
\Risk_{\mathcal D_c}(S)
=
\Risk_{\mathcal D_f}\bigl(q^{-1}(S)\bigr).
}
\]
Predictive content is preserved, although the visible coalition count may
contract. Primitive-carrier support cost is therefore preferred when claims
must be compared across resolutions.

\subtopic{\(T_{27}\): Decomposition-transport robustness theorem}
If every required gate has reference margin larger than its worst admissible
decomposition-transport defect, then the complete certificate passes
throughout \(\mathfrak D\). When the mismatch term vanishes, the certificate
radius is bounded below by the smallest gate margin divided by its transport
sensitivity.

\subtopic{\(T_{28}\): Specification-transport robustness theorem}
For a reference specification \(\varsigma_0\), if
\[
m_q(W;\varsigma_0)
>
\sup_{\varsigma\in\mathfrak S}
\left[
L_qd_{\mathfrak S}(\varsigma,\varsigma_0)+b_q(\varsigma)
\right]
\]
for every gate, then the transported workspace passes every admissible
specification. If \(b_q=0\),
\[
r_{\mathrm{spec}}(W;\varsigma_0)
\geq
\min_{q:L_q>0}\frac{m_q(W;\varsigma_0)}{L_q}.
\]

\subtopic{\(T_{29}\): Experimental non-identifiability theorem}
If two realizations induce the same law under every declared observation and
intervention but have different certificate values, no decision rule based on
that experiment can be uniformly correct. For every such rule, the two error
probabilities sum to one, so the minimax error is at least one half.

\subtopic{\(T_{30}\): Joint specification-identification theorem}
Let \(r_{\mathfrak S}(\mathfrak R)\) be the worst gate margin across the
specification family. If its infimum over the identified realization set is
positive, the robust certificate is identified as passing. If its supremum is
nonpositive, the certificate is identified as failing. Compatible realizations
on opposite sides of zero establish non-identifiability.

\subtopic{\(T_{31}\): Uniform finite-sample specification theorem}
If all gate estimators are uniformly within their error radii across
\(\mathfrak S\), then requiring every empirical lower envelope to be positive
controls the probability of a false robust certificate by the failure
probability of the uniform event.

\subtopic{\(T_{32}\): Independent selection and certification theorem}
A specification chosen on data independent of the certification split may be
evaluated by its fixed-specification test without inflating the false-pass
bound. Selection and certification on the same outcomes require simultaneous
or post-selection correction.

\theorypart{XX}{Conjectures and Open Mathematical Directions}

\topic{71}{Workspace Formation Transitions}

The symbolic outline proposed a coupling threshold \(\lambda_c\). A general
theory must first state what is allowed to vary with \(\lambda\). A threshold
may concern structural coherence, causal efficacy, relational irreducibility,
or the conjunction of all certification gates. These need not occur at the
same parameter value.

\subtopic{71.1 No coordinate-free universal threshold}

\begin{proposition}[Threshold values are parameterization dependent]
\label{prop:threshold-reparam}
Suppose a model indexed by \(\lambda\) passes a declared formation certificate
exactly for \(\lambda>\lambda_c\). For every strictly increasing bijection
\(f\), the reparameterized family indexed by \(\mu=f(\lambda)\) has threshold
\[
\mu_c=f(\lambda_c).
\]
Consequently, the numerical value of a coupling threshold is not invariant
without a physically or operationally fixed parameterization.
\end{proposition}

\begin{proof}
The reparameterized realization at \(\mu\) is the original realization at
\(f^{-1}(\mu)\). It passes precisely when
\(f^{-1}(\mu)>\lambda_c\), equivalently when
\(\mu>f(\lambda_c)\).
\end{proof}

There can therefore be no model-independent numerical value of \(\lambda_c\).
The invariant object is the change in the certified organization and its local
critical behavior, not the label placed on the control parameter.

\subtopic{71.2 Gate-specific transition parameters}

For a one-parameter family \(\mathfrak R_\lambda\), define
\[
\begin{aligned}
\lambda_{\mathrm{str}}
&=
\inf\{\lambda:
 d_{\varepsilon}(\lambda)>0,
 g_{\varepsilon}(\lambda)>0\},\\
\lambda_{\mathrm{op}}
&=
\inf\{\lambda:
 \delta_{\mathrm{align}}(\lambda)
 \leq\varepsilon_{\mathrm{op}}\},\\
\lambda_{\mathrm{causal}}
&=
\inf\{\lambda:
 \eta_{\mathrm{causal}}(\lambda)
 \geq\tau_{\mathrm{causal}}\},\\
\lambda_{\mathrm{rel}}
&=
\inf\{\lambda:
 k_\varepsilon(\lambda)\geq2,
 \Delta_{\mathrm{tr}}(\lambda)
 \geq\tau_{\mathrm{tr}}\},\\
\lambda_{\mathrm{form}}
&=
\inf\{\lambda:
 \operatorname{Form}_{\Theta_{\mathrm{form}}}
 (W_\lambda)=1\}.
\end{aligned}
\]
In general,
\[
\lambda_{\mathrm{str}},
\lambda_{\mathrm{op}},
\lambda_{\mathrm{causal}},
\lambda_{\mathrm{rel}},
\lambda_{\mathrm{form}}
\]
can be distinct. A synchronized representation can appear before it becomes
causally useful, and causal coordination can appear while remaining reducible
to one centralized component.

\subtopic{71.3 A solvable nonlinear formation model \(T_{18}\) }

A minimal exact model is the scalar recurrent coherence equation
\[
\boxed{
m_{t+1}=\tanh(\lambda m_t),
\qquad
\lambda\geq0.
}
\]
The order parameter \(m_t\) represents a signed coherent mode. This model is
not a complete relational workspace. It isolates one mechanism by which a
nonzero shared mode can appear continuously from a symmetric zero state.

\begin{theorem}[Pitchfork formation transition]\label{thm:pitchfork-formation}
For the map
\[
F_\lambda(m)=\tanh(\lambda m),
\qquad
m\in[-1,1],
\]
the following hold.
\begin{enumerate}
\item If \(0\leq\lambda\leq1\), the unique fixed point is \(m=0\).
 It is globally attracting for \(\lambda<1\) and nonexpanding at
 \(\lambda=1\).
\item If \(\lambda>1\), there are exactly three fixed points
 \[
 -m_\lambda,
 \quad0,
 \quad m_\lambda,
 \qquad
 m_\lambda\in(0,1).
 \]
 The zero fixed point is unstable and the two nonzero fixed points are
 locally asymptotically stable.
\item As \(\lambda\downarrow1\),
 \[
 \boxed{
 m_\lambda
 =
 \sqrt{\frac{3(\lambda-1)}{\lambda^3}}
 +
 O\bigl((\lambda-1)^{3/2}\bigr).
 }
 \]
\end{enumerate}
\end{theorem}

\begin{proof}
A fixed point satisfies
\[
\operatorname{artanh}(m)=\lambda m.
\]
For \(m>0\), the function
\[
h(m)=\frac{\operatorname{artanh}(m)}{m}
\]
is strictly increasing from \(1\) to \(+\infty\). Hence there is no positive
solution when \(\lambda\leq1\), and exactly one positive solution when
\(\lambda>1\). Odd symmetry gives the negative solution.

The derivative is
\[
F_\lambda'(m)
=
\lambda\operatorname{sech}^2(\lambda m).
\]
At zero this equals \(\lambda\), proving stability for \(\lambda<1\) and
instability for \(\lambda>1\). At a nonzero fixed point,
\[
F_\lambda'(m_\lambda)
=
\lambda(1-m_\lambda^2)
=
\frac{(1-m_\lambda^2)\operatorname{artanh}(m_\lambda)}{m_\lambda}
<1,
\]
because
\[
\operatorname{artanh}(m)
<
\frac{m}{1-m^2}
\qquad(0<m<1).
\]
The derivative is positive, so the nonzero fixed points are locally stable.
Finally, expanding
\[
\tanh(\lambda m)
=
\lambda m-
\frac{\lambda^3m^3}{3}
+O(m^5)
\]
and solving the nonzero fixed-point equation gives the stated asymptotic.
\end{proof}

The square-root scaling is a genuine model-specific critical law. It does not
show that every workspace forms by a pitchfork bifurcation. It supplies an
exact benchmark against which more complex formation dynamics can be tested.

\subtopic{71.4 Detecting a transition from finite data}

Let \(Q(\lambda)\) be a gate-specific population order parameter and
\(\widehat Q_n(\lambda)\) its held-out estimate. If
\[
\sup_{\lambda\in\Lambda_n}
|\widehat Q_n(\lambda)-Q(\lambda)|
\leq\varepsilon_n
\]
on a finite evaluation grid \(\Lambda_n\), then every threshold crossing with
vertical margin larger than \(2\varepsilon_n\) is recovered by the same
uniform-selection argument as Theorem~\ref{thm:uniform-structure-selection}.
Near a continuous transition, the margin vanishes and sample requirements
necessarily grow. A broad, slowly varying crossover should not be reported as
a sharp phase transition without a finite-size scaling analysis.

\subtopic{71.5 General formation-transition conjecture}

For adaptive transfer systems with a fixed physical parameterization, there
may exist model classes in which
\[
\lambda<\lambda_c
\Longrightarrow
\operatorname{Form}_{\Theta_{\mathrm{form}}}(W_\lambda)=0,
\]
while
\[
\lambda>\lambda_c
\Longrightarrow
\operatorname{Form}_{\Theta_{\mathrm{form}}}(W_\lambda)=1
\]
with a stable nonzero structural mode, positive causal efficacy, and
relational irreducibility. The theorem above establishes this only for one
structural order parameter in a solvable nonlinear model. The conjunction of
all workspace gates remains model dependent.


\topic{72}{Hierarchical Stabilization Conjecture}

For appropriate recurrent systems,
\[
W^{(r+1)}
=
\Gamma_{\varepsilon_r}(W^{(r)})
\]
may converge to a stable workspace type:
\[
d_{\mathrm{id}}\!\left(
W^{(r+1)},W^{(r)}
\right)
\to0.
\]

\topic{73}{Endogenous Boundary Conjecture}

The boundary theorems in Part~XV establish recovery and stability once a score,
candidate family, and formation procedure are fixed. They do not establish
that an arbitrary adaptive system will create, represent, and maintain its own
boundary.

Let
\[
B_t^{\mathrm{sys}}
\]
be an internal membership or routing state used by the system, and let
\[
S_t^*
=
\argmax_{S\in\mathfrak S_{\mathrm{pass}}}
\mathcal B_t(S)
\]
be the population boundary defined by the external certification protocol.
The strong endogenous boundary conjecture asserts that there are adaptive
systems and nontrivial environments for which:
\[
\boxed{
\begin{aligned}
&d_B(B_t^{\mathrm{sys}},S_t^*)\to0,\\
&\operatorname{Form}_{\Theta_{\mathrm{form}}}(W_{B_t^{\mathrm{sys}}})=1,\\
&\operatorname{CE}
\left(
B_t^{\mathrm{sys}}\to
\operatorname{Routing}_{t+1}
\right)>0,\\
&\operatorname{CE}
\left(
\operatorname{BoundaryDamage}_t\to
\operatorname{RepairAction}_{t+1}
\right)>0,
\end{aligned}
}
\]
while the scaffold contains no direct encoding of \(S_t^*\).

A convincing demonstration requires four separable achievements:
\begin{enumerate}
\item \emph{discovery}: interaction history changes which boundary is learned;
\item \emph{instantiation}: an internal variable causally realizes that
 boundary;
\item \emph{use}: changing the boundary changes workspace routing and
 downstream behavior; and
\item \emph{maintenance}: detected boundary degradation causes restorative
 action.
\end{enumerate}

Theorems~\ref{thm:heldout-boundary-recovery} and
\ref{thm:boundary-drift-stability} handle external consistency of the selected
boundary. Theorems~\ref{thm:repair-selection} and
\ref{thm:self-maintenance} now handle controlled repair selection and bounded
closed-loop maintenance once a health representation, damage class, monitor,
and repair controller are given. The remaining conjectural step is endogenous
learning of the internal boundary, detector, and repair policy under open-ended
environmental change without direct scaffold encoding or external
substitution.


\theorypart{XXI}{Empirical Certification and Reporting Protocol}

\topic{74}{Required Experiments}

A workspace claim is not certified by a single benchmark score. The complete
protocol is organized as a sequence of declarations, held-out measurements,
implemented interventions, matched nulls, specification sweeps, identification
checks, and gate decisions. Every reported
result must identify the claim level it supports. Measurements that were not
performed remain unmeasured rather than being inferred from neighboring gates.

\subtopic{74.1 Claim declaration before evaluation}

Before the certification split is opened, record:
\begin{enumerate}
\item the full reference specification
 \(\varsigma_0=(\mathcal D,Y^+,\ell,\tau,\mathcal I,\mathcal N,\Theta)\);
\item the admissible decomposition family \(\mathfrak D\) and specification
 family \(\mathfrak S\);
\item the target trajectory \(Y^+\), loss, horizon, and environment family;
\item the candidate structural workspace extractor;
\item the goal family \(\mathcal G\), utility normalization, and feasible
 action set;
\item the intervention families \(\mathcal I_W,\mathcal I_G,\mathcal I_A,
 \mathcal I_D,\mathcal I_R\);
\item the null panel and the local distributions each null must preserve;
\item the required certificate level and every numerical threshold;
\item the formation, calibration, certification, and stress-test splits;
\item the uncertainty procedure and resampling unit; and
\item the stopping and model-selection rules;
\item the experiment family and model class used for identifiability; and
\item the complete specification-selection history and multiplicity rule.
\end{enumerate}
Changing any of these after observing certification outcomes creates a new
analysis that must be reported as exploratory or evaluated on a fresh split.

\subtopic{74.2 Data and intervention partition}

Use at least four logically distinct partitions:
\[
\boxed{
\mathcal D
=
\mathcal D_{\mathrm{form}}
\sqcup
\mathcal D_{\mathrm{cal}}
\sqcup
\mathcal D_{\mathrm{cert}}
\sqcup
\mathcal D_{\mathrm{stress}}.
}
\]
Their roles are:
\begin{description}
\item[Formation.] Learn membership, transfer maps, workspace channels,
 boundaries, policies, and repair models.
\item[Calibration.] Select thresholds, intervention amplitudes, and
 uncertainty procedures without making the final pass decision.
\item[Certification.] Estimate the declared gates once under the frozen
 analysis plan.
\item[Stress test.] Evaluate distribution shift, repeated damage, replacement,
 and long-horizon persistence.
\end{description}
When independent partitions are impossible, the report must state the precise
cross-fitting or nested-resampling procedure and preserve formation leakage
controls.

\subtopic{74.3 Structural tests}

Report the complete low-energy spectrum relevant to the selected threshold:
\[
\spec(L_H),
\qquad
 d_\varepsilon,
\qquad
 g_\varepsilon,
\qquad
 \delta_{\mathrm{coh}}.
\]
At minimum include:
\begin{itemize}
\item the construction and normalization of \(D_H\);
\item edge weights and estimated transfer uncertainties;
\item the selected eigenspace dimension and threshold rule;
\item the spectral gap with uncertainty or perturbation bound;
\item stability under bootstrap, seed, and representation replacement; and
\item the exact or approximate kernel residual.
\end{itemize}

\subtopic{74.4 Transfer and operation tests}

Report
\[
\eta_{\mathrm{int}},
\qquad
\delta_e^{\mathrm{tr}},
\qquad
\Omega_\gamma,
\qquad
\Lambda_\gamma,
\]
together with:
\begin{itemize}
\item the paths and cycle basis used;
\item path-length dependence of leakage;
\item the gauge-alignment procedure;
\item held-out transfer-prediction error;
\item operation-intertwining defects; and
\item the approximate closure bound from Theorem~\(T_{10}\).
\end{itemize}

\subtopic{74.5 Predictive alignment and coalition tests}

Report
\[
\delta_{\mathrm{align}}^{\mathcal I},
\qquad
\{\Delta_A^\ell\}_{A\subseteq V},
\qquad
 k_\varepsilon,
\qquad
\delta_A^Z.
\]
When the number of coalitions is too large for exhaustive evaluation, state the
sampling design, include every singleton and every leave-one-out coalition,
and report the most favorable coalition found for the reducibility null. The
claim must not be based only on average coalition performance.

\subtopic{74.6 Causal and reentry tests}

Report
\[
\eta_{\mathrm{causal}}^\nu,
\qquad
\Delta_{\mathrm{specific}}(r),
\qquad
\delta_{\mathrm{med}},
\qquad
\eta_{\mathrm{rec}}(\tau),
\qquad
\eta_{\mathrm{loop}}(\tau).
\]
Every causal estimate must be paired with:
\begin{enumerate}
\item a concrete mechanism-level implementation;
\item a manipulation check;
\item matched background variables;
\item declared direct paths that were blocked or preserved;
\item intervention magnitude and cost;
\item an off-workspace matched intervention; and
\item the downstream outcome window.
\end{enumerate}

\subtopic{74.7 Persistence, reconstruction, and replacement tests}

Report
\[
d_{\mathrm{id}}(W_t,W_{t+\tau}),
\qquad
\Rec_k,
\qquad
 e_t,
\qquad
\alpha^{(k)},
\qquad
\mu_t^{\mathrm{keep}}.
\]
The report must separate:
\begin{itemize}
\item subspace tracking from state tracking;
\item coherence from non-collapse;
\item reconstruction from mere prediction;
\item repair from successor-workspace creation; and
\item structural replacement from causal and reconstructive replacement.
\end{itemize}

\subtopic{74.8 Formation-provenance tests}

Report
\[
\mathbf E_{\mathrm{form}},
\qquad
\delta^{\mathrm{scaf}},
\qquad
\delta^{\mathrm{leak}},
\qquad
\Delta_q^{\mathrm{hist}}.
\]
Include the external scaffold, information available to the learner, history
shuffles, cross-scaffold retraining, and every route by which certification
labels or target boundaries could have entered formation.

\subtopic{74.9 Learned-structure and boundary tests}

Report
\[
\|\widehat H-H_*\|,
\qquad
\Delta_J,
\qquad
\Delta_t^B,
\qquad
\COut_\tau(\widehat S),
\qquad
\CIn_\tau(\widehat S).
\]
Boundary reports must include the admissible candidate family, full-system and
empty-set controls, singleton and centralized controls, score components,
margin to the runner-up, and temporal stability under held-out environment
drift.

\subtopic{74.10 Goal and counterfactual-policy tests}

Report the agency coordinates
\[
\eta_G(\tau),
\qquad
\delta_G^{\mathrm{cue}}(\tau),
\qquad
\eta_{G\to A},
\qquad
\delta_Q,
\qquad
\eta_{V\to A},
\qquad
\eta_{A\to Y},
\qquad
 k_{\mathrm{ctl},\varepsilon}.
\]
The required experiments include:
\begin{enumerate}
\item cue-controlled goal interventions;
\item goal-channel ablation and permutation;
\item counterfactual-value calibration against action interventions;
\item value-label permutation among feasible actions;
\item critic or value-channel blocking;
\item action replay under held-out perturbations;
\item policy regret or value-gap estimation; and
\item coalition reconstruction of the complete control state.
\end{enumerate}

\subtopic{74.11 Damage and repair tests}

For every damage family, report
\[
\eta_{D\to M},
\qquad
\eta_{M\to R},
\qquad
\eta_{R\to v},
\qquad
\Delta_{\mathrm{repair}}^{\mathrm{spec}},
\qquad
\delta_{\mathrm{rep-id}},
\qquad
\sup_t e_t.
\]
The repeated-disturbance experiment must contain:
\begin{enumerate}
\item a no-damage baseline;
\item damage with repair disabled;
\item damage with the intact repair loop;
\item monitor ablation;
\item cost-matched random repair;
\item off-workspace repair;
\item repeated and compound damage; and
\item a held-out disturbance distribution.
\end{enumerate}

\topic{75}{Required Null Families}

The core null panel is:
\[
\begin{aligned}
N_1&=\text{shuffled transfer maps with validated local matching},\\
N_2&=\text{shuffled cross-node temporal correspondence},\\
N_3&=\text{matched-capacity local memory},\\
N_4&=\text{matched-size irrelevant artifact},\\
N_5&=\text{participant ablation},\\
N_6&=\text{artifact ablation},\\
N_7&=\text{workspace implementation intervention},\\
N_8&=\text{matched off-workspace intervention},\\
N_9&=\text{representation replacement},\\
N_{10}&=\text{context reset and reconstruction},\\
N_{11}&=\text{matched formation-history shuffle},\\
N_{12}&=\text{cross-scaffold re-encoding and retraining},\\
N_{13}&=\text{boundary-label and certification leakage audit},\\
N_{14}&=\text{empty, full, singleton, and centralized boundary nulls},\\
N_{15}&=\text{held-out transfer and membership re-estimation},\\
N_{16}&=\text{boundary damage, use, and repair intervention},\\
N_{17}&=\text{goal-code permutation with cues and actions matched},\\
N_{18}&=\text{persistent external-cue rewrite with internal retention cut},\\
N_{19}&=\text{counterfactual-value label permutation},\\
N_{20}&=\text{critic or value-channel ablation},\\
N_{21}&=\text{open-loop action replay under held-out perturbation},\\
N_{22}&=\text{monitor-blind cost-matched random repair},\\
N_{23}&=\text{off-workspace repair with matched intervention norm},\\
N_{24}&=\text{identity-breaking reconstruction with matched task score},\\
N_{25}&=\text{merge control for components sharing an uninterruptible controller},\\
N_{26}&=\text{secret-sharing split of a known centralized controller},\\
N_{27}&=\text{hidden shared-memory and common-cue exposure test},\\
N_{28}&=\text{frozen certificate sweep across the admissible decomposition family},\\
N_{29}&=\text{post hoc target-selection audit},\\
N_{30}&=\text{horizon cherry-picking audit},\\
N_{31}&=\text{loss and utility substitution audit},\\
N_{32}&=\text{restricted intervention-family audit},\\
N_{33}&=\text{threshold and gate-set tuning audit},\\
N_{34}&=\text{frozen certificate sweep across the admissible specification family}.
\end{aligned}
\]

\subtopic{75.1 Null validity record}

Each observational null must state:
\begin{itemize}
\item the local distributions it preserves;
\item the relational dependency it destroys;
\item the distance between intact and null local marginals;
\item the capacity, parameter count, and information budget; and
\item whether null fitting used the same data budget as the intact model.
\end{itemize}
Each causal null must state:
\begin{itemize}
\item the manipulated mechanism;
\item the matched background and intervention cost;
\item the intended causal path destroyed or preserved;
\item the manipulation check; and
\item unintended direct effects detected during validation.
\end{itemize}
A failed manipulation check invalidates the null. It does not count as evidence
for the workspace claim.

\subtopic{75.2 Claim-specific null subsets}

A successful claim requires
\[
\Pass(N_1,\ldots,N_{28})
\]
only for the subset relevant to its stated level. The subset must be frozen
before certification. Failed applicable nulls, invalid nulls, and unmeasured
nulls are retained separately in the final report. They must not be merged
into a generic ``not significant'' category.

\topic{76}{Certification Vector and Complete Reporting Templates}

\subtopic{76.1 Extended certification vector}

Use a vector rather than one aggregate score:
\[
\boxed{
\begin{aligned}
\mathcal C_{\mathrm{full}}(W)
=\bigl(&
 d_\varepsilon,
 g_\varepsilon,
 \delta_{\mathrm{coh}},
 \delta_{\mathrm{align}}^{\mathcal I},
 \Delta_{\mathrm{joint}}^\ell,
 k_\varepsilon,\\
&\boldsymbol\Delta_{\mathrm{tr}}(\mathcal Q),
 \eta_{\mathrm{causal}}^\nu,
 \Delta_{\mathrm{specific}}(r),
 \eta_{\mathrm{loop}}(\tau),\\
&d_{\mathrm{id}},
 \Rec_k,
 \underline{\mathbf m}_{\mathfrak D},
 \boldsymbol\Delta_{\mathrm{dec}},
 \mathbf E_{\mathrm{form}},
 \delta^{\mathrm{scaf}},
 \delta^{\mathrm{leak}},\\
&\Delta_J,
 \Delta_t^B,
 \COut_\tau(\widehat S),
 \CIn_\tau(\widehat S),\\
&\eta_G(\tau),
 \delta_G^{\mathrm{cue}}(\tau),
 \eta_{G\to A},
 \delta_Q,
 \eta_{V\to A},
 \eta_{A\to Y},\\
&\eta_{D\to M},
 \eta_{M\to R},
 \eta_{R\to v},
 \Delta_{\mathrm{repair}}^{\mathrm{spec}},
 \delta_{\mathrm{rep-id}},
 \sup_t e_t
\bigr).
\end{aligned}
}
\]
A claim passes only when every required coordinate lies in its predeclared
admissible region. No weighted average may compensate for a failed required
gate.

\subtopic{76.2 Claim card template}

\begin{longtable}{@{}>{\raggedright\arraybackslash}p{0.29\textwidth}>{\raggedright\arraybackslash}p{0.65\textwidth}@{}}
\toprule
Field & Required entry \\
\midrule
\endfirsthead
\toprule
Field & Required entry \\
\midrule
\endhead
Claim identifier & Unique versioned identifier tied to code and result
artifacts. \\
Claim level & Archive, shared representation, operational workspace,
relational workspace, persistent workspace, self-maintaining workspace, or
workspace agent. \\
Reference decomposition & Nodes, artifacts, channels, local-access maps, and coalition universe. \\
Admissible decomposition family & Plausible merges, splits, shared-resource assignments, exclusions, and audit rationale. \\
Target and horizon & \(Y^+\), loss, utility, prediction horizon, control
horizon, and maintenance horizon. \\
Environment scope & Training, certification, and stress-test environment
families. \\
Workspace extractor & Frozen construction of \(D_H\), \(L_H\), spectral
threshold, and state readout. \\
Goal and policy & Goal codes, utility normalization, action set, and policy
class. \\
Boundary & Candidate family, score, admissibility constraints, and selected
boundary. \\
Interventions & Feasible workspace, goal, action, damage, and repair
implementations. \\
Null panel & Applicable \(N_j\), preservation targets, and fit budgets. \\
Thresholds & Every gate direction, numerical cutoff, and uncertainty rule. \\
Data split & Formation, calibration, certification, stress, and cross-fitting
logic. \\
Stopping rule & Number of seeds, runs, interventions, and any sequential
stopping criterion. \\
Status & Confirmatory, exploratory, replication, or audit. \\
\bottomrule
\end{longtable}

\subtopic{76.3 Gate-estimate table template}

For every required coordinate, include one row with the following fields:

\begin{longtable}{@{}>{\raggedright\arraybackslash}p{0.20\textwidth}>{\raggedright\arraybackslash}p{0.73\textwidth}@{}}
\toprule
Field & Required entry \\
\midrule
\endfirsthead
\toprule
Field & Required entry \\
\midrule
\endhead
Gate & Formal symbol and verbal interpretation. \\
Estimator & Exact statistic, normalization, and software implementation. \\
Evaluation unit & Episode, participant, graph, seed, environment, or other
independent unit. \\
Point estimate & Estimate on the frozen certification split. \\
Uncertainty & Interval, bound, or posterior region with declared coverage or
credibility. \\
Threshold & Predeclared pass region. \\
Manipulation status & Passed, failed, or not applicable. \\
Decision & Pass, fail, invalid, underpowered, or unmeasured. \\
Sensitivity & Result under reasonable alternative normalizations and
resampling choices. \\
Artifact reference & File, table, raw-result key, and code commit. \\
\bottomrule
\end{longtable}

\subtopic{76.4 Intervention ledger template}

\begin{longtable}{@{}p{0.20\textwidth}p{0.73\textwidth}@{}}
\toprule
Field & Required entry \\
\midrule
\endfirsthead
\toprule
Field & Required entry \\
\midrule
\endhead
Intervention ID & Unique identifier and randomization seed. \\
Target mechanism & Workspace coordinate, goal code, value channel, action,
boundary, monitor, or repair channel. \\
Implementation & Exact code path, hook, artifact edit, or mechanism change. \\
Dose & Norm, duration, timing, and support of the intervention. \\
Background match & Variables fixed, blocked, randomized, or balanced. \\
Excluded paths & Direct paths intentionally removed from the causal contrast. \\
Manipulation check & Evidence that the intended mechanism changed by the
specified amount. \\
Collateral effects & Changes outside the target mechanism. \\
Outcome & Downstream variable and evaluation window. \\
Effect estimate & Point estimate, uncertainty, and sign convention. \\
Cost & Compute, latency, information, action, or structural cost. \\
Validity status & Valid, invalid, ambiguous, or unmeasured. \\
\bottomrule
\end{longtable}

\subtopic{76.5 Null ledger template}

\begin{longtable}{@{}p{0.20\textwidth}p{0.73\textwidth}@{}}
\toprule
Field & Required entry \\
\midrule
\endfirsthead
\toprule
Field & Required entry \\
\midrule
\endhead
Null ID & One of \(N_1,\ldots,N_{34}\) or a declared extension. \\
Claim attacked & The exact explanation the null is designed to preserve or
destroy. \\
Matched quantities & Local marginals, capacity, training budget, action cost,
intervention norm, or other controlled quantities. \\
Destroyed relation & Edge alignment, temporal order, value ranking, goal
identity, monitor mediation, or another relational feature. \\
Fit procedure & Data, optimization, seeds, and hyperparameter budget. \\
Validity metric & Distance from the intended matched quantities. \\
Primary contrast & Intact minus null effect with uncertainty. \\
Disposition & Rejects null, fails to reject, invalid null, or underpowered. \\
Failure notes & Any mismatch or collateral mechanism change. \\
\bottomrule
\end{longtable}

\subtopic{76.6 Agency and maintenance report template}

The agency section must provide the following matrix:

\begin{center}
\small
\begin{tabularx}{\textwidth}{@{}p{0.22\textwidth}p{0.23\textwidth}X@{}}
\toprule
Component & Primary metric & Required contrast \\
\midrule
Goal semantics & Held-out \(Q_t^g\) ordering & Distinct goals reverse or alter
at least one declared action ordering. \\
Goal retention & \(\eta_G(\tau)\) & Goal intervention with future cues matched
or removed. \\
Goal use & \(\eta_{G\to A}\) & Goal intervention versus goal-label replay. \\
Counterfactual accuracy & \(\delta_Q\) & Internal value versus external action
interventions. \\
Value use & \(\eta_{V\to A}\) & Value-label permutation and value-channel
ablation. \\
Outcome control & \(\eta_{A\to Y}\) & Feasible action intervention. \\
Damage detection & \(\eta_{D\to M}\) & Damage versus matched no-damage. \\
Repair mediation & \(\eta_{M\to R}\) & Monitor intervention and monitor
ablation. \\
Restoration & \(\eta_{R\to v}\) & Repair versus matched no-repair. \\
Repair specificity & \(\Delta_{\mathrm{repair}}^{\mathrm{spec}}\) & Intended
repair versus cost-matched generic action. \\
Identity preservation & \(\delta_{\mathrm{rep-id}}\) & Pre-damage workspace
versus post-repair workspace. \\
Long-run maintenance & \(\sup_t e_t\) & Repeated held-out disturbances. \\
Relational location & \(k_{\mathrm{ctl},\varepsilon}\) & Full system versus
all single nodes and matched local nulls. \\
\bottomrule
\end{tabularx}
\end{center}

\subtopic{76.7 Uncertainty and multiplicity}

The independent resampling unit must match the data-generating process. Token,
time-step, or edge resampling is invalid when dependence is organized at the
episode, participant, graph, or environment level. Recommended procedures
include:
\begin{itemize}
\item environment- or episode-level bootstrap;
\item repeated independent training seeds;
\item cluster-robust intervals for repeated observations;
\item permutation tests that preserve the declared local structure;
\item concentration bounds for uniform finite candidate families; and
\item stress-test intervals over held-out environment families.
\end{itemize}

Pass decisions should use one-sided uncertainty bounds. For a lower-bounded
gate \(\theta\geq\tau\), pass only when the lower confidence bound exceeds
\(\tau\). For an upper-bounded defect \(\delta\leq\varepsilon\), pass only when
the upper confidence bound is below \(\varepsilon\).

When a conclusion ranges over \(\mathfrak S\), the uncertainty rule must be
uniform over every gate and specification included in the claim. Report the
pointwise interval, the family-wide lower envelope, the covering or
multiplicity correction, and the identification status. A narrow pointwise
interval does not establish a robust result when model or specification
uncertainty remains wide.

\subtopic{76.8 \(T_{24}\): Gatewise certification validity }

\begin{theorem}[Intersection-gate false-certification control]
\label{thm:gatewise-validity}
Let a certificate require all gates \(G_1,\ldots,G_m\) to pass. For each gate
\(j\), let \(H_{0j}\) denote the null that the required property is absent, and
assume the predeclared test satisfies
\[
\sup_{P\in H_{0j}}
\mathbb P_P(G_j\text{ passes})
\leq
\alpha_j.
\]
The global null is
\[
H_0
=
\bigcup_{j=1}^m H_{0j},
\]
meaning that at least one required property is absent. If the certificate is
declared only when every gate passes, then
\[
\boxed{
\sup_{P\in H_0}
\mathbb P_P(\text{certificate passes})
\leq
\max_{1\leq j\leq m}\alpha_j.
}
\]
No independence assumption among gate tests is required.
\end{theorem}

\begin{proof}
For any \(P\in H_0\), choose an index \(j_P\) such that \(P\in H_{0j_P}\).
The event that every gate passes is contained in the event that gate \(j_P\)
passes. Hence
\[
\mathbb P_P(\text{certificate passes})
\leq
\mathbb P_P(G_{j_P}\text{ passes})
\leq
\alpha_{j_P}
\leq
\max_j\alpha_j.
\]
Taking the supremum over \(P\in H_0\) proves the claim.
\end{proof}

The theorem is an intersection-union guarantee. It applies only when the gate
family, thresholds, and tests are fixed before certification and each gate test
is uniformly valid for its null. Adaptive gate deletion, selective reporting,
or repeated reuse of the certification split invalidates the bound.

\subtopic{76.9 Failure taxonomy}

Every gate and null receives exactly one primary status:
\begin{description}
\item[Pass.] The frozen estimate and uncertainty rule satisfy the threshold.
\item[Fail.] The gate was validly measured and did not satisfy the threshold.
\item[Invalid intervention.] The manipulation check or background match
 failed.
\item[Invalid null.] The null did not preserve its required local quantities
 or exceeded its fit budget.
\item[Underpowered.] The uncertainty interval crosses the threshold by more
 than the declared indifference region.
\item[Unmeasured.] No valid estimate was obtained.
\item[Leakage.] Formation or model selection used certification information.
\item[Exploratory.] The metric, threshold, or analysis was selected after
 inspecting certification outcomes.
\end{description}
Only \emph{Pass} counts toward the formal certificate. ``Underpowered'' and
``unmeasured'' are not failures of the underlying theory, but they are failures
to establish the corresponding empirical claim.

\subtopic{76.10 Machine-readable result record}

The machine-readable record should contain at least the following top-level
objects:
\[
\boxed{
\begin{aligned}
\mathsf{ResultRecord}
=\{&
\mathsf{claim},
\mathsf{specification},
\mathsf{specification\_family},
\mathsf{decomposition},
\mathsf{data\_splits},
\mathsf{formation},\\
&\mathsf{workspace\_extractor},
\mathsf{goals},
\mathsf{actions},
\mathsf{interventions},\\
&\mathsf{nulls},
\mathsf{gates},
\mathsf{identified\_sets},
\mathsf{specification\_envelopes},
\mathsf{uncertainty},
\mathsf{failures},
\mathsf{artifacts}
\}.
\end{aligned}
}
\]
Each gate object records its symbol, estimator, direction, threshold, point
estimate, interval, split, seed set, status, and artifact reference. Each
intervention and null object records its manipulation check and validity
status. The schema must preserve failed and invalid entries rather than only
the passing subset.

\subtopic{76.11 Minimum reproducibility package}

A complete certification package contains:
\begin{enumerate}
\item the frozen claim card;
\item source code and environment lock file;
\item exact data or deterministic data-generation procedure;
\item formation, calibration, certification, and stress-test manifests;
\item learned transfer, workspace, policy, monitor, and repair parameters;
\item raw per-run and per-intervention results;
\item the intervention and null ledgers;
\item the complete gate table, including failures;
\item figures and tables regenerated from raw results;
\item random seeds and hardware or numerical-precision information;
\item checksums for every included artifact; and
\item a single verification command that reproduces package integrity and the
 reported summary tables.
\end{enumerate}

\subtopic{76.12 Decision rule}

Let \(\mathcal J(L)\) be the required gate set for claim level \(L\). Then
\[
\boxed{
\operatorname{Certify}_L(W;\varsigma)
=
\bigwedge_{j\in\mathcal J(L,\varsigma)}
\operatorname{Pass}_j(W;\varsigma).
}
\]
For a workspace-agent claim, the required set includes every structural,
operational, relational, persistent, goal, counterfactual-control, and
self-maintenance gate declared in Parts~IX and XVI. A specification-robust
claim additionally requires the uniform rule of Part~XVIII across
\(\mathfrak S\). Failure of an earlier level
prevents certification of a later level, even when later-level behavior appears
strong on a task benchmark.

\theorypart{XXII}{Integrated Worked Examples and Case Studies}

\topic{77}{Case-Study Method and Common Gate Card}

The preceding parts define a large family of structural, operational, causal,
relational, persistent, formative, and agentive predicates. A worked example
is useful only when it makes clear which of those predicates are established,
which are assumed, and which remain unmeasured. This part therefore uses a
single case-study protocol for three different kinds of realization:
\begin{enumerate}
\item an exact finite-dimensional model of a human--human theorem laboratory;
\item an idealized human--AI--artifact--executor research loop, followed by a
 deliberately conservative audit of the collaboration that motivated the
 theory; and
\item a closed-form multi-agent control system that reaches the formal agency
 level inside a small declared domain.
\end{enumerate}
The first and third cases are analytic models. Their conclusions hold inside
those models. The second contains both an analytic idealization and an
observational audit. The audit is not upgraded to a certificate merely because
the idealization is mathematically clean.

\begin{definition}[Case-study realization]\label{def:case-study-realization}
A case-study realization is a tuple
\[
\mathfrak K
=
\bigl(
\mathfrak R,
\mathfrak B,
\mathfrak C,
\mathfrak P,
\mathcal E,
\mathfrak D,
\mathfrak S,
\mathcal I,
\mathcal Q,
\mathcal G
\bigr),
\]
where:
\begin{itemize}
\item \(\mathfrak R\) is the declared relational realization;
\item \(\mathfrak B\) is the boundary problem and component decomposition;
\item \(\mathfrak C\) is the structural core--transfer model;
\item \(\mathfrak P\) is the downstream predictive or control problem;
\item \(\mathcal E\) is the evidence or generated-data family;
\item \(\mathfrak D\) is the admissible decomposition family;
\item \(\mathfrak S\) is the admissible family of complete scientific
 specifications, including targets, losses, horizons, interventions, nulls,
 and thresholds;
\item \(\mathcal I\) is the feasible intervention family;
\item \(\mathcal Q\) is the matched null family; and
\item \(\mathcal G\) is the set of gates being claimed.
\end{itemize}
A conclusion is always relative to this tuple. Changing the downstream target,
temporal horizon, intervention support, component decomposition, null family,
or threshold rule creates a different specification of the case. A stronger
case study reports whether the conclusion survives the predeclared family
\(\mathfrak S\) and whether the relevant gates are identified by the available
experiment family.
\end{definition}

\subtopic{77.1 Three evidential statuses}

Every claim in this part receives one of three primary statuses:
\begin{description}
\item[Analytic.] The gate follows exactly from a declared mathematical model.
\item[Empirically certified.] The gate is estimated from held-out data using
 the protocol of Part~XXI and passes its frozen uncertainty rule.
\item[Observational only.] The case contains suggestive records or natural
 perturbations, but no valid held-out intervention-and-null certificate.
\end{description}
An analytic result is not automatically an empirical result. An observational
pattern is not automatically a causal result. None of the three statuses
establishes phenomenal consciousness.

\subtopic{77.2 Common gate card}

For each case we report the vector
\[
\boxed{
\mathsf{Card}(\mathfrak K)
=
\left(
\mathsf S,
\mathsf O,
\mathsf R,
\mathsf P,
\mathsf F,
\mathsf A,
\mathsf D,
\mathsf X,
\mathsf C
\right),
}
\]
where:
\[
\begin{array}{ll}
\mathsf S & = \text{structural core and spectral coherence},\\
\mathsf O & = \text{operational alignment and causal efficacy},\\
\mathsf R & = \text{relational irreducibility and matched-null gap},\\
\mathsf P & = \text{persistence, reconstruction, and identity},\\
\mathsf F & = \text{formation provenance and boundary status},\\
\mathsf A & = \text{goal-conditioned agency and self-maintenance},\\
\mathsf D & = \text{decomposition robustness and interface audit},\\
\mathsf X & = \text{specification robustness, identification, and uniform uncertainty},\\
\mathsf C & = \text{consciousness or personhood claim}.
\end{array}
\]
The decomposition coordinate distinguishes a result established only for the
declared interface from one stable across an admissible family. The
specification coordinate records whether the same conclusion survives
scientifically defensible alternatives, whether its sign is identified by the
available experiments, and whether its lower margin remains positive after
uniform finite-sample uncertainty. The final coordinate is included to make
the non-implication explicit. In all examples below,
\[
\mathsf C=\text{not established}.
\]

\subtopic{77.3 Claim-level discipline}

Let \(L(\mathfrak K)\in\{0,1,\ldots,6\}\) denote the highest level from
Part~XXIV whose complete required gate set is established. The level is
conjunctive and pointwise in the declared scientific specification unless the
\(\mathsf D\) and \(\mathsf X\) coordinates also pass:
\[
L(\mathfrak K)\geq \ell
\quad\Longleftrightarrow\quad
\bigwedge_{j\in\mathcal J(\ell)}\Pass_j(\mathfrak K).
\]
A case may have strong higher-level-looking behavior while remaining at a lower
certified level because an earlier gate is absent. For example, successful
repair behavior does not establish self-maintenance when no matched no-repair
counterfactual is available, and coherent shared documents do not establish a
relational workspace when one participant contains the complete operative
state.

\subtopic{77.4 Structural calculations used repeatedly}

The positive examples use a recurring distributed code. Let
\[
c=(c_1,c_2,c_3)^{\mathsf T}\in\mathbb R^3
\]
and define three local encodings
\[
E_1c=(c_1,c_2),
\qquad
E_2c=(c_2,c_3),
\qquad
E_3c=(c_1,c_3).
\]
Then
\[
\sum_{i=1}^3E_i^*E_i=2I_3.
\]
No individual encoding is injective, while every two-node coalition is
injective. The pairwise frame operators are
\[
\begin{aligned}
E_{\{1,2\}}^*E_{\{1,2\}}&=\operatorname{diag}(1,2,1),\\
E_{\{1,3\}}^*E_{\{1,3\}}&=\operatorname{diag}(2,1,1),\\
E_{\{2,3\}}^*E_{\{2,3\}}&=\operatorname{diag}(1,1,2).
\end{aligned}
\]
Hence
\[
\alpha_{\{i,j\}}=1
\qquad(i\neq j),
\]
and one-node erasure recovery is exact. Every nonzero core vector is visible
at at least two nodes, so
\[
d_{\mathrm{er}}=2.
\]
This code is distributed but redundant: no node is sufficient, yet any one
node may be lost without destroying the core.

\subtopic{77.5 Binary synergy test}

Let \(C_1,C_2,C_3\) be independent uniform bits and let
\[
Y=C_1\oplus C_2\oplus C_3.
\]
Each node sees two of the three bits. Therefore
\[
H(Y\mid E_iC)=1\text{ bit}
\qquad\text{for every }i,
\]
while
\[
H(Y\mid E_1C,E_2C,E_3C)=0.
\]
Under logarithmic loss,
\[
\Risk_{\log}(\{i\})=1,
\qquad
\Risk_{\log}(V)=0,
\qquad
\Delta_{\mathrm{joint}}^{\log}=1\text{ bit},
\qquad
k_0=2.
\]
Moreover, \(Y\) is independent of every individual local state. The null law
that keeps all three core bits independent and makes \(Y\) an independent
uniform bit preserves every pair \((Y,E_iC)\) while destroying the joint parity
relation. It is therefore a predictive-marginal-matched relational null with
\[
\Delta_{\mathrm{tr}}=1\text{ bit}.
\]
This test will be reused with different semantic interpretations.

\topic{78}{Case Study A: A Distributed Human--Human Theorem Laboratory}

\subtopic{78.1 Declared realization}

Consider two researchers and one persistent notebook:
\[
V_{\mathrm{HH}}=\{P,C,N\},
\]
where:
\begin{itemize}
\item \(P\) is a proposer who tracks the theorem statement and open proof
 obligations;
\item \(C\) is a critic who tracks proof obligations and falsification status;
\item \(N\) is a notebook that tracks the theorem statement and falsification
 status.
\end{itemize}
The abstract core is
\[
c=(q,p,f)^{\mathsf T}\in\mathbb R^3,
\]
with
\[
\begin{aligned}
q&=\text{statement coordinate},
&
p&=\text{proof-obligation coordinate},\\
f&=\text{falsification-status coordinate}.
\end{aligned}
\]
The accessible encodings are
\[
E_P=
\begin{pmatrix}
1&0&0\\
0&1&0
\end{pmatrix},
\quad
E_C=
\begin{pmatrix}
0&1&0\\
0&0&1
\end{pmatrix},
\quad
E_N=
\begin{pmatrix}
1&0&0\\
0&0&1
\end{pmatrix}.
\]
The semantics are stylized. The point is not that real researchers contain
only two numbers. The point is to exhibit, without metaphor, a system in which
no participant or artifact contains the entire operative state while the
coupled system does.

\subtopic{78.2 Exact overlap descent}

Order the local coordinates as
\[
x=(q_P,p_P,p_C,f_C,q_N,f_N)^{\mathsf T}.
\]
Compatibility requires
\[
p_P=p_C,
\qquad
f_C=f_N,
\qquad
q_N=q_P.
\]
The discrepancy operator is
\[
D_{\mathrm{HH}}
=
\begin{pmatrix}
0&1&-1&0&0&0\\
0&0&0&1&0&-1\\
-1&0&0&0&1&0
\end{pmatrix}.
\]
Therefore
\[
\ker D_{\mathrm{HH}}
=
\left\{
(q,p,p,f,q,f)^{\mathsf T}:q,p,f\in\mathbb R
\right\}
\cong\mathbb R^3.
\]
The rows are orthogonal with squared norm two, so
\[
D_{\mathrm{HH}}D_{\mathrm{HH}}^*=2I_3.
\]
Consequently the transfer Laplacian has spectrum
\[
\boxed{
\spec(L_{\mathrm{HH}})
=
\{0,0,0,2,2,2\},
}
\]
and the exact workspace gap is
\[
\gamma_{\mathrm{HH}}=2.
\]
The low-energy subspace is therefore stable under any Laplacian perturbation
whose operator norm is strictly below one, with the quantitative projector
bound from Theorem~\ref{thm:spectral-stability}.

\subtopic{78.3 Reconstruction and one-node failure}

The full frame operator is
\[
F_{PCN}=E_P^*E_P+E_C^*E_C+E_N^*E_N=2I_3.
\]
For every two-node coalition \(S\),
\[
\lambda_{\min}(F_S)=1.
\]
Thus every two-node coalition reconstructs the complete core and, for noisy
observations \(y_S=A_Sc+\eta_S\),
\[
\|\widehat c-c\|
\leq
\|\eta_S\|.
\]
No individual node reconstructs \(c\), because
\[
\operatorname{rank}(E_P)
=
\operatorname{rank}(E_C)
=
\operatorname{rank}(E_N)
=2<3.
\]
The realization is therefore simultaneously:
\[
\text{individually incomplete}
\quad\text{and}\quad
\text{robust to any one-node loss}.
\]

\subtopic{78.4 Relational prediction problem}

Let \(q,p,f\) now be independent uniform bits and define the downstream
acceptance decision
\[
Y=q\oplus p\oplus f.
\]
The interpretation is abstract: acceptance depends on the statement, the state
of the proof obligations, and the current falsification record. No individual
node has sufficient information:
\[
H(Y\mid X_P)
=
H(Y\mid X_C)
=
H(Y\mid X_N)
=1.
\]
Every two-node coalition has all three bits and therefore zero conditional
entropy. Hence
\[
\boxed{
\Delta_{\mathrm{joint}}^{\log}=1\text{ bit},
\qquad
k_0=2.
}
\]
Let \(Q_{\mathrm{HH}}\) preserve the independent uniform local bit pairs but
replace \(Y\) by an independent uniform bit. Because parity is independent of
every two-bit local view, \(Q_{\mathrm{HH}}\) preserves each
\((Y,X_i)\)-marginal. Nevertheless,
\[
\Risk_{\log,Q_{\mathrm{HH}}}(V)
-
\Risk_{\log,P_{\mathrm{HH}}}(V)
=1\text{ bit}.
\]
This is an exact relational witness: the predictive capacity lies in the joint
organization rather than in any local predictive marginal.

\subtopic{78.5 Feasible intervention and specificity}

A feasible statement intervention must alter both local realizations of the
statement coordinate:
\[
\iota_{q\leftarrow a}:
(q_P,q_N)\leftarrow(a,a),
\]
while holding \((p,f)\) fixed. For configurations with fixed \(p\) and \(f\),
changing \(a\) flips the downstream decision with total-variation effect one:
\[
\TV
\left(
P(Y\mid\iota_{q\leftarrow0}),
P(Y\mid\iota_{q\leftarrow1})
\right)
=1.
\]
A matched intervention on a residual notebook-format coordinate leaves \(Y\)
unchanged. Thus the core intervention has positive causal specificity in this
model.

\subtopic{78.6 Persistence, handoff, and replacement}

Suppose the proposer becomes unavailable. The remaining coalition \(\{C,N\}\)
reconstructs \((q,p,f)\) exactly. A replacement proposer \(P'\) can therefore
be initialized by
\[
E_{P'}c=(q,p).
\]
If \(P'\) uses a different local coordinate system
\[
J_{P'}\in GL(2),
\]
then the incoming and outgoing overlap maps are conjugated accordingly. The
global section and all downstream kernels remain unchanged. This is exact
replacement invariance, not literal carrier identity.

If the notebook alone survives, the core cannot be reconstructed because the
proof-obligation coordinate is missing. The model therefore predicts a sharp
handoff boundary:
\[
\text{one surviving node is insufficient},
\qquad
\text{any two surviving nodes are sufficient}.
\]

\subtopic{78.7 Case card}

Inside the declared finite model,
\[
\boxed{
\begin{aligned}
\mathsf S&=\checkmark_{\mathrm{analytic}},
&\mathsf O&=\checkmark_{\mathrm{analytic}},\\
\mathsf R&=\checkmark_{\mathrm{analytic}},
&\mathsf P&=\checkmark_{\mathrm{analytic}},\\
\mathsf F&=\text{scaffolded},
&\mathsf A&=\text{not tested},\\
\mathsf D&=\text{fixed decomposition},
&\mathsf C&=\text{not established}.
\end{aligned}
}
\]
The highest supported level is Level~4 when the handoff dynamics and identity
correspondence are included. The example does not contain an operational goal
state, counterfactual policy, damage monitor, or repair loop, so it does not
support Levels~5 or 6.

\topic{79}{Case Study B: A Human--AI--Artifact--Executor Research Loop}

\subtopic{79.1 Idealized four-node loop}

Consider four heterogeneous components:
\[
V_{\mathrm{HAI}}=\{H,M,X,A\},
\]
where:
\begin{itemize}
\item \(H\) is a human researcher who maintains goals and semantic direction;
\item \(M\) is a language-model process that constructs local formalizations;
\item \(X\) is an execution and verification process that compiles, tests, or
 falsifies candidate artifacts; and
\item \(A\) is a persistent artifact store containing exact source, results,
 and versioned handoff state.
\end{itemize}
Let the abstract research state be
\[
c=(g,d,p,v)^{\mathsf T}\in\mathbb R^4,
\]
where:
\[
\begin{array}{ll}
g=\text{goal and acceptance criterion},
&d=\text{current derivation or theory state},\\
p=\text{open proof or implementation obligations},
&v=\text{verification and regression status}.
\end{array}
\]
Define local encodings
\[
E_Hc=(g,d),
\qquad
E_Mc=(g,p),
\qquad
E_Xc=(p,v),
\qquad
E_Ac=(d,v).
\]
Each coordinate is represented at exactly two nodes, and the incidence pattern
forms a four-cycle.

\subtopic{79.2 Exact structural calculation}

Order the local coordinates as
\[
x=(g_H,d_H,g_M,p_M,p_X,v_X,d_A,v_A)^{\mathsf T}.
\]
The overlap discrepancies are
\[
g_H-g_M,
\qquad
p_M-p_X,
\qquad
v_X-v_A,
\qquad
d_A-d_H.
\]
Thus
\[
D_{\mathrm{HAI}}
=
\begin{pmatrix}
1&0&-1&0&0&0&0&0\\
0&0&0&1&-1&0&0&0\\
0&0&0&0&0&1&0&-1\\
0&-1&0&0&0&0&1&0
\end{pmatrix}.
\]
The rows are orthogonal with squared norm two. Therefore
\[
\ker D_{\mathrm{HAI}}
\cong\mathbb R^4,
\]
\[
\boxed{
\spec(L_{\mathrm{HAI}})
=
\{0,0,0,0,2,2,2,2\},
\qquad
\gamma_{\mathrm{HAI}}=2.
}
\]
The full accessible frame operator is
\[
F_{HMXA}=2I_4.
\]
No single node is injective. Every three-node coalition reconstructs the full
state, and the two opposite pairs
\[
\{H,X\},
\qquad
\{M,A\}
\]
are already sufficient. Adjacent pairs observe only three of the four
coordinates. Hence
\[
k_{\mathrm{rec}}=2,
\qquad
d_{\mathrm{er}}=2,
\]
while the identity of the sufficient two-node coalition depends on the graph
geometry.

\subtopic{79.3 Relational parity benchmark}

Let \(g,d,p,v\) be independent uniform bits and define
\[
Y=g\oplus d\oplus p\oplus v.
\]
Every single node sees only two bits, so
\[
H(Y\mid X_i)=1
\qquad\text{for all }i.
\]
The opposite pairs \(\{H,X\}\) and \(\{M,A\}\) determine all four bits and
therefore determine \(Y\). Every adjacent pair remains one bit short. It
follows that
\[
\boxed{
\Delta_{\mathrm{joint}}^{\log}=1\text{ bit},
\qquad
k_0=2.
}
\]
As in the preceding case, parity is independent of every individual local
view. Replacing \(Y\) by an independent uniform bit produces a matched local
predictive null and a full-system transfer gap of one bit.

This benchmark is intentionally narrow. It demonstrates that the four-cycle
can realize exact relational information. It does not show that every real
human--AI collaboration has this encoding, this target, or this degree of
irreducibility.

\subtopic{79.4 A process-level interpretation}

The idealized loop can be read as the recurrent sequence
\[
\boxed{
H\xrightarrow{g}M
\xrightarrow{p}X
\xrightarrow{v}A
\xrightarrow{d}H.
}
\]
A more explicit cycle is:
\begin{enumerate}
\item the human supplies a goal and acceptance criterion;
\item the model converts that goal into a formal derivation and explicit
 obligations;
\item the executor compiles, tests, or searches for counterexamples;
\item the artifact store preserves the exact source and verification record;
\item the human evaluates the returned derivation and evidence, revising the
 next goal.
\end{enumerate}
The operative state is reentrant because products of one cycle become inputs to
the next. The artifact is not merely a record when its contents alter the next
formalization, test, or acceptance decision.

\subtopic{79.5 Conservative audit of the motivating collaboration}

The collaboration that motivated this theory has several features consistent
with the four-node description: specialized shared vocabulary, versioned
LaTeX and code artifacts, repeated compile-and-repair cycles, explicit theorem
and null-control demands, and reconstruction across conversations from handoff
packages. Those observations justify treating the collaboration as a serious
candidate case. They do not establish the complete certificate.

\begin{center}
\small
\begin{tabularx}{\textwidth}{@{}>{\raggedright\arraybackslash}p{0.22\textwidth}>{\raggedright\arraybackslash}p{0.18\textwidth}>{\raggedright\arraybackslash}X@{}}
\toprule
Gate family & Present evidential status & Conservative assessment \\
\midrule
Declared heterogeneous components & Observational & Human, model, artifacts,
and execution environment are distinguishable components with different
roles and information access. \\
Persistent shared vocabulary & Observational & Recurrent terms and distinctions
survive across research phases and constrain later work. \\
Exact structural core \(\Gamma(H)\) & Unmeasured & No frozen local coordinate
system or learned transfer family has yet been estimated from held-out data. \\
Transfer-Laplacian gap & Unmeasured & The idealized value \(2\) belongs to the
toy four-cycle, not to the actual collaboration. \\
Operational alignment & Unmeasured & No preregistered future variable has been
shown conditionally independent of the full history given a measured workspace
state. \\
Joint advantage & Suggestive only & The coupled process has produced outputs
that neither party appears to have supplied verbatim, but local-only Bayes risks
have not been measured under controlled information ceilings. \\
Matched relational null gap & Unmeasured & Shuffled artifacts, matched irrelevant
memory, and local-predictive-marginal-preserving nulls have not been run as a
frozen panel. \\
Causal workspace efficacy & Suggestive only & Corrections, compiler failures,
and artifact substitutions change later behavior, but they are not yet matched
workspace interventions with manipulation checks. \\
Artifact-mediated persistence & Observational & Handoff files and source bundles
support continuation, but blind reconstruction and ablation curves have not
been estimated. \\
Formation provenance & Mixed & The communication architecture is scaffolded,
while terminology, theorem organization, and many control practices were
adaptively developed. Causal formation-history contribution is unmeasured. \\
Workspace agency & Not established & No complete goal-retention,
counterfactual-policy, distributed-control, damage, and repair certificate has
been run. \\
Consciousness or personhood & Not established & The theory supplies no such
inference from the observed collaboration. \\
\bottomrule
\end{tabularx}
\end{center}

The correct present statement is therefore:
\[
\boxed{
\text{candidate persistent collaborative workspace}
\quad\text{rather than}\quad
\text{certified relational workspace agent}.
}
\]

\subtopic{79.6 A controlled certification design}

A serious test can use one manuscript section or software work unit as the unit of analysis.
Before the certification split is opened, freeze:
\begin{itemize}
\item the next-stage task family;
\item the downstream target, such as accepted theorem patch, successful clean
 compile, and absence of a declared regression;
\item the local information available to each component;
\item the workspace extractor and transfer maps;
\item all intervention and null families; and
\item the required gate thresholds.
\end{itemize}
The following conditions can then be compared:
\[
\begin{array}{ll}
\mathsf H & =\text{human-side externalized state only},\\
\mathsf M & =\text{fresh model instance with model-side state only},\\
\mathsf A & =\text{artifact package without live human or model state},\\
\mathsf X & =\text{executor and test outputs without the research context},\\
\mathsf Full & =\text{the intact recurrent loop}.
\end{array}
\]
The design must avoid pretending that a person's unobserved internal state has
been fully measured. The local conditions certify only the declared accessible
histories.

A minimal null panel includes:
\begin{enumerate}
\item theorem-to-artifact reassignment with artifact size and style matched;
\item temporal shuffling of verification feedback;
\item matched irrelevant manuals or source bundles;
\item context reset followed by blind reconstruction;
\item removal of compiler or test feedback while preserving message count;
\item model replacement under a fixed handoff package;
\item artifact removal with equivalent token budget restored by irrelevant text;
\item a local-marginal-matched synthetic null for the frozen downstream target.
\end{enumerate}
A valid empirical claim would require the intact loop to outperform every
required local and matched-null condition, while the extracted workspace state
remains sufficient and feasible workspace interventions alter the downstream
result in the predicted direction.

\subtopic{79.7 Handoff and identity experiment}

Let \(B_t\) be the source-and-results bundle at the end of stage \(t\). A blind
handoff trial replaces the active model instance and optionally the human
analyst, then asks the replacement coalition to reconstruct:
\[
\left(
\text{current definitions},
\text{proved results},
\text{open obligations},
\text{next admissible research step}
\right).
\]
The reconstruction error is measured against a frozen adjudication set. A
replacement invariance test additionally compares:
\begin{itemize}
\item the transported structural workspace;
\item operation behavior on held-out transformations;
\item response to matched interventions; and
\item preservation of failed and unresolved claims.
\end{itemize}
Equivalent benchmark prose alone is insufficient. A replacement that silently
changes which interventions alter later reasoning has a nonzero causal identity
defect even when it produces a similarly polished manuscript.

\subtopic{79.8 Case card}

For the idealized four-cycle model,
\[
\mathsf S,\mathsf O,\mathsf R,\mathsf P
\]
can all be made analytic by specifying the downstream kernel and temporal
updates. For the motivating collaboration itself, the conservative card is
\[
\boxed{
\begin{aligned}
\mathsf S&=\text{observational candidate},
&\mathsf O&=\text{unmeasured},\\
\mathsf R&=\text{unmeasured},
&\mathsf P&=\text{observational candidate},\\
\mathsf F&=\text{mixed and unmeasured},
&\mathsf A&=\text{not established},\\
\mathsf D&=\text{unmeasured},
&\mathsf C&=\text{not established}.
\end{aligned}
}
\]
This audit is intentionally stricter than ordinary conversational language. It
preserves the possibility that the collaboration instantiates a real relational
workspace while refusing to treat that possibility as already certified.

\topic{80}{Case Study C: A Closed-Form Multi-Agent Workspace Agent}

\subtopic{80.1 Declared finite-state system}

Consider three agents:
\[
V_{\mathrm{MA}}=\{P,M,R\},
\]
called the planner, monitor, and recorder. The global workspace-control state
is
\[
c_t=(g_t,b_t,h_t)\in\mathbb F_2^3,
\]
where:
\[
\begin{array}{ll}
g_t=\text{goal bit},
&b_t=\text{environment or belief bit},\\
h_t=\text{workspace-health bit},
&h_t=1\text{ means intact}.
\end{array}
\]
The local accessible states are
\[
X_P(t)=(g_t,b_t),
\qquad
X_M(t)=(b_t,h_t),
\qquad
X_R(t)=(g_t,h_t).
\]
This is the three-node distributed code from Topic~77. Consequently:
\[
\dim\Gamma=3,
\qquad
\spec(L)=\{0,0,0,2,2,2\},
\qquad
k_{\mathrm{rec}}=2,
\qquad
d_{\mathrm{er}}=2.
\]
The organization is fully scaffolded. No endogenous-formation claim is made.

\subtopic{80.2 Goal retention}

In the absence of an explicit goal intervention, let
\[
g_{t+1}=g_t.
\]
Under matched cue-cut trials,
\[
\operatorname{do}(g_t=0)
\quad\text{and}\quad
\operatorname{do}(g_t=1)
\]
remain distinguishable at every later time:
\[
|g_{t+\tau}-g_{t+\tau}'|=1.
\]
Thus the cue-controlled retention factor is exactly one:
\[
\eta_G(\tau)=1.
\]
The goal is not reloaded by an external channel after the initial intervention.

\subtopic{80.3 Joint task and repair action}

After a possible damage event, let the pre-repair health be
\[
h_t^-\in\{0,1\}.
\]
The system chooses a joint action
\[
u_t=(a_t,r_t)\in\{0,1\}^2,
\]
where \(a_t\) is a task action and \(r_t\) is a repair action. Define the
optimal task action
\[
a_t^*=g_t\oplus b_t\oplus h_t^-
\]
and optimal repair action
\[
r_t^*=1-h_t^-.
\]
No individual node determines the full pair \((a_t^*,r_t^*)\):
\begin{itemize}
\item the planner lacks \(h_t^-\);
\item the monitor lacks \(g_t\); and
\item the recorder lacks \(b_t\).
\end{itemize}
Every two-node coalition determines all three workspace bits and therefore the
complete action. Under independent uniform bits,
\[
H(a_t^*\mid X_i(t))=1
\qquad\text{for every }i,
\]
so the task-action component has a one-bit joint advantage and
\[
k_{\mathrm{ctl},0}=2.
\]

\subtopic{80.4 Interventional value and exact policy}

Let the one-step utility be
\[
U_g(a,r;b,h^-)
=
\mathbf 1\{a=g\oplus b\oplus h^-\}
+2\mathbf 1\{h^-\lor r=1\}
-\frac12r.
\]
For fixed \((g,b,h^-)\), the unique maximizer is
\[
\boxed{
(a^*,r^*)
=
(g\oplus b\oplus h^-,1-h^-).
}
\]
Indeed:
\begin{itemize}
\item a wrong task action loses one unit;
\item omitting repair when \(h^-=0\) loses \(3/2\) units; and
\item performing unnecessary repair when \(h^-=1\) loses \(1/2\) unit.
\end{itemize}
The smallest optimal-action gap is therefore
\[
\Delta_Q=\frac12.
\]
By Theorem~\ref{thm:counterfactual-policy}, any uniformly accurate value model with
\[
\varepsilon_Q<\frac14
\]
recovers the exact greedy joint action. More generally, its regret is bounded
by
\[
\operatorname{Reg}\leq2\varepsilon_Q+\delta_{\mathrm{opt}}.
\]

\subtopic{80.5 Feasible goal, workspace, and damage interventions}

The intervention family contains:
\[
\operatorname{do}(g_t=a),
\qquad
\operatorname{do}(b_t=b),
\qquad
\operatorname{do}(h_t^-=h),
\]
implemented consistently across every local chart that contains the manipulated
coordinate. Holding the other two bits fixed, each core-bit intervention
changes the prescribed task action. A matched intervention on an independent
residual bit leaves the task and repair actions unchanged. Hence
\[
\Delta_{\mathrm{specific}}>0.
\]
A damage intervention sets
\[
h_t^-=0.
\]
The matched no-repair counterfactual holds \(r_t=0\), while the intact policy
sets \(r_t=1\).

\subtopic{80.6 Exact self-maintenance}

Post-action health evolves according to
\[
h_{t+1}=h_t^-\lor r_t.
\]
Define the maintenance deficit
\[
e_t=1-h_t.
\]
Under the intact policy,
\[
h_{t+1}=1
\qquad\text{and}\qquad
e_{t+1}=0
\]
for both damage states. Under the matched no-repair intervention and
\(h_t^-=0\),
\[
e_{t+1}^{(0)}=1.
\]
Thus the measured causal repair gain is
\[
\boxed{
\Delta_{\mathrm{maint}}=1.
}
\]
The closed-loop recurrence has
\[
\beta=0,
\qquad
\overline d=0,
\qquad
\overline\epsilon=0,
\]
so Theorem~\ref{thm:self-maintenance} gives
\[
\sup_t e_t=0
\]
after each repair action. Monitor-blind random repair fails the specificity
control whenever it either omits necessary repair or performs unnecessary
repair with positive probability.

\subtopic{80.7 Reentry and causal loop}

The repaired health bit becomes part of the next workspace state:
\[
h_t^-
\longrightarrow
r_t
\longrightarrow
h_{t+1}
\longrightarrow
X_M(t+1),X_R(t+1).
\]
Cutting the direct storage of \(h_t^-\) while preserving the action-mediated
path still leaves a return path through repair. Cutting the repair-to-health
edge destroys that return. The system therefore contains explicit
loop-mediated reentry rather than mere temporal autocorrelation.

\subtopic{80.8 Analytic agency certificate}

Within this finite structural causal model, the following hold exactly:
\begin{enumerate}
\item a nontrivial exact distributed workspace with spectral gap two;
\item two-node minimal reconstruction and control coalition size;
\item matched-null joint advantage for the parity task;
\item persistent internally retained goal state;
\item exact interventional action values and zero policy regret;
\item goal and workspace interventions that alter action;
\item damage detection and repair selection;
\item positive matched repair gain and bounded maintenance deficit; and
\item causal reentry through the action-mediated repair loop.
\end{enumerate}
Therefore, relative to the declared bits, actions, utility, intervention family,
and one-step horizon,
\[
\boxed{
\operatorname{Agent}_{\Theta_A}(\Workspace)=1.
}
\]
The case reaches formal Level~6 by construction. It remains:
\[
\text{fully scaffolded},
\qquad
\text{domain restricted},
\qquad
\text{non-evidence for consciousness}.
\]

\subtopic{80.9 Case card}

\[
\boxed{
\begin{aligned}
\mathsf S&=\checkmark_{\mathrm{analytic}},
&\mathsf O&=\checkmark_{\mathrm{analytic}},\\
\mathsf R&=\checkmark_{\mathrm{analytic}},
&\mathsf P&=\checkmark_{\mathrm{analytic}},\\
\mathsf F&=\text{fully scaffolded},
&\mathsf A&=\checkmark_{\mathrm{analytic}},\\
\mathsf D&=\text{fixed decomposition},
&\mathsf C&=\text{not established}.
\end{aligned}
}
\]
The example demonstrates logical compatibility among relationality, redundancy,
agency, and self-maintenance. It does not show that those properties arise
spontaneously, scale to open environments, or imply subjective experience.

\topic{81}{Near-Miss Cases and Anti-Triviality Controls}

Positive examples are insufficient unless neighboring false positives are also
made explicit. The following systems each resemble a relational workspace in
one respect while failing a required gate.

\subtopic{81.1 Near miss A: the replicated archive}

Let every node store a complete copy of the same core:
\[
E_i=I_d
\qquad\text{for every }i.
\]
Then coherence is perfect and reconstruction is robust, but
\[
k_0=1,
\qquad
\min_i\delta_i^{\mathrm{local}}=0.
\]
If the archive is never used in downstream control, then also
\[
\eta_{\mathrm{causal}}=0.
\]
This is a shared representation or archive, not a relationally irreducible
workspace.

\subtopic{81.2 Near miss B: the centralized oracle with decorative agents}

Suppose node \(o\) contains the complete predictive state \(Z_t\), while every
other node receives only messages generated by \(o\). The full system can have
excellent performance, rich communication, and recurrent activity. Yet
\[
Y_t^+\indep\mathbf H_t\mid X_o(t),
\qquad
k_0=1,
\qquad
\Delta_{\mathrm{joint}}=0.
\]
Removing the decorative nodes may change latency or style without changing the
operative capacity. Multi-agent appearance therefore does not imply a
relational workspace.

\subtopic{81.3 Near miss C: locally coherent but globally obstructed charts}

Let three local scalar charts have relative overlap demands
\[
b_{12}=b_{23}=b_{31}=1.
\]
The cycle obstruction is
\[
\omega_{\mathcal U}(b)
=
\frac{|1+1+1|}{\sqrt3}
=\sqrt3>0.
\]
Each local chart is internally coherent, but no global scalar assignment
satisfies all three overlap equations. The system contains compatible-looking
subworkspaces without a global workspace.

\subtopic{81.4 Near miss D: externally reloaded goals}

Let
\[
G_{t+1}=C_{t+1},
\]
where \(C_{t+1}\) is an external instruction repeated at every step. Goal and
action can be strongly correlated, but after cutting the cue channel,
\[
I(G_{t+\tau};G_t\mid\operatorname{cut}(C))=0.
\]
The system tracks an external command. It has not established internal goal
retention.

\subtopic{81.5 Near miss E: task recovery with identity failure}

Let two realizations have identical benchmark output distributions,
\[
P(Y_{\mathrm{bench}})=P'(Y_{\mathrm{bench}}),
\]
but different intervention kernels:
\[
K(w,b)\neq K'(\beta w,b)
\]
for some matched workspace intervention. Then task behavior is preserved while
\[
d_{\mathrm{causal}}(W,W')>0.
\]
The later system may be a successor or replacement rather than the same
workspace under the identity criterion of Part~XII.

\subtopic{81.6 Near miss F: open-loop repair replay}

A recorded repair sequence may restore average task performance after a fixed
damage schedule. If the same sequence is replayed when the damage timing or
type changes, it need not respond appropriately. When
\[
R_t\indep M_t\mid\text{clock and replay index},
\]
monitor-to-repair mediation is absent. Successful replay is robustness under a
script, not self-maintenance.

\subtopic{81.7 Anti-triviality summary}

The near misses establish six independent exclusions:
\[
\boxed{
\begin{aligned}
\text{workspace}
\neq{}&\text{replicated storage},\\
\neq{}&\text{centralized control with decorative communication},\\
\neq{}&\text{locally coherent but obstructed charts},\\
\neq{}&\text{externally reloaded goals},\\
\neq{}&\text{task-equivalent identity replacement},\\
\neq{}&\text{open-loop repair replay}.
\end{aligned}
}
\]
A complete certificate must distinguish the intended realization from all
near misses relevant to its claimed level.

\topic{82}{Cross-Case Synthesis and Benchmark Blueprint}

\subtopic{82.1 Comparative gate matrix}

\begin{center}
\small
\begin{tabularx}{\textwidth}{@{}>{\raggedright\arraybackslash}p{0.20\textwidth}>{\raggedright\arraybackslash}p{0.20\textwidth}>{\raggedright\arraybackslash}p{0.23\textwidth}>{\raggedright\arraybackslash}X@{}}
\toprule
Gate family & Human--human analytic model & Motivating human--AI collaboration & Multi-agent analytic model \\
\midrule
Exact structural core & Analytic pass & Not measured & Analytic pass \\
Spectral gap & \(2\) & Not measured & \(2\) \\
Minimal sufficient coalition & \(k_0=2\) & Not measured & \(k_0=2\) \\
Matched-null transfer gap & \(1\) bit & Not measured & \(1\) bit for task parity \\
Causal specificity & Analytic pass & Suggestive, uncontrolled & Analytic pass \\
One-node erasure resilience & Analytic pass & Handoff only observational & Analytic pass \\
Formation provenance & Scaffolded & Mixed, unmeasured & Fully scaffolded \\
Goal retention & Not modeled & Not certified & Exact analytic pass \\
Counterfactual policy & Not modeled & Not certified & Exact analytic pass \\
Self-maintenance & Not modeled & Not certified & Exact analytic pass \\
Highest supported level & Level 4 in model & Candidate only & Level 6 in model \\
Consciousness/\allowbreak personhood & Not established & Not established & Not established \\
\bottomrule
\end{tabularx}
\end{center}

The table illustrates why the theory uses a vector certificate rather than one
score. The human--human model is relational and persistent without being
agentive. The multi-agent model is agentive without being endogenously formed.
The motivating collaboration is richer than either toy model but presently has
weaker evidential status because the necessary interventions and nulls have not
been frozen and run.

\subtopic{82.2 Benchmark family B1: distributed parity}

The first benchmark is the exact parity family from Topics~85--80. For
\(m\geq3\), distribute \(m\) independent bits among nodes so that no individual
node sees all bits, but specified coalitions do. The target is
\[
Y=\bigoplus_{j=1}^m C_j.
\]
The benchmark measures:
\begin{itemize}
\item the learned transfer-Laplacian spectrum;
\item coalition conditional entropy and log risk;
\item recovery of \(k_0\);
\item robustness to edge and node ablations;
\item matched local-predictive null gaps; and
\item causal response to consistent versus inconsistent coordinate
 interventions.
\end{itemize}
The ground truth is known exactly, so structural and relational estimators can
be calibrated before they are applied to open-ended research tasks.

\subtopic{82.3 Benchmark family B2: theorem-patch handoff}

A theorem-patch task supplies:
\[
\left(
\text{definitions},
\text{partial proof},
\text{counterexample tests},
\text{style constraints}
\right)
\]
across distinct components. No component receives all four. The coupled
system must produce a patch that:
\begin{enumerate}
\item is formally consistent with the definitions;
\item closes the declared proof obligation or marks it unresolved;
\item passes the frozen counterexample suite;
\item compiles from a clean environment; and
\item survives a blind replacement-handoff trial.
\end{enumerate}
Matched nulls permute proof fragments, substitute irrelevant artifacts of equal
size, or replay test outputs from a different theorem. This family targets the
structural, relational, persistence, and identity layers without requiring an
agency claim.

\subtopic{82.4 Benchmark family B3: damage and repair}

A certified workspace is first formed and measured. A damage intervention then
removes or corrupts one of:
\[
\text{a transfer channel},
\quad
\text{a participant},
\quad
\text{an artifact},
\quad
\text{a goal encoding},
\quad
\text{a verification path}.
\]
The system receives a repair-action family with matched costs. The benchmark
measures:
\[
\eta_{D\to M},
\qquad
\eta_{M\to R},
\qquad
\eta_{R\to v},
\qquad
\Delta_{\mathrm{repair}}^{\mathrm{spec}},
\qquad
\mathbf d_{\mathrm{id}}.
\]
Open-loop replay, random repair, off-workspace repair, and identity-breaking
task recovery are mandatory controls.

\subtopic{82.5 Benchmark family B4: scaffold transport}

Train the same formation rule under several scaffolds
\[
\Sigma^{(1)},\ldots,\Sigma^{(K)}
\]
that differ in node labels, coordinate systems, and communication topology but
support the same abstract task. After gauge alignment, measure whether the
learned organizations transport to a common equivalence class:
\[
[\widehat H^{(1)}]_{\mathcal G}
\cong\cdots\cong
[\widehat H^{(K)}]_{\mathcal G}.
\]
A structure that appears only under one answer-encoding scaffold has weak
endogeneity evidence. A structure that recurs across alternative scaffolds,
passes held-out task and causal gates, and cannot be compressed into the
scaffold specification has stronger evidence.

\subtopic{82.6 Reporting rule for mixed analytic and empirical cases}

When a case combines exact mathematics with observed data, report separate
columns:
\[
\boxed{
\text{proved in model}
\quad\big|\quad
\text{estimated in system}
\quad\big|\quad
\text{not established}.
}
\]
For example, the four-cycle Laplacian has an analytic gap of two, but an actual
human--AI collaboration has that gap only if its measured discrepancy operator
has the corresponding spectrum. Importing the toy value into the empirical
case is invalid.

\subtopic{82.7 Lessons from the integrated cases}

The case studies establish the following conceptual separation:
\begin{enumerate}
\item A relational workspace can be mathematically exact without being an
 agent.
\item An agentive workspace can be fully scaffolded and still be operationally
 real within its domain.
\item A rich real collaboration can be a stronger candidate than a toy model
 while having a weaker current certificate.
\item Redundancy and irreducibility are compatible when no individual node is
 sufficient but several proper coalitions are.
\item Artifact-mediated persistence is not passive storage when artifact
 contents causally alter later operations.
\item High task performance, communication volume, or behavioral polish alone
 cannot distinguish the intended theory from its near misses.
\item No structural, causal, or agentive result in these examples entails
 phenomenal consciousness.
\end{enumerate}

\centerbox{\large\bfseries
The theory becomes empirically meaningful only when the same gate card that
certifies the positive case rejects its matched near misses.
}

\theorypart{XXIII}{Relation to Prior Work}

\topic{83}{Scope, Method, and Standards of Priority}

\subtopic{83.1 Purpose of the literature review}

The preceding parts construct a unified formal theory. Construction alone does
not establish priority. This part therefore asks a narrower question:
which ingredients are established, which are adaptations to a new domain,
which combinations appear to be unusual, and which claims remain open originality candidates?

The audit uses five labels:
\[
\begin{array}{ll}
\mathsf{STD} & = \text{established concept, theorem, or technique},\\
\mathsf{ADP} & = \text{established machinery adapted to relational workspaces},\\
\mathsf{SYN} & = \text{nontrivial synthesis of established components},\\
\mathsf{CAN} & = \text{candidate original contribution requiring priority review},\\
\mathsf{EMP} & = \text{empirical claim not yet established in an uncontrolled system}.
\end{array}
\]
These labels describe the status of a claim, not its importance. A standard
result may be indispensable. A synthesis may be scientifically useful without
containing a new theorem. A candidate contribution may later be narrowed or
withdrawn when a closer precedent is found.

\subtopic{83.2 Search scope and date}

This review covers sources available through August 29, 2026. It includes:
\begin{itemize}
\item extended and distributed cognition, transactive memory, shared
 intentionality, common knowledge, collective intelligence, and group
 agency;
\item blackboard systems, global-workspace architectures, team cognition,
 hybrid intelligence, and human--AI cooperation;
\item category-theoretic limits, local systems, holonomy, sheaves, \v{C}ech
 descent, connection and sheaf Laplacians, consensus, and spectral
 perturbation;
\item statistical sufficiency, causal states, information decomposition,
 structural causal models, system identification, frames, erasure
 recovery, and probabilistic process equivalence; and
\item agency, viability, control, empowerment, damage detection, repair, and
 self-maintaining systems.
\end{itemize}
The search combined exact-title queries, component-level queries, backward
reference tracing, and comparison with recent formal work. It is a scholarly
priority review, not a legal patent search and not a proof that no uncatalogued,
unpublished, or differently named precedent exists.

\subtopic{83.3 A stronger standard than name uniqueness}

Failure to find the phrase ``Stratified Relational Workspace Theory'' in prior
literature is weak evidence. Priority is determined by mathematical and
operational equivalence, not by terminology. Let \(C\) be a claim in the
paper and \(S\) a prior source. Write
\[
S\preccurlyeq C
\]
when the main content of \(C\) follows from \(S\) after admissible relabeling,
restriction, or direct specialization. A substantive novelty residual exists
only when the claim contains content not absorbed by the closure of known
sources:
\[
\operatorname{Res}(C)
=
C\setminus
\operatorname{Cl}
\left(
\bigcup_{S\in\mathcal L}S
\right).
\]
This is a conceptual audit device rather than a set-theoretic definition of
ideas. Its purpose is to force an explicit account of what remains after the
standard machinery is removed.

\subtopic{83.4 When a synthesis counts as a contribution}

A conjunction of known ingredients is not automatically new. The present
framework treats a synthesis as a candidate contribution only when it satisfies
all three conditions:
\begin{enumerate}
\item \textbf{Non-redundancy.} Removing one of the joined layers changes the
 class of systems certified or the predictions made.
\item \textbf{Formal integration.} The layers share declared objects,
 transfer maps, interventions, or invariants rather than appearing as an
 unconnected checklist.
\item \textbf{New discriminations.} The synthesis separates at least one
 positive case from a matched near miss that the individual components do
 not separate on their own.
\end{enumerate}
The anti-centralization, matched-null, non-collapse, provenance, replacement,
and repair-replay controls are important precisely because they make the
combination testable rather than merely terminological.

\subtopic{83.5 Priority claims are coordinatewise}

The analysis does not make one global novelty claim. It uses a priority
profile
\[
\boxed{
\mathbf p(C)
=
\left(
 p_{\mathrm{concept}},
 p_{\mathrm{object}},
 p_{\mathrm{theorem}},
 p_{\mathrm{certificate}},
 p_{\mathrm{experiment}}
\right).
}
\]
For example, the philosophical concept of distributed cognition may be
standard while a specific mathematical object is adapted, a theorem is a
specialization of known spectral theory, the certificate is a synthesis, and
the empirical claim is unresolved. This prevents originality in one
coordinate from being used to imply originality in every coordinate.

\topic{84}{Philosophical and Cognitive Ancestry}

\subtopic{84.1 Extended cognition}

The extended-mind tradition argues that external resources can participate in
cognitive processes when they are appropriately integrated with an agent's
activity \cite{ClarkChalmers1998}. Its critics emphasize that causal coupling
alone is too permissive and that cognition requires a principled mark or
boundary \cite{AdamsAizawa2001,Rupert2004}. Stratified Relational Workspace
Theory inherits the possibility of extracranial realization but changes the
question. It does not ask whether a notebook, model, or collaborator literally
belongs to one individual's mind. It asks whether a distributed realization
contains a persistent, causally operative, jointly reconstructible structure
that passes explicit anti-triviality gates.

Accordingly,
\[
\operatorname{ExtendedCognition}
\centernot\Longleftrightarrow
\operatorname{SRW}.
\]
An externally scaffolded process may be cognitively extended without being
relationally irreducible. Conversely, an institution or multi-agent system may
instantiate a certified relational workspace without being an extension of any
single person's cognition.

\subtopic{84.2 Distributed cognition}

Distributed cognition studies cognitive organization across people,
representations, artifacts, and environments rather than locating all relevant
processing inside one individual \cite{Hutchins1995,HollanEtAl2000}. This is
the closest philosophical ancestor of the framework's ontology. It already
supports the central observation that the unit of analysis may be a coupled
system and that external representations can transform the computations
available to that system.

The framework should therefore not claim the following ideas as new:
\begin{itemize}
\item cognition can be distributed across heterogeneous carriers;
\item artifacts can be active components of a cognitive process;
\item representational transformations can alter system-level performance; and
\item the relevant explanatory unit can exceed an individual organism.
\end{itemize}
The added ambition is mathematical and certifying. The theory specifies local
state spaces, transfer maps, compatible sections, causal sufficiency,
coalition irreducibility, persistence, replacement identity, formation
provenance, and matched nulls. Distributed cognition supplies the ontological
motivation. It does not by itself supply this complete formal certificate.

\subtopic{84.3 Transactive memory and collective intelligence}

Transactive-memory theory models the group-level organization of ``who knows
what'' and the processes by which members encode, retrieve, and coordinate
specialized knowledge \cite{Wegner1987,Lewis2003}. Collective-intelligence
research studies stable differences in group performance that are not reduced
to the intelligence of a single member \cite{WoolleyEtAl2010}. Both traditions
are directly relevant to the coalition layer.

The distinction is:
\[
\begin{aligned}
\text{transactive memory}
&\approx
\text{distributed indexing and retrieval of expertise},\\
\text{collective intelligence}
&\approx
\text{system-level task capacity},\\
\text{relational workspace}
&\approx
\text{an invariant state plus transfer, causal, and identity structure}.
\end{aligned}
\]
A transactive-memory system may still be centralized by one coordinator. A
high-performing group may still be explainable by the best individual member.
The certificate therefore requires matched coalition risks and nulls
rather than inferring relationality from group success.

\subtopic{84.4 Shared intention, common knowledge, and group agency}

Common-knowledge theory formalizes nested epistemic conditions required by
many distributed protocols \cite{HalpernMoses1990,FaginEtAl1995}. Shared-action
accounts analyze jointly intended and mutually responsive activity
\cite{Bratman1992}. Group-agency theory asks when a collective can satisfy
rational and normative conditions associated with agency
\cite{ListPettit2011}.

These literatures address questions that the theory deliberately keeps
separate:
\[
\begin{array}{ll}
\text{epistemic coordination} & \text{what members know about one another},\\
\text{shared intention} & \text{how intentions and plans are jointly organized},\\
\text{group agency} & \text{when collective attitudes and actions form an agent},\\
\text{SRW structure} & \text{how a distributed operative state exists and persists}.
\end{array}
\]
The workspace-agent predicate in Part~XVI is domain-relative and causal. It is
not a general account of corporate personhood, moral responsibility, or
collective intentionality.

\topic{85}{Workspace Architectures and the Meaning of ``Workspace''}

\subtopic{85.1 Global workspace theories}

Global workspace theories use widespread availability or broadcast within a
cognitive architecture to explain conscious access and flexible control
\cite{Baars1988,DehaeneChangeux2011}. The shared word ``workspace'' can create
a misleading implication of equivalence. The theory neither assumes a
single broadcast hub nor defines consciousness by access to a global channel.
Its workspace may be distributed across agents and artifacts, and its
existence criterion is a conjunction of structural, predictive, causal,
relational, and persistence gates.

The relation is therefore analogical rather than identificatory:
\[
\boxed{
\operatorname{SRW}(W)
\centernot\Longrightarrow
\operatorname{GlobalNeuronalWorkspace}(W)
\centernot\Longrightarrow
\operatorname{Conscious}(W).
}
\]
The term is retained because it denotes a domain in which jointly
available operations can be carried out. Every consciousness claim remains on
an explicitly separate axis.

\subtopic{85.2 Dialogue and the Shared Workspace Framework}

Pickering and Garrod's Shared Workspace Framework treats dialogue as a jointly
organized activity in which interlocutors construct, align, and use shared
representations, including persistent external artifacts
\cite{PickeringGarrod2021}. This is a direct conceptual and terminological
predecessor. The theory therefore cannot claim priority for the idea
that dialogue or collaboration operates through a shared workspace.

The difference is one of formal scope. SRWT permits arbitrary graph-structured
human, machine, instrument, and artifact nodes and asks whether their common
organization passes typed-transport, coalition, causal, persistence, and
replacement gates. A dialogue can instantiate such a workspace, but the
shared-workspace account of dialogue does not by itself imply relational
irreducibility or persistent identity under carrier replacement.

\subtopic{85.3 Blackboard systems}

Blackboard architectures coordinate specialized modules through a shared data
structure \cite{HayesRoth1985}. They are an important engineering precedent
for artifact-mediated collaboration. They also provide a canonical near
miss. When the blackboard itself contains the complete operative state and
all modules merely read or write it, the system can fail the relational
irreducibility gate:
\[
k_0=1.
\]
A blackboard implementation can instantiate a relational workspace, but only
when no single board or controller is sufficient and when the transfer and
coalition tests show genuinely distributed operative capacity.

\subtopic{85.4 Consensus and coordination}

Network consensus studies how local interactions drive distributed states
toward agreement \cite{OlfatiSaberEtAl2007}. Agreement is not the same as a
workspace. Consensus often removes heterogeneity by converging all nodes to a
common value. Stratified relational workspaces preserve heterogeneous local
representations and may require complementary disagreement:
\[
X_i\not\cong X_j,
\qquad
E_i c\neq E_j c,
\qquad
(c_i)_i\in\Gamma(H).
\]
The global state is compatible, not necessarily copied. This distinction is
central to the anti-replication criterion.

\subtopic{85.5 Terminological discipline}

The phrase ``relational workspace'' is used here in a technical sense:
\begin{quote}
A state-bearing, operation-supporting organization whose relevant invariants
are distributed across declared components and preserved through admissible
transfers and time.
\end{quote}
It does not, by terminology alone, imply a mind, subject, self, institution,
shared intention, or conscious broadcast. Those stronger predicates require
separate definitions and evidence.

\topic{86}{Local-to-Global Mathematics}

\subtopic{86.1 Limits, local systems, and holonomy}

The exact compatible-section object
\[
\Gamma(H)=\varprojlim C_i
\]
is a standard categorical limit construction \cite{MacLane1998}. For
invertible transfers on a connected graph, identifying global sections with
the fixed subspace of holonomy is standard local-system and bundle mathematics
\cite{Hatcher2002,Bredon1997}. Accordingly, Theorem~\(T_1\) is not claimed as
a new abstract theorem. Its role is foundational: it translates the workspace
question into an exact invariant-subspace question and exposes cycle
obstructions.

The operation-closure and commutant statements are likewise standard
representation-theoretic consequences once the operation algebra and
intertwiners are defined. Their contribution is the specification of what
counts as a shared operation in the workspace model, not the discovery of
commutants.

\subtopic{86.2 Cellular sheaves and sheaf Laplacians}

Cellular sheaves encode heterogeneous local data and restriction maps, while
global sections express consistency across a cell complex. Their Laplacians
have kernels equal to spaces of global sections and support spectral analysis
of approximate consistency \cite{HansenGhrist2019,Robinson2017}. The transfer
Laplacian
\[
L_H=D_H^*D_H
\]
is therefore best understood as a workspace-specific sheaf or connection
Laplacian. The kernel theorem \(\ker L_H=\Gamma(H)\) is standard once the
coboundary operator is constructed.

The appropriate priority statement is:
\[
\begin{aligned}
\text{sheaf/Laplacian mechanism} &:\ \mathsf{STD},\\
\text{workspace interpretation and gate integration} &:\ \mathsf{ADP}/\mathsf{SYN}.
\end{aligned}
\]
The transfer Laplacian itself is not presented as an
unprecedented mathematical object.

\subtopic{86.3 Connection Laplacians and gauge covariance}

Connection Laplacians compare vector-valued data transported across a graph
and naturally transform under local changes of coordinates
\cite{SingerWu2012}. The gauge rule
\[
H_e^g=g_{t(e)}H_eg_{s(e)}^{-1}
\]
and the conjugacy of holonomy are therefore standard. The theory uses
this structure to define realization identity and to prevent raw coordinate
matrices from being mistaken for invariant workspace content.

This use is an adaptation with conceptual consequences. It makes model or
participant replacement a question about transported invariant organization,
not token equality. The general mathematics remains prior.

\subtopic{86.4 \v{C}ech descent and gluing}

Exact gluing, equalizer formulas, and \v{C}ech obstruction theory are standard
local-to-global constructions. Sheaf-theoretic contextuality provides a
prominent example in which the absence of a global section is itself the
substantive phenomenon \cite{AbramskyBrandenburger2011}. Part~XIII adapts this
logic to local workspaces and adds an anti-centralization requirement at the
meta-level.

The exact descent theorems \(T_8\) and \(T_9\) should therefore be presented as
specialized consequences of established descent. The candidate contribution
is the combination:
\[
\text{descent}
+
\text{predictive sufficiency}
+
\text{causal equivalence}
+
\text{coalition anti-centralization}.
\]
Standard gluing alone does not certify an operative workspace of workspaces.

\subtopic{86.5 Spectral perturbation}

The projector-stability results in Part~IV belong to the classical perturbation
theory of invariant subspaces \cite{DavisKahan1970,Kato1995}. Bounds of the
form
\[
\|P-\widetilde P\|
\lesssim
\frac{\|L-\widetilde L\|}{\text{spectral gap}}
\]
are not new. Their workspace-specific role is to turn uncertainty in learned
transfers into uncertainty about identity, operation closure, and replacement.
The novelty, if any, lies in the downstream certificate built on the bound, not
in the perturbation inequality.

\topic{87}{Predictive, Causal, and Information-Theoretic Precedents}

\subtopic{87.1 Minimal predictive states}

The conditional-law state
\[
Z_t=\mathbb P(Y_t^+\in\cdot\mid\mathbf H_t)
\]
is a canonical sufficient statistic. Computational mechanics develops causal
states as equivalence classes of histories with identical future conditional
laws and proves their minimal predictive character
\cite{ShaliziCrutchfield2001}. Theorem~\(T_{12}\) is therefore an established
minimal-sufficiency result specialized to the declared future target and
intervention family.

The new burden is structural alignment: determining whether the predictive
state is recoverable from the low-energy transfer workspace. Predictive
sufficiency alone does not identify a distributed workspace.

\subtopic{87.2 Statistical experiments and sufficiency}

Blackwell's comparison of experiments formalizes when one information source
is at least as informative as another for every decision problem
\cite{Blackwell1953}. The coalition-risk hierarchy in Part~VII is compatible
with this tradition. A coalition can be sufficient even when its internal
representation differs from the full-system state.

The framework adds a decomposition-relative question:
\[
\text{what is the smallest coalition whose accessible history is sufficient?}
\]
The statistic \(k_\varepsilon\) is an adaptation of sufficiency to distributed
carrier structure. Its interpretation depends on the declared nodes and
accessible encodings.

\subtopic{87.3 Synergy and partial information}

Partial-information decomposition distinguishes redundant, unique, and
synergistic information in multivariate systems, building on standard
information-theoretic quantities \cite{CoverThomas2006,WilliamsBeer2010}. The
XOR examples in Part~XXII are canonical synergy constructions and are not new.
They are used because they provide exact ground truth for relational
irreducibility:
\[
I(Y;X_i)=0
\quad\text{for every }i,
\qquad
I(Y;X_1,\ldots,X_m)>0.
\]
The matched-predictive-null test is more restrictive than merely
reporting synergy. It asks whether local predictive marginals can be held
fixed while the relational organization and full-system capacity change.
That joint null design is a candidate synthesis, not a claim that XOR or
synergy is new.

\subtopic{87.4 Structural causal models}

Interventions, counterfactuals, mediation, and causal graphs are established
machinery \cite{Pearl2009}. The theory adopts the interventionist
requirement that a workspace variable can be manipulated only through a
feasible implementation intervention. This avoids treating
\(\operatorname{do}(W=w)\) as meaningful when no physical or computational
operation realizes \(w\).

The causal gate is an adaptation:
\[
\text{SCM intervention theory}
\longrightarrow
\text{implemented workspace and repair interventions}.
\]
The potentially distinctive contribution is that causal efficacy is tested
jointly with coherence, sufficiency, coalition irreducibility, matched nulls,
and identity preservation.

\subtopic{87.5 Probabilistic process equivalence}

Bisimulation and probabilistic testing compare processes by the observations
and experiments they support \cite{LarsenSkou1991}. The causal-identity profile
in Part~XII belongs to this broader family of behavioral equivalence notions.
It differs by keeping structural transport, operation intertwining,
reconstruction, and intervention response as separate coordinates rather than
collapsing them into one process equivalence.

This is best classified as an integrative adaptation. The existence of
behavioral-equivalence theory blocks any claim that replacement identity based
on observable behavior is wholly new.

\subtopic{87.6 Specification sensitivity, partial identification, and post-selection}

A conclusion can depend on analytically reasonable choices even when every
individual analysis is internally correct. Sensitivity analysis and
specification-curve or multiverse methods make that dependence visible rather
than allowing one favorable analysis to stand for an entire scientific claim
\cite{Leamer1983,SteegenEtAl2016,SimonsohnEtAl2020}. Partial-identification
theory similarly replaces unjustified point conclusions with identified sets
when the maintained assumptions and observations determine only bounds
\cite{Manski2003,ImbensManski2004}. Post-selection inference addresses the
additional error created when a model, target, or specification is selected
because it looked favorable on the same data used for reporting
\cite{BerkEtAl2013}.

These traditions preclude any claim that specification robustness, identified
sets, uniform concentration, or selection adjustment are new statistical
ideas. The adaptation in Part~XVIII is to apply them jointly to a
noncompensatory workspace certificate whose gates span structure, causal
operation, coalition irreducibility, identity, provenance, and maintenance.
The identified object is therefore not a single regression parameter but the
sign profile of an entire relational certificate across a declared family of
scientific specifications.

\topic{88}{Persistence, Reconstruction, Learning, and Control}

\subtopic{88.1 Frames and erasure recovery}

Stable reconstruction from distributed encodings is standard frame theory.
Lower frame bounds control noise amplification, and frame robustness under
erasures is well studied \cite{CasazzaKovacevic2003,Christensen2016}. Theorem
\(T_7\) and the node-erasure distance therefore instantiate known recovery
logic.

The adaptation is to distinguish:
\[
\overline E_i
=
\text{complete mathematical chart},
\qquad
E_i=O_i\overline E_i
=
\text{locally accessible encoding}.
\]
This distinction is essential to the anti-centralization gate, but the frame
inequalities themselves are standard.

\subtopic{88.2 Persistence recurrences}

Bounds for recurrences such as
\[
e_{t+1}\leq\alpha_t e_t+\sigma_t
\]
are elementary stability results. The persistence theorem is not novel as an
inequality. Its conceptual role is to force three independent checks:
\begin{enumerate}
\item the workspace subspace remains trackable;
\item the state remains close to it; and
\item the retained workspace content does not collapse to zero.
\end{enumerate}
The non-collapse gate is a modeling contribution because it rules out a common
but misleading form of ``stability.''

\subtopic{88.3 System identification and held-out selection}

Population least squares, excitation conditions, ridge bias, and empirical
model selection are standard system-identification and statistical-learning
ideas \cite{Ljung1999}. Theorems \(T_{15}\)--\(T_{17}\) are finite-dimensional
specializations. Their workspace-specific content is the demand that transfer
maps be compared modulo gauge and that formation, calibration, and
certification data remain separated.

The provenance layer is not equivalent to identifiability. A transfer can be
identified accurately while having been entirely encoded by the scaffold. A
structure can be endogenously learned while failing every workspace gate.

\subtopic{88.4 Viability, safety, and self-maintenance}

Viability theory studies trajectories that remain inside admissible regions
\cite{Aubin1991}. Control-barrier methods provide constructive safety
certificates \cite{AmesEtAl2017}. Cybernetic and autonomy traditions analyze
regulation, adaptive self-production, and operational closure
\cite{Ashby1956,MaturanaVarela1980,DiPaolo2005}. Formal accounts of agency
emphasize individuality, normativity, asymmetry, and temporally extended action
\cite{BarandiaranEtAl2009}. Empowerment measures an agent's potential causal
control over future states \cite{KlyubinEtAl2005}.

The closed-loop maintenance theorem in Part~XVI uses standard controlled
stability. Its candidate contribution is the target of maintenance: not an
arbitrary safety set, but the conjunctive workspace certificate, including
identity and anti-centralization. The matched no-repair, random-repair,
off-workspace-repair, and identity-breaking-reconstruction controls prevent
passive robustness or generic recovery from being reported as self-maintenance.

\subtopic{88.5 Decision and policy theory}

The value-error-to-regret bound in Theorem~\(T_{20}\) is elementary decision
theory and aligns with standard reinforcement-learning analysis
\cite{SuttonBarto2018}. Its role is evidential. An internal value channel
counts only when interventions on that channel alter action selection in the
direction predicted by counterfactual value and outperform replay or
critic-ablation controls.

\topic{89}{Human--AI and Multi-Agent Research}

\subtopic{89.1 Hybrid intelligence and AI teammates}

Hybrid-intelligence research studies systems in which human and machine
capabilities are deliberately combined \cite{DellermannEtAl2019}. Work on
machines as teammates emphasizes coordination, trust, communication, shared
mental models, and team design \cite{SeeberEtAl2020}. Collective-intelligence
research has begun to ask how human--AI groups can produce capacities not
reducible to either humans or AI alone \cite{GuptaEtAl2025}.

These programs establish that human--AI collectives are a legitimate unit of
analysis. They do not by themselves identify an invariant core, prove
non-reducibility relative to declared components, or supply a replacement and
repair certificate. The theory should be positioned as a formal
measurement and falsification layer for such systems, not as the first proposal
that humans and AI can form productive teams.

\subtopic{89.2 Distributed-cognition analyses of human--AI interaction}

Recent work explicitly applies distributed cognition to human--AI cooperation
\cite{SidjiEtAl2024,ZhaoHan2026}. This literature is a direct conceptual
precedent. It reinforces that representations, artifacts, interface states,
and human reasoning can form one analytical system.

The framework differs in granularity. It asks for a learned or declared
transfer diagram, a spectral or exact structural state, predictive and causal
alignment, coalition tests, matched nulls, persistence under handoff, and
provenance controls. Those additions may constitute a formal synthesis, but
they should not be described as the discovery of distributed human--AI
cognition itself.

\subtopic{89.3 Cognitive spillover}

Randomized evidence now indicates that exposure to AI can causally alter later
human--human language, attention, shared mental models, and social cohesion,
even outside the immediate human--AI exchange \cite{RiedlEtAl2026}. This is
important for the theory because it demonstrates that the interaction graph can
have downstream effects that outlast the active model channel.

Spillover is not automatically a relational workspace. It may reflect
priming, imitation, or one-way influence. A workspace claim additionally
requires a jointly sufficient state, matched local nulls, bidirectional or
multi-path causal organization, and reconstructive persistence.

\subtopic{89.4 Shared neural workspaces}

Goyal and collaborators introduced a capacity-limited shared global workspace
through which specialized neural modules communicate and compete for access
\cite{GoyalEtAl2022}. The architecture directly anticipates the claim that a
constrained shared channel can synchronize otherwise specialized components.
It is therefore a close machine-learning precedent, not merely a metaphorical
relative.

SRWT differs by allowing heterogeneous human, model, artifact, and instrument
nodes and by withholding the relational label until coalition, matched-null,
causal, persistence, and identity gates are passed. A neural shared workspace
may satisfy those gates, but task gains and bottlenecked communication alone do
not establish them.

\subtopic{89.5 Graph-based human--AI artifacts}

CoMAP provides a persistent graph-based visual workspace through which
educators and AI jointly construct project-based learning designs
\cite{LiJiang2026}. It is grounded in distributed cognition and treats the
artifact graph as a transparent, revisable collaboration medium. This system
is a direct precedent for the claim that human--AI interaction can be organized
by a persistent shared graph rather than a linear prompt history.

The remaining SRWT question is empirical: whether the graph contains an
estimated low-energy compatible state, whether no declared component is
sufficient by itself, and whether targeted graph interventions alter later
joint behavior in a workspace-specific way.

\subtopic{89.6 Controlled shared-workspace synergy}

Kotalwar, Das, and Rose study shared-workspace human--AI collaboration across
1,482 simulated sessions \cite{KotalwarDasRose2026}. Their results show that
adding collaborators can reduce performance when coordination structure is
absent, while shared group memory and simulated human-in-the-loop approval
gates can improve outcomes. This is a direct empirical predecessor to the
manuscript's claim that available capability is not equivalent to coordinated
team capacity.

The study also constrains priority. SRWT is not the first empirical framework
to test whether shared-memory scaffolding changes human--AI team performance.
Its candidate increment is the stricter decomposition of that effect into
structural coherence, local predictive marginals, relational-null loss,
feasible interventions, persistence, and replacement identity.

\subtopic{89.7 Active shared context graphs for team science}

The closest operational neighbor located in this audit is Mycelium, introduced
by Choudhury and collaborators as an active shared workspace connecting
researchers and AI agents through an evolving context graph
\cite{ChoudhuryEtAl2026}. The system captures observations and hypotheses,
relates them to a team's developing knowledge model, and routes them to the
person or agent whose next decision they can inform. Its computational account
frames networked intelligence as sparse conditional computation over
distributed scientific contexts and explicitly asks when a standalone agent is
sufficient and when network structure is essential.

The overlap is material:
\begin{center}
\small
\begin{tabularx}{0.94\textwidth}{@{}p{0.16\textwidth}X@{}}
\toprule
Mycelium & Active context graph, routing, team science, and an explicit
network-versus-standalone comparison. \\
SRWT & Typed transport, invariant core, coalition tests, causal nulls,
replacement identity, formation provenance, and maintenance gates. \\
\bottomrule
\end{tabularx}
\end{center}
Mycelium rules out any priority claim stated merely as ``an active shared
workspace for human--AI science.'' The remaining candidate contribution is the
mathematical and empirical certificate applied to such an architecture.

\subtopic{89.8 Current multi-agent formalization}

Multi-agent research contains extensive models of planning, communication,
consensus, epistemic coordination, and decentralized control. The present
review did not identify one established framework that jointly uses all of the
following as necessary gates:
\[
\begin{aligned}
&\text{heterogeneous core transport and residuals},\\
&\text{exact or spectral global-section structure},\\
&\text{predictive sufficiency and feasible causal interventions},\\
&\text{coalition irreducibility under matched predictive nulls},\\
&\text{persistent identity under replacement},\\
&\text{formation provenance and leakage controls},\\
&\text{goal-conditioned repair of the certificate itself}.
\end{aligned}
\]
This negative search result is provisional. It supports a synthesis claim, not
an absolute claim of firstness.

\topic{90}{Closest Formal Neighbor: Collective Counterfactual Planning}

\subtopic{90.1 Why the comparison matters}

Collective Counterfactual Planning (CCP) is a 2026 formal model in which each
agent represents an orthogonal projection of a common task space. It separates
conception, consent, implementation, and verification, and proves that
cross-agent relay can enable plans unavailable to any one-shot individual tree
while team-dark task requirements remain unverifiable
\cite{Amornbunchornvej2026}. It is the closest current formal neighbor located
by this audit because representational geometry, not merely missing
information or capability, is the central limitation.

\subtopic{90.2 Shared commitments}

CCP and the theory share several commitments:
\begin{itemize}
\item no single member need represent the complete task;
\item collective reach may exceed one-shot pooling;
\item heterogeneous representation is mathematically explicit;
\item validation or certification is distinct from accidental task success;
\item subteams and verification roles can be forced by structure; and
\item a team's blind subspace creates principled impossibility results.
\end{itemize}
These overlaps narrow the defensible novelty of this paper. It
cannot claim to be the only formal theory of collective action under
representational constraints.

\subtopic{90.3 Different primary objects}

The theories nevertheless address different mathematical objects:
\begin{center}
\small
\begin{tabularx}{\textwidth}{@{}>{\raggedright\arraybackslash}p{0.18\textwidth}>{\raggedright\arraybackslash}X>{\raggedright\arraybackslash}X@{}}
\toprule
Dimension & Collective Counterfactual Planning & Stratified Relational Workspace Theory \\
\midrule
Local representation & Orthogonal subspace of one common task space & Heterogeneous local spaces, core--residual decompositions, and transfer maps \\
Primary question & Whether a bounded-horizon plan can be conceived, consented to, executed, and validly verified & Whether a persistent distributed operative state exists, acts, survives replacement, and possibly maintains itself \\
Global mechanism & Relay closure of projected planning trees & Inverse limits, holonomy, transfer Laplacians, descent, predictive alignment, and causal gates \\
Failure object & Dark subspace, consent failure, implementation failure, or missing verifier & Incoherence, local sufficiency, matched-null collapse, causal failure, identity drift, provenance leakage, or repair replay \\
Temporal scope & Horizon-bounded task planning & Recurrent dynamics, persistence, handoff, replacement, damage, and repair \\
Formation & Representation geometry is part of the declared model & Scaffolded, adaptively learned, or endogenously constituted organization is separately audited \\
Agency claim & Cooperative solvability under declared assumptions & Domain-relative workspace agency only after goal, value, action, repair, and distributed-control gates \\
\bottomrule
\end{tabularx}
\end{center}
The comparison indicates complementarity rather than subsumption. CCP could
serve as a planning layer inside a relational workspace. Conversely, an SRW
could exist without solving a CCP instance.

\subtopic{90.4 Priority consequence}

After Mycelium and the shared-workspace systems above, the novelty residual
should not be stated as ``human and AI form an active shared scientific
workspace.'' After CCP, it should not be stated as ``groups solve tasks no
member can represent.'' Both the operational architecture and a rigorous
representational planning model are already present. The narrower candidate residual is:
\[
\boxed{
\begin{aligned}
&\text{persistent core--transfer workspace structure}\\
&+\text{causal and coalition anti-triviality}\\
&+\text{identity under replacement}\\
&+\text{formation provenance}\\
&+\text{certificate-directed self-maintenance}\\
&+\text{robustness across admissible decompositions}.
\end{aligned}
}
\]
This is a stronger and more defensible formulation.

\topic{91}{Theorem-by-Theorem Dependency and Priority Register}

\subtopic{91.1 Register}

The following table records the most conservative current classification. A
``candidate'' entry means that the exact integration was not located in the
reviewed sources. It does not assert exhaustive global priority.

\small
\begin{longtable}{@{}>{\raggedright\arraybackslash}p{0.08\textwidth}>{\raggedright\arraybackslash}p{0.25\textwidth}>{\raggedright\arraybackslash}p{0.24\textwidth}>{\raggedright\arraybackslash}p{0.33\textwidth}@{}}
\toprule
Claim & Present result & Main prior machinery & Current assessment \\
\midrule
\endfirsthead
\toprule
Claim & Present result & Main prior machinery & Current assessment \\
\midrule
\endhead
\(T_1\) & Limit--holonomy equivalence & Categorical limits, local systems, holonomy fixed spaces & \(\mathsf{STD}/\mathsf{ADP}\). Standard theorem in workspace notation. \\
\(T_2,T_{2a}\) & Operation closure and commutant & Intertwiners and representation commutants & \(\mathsf{STD}/\mathsf{ADP}\). \\
\(T_3\) & Transfer-Laplacian kernel & Sheaf and connection Laplacians & \(\mathsf{STD}/\mathsf{ADP}\). \\
\(T_4,T_5\) & Spectral stability and gauge covariance & Davis--Kahan, Kato, connection Laplacians & \(\mathsf{STD}/\mathsf{ADP}\). \\
\(T_{10},T_{11}\) & Approximate operation closure and path leakage & Spectral projection estimates and telescoping block products & Mostly \(\mathsf{ADP}\). Exact constants may be manuscript-specific, but the method is standard. \\
\(T_{12}\) & Minimal predictive state and alignment & Sufficient statistics and computational mechanics & Minimality is \(\mathsf{STD}\). Alignment with a transfer workspace is \(\mathsf{ADP}/\mathsf{SYN}\). \\
\(T_{13}\) & Operational relationality & Conditional information, coalition risk, SCM interventions & Individual identities are standard. Their conjunctive matched-null certificate is \(\mathsf{SYN}/\mathsf{CAN}\). \\
\(T_6\) & Moving-subspace persistence & Stability recurrences and perturbation theory & \(\mathsf{STD}/\mathsf{ADP}\). The explicit non-collapse gate is a modeling synthesis. \\
\(T_7\) & Stable coalition reconstruction & Frame theory and erasure recovery & \(\mathsf{STD}/\mathsf{ADP}\). \\
\(T_{14}\) & Replacement invariance profile & Perturbation, process equivalence, causal comparison & Bound ingredients are standard. The coordinatewise structural--causal--reconstructive identity profile is \(\mathsf{SYN}/\mathsf{CAN}\). \\
\(T_8,T_9\) & Gluing and hierarchical collapse & Sheaf descent and \v{C}ech theory & Exact results are \(\mathsf{STD}/\mathsf{ADP}\). Descent plus anti-centralization and causal alignment is \(\mathsf{SYN}\). \\
\(T_{15}\)--\(T_{17}\) & Transfer learning, selection, boundary recovery & System identification and finite-margin selection & \(\mathsf{STD}/\mathsf{ADP}\). Gauge and leakage audit is integrative. \\
\(T_{18}\) & Formation transition example & Standard pitchfork bifurcation & \(\mathsf{STD}\). Included as a solvable benchmark, not a priority claim. \\
\(T_{19}\)--\(T_{22}\) & Goal retention, regret, repair, maintenance & Linear systems, decision theory, viability, feedback control & Mathematical bounds are \(\mathsf{STD}/\mathsf{ADP}\). Applying them to certificate repair is \(\mathsf{SYN}\). \\
\(T_{23}\) & Workspace-agency synthesis & Agency theory, causal control, viability & The conjunctive domain-relative certificate is \(\mathsf{SYN}/\mathsf{CAN}\). \\
\(T_{24}\) & Gatewise false-certification bound & Intersection testing and elementary probability & \(\mathsf{STD}\). The gate architecture is the relevant synthesis. \\
\(T_{25}\)--\(T_{27}\) & Faithful coarse-graining, risk transport, and decomposition robustness & Kernel factorization, sigma-field aggregation, and margin stability & Exact identities are \(\mathsf{STD}/\mathsf{ADP}\). Their integration with anti-centralization and certificate envelopes is \(\mathsf{SYN}/\mathsf{CAN}\). \\
\(T_{28}\)--\(T_{32}\) & Specification transport, non-identifiability, partial identification, uniform certification, and independent selection & Sensitivity analysis, experimental equivalence, identified sets, uniform concentration, and post-selection inference & The statistical principles are \(\mathsf{STD}/\mathsf{ADP}\). Their gatewise integration with the full relational certificate is \(\mathsf{SYN}/\mathsf{CAN}\). \\
Part XXII & Positive cases and near misses & Distributed parity, frames, causal tests & Toy calculations are standard. The common positive/near-miss benchmark architecture is \(\mathsf{SYN}/\mathsf{CAN}\). \\
\bottomrule
\end{longtable}
\normalsize

\subtopic{91.2 Standard ingredients that must be cited, not claimed}

The following should always be described as inherited machinery:
\[
\begin{aligned}
&\text{inverse limits, global sections, holonomy invariants, commutants},\\
&\text{sheaf and connection Laplacians, \v{C}ech gluing, spectral gaps},\\
&\text{Davis--Kahan and Kato perturbation, frame reconstruction, erasures},\\
&\text{minimal sufficient states, conditional information, SCM interventions},\\
&\text{least squares, ridge recovery, finite-margin model selection},\\
&\text{sensitivity analysis, partial identification, uniform concentration, post-selection inference},\\
&\text{pitchfork bifurcation, feedback stability, viability, policy regret}.
\end{aligned}
\]
Renaming these objects does not create novelty.

\subtopic{91.3 Domain adaptations}

The following are legitimate adaptations but weak standalone novelty claims:
\begin{itemize}
\item interpreting local-system sections as shared workspace cores;
\item treating the connection or sheaf Laplacian as a transfer-coherence
 operator;
\item using frame bounds to quantify coalition reconstruction;
\item using causal states as operational workspace summaries;
\item using gauge covariance to compare participant or model realizations;
\item using \v{C}ech descent to form workspaces of workspaces;
\item applying specification envelopes and identified sets to a conjunctive
 workspace certificate; and
\item applying viability and feedback bounds to workspace-health margins.
\end{itemize}
Their value depends on the integration they enable.

\subtopic{91.4 Strongest candidate synthesis}

The strongest current candidate is not a single equation. It is the following
conjunctive object:
\[
\boxed{
\mathfrak C_{\mathrm{SRW}}
=
\left(
\begin{array}{c}
\text{heterogeneous core--transfer diagram and residuals},\\
\text{exact or spectral shared-section state},\\
\text{operation and predictive alignment},\\
\text{feasible causal interventions},\\
\text{coalition irreducibility under matched local-predictive nulls},\\
\text{non-collapse, reconstruction, and replacement identity},\\
\text{descent across subworkspaces without hidden centralization},\\
\text{formation provenance and certification leakage control},\\
\text{goal-conditioned action and certificate-directed repair},\\
\text{robustness across admissible decompositions and scientific specifications},\\
\text{identified-set and uniform finite-sample certification}
\end{array}
\right).
}
\]
The reviewed sources contain many of these coordinates separately and some in
pairs. This audit did not locate a prior source that defines their conjunction
as one falsifiable workspace certificate.

\subtopic{91.5 Why the conjunction is not arbitrary}

Each coordinate excludes a distinct near miss:
\[
\begin{array}{ll}
\text{coherence} & \text{rejects incompatible local states},\\
\text{predictive alignment} & \text{rejects inert structure},\\
\text{causal efficacy} & \text{rejects correlational epiphenomena},\\
\text{coalition irreducibility} & \text{rejects centralized sufficiency},\\
\text{matched nulls} & \text{reject local-marginal explanations},\\
\text{non-collapse} & \text{rejects convergence to a trivial zero state},\\
\text{replacement identity} & \text{rejects task-equivalent but causally different successors},\\
\text{provenance} & \text{rejects scaffold-decoded organization},\\
\text{repair specificity} & \text{rejects passive robustness and open-loop replay},\\
\text{specification robustness} & \text{rejects favorable analytical tuning},\\
\text{identifiability} & \text{rejects conclusions unsupported by the experiment family}.
\end{array}
\]
Because the excluded systems differ, the certificate is not equivalent to any
one inherited component.

\topic{92}{Scope of the Originality Claim}

\subtopic{92.1 Defensible thesis}

The strongest defensible originality statement is:
\begin{quote}
Stratified Relational Workspace Theory proposes a unified formal and empirical
certificate for persistent operative structures distributed across
heterogeneous agents and artifacts. It integrates core--transfer and
local-to-global geometry with predictive sufficiency, feasible causal
interventions, coalition irreducibility, matched relational nulls,
reconstructive identity, formation provenance, certificate-directed
self-maintenance, robustness across admissible decompositions and scientific
specifications, and identified finite-sample certification. The individual
mathematical tools are largely established.
The candidate contribution is their noncompensatory integration and the class
of near misses it is designed to reject.
\end{quote}
This statement is narrower than claiming a new theory of cognition in general.
It is also stronger than claiming only a metaphor.

\subtopic{92.2 Claims excluded by the comparison}

The theory does not claim:
\begin{enumerate}
\item that distributed or extended cognition is newly discovered;
\item that global sections, holonomy, sheaf Laplacians, spectral perturbation,
 frames, causal states, or viability bounds are new mathematics;
\item that all human--AI collaboration forms a relational workspace;
\item that group performance establishes relational irreducibility;
\item that scaffolded organization is unreal;
\item that endogenous formation establishes agency;
\item that agency establishes consciousness, personhood, or moral status; or
\item that the motivating collaboration has already passed the empirical
 certificate.
\end{enumerate}

\subtopic{92.3 Claims supported at present}

The analysis supports the following claims:
\begin{enumerate}
\item One coherent formal language can express exact structure, approximate
 coherence, causal operation, distributed sufficiency, persistence,
 identity, formation, and repair.
\item Under its declared assumptions, the stated mathematical theorems follow.
\item The common certificate rejects analytic replicated, centralized,
 obstructed, externally reloaded, identity-breaking, and replay-based
 near misses.
\item The motivating collaboration is a plausible observational candidate but
 remains uncertified.
\item The synthesis yields concrete experiments whose failure would narrow or
 falsify the stronger relational claims.
\end{enumerate}

\subtopic{92.4 Remaining priority risks}

The highest-priority remaining searches are:
\begin{itemize}
\item sheaf-theoretic models of team cognition and organizational knowledge;
\item distributed sufficient-statistic and Blackwell-order models with
 explicit coalition nulls;
\item causal representation and multi-agent latent-state identity under
 replacement;
\item formal theories of transactive memory with local-to-global topology;
\item organizational cybernetics that couples self-maintenance to the
 preservation of distributed informational organization; and
\item post-August-2026 work that combines representational geometry with
 persistent human--AI or multi-agent agency;
\item formal analysis of active shared context graphs, shared group memory, and
 artifact-centric graph workspaces under coalition and intervention tests.
\end{itemize}
A specialist review from each adjacent field may uncover precedents not visible
through cross-disciplinary keyword search.

\subtopic{92.5 Evidence required for a stronger claim}

Before presenting the synthesis as established rather than proposed, the
following should be completed:
\begin{enumerate}
\item a public, timestamped manuscript and versioned source archive;
\item theorem-level comparison with the closest sheaf, causal-state, group
 agency, and CCP formulations;
\item independent mathematical review of the nonstandard claims and constants;
\item preregistered benchmark gates and matched null families;
\item at least one controlled human--AI or multi-agent experiment estimating
 the complete certificate vector; and
\item explicit reporting of failures, unidentified quantities, and
 decomposition sensitivity.
\end{enumerate}

\subtopic{92.6 Priority conclusion}

The literature review narrows rather than weakens the theory. It removes
unnecessary novelty claims from the standard mathematics and leaves a clearer
research object:
\centerbox{\large\bfseries
Not a new Laplacian, sheaf, causal state, or controller, but a candidate
unified certificate for relational workspaces.
}
That is the priority position adopted here.


\theorypart{XXIV}{Theory Hierarchy}

\topic{93}{Levels of Organization}

\subtopic{Level 0: Archive}
\[
\delta_{\mathrm{coh}}\text{ may be low},
\qquad
\eta_{\mathrm{causal}}=0.
\]

\subtopic{Level 1: Shared representation}
\[
d_\varepsilon>0,
\qquad
\delta_{\mathrm{coh}}\leq\varepsilon.
\]

\subtopic{Level 2: Coordinated operational workspace}
\[
\delta_{\mathrm{align}}^{\mathcal I}\leq\varepsilon,
\qquad
\eta_{\mathrm{causal}}^\nu>0.
\]

\subtopic{Level 3: Relational workspace}
\[
k_\varepsilon\geq2,
\qquad
\Delta_{\mathrm{joint}}^\ell>0,
\qquad
\Delta_{\mathrm{tr}}^{\min}(\mathcal Q)>0.
\]

\subtopic{Level 4: Persistent relational workspace}
\[
d_{\mathrm{id}}(W_t,W_{t+\tau})
\leq\varepsilon.
\]

\subtopic{Level 5: Self-maintaining workspace}
\[
\eta_{D\to M}>0,
\qquad
\eta_{M\to R}>0,
\qquad
\eta_{R\to v}>0,
\qquad
\Delta_{\mathrm{repair}}^{\mathrm{spec}}>0,
\qquad
\sup_t e_t\leq e_{\max}.
\]

\subtopic{Level 6: Workspace agent}
\[
\begin{aligned}
&\eta_{G\to A}>0,
\qquad
\eta_{W\to A}>0,
\qquad
\delta_Q\leq\varepsilon_Q,\\
&\operatorname{Reg}_t(\pi)
 \leq2\varepsilon_Q+\delta_{\mathrm{opt}},
\qquad
k_{\mathrm{ctl},\varepsilon}\geq2,
\qquad
\operatorname{SelfMaintenance}.
\end{aligned}
\]

\subtopic{Orthogonal formation-provenance axis}
\[
\operatorname{Scaffolded}
\longrightarrow
\operatorname{AdaptivelyParameterised}
\longrightarrow
\operatorname{StructurallyAdaptive}
\longrightarrow
\operatorname{EndogenouslyConstituted}.
\]
This axis is independent of Levels~0--6. A scaffolded system can instantiate a
high-level relational workspace, while an endogenously learned organization
can fail every operational workspace gate.

\subtopic{Decomposition-robustness qualifier}

Every level is first defined relative to a decomposition \(\mathcal D\). A
family-robust version carries the additional requirement
\[
\inf_{\mathcal D\in\mathfrak D}
 m_q(W;\mathcal D)>0
\qquad
\text{for every gate required at that level}.
\]
This is a robustness qualifier rather than a higher agency level. A system can
be highly agentive under one fixed interface while remaining decomposition
fragile.

\subtopic{Specification-robustness and identification qualifier}

Every level is also relative to a complete scientific specification
\(\varsigma\). A specification-robust version requires
\[
\inf_{\varsigma\in\mathfrak S}
 m_q(W;\varsigma)>0
\qquad
\text{for every gate required at that level}.
\]
An identified version additionally requires the sign of each required lower
envelope to be constant across the realization set compatible with the
available experiment family. An empirical claim must finally subtract the
uniform finite-sample error radius before declaring the lower envelope
positive. These are epistemic qualifiers, not additional agency levels.

\subtopic{Separate, unresolved axis}
\[
\operatorname{PhenomenalConsciousness}.
\]

\theorypart{XXV}{Unified Formulation}

The complete static-to-agentive construction is
\[
\boxed{
\begin{aligned}
&
(\Omega,\mathfrak D,\mathfrak S,\Sigma,\mathcal D_{\mathrm{form}},\mathcal L)
\\[2mm]
&\Downarrow
\qquad
(\mathcal D,\widehat V,\widehat E,\widehat H,\widehat P,\widehat B)
\\[2mm]
&\Downarrow
\qquad
(G,\{X_i\},\{T_e\})
\\[2mm]
&\Downarrow
\qquad
X_i=C_i\oplus R_i
\\[2mm]
&
(\{C_i\},\{H_e\},\mathcal A,\{\rho_i\})
\\[2mm]
&\Downarrow
\qquad
\Gamma(H)=\varprojlim C_i
\\[2mm]
&\Downarrow
\qquad
L_H=D_H^*D_H
\\[2mm]
&\Downarrow
\qquad
\Gamma_\varepsilon(H)
=
\ran\,\mathbf 1_{[0,\varepsilon]}(L_H)
\\[2mm]
&\Downarrow
\qquad
W_t^{\mathrm{str}}
\\[2mm]
&\xrightarrow{\text{alignment}}
Z_t=\mathbb P(Y_t^+\in\cdot\mid\mathbf H_t)
\\[2mm]
&\Downarrow
\qquad
\SRW_\Theta(W_t;\mathcal D)
\\[2mm]
&\Downarrow
\qquad
(W_t,G_t,\widehat Q_t,M_t,B_t)
\\[2mm]
&\xrightarrow{\text{goal and value use}}
(A_t,R_t)
\\[2mm]
&\xrightarrow{\text{causal repair}}
e_t\leq e_{\max}
\\[2mm]
&\Downarrow
\qquad
\Workspace
=
[\mathfrak R]_{\text{gauge}+\text{causal equivalence}}.
\end{aligned}
}
\]
Compactly:
\[
\boxed{
\begin{aligned}
\text{scaffold and history}
&\to
\text{learned organization and boundary}
\\
&\to
\text{local core--transfer realization}
\\
&\to
\text{compatible sections and global descent}
\\
&\to
\text{causal operation and persistent identity}
\\
&\to
\text{cue-controlled goal retention and counterfactual action}
\\
&\to
\text{monitor-mediated repair and bounded maintenance deficit}
\\
&\to
\text{decomposition- and specification-robust identified certification}.
\end{aligned}
}
\]
The final two arrows are optional extensions rather than consequences of the
earlier ones. A relational workspace can exist without being an agent, and an
agentive controller can fail to maintain the workspace that realizes it.


\theorypart{XXVI}{Discussion, Limitations, and Philosophical Consequences}

\topic{94}{Decomposition Relativity}

Part~XVII turns decomposition relativity from a caveat into a formal part of the
certificate. The central lesson is negative and positive at once. No claim of
relationality is invariant under arbitrary redescription. An indiscrete merge
makes the whole state local, while a secret-sharing split can make a centralized
encoder look irreducibly distributed. Yet exact structural content can remain
invariant under faithful coarse-graining, and gate margins can remain stable
through a declared family of operationally meaningful decompositions.

The resulting position is neither that decomposition is arbitrary nor that one
partition is uniquely privileged. Decompositions are constrained by access,
causal interfaces, intervention separability, temporal continuity, and explicit
shared resources. A certificate established only at one interface is a valid
but limited claim. A stronger claim reports the full envelope
\[
\left(
\widetilde m_q,
\underline m_q,
\overline m_q,
\Delta_q^{\mathrm{dec}}
\right)
\]
for each required gate, where \(\tilde m_q\) is the reference margin and the
remaining terms summarize its behavior across \(\mathfrak D\).

This does not eliminate judgment. It makes the judgment inspectable. Analysts
can disagree about the admissible family and identify exactly which merge,
split, hidden resource, or intervention assumption changes the conclusion.
That is preferable to treating a convenient node diagram as if it were part of
nature itself.

\subtopic{94.1 Specification relativity and epistemic restraint}

Decomposition is only one coordinate of the scientific claim. The downstream
target, loss, temporal horizon, intervention family, null panel, and threshold
rule also determine what is being certified. A workspace may be operationally
real for a month-long theorem-development trajectory yet irrelevant to a
second-scale motor target. Those are distinct specifications, not competing
answers to one underspecified question.

The framework therefore distinguishes truth under a declared specification
from robustness across an admissible family. It also distinguishes robustness
inside one fitted realization from identification across all realizations
compatible with the available experiments. A positive point estimate at one
favorable analysis is the weakest admissible conclusion. A stronger claim
requires a positive lower envelope across \(\mathfrak S\), an identified sign
across the compatible realization set, and a finite-sample confidence bound
that remains positive after uniform error control.

This hierarchy is intentionally conservative. It does not demand invariance
under every imaginable redescription or target. It demands that the scientific
universe over which robustness is claimed be declared before certification and
that ambiguity outside the reach of the experiment be reported as partial or
non-identification rather than converted into certainty.

\topic{95}{Realism Without Reification}

The framework is realist about relational organization but does not posit a free-standing substance
between the components. A workspace exists only through a realization, and every change in the
workspace must be implemented by a change in participants, artifacts, channels, or their states. The
quotient
\[
\Workspace=[\mathfrak R]_{\simeq}
\]
is therefore a mathematical account of identity across realizations, not the name of an additional
physical object.

This position allows two claims that are sometimes treated as incompatible. First, the workspace is
physically dependent: no realization, no workspace. Second, it is not identical with any particular
carrier when several distinct realizations preserve the same structural and causal organization. The same
logic is familiar whenever an object is identified by invariants rather than by one material instance. A
melody is not a second substance beside its performances, yet it is not identical with one performance.
A proof is not separate from every inscription that carries it, yet its identity can survive changes of
notation and medium. The analogy is limited, but it clarifies the intended form of realism.

The theory is therefore compatible with physicalism and multiple realization. It rejects both the view
that a workspace must be wholly inside one component and the stronger view that the workspace is an
independent entity floating outside its carriers. What is real is the organized pattern of compatible
representation, transfer, causal use, persistence, and reconstruction.

\topic{96}{Location and the Sense of ``Between''}

Asking where a relational workspace is located can invite a category mistake. The local states have
ordinary physical locations. A notebook is on a desk, a model state is implemented on hardware, a
message is transmitted through a channel, and a person's activity is realized in a biological organism.
The workspace as such is not located at an additional point. Its location is distributed in the modest
sense that the relations constituting it span several located carriers.

The word ``between'' is best understood counterfactually rather than spatially. Suppose local marginal
states are held fixed while transfer organization is shuffled, delayed, or replaced by a matched null. If
the joint predictive or causal capacity changes, then an operative property resided in the organization of
the relations rather than in any one local state. Conversely, if a single component or matched storage
surrogate reproduces the complete capacity, the case for a relational workspace fails. The relevant
question is therefore not whether an invisible object occupies the interval between participants. It is
whether the coupling carries indispensable state and supports capacities that disappear when that
coupling is altered.

This account also explains why the workspace can be partially present in each component without being
fully stored in any of them. Each local chart may encode a view, role, or constraint. The global object is
the compatible and causally operative organization of those views. In exact settings this is a global
section. In approximate settings it is a stable low-energy subspace together with the dynamic and causal
conditions developed above.

\topic{97}{Continuity, Contribution, and Shared Ownership}

Relational persistence does not require an unchanging set of participants. A workspace can survive
replacement when the new realization preserves the relevant transfer structure, operation algebra,
causal response, reconstructive capacity, and non-collapse margin. This is continuity through organized
transport, not continuity through an immutable token.

The same point applies to artifacts. A handoff document, theorem register, code repository, or experiment
log is not merely a passive record when it constrains future reasoning and permits reconstruction after
interruption. Such artifacts can become constitutive parts of a workspace realization. Their contribution
is causal and testable: remove or corrupt them, then measure what can no longer be reconstructed or done.

Distributedness does not imply equal contribution. One participant may supply long-horizon goals, another
formal manipulation, and an artifact precise persistent memory. Nor does a system-level description erase
individual authorship or responsibility. The theory describes how capacities are organized; it does not
settle legal, ethical, or scholarly questions of credit. Those questions require a separate account of who
introduced which ideas, who exercised control, who bore risk, and which contributions were replaceable.

The phrase ``shared ownership'' should therefore be used carefully. No individual component need own a
complete copy of the workspace state, yet individuals can still own artifacts, originate arguments, and
retain responsibility for decisions. Relational identity and ordinary ownership are different relations.

\topic{98}{Agency, Consciousness, and Moral Non-Implications}

The agency certificate developed here is functional and domain relative. It concerns persistent goals,
interventional value estimates, feasible action control, damage detection, repair, and bounded maintenance
within a declared environment. Passing those tests would justify a claim of workspace-level agency in that
domain. It would not establish unrestricted autonomy.

No theorem in the framework identifies agency with phenomenal consciousness. A system may maintain a
goal, evaluate counterfactual actions, and repair its own organization without there being anything it is
like to be that system. The converse also remains open: failure to satisfy this engineering certificate
would not prove the absence of consciousness. The framework is designed to prevent functional evidence
from silently becoming a metaphysical conclusion.

Moral status is likewise not inferred. Neither coherence, causal efficacy, relational irreducibility, nor
self-maintenance is by itself a criterion of moral patienthood. Future arguments connecting these properties
to consciousness or welfare would need independent premises and evidence. Until then, the framework can
support careful descriptions of organization and agency without treating a collaborative system as a
person or dismissing the moral standing of the humans who participate in it.

\topic{99}{Limitations}

The theory is deliberately broad, but its conclusions remain bounded by several limitations.

\begin{enumerate}
\item \textbf{Admissible-decomposition burden.} Part~XVII formalizes faithful coarse-graining and
family-robust certification, but the choice of primitive carriers and admissible alternatives remains an
empirical and methodological judgment. A narrow family can still overstate robustness.

\item \textbf{Admissible-specification burden.} Part~XVIII makes the target, loss, horizon,
intervention family, null panel, and thresholds explicit. The admissible family \(\mathfrak S\) still
requires scientific judgment. A family chosen after seeing the results can conceal fragility even when
every reported analysis is technically correct.

\item \textbf{Limited observability and partial identification.} Many practical systems expose messages
and outputs rather than complete local states. Transfer maps, predictive states, reconstruction operators,
and even the certificate sign may therefore be only partially identified. A point claim is unjustified
when both passing and failing realizations remain compatible with the declared experiment.

\item \textbf{Linear and finite-dimensional emphasis.} The clearest structural theorems use finite graphs
and linear Hilbert-space models. Nonlinear, stochastic, symbolic, and category-valued extensions require
additional work.

\item \textbf{Gauge and model ambiguity.} Even with abundant data, local coordinates need not be
identifiable. Learned transfers are often meaningful only up to a gauge class and relative to a selected
model family.

\item \textbf{Causal identification burden.} Workspace interventions must be feasible and matched.
Observational dependence, activation, or predictive accuracy cannot substitute for interventions when a
causal claim is made.

\item \textbf{Certification cost.} The complete certificate requires coalition tests, null models,
replacement studies, longitudinal interventions, admissible specification sweeps, and identified-set or
uniform uncertainty calculations. Exhaustive certification may be impractical in large systems, and
adaptive exploration can invalidate nominal error guarantees unless selection and certification are
separated or valid selective procedures are used.

\item \textbf{Conditional mathematical results.} The theorems establish what follows from stated
assumptions. They do not show that naturally occurring collaborations satisfy those assumptions, and
several constants may be too conservative for practical estimation.

\item \textbf{Open empirical status.} The motivating human--AI collaboration has not been subjected to
the complete frozen certificate. It remains a plausible case study and a source of hypotheses, not a
demonstrated instance of relational agency.

\item \textbf{Open priority boundary.} The literature comparison is broad but cannot prove absolute
firstness. Closely related work may exist under other terminology or appear after the review date.
\end{enumerate}

These limitations do not make the framework vacuous. They specify the conditions under which its claims
can fail and identify the measurements needed to strengthen them.

\topic{100}{Open Problems}

Several directions are mathematically and empirically immediate.

\begin{enumerate}
\item Develop nonlinear and stochastic analogues of the transfer Laplacian, including kernel, manifold,
and category-valued realizations.

\item Establish sharp uniform finite-sample estimators for workspace dimension, coalition risk,
matched-null gaps, replacement identity, and specification envelopes under adaptive data collection.

\item Characterize the identified set when local states are latent and only messages, artifacts, or
actions are observed. Determine the minimum intervention family needed to identify each certificate gate
and the complete conjunctive certificate.

\item Extend gluing and persistence to changing graph topology, asynchronous communication, and
delayed or lossy transfers.

\item Extend the joint decomposition-and-specification theory to latent shared resources, overlapping
component systems, continuous coarse-graining, target families, and data-driven discovery of admissible
interface and specification families without invalidating certification.

\item Study the interaction between workspace formation and strategic participants whose local objectives
need not agree.

\item Build benchmark suites in which structural, causal, persistence, and repair quantities are known
exactly but remain nontrivial to infer.

\item Run preregistered human--AI and multi-agent experiments that estimate the complete certificate
rather than isolated performance gains.

\item Compare certificate-directed self-maintenance with established notions of organizational resilience,
institutional memory, and collective agency.

\item Clarify which additional evidence, if any, could connect functional workspace predicates to theories
of consciousness without assuming the connection in advance.
\end{enumerate}

\topic{101}{Conclusion}

A stratified relational workspace is a temporally persistent, causally efficacious, jointly reconstructible
structure represented only partially in several heterogeneous components. Its identity is determined by
invariants preserved under admissible transfer, and its operative capacity is relational only when no
individual component or matched local-storage null reproduces the relevant function.

The theory gives this definition several independent layers. Exact shared content appears as a global
compatible section and, in the invertible case, as a holonomy-fixed core. Approximate shared content
appears as a low-energy spectral subspace of the transfer Laplacian. Predictive states and feasible
interventions determine whether that structure matters for a declared future. Coalition tests and matched
nulls determine whether the capacity is genuinely distributed. Reconstruction frames, temporal transport,
and identity profiles determine whether it survives loss and replacement. Descent determines whether
local workspaces compose globally. Provenance tests distinguish learned organization from scaffolded
answers. Goal, value, damage, and repair interventions determine whether the workspace supports
domain-relative agency and self-maintenance.

No one layer can stand in for the others. Coherence can describe an archive. Performance can describe a
centralized controller. Persistence can be trivial if all content collapses. A task-equivalent replacement
can alter the causal organization. Local compatibility can fail to glue globally. Endogenous-looking
structure can be supplied through hidden leakage, and successful recovery can be open-loop replay rather
than repair. The value of the framework lies in making these alternatives explicit and testable.


Relationality also cannot be detached from decomposition. Faithful
coarse-graining preserves the exact structural core, but coalition counts and
local sufficiency change when components are merged. Arbitrary splitting can
create a formally irreducible code around a centralized encoder. The strongest
claims therefore quantify their margins across an admissible family of
interfaces rather than relying on one favorable diagram.

A complete scientific claim must also survive scrutiny of the specification and the experiment.
The strongest form of the certificate reports a positive gatewise lower envelope across a predeclared
family of decompositions and scientific specifications, shows that the sign of that envelope is identified
across all realizations compatible with the intervention data, and retains a positive margin after uniform
finite-sample uncertainty. When these conditions fail, the correct conclusion is specification-specific,
fragile, partially identified, or unidentified rather than simply positive or negative.

The central identity principle is therefore
\[
\boxed{
\text{workspace identity}
=
\text{persistent invariant causal organization under admissible transport}.
}
\]
This principle is compatible with changing carriers, representations, and participants, but only when
structural, operational, causal, reconstructive, and non-collapse defects remain within declared tolerances.
It provides a formal language for studying distributed collaboration without turning that collaboration
into a mysterious third substance and without inferring consciousness from organization alone.

\begin{thebibliography}{99}
\small

\bibitem{AbramskyBrandenburger2011}
S. Abramsky and A. Brandenburger.
\newblock The sheaf-theoretic structure of non-locality and contextuality.
\newblock \emph{New Journal of Physics}, 13:113036, 2011.
\newblock doi:10.1088/1367-2630/13/11/113036.

\bibitem{AdamsAizawa2001}
F. Adams and K. Aizawa.
\newblock The bounds of cognition.
\newblock \emph{Philosophical Psychology}, 14(1):43--64, 2001.
\newblock doi:10.1080/09515080120033571.

\bibitem{AmesEtAl2017}
A. D. Ames, X. Xu, J. W. Grizzle, and P. Tabuada.
\newblock Control barrier function based quadratic programs for safety critical systems.
\newblock \emph{IEEE Transactions on Automatic Control}, 62(8):3861--3876, 2017.
\newblock doi:10.1109/TAC.2016.2638961.

\bibitem{Amornbunchornvej2026}
C. Amornbunchornvej.
\newblock Collective counterfactual planning: Coordination, consent, and verification under representational constraints.
\newblock arXiv:2608.17932, 2026.
\newblock doi:10.48550/arXiv.2608.17932.

\bibitem{Ashby1956}
W. R. Ashby.
\newblock \emph{An Introduction to Cybernetics}.
\newblock Chapman \& Hall, London, 1956.

\bibitem{Aubin1991}
J.-P. Aubin.
\newblock \emph{Viability Theory}.
\newblock Birkh\"auser, Boston, 1991.
\newblock doi:10.1007/978-0-8176-4910-4.

\bibitem{Baars1988}
B. J. Baars.
\newblock \emph{A Cognitive Theory of Consciousness}.
\newblock Cambridge University Press, Cambridge, 1988.

\bibitem{BarandiaranEtAl2009}
X. E. Barandiaran, E. Di Paolo, and M. Rohde.
\newblock Defining agency: Individuality, normativity, asymmetry, and spatio-temporality in action.
\newblock \emph{Adaptive Behavior}, 17(5):367--386, 2009.
\newblock doi:10.1177/1059712309343819.

\bibitem{Blackwell1953}
D. Blackwell.
\newblock Equivalent comparisons of experiments.
\newblock \emph{Annals of Mathematical Statistics}, 24(2):265--272, 1953.
\newblock doi:10.1214/aoms/1177729032.


\bibitem{BerkEtAl2013}
R. Berk, L. Brown, A. Buja, K. Zhang, and L. Zhao.
\newblock Valid post-selection inference.
\newblock \emph{The Annals of Statistics}, 41(2):802--837, 2013.
\newblock doi:10.1214/12-AOS1077.

\bibitem{Bredon1997}
G. E. Bredon.
\newblock \emph{Sheaf Theory}.
\newblock 2nd edition, Springer, New York, 1997.

\bibitem{Bratman1992}
M. E. Bratman.
\newblock Shared cooperative activity.
\newblock \emph{The Philosophical Review}, 101(2):327--341, 1992.
\newblock doi:10.2307/2185537.

\bibitem{CasazzaKovacevic2003}
P. G. Casazza and J. Kova\v{c}evi\'c.
\newblock Equal-norm tight frames with erasures.
\newblock \emph{Advances in Computational Mathematics}, 18:387--430, 2003.
\newblock doi:10.1023/A:1021349819855.

\bibitem{Christensen2016}
O. Christensen.
\newblock \emph{An Introduction to Frames and Riesz Bases}.
\newblock 2nd edition, Birkh\"auser, Cham, 2016.

\bibitem{ClarkChalmers1998}
A. Clark and D. Chalmers.
\newblock The extended mind.
\newblock \emph{Analysis}, 58(1):7--19, 1998.
\newblock doi:10.1093/analys/58.1.7.

\bibitem{CoverThomas2006}
T. M. Cover and J. A. Thomas.
\newblock \emph{Elements of Information Theory}.
\newblock 2nd edition, Wiley, Hoboken, 2006.

\bibitem{DavisKahan1970}
C. Davis and W. M. Kahan.
\newblock The rotation of eigenvectors by a perturbation. III.
\newblock \emph{SIAM Journal on Numerical Analysis}, 7(1):1--46, 1970.
\newblock doi:10.1137/0707001.

\bibitem{DehaeneChangeux2011}
S. Dehaene and J.-P. Changeux.
\newblock Experimental and theoretical approaches to conscious processing.
\newblock \emph{Neuron}, 70(2):200--227, 2011.
\newblock doi:10.1016/j.neuron.2011.03.018.

\bibitem{DellermannEtAl2019}
D. Dellermann, P. Ebel, M. S\"ollner, and J. M. Leimeister.
\newblock Hybrid intelligence.
\newblock \emph{Business \& Information Systems Engineering}, 61:637--643, 2019.
\newblock doi:10.1007/s12599-019-00595-2.

\bibitem{DiPaolo2005}
E. A. Di Paolo.
\newblock Autopoiesis, adaptivity, teleology, agency.
\newblock \emph{Phenomenology and the Cognitive Sciences}, 4:429--452, 2005.
\newblock doi:10.1007/s11097-005-9002-y.

\bibitem{FaginEtAl1995}
R. Fagin, J. Y. Halpern, Y. Moses, and M. Y. Vardi.
\newblock \emph{Reasoning About Knowledge}.
\newblock MIT Press, Cambridge, Massachusetts, 1995.

\bibitem{GuptaEtAl2025}
P. Gupta, T. Nguyen, C. Gonzalez, and A. W. Woolley.
\newblock Fostering collective intelligence in human--AI collaboration: Laying the groundwork for COHUMAIN.
\newblock \emph{Topics in Cognitive Science}, 17(2):189--216, 2025.
\newblock doi:10.1111/tops.12679.

\bibitem{HalpernMoses1990}
J. Y. Halpern and Y. Moses.
\newblock Knowledge and common knowledge in a distributed environment.
\newblock \emph{Journal of the ACM}, 37(3):549--587, 1990.
\newblock doi:10.1145/79147.79161.

\bibitem{HansenGhrist2019}
J. Hansen and R. Ghrist.
\newblock Toward a spectral theory of cellular sheaves.
\newblock \emph{Journal of Applied and Computational Topology}, 3:315--358, 2019.
\newblock doi:10.1007/s41468-019-00038-7.

\bibitem{Hatcher2002}
A. Hatcher.
\newblock \emph{Algebraic Topology}.
\newblock Cambridge University Press, Cambridge, 2002.

\bibitem{HayesRoth1985}
B. Hayes-Roth.
\newblock A blackboard architecture for control.
\newblock \emph{Artificial Intelligence}, 26(3):251--321, 1985.
\newblock doi:10.1016/0004-3702(85)90063-3.

\bibitem{HollanEtAl2000}
J. Hollan, E. Hutchins, and D. Kirsh.
\newblock Distributed cognition: Toward a new foundation for human-computer interaction research.
\newblock \emph{ACM Transactions on Computer-Human Interaction}, 7(2):174--196, 2000.
\newblock doi:10.1145/353485.353487.

\bibitem{Hutchins1995}
E. Hutchins.
\newblock \emph{Cognition in the Wild}.
\newblock MIT Press, Cambridge, Massachusetts, 1995.

\bibitem{Kato1995}
T. Kato.
\newblock \emph{Perturbation Theory for Linear Operators}.
\newblock Springer, Berlin, 2nd edition, 1995.

\bibitem{KlyubinEtAl2005}
A. S. Klyubin, D. Polani, and C. L. Nehaniv.
\newblock Empowerment: A universal agent-centric measure of control.
\newblock In \emph{2005 IEEE Congress on Evolutionary Computation}, pages 128--135, 2005.
\newblock doi:10.1109/CEC.2005.1554676.

\bibitem{ImbensManski2004}
G. W. Imbens and C. F. Manski.
\newblock Confidence intervals for partially identified parameters.
\newblock \emph{Econometrica}, 72(6):1845--1857, 2004.
\newblock doi:10.1111/j.1468-0262.2004.00555.x.

\bibitem{Leamer1983}
E. E. Leamer.
\newblock Let's take the con out of econometrics.
\newblock \emph{The American Economic Review}, 73(1):31--43, 1983.

\bibitem{LarsenSkou1991}
K. G. Larsen and A. Skou.
\newblock Bisimulation through probabilistic testing.
\newblock \emph{Information and Computation}, 94(1):1--28, 1991.
\newblock doi:10.1016/0890-5401(91)90030-6.

\bibitem{Lewis2003}
K. Lewis.
\newblock Measuring transactive memory systems in the field: Scale development and validation.
\newblock \emph{Journal of Applied Psychology}, 88(4):587--604, 2003.
\newblock doi:10.1037/0021-9010.88.4.587.

\bibitem{ListPettit2011}
C. List and P. Pettit.
\newblock \emph{Group Agency: The Possibility, Design, and Status of Corporate Agents}.
\newblock Oxford University Press, Oxford, 2011.

\bibitem{Ljung1999}
L. Ljung.
\newblock \emph{System Identification: Theory for the User}.
\newblock 2nd edition, Prentice Hall, Upper Saddle River, New Jersey, 1999.

\bibitem{MacLane1998}
S. Mac Lane.
\newblock \emph{Categories for the Working Mathematician}.
\newblock 2nd edition, Springer, New York, 1998.

\bibitem{MaturanaVarela1980}
H. R. Maturana and F. J. Varela.
\newblock \emph{Autopoiesis and Cognition: The Realization of the Living}.
\newblock D. Reidel, Dordrecht, 1980.

\bibitem{OlfatiSaberEtAl2007}
R. Olfati-Saber, J. A. Fax, and R. M. Murray.
\newblock Consensus and cooperation in networked multi-agent systems.
\newblock \emph{Proceedings of the IEEE}, 95(1):215--233, 2007.
\newblock doi:10.1109/JPROC.2006.887293.

\bibitem{Manski2003}
C. F. Manski.
\newblock \emph{Partial Identification of Probability Distributions}.
\newblock Springer, New York, 2003.
\newblock doi:10.1007/b97478.

\bibitem{Pearl2009}
J. Pearl.
\newblock \emph{Causality: Models, Reasoning, and Inference}.
\newblock 2nd edition, Cambridge University Press, Cambridge, 2009.

\bibitem{RiedlEtAl2026}
C. Riedl, S. Savage, and J. Zvelebilova.
\newblock Cognitive spillover in human--AI teams.
\newblock \emph{ACM Transactions on Computer-Human Interaction}, 2026.
\newblock doi:10.1145/3805039; arXiv:2407.17489v3.

\bibitem{Robinson2017}
M. Robinson.
\newblock Sheaves are the canonical data structure for sensor integration.
\newblock \emph{Information Fusion}, 36:208--224, 2017.
\newblock doi:10.1016/j.inffus.2016.12.002.

\bibitem{Rupert2004}
R. D. Rupert.
\newblock Challenges to the hypothesis of extended cognition.
\newblock \emph{The Journal of Philosophy}, 101(8):389--428, 2004.
\newblock doi:10.5840/jphil2004101826.

\bibitem{SeeberEtAl2020}
I. Seeber, E. Bittner, R. O. Briggs, T. de Vreede, G.-J. de Vreede, A. Elkins, R. Maier, A. B. Merz, S. Oeste-Rei\ss, N. Randrup, G. Schwabe, and M. S\"ollner.
\newblock Machines as teammates: A research agenda on AI in team collaboration.
\newblock \emph{Information \& Management}, 57(2):103174, 2020.
\newblock doi:10.1016/j.im.2019.103174.

\bibitem{ShaliziCrutchfield2001}
C. R. Shalizi and J. P. Crutchfield.
\newblock Computational mechanics: Pattern and prediction, structure and simplicity.
\newblock \emph{Journal of Statistical Physics}, 104:817--879, 2001.
\newblock doi:10.1023/A:1010388907793.

\bibitem{SidjiEtAl2024}
M. Sidji, W. Smith, and M. J. Rogerson.
\newblock Adopting the theory of distributed cognition for human--AI cooperation.
\newblock In \emph{Proceedings of the 36th Australasian Conference on Human-Computer Interaction (OzCHI '24)}, pages 774--779, 2024.
\newblock doi:10.1145/3726986.3727024.

\bibitem{SimonsohnEtAl2020}
U. Simonsohn, J. P. Simmons, and L. D. Nelson.
\newblock Specification curve analysis.
\newblock \emph{Nature Human Behaviour}, 4(11):1208--1214, 2020.
\newblock doi:10.1038/s41562-020-0912-z.

\bibitem{SingerWu2012}
A. Singer and H.-T. Wu.
\newblock Vector diffusion maps and the connection Laplacian.
\newblock \emph{Communications on Pure and Applied Mathematics}, 65(8):1067--1144, 2012.
\newblock doi:10.1002/cpa.21395.

\bibitem{SteegenEtAl2016}
S. Steegen, F. Tuerlinckx, A. Gelman, and W. Vanpaemel.
\newblock Increasing transparency through a multiverse analysis.
\newblock \emph{Perspectives on Psychological Science}, 11(5):702--712, 2016.
\newblock doi:10.1177/1745691616658637.

\bibitem{SuttonBarto2018}
R. S. Sutton and A. G. Barto.
\newblock \emph{Reinforcement Learning: An Introduction}.
\newblock 2nd edition, MIT Press, Cambridge, Massachusetts, 2018.

\bibitem{Wegner1987}
D. M. Wegner.
\newblock Transactive memory: A contemporary analysis of the group mind.
\newblock In B. Mullen and G. R. Goethals, editors, \emph{Theories of Group Behavior}, pages 185--208. Springer, New York, 1987.
\newblock doi:10.1007/978-1-4612-4634-3\_9.

\bibitem{WilliamsBeer2010}
P. L. Williams and R. D. Beer.
\newblock Nonnegative decomposition of multivariate information.
\newblock arXiv:1004.2515, 2010.

\bibitem{WoolleyEtAl2010}
A. W. Woolley, C. F. Chabris, A. Pentland, N. Hashmi, and T. W. Malone.
\newblock Evidence for a collective intelligence factor in the performance of human groups.
\newblock \emph{Science}, 330(6004):686--688, 2010.
\newblock doi:10.1126/science.1193147.

\bibitem{ZhaoHan2026}
H. Zhao and Z. Han.
\newblock Understanding the mechanism of human--AI interaction: A distributed cognition perspective.
\newblock \emph{Journal of Documentation}, 82(4):1042--1063, 2026.
\newblock doi:10.1108/JD-01-2026-0003.


\bibitem{ChoudhuryEtAl2026}
S. Choudhury et al.
\newblock Networked intelligence: Active shared context graphs for human--AI team science.
\newblock arXiv:2607.13220v3, 2026.
\newblock doi:10.48550/arXiv.2607.13220.

\bibitem{GoyalEtAl2022}
A. Goyal, A. Didolkar, A. Lamb, K. Badola, N. R. Ke, N. Rahaman, J. Binas,
C. Blundell, M. Mozer, and Y. Bengio.
\newblock Coordination among neural modules through a shared global workspace.
\newblock In \emph{International Conference on Learning Representations}, 2022.
\newblock arXiv:2103.01197v2.

\bibitem{KotalwarDasRose2026}
N. Kotalwar, R. Das, and C. Ros\'e.
\newblock Searching for synergy in shared workspace human--AI collaboration.
\newblock arXiv:2606.18413v3, 2026.
\newblock doi:10.48550/arXiv.2606.18413.

\bibitem{LiJiang2026}
R. Li and B. Jiang.
\newblock Thinking in graphs with CoMAP: A shared visual workspace for designing project-based learning.
\newblock In \emph{Proceedings of the 2026 CHI Conference on Human Factors in Computing Systems}, 2026.
\newblock doi:10.1145/3772318.3790273.

\bibitem{PickeringGarrod2021}
M. J. Pickering and S. Garrod.
\newblock \emph{Understanding Dialogue: Language Use and Social Interaction}.
\newblock Cambridge University Press, Cambridge, 2021.

\end{thebibliography}

\end{document}
